% concatenated from main.tex, math_commands.tex
\documentclass[twoside]{article} % AISTATS 2027 template
%\usepackage{aistats2027}
\usepackage[preprint]{aistats2027}
%%%%% NEW MATH DEFINITIONS %%%%%

\usepackage{amsmath,amsfonts,bm}

% Mark sections of captions for referring to divisions of figures
\newcommand{\figleft}{{\em (Left)}}
\newcommand{\figcenter}{{\em (Center)}}
\newcommand{\figright}{{\em (Right)}}
\newcommand{\figtop}{{\em (Top)}}
\newcommand{\figbottom}{{\em (Bottom)}}
\newcommand{\captiona}{{\em (a)}}
\newcommand{\captionb}{{\em (b)}}
\newcommand{\captionc}{{\em (c)}}
\newcommand{\captiond}{{\em (d)}}

% Highlight a newly defined term
\newcommand{\newterm}[1]{{\bf #1}}


% Figure reference, lower-case.
\def\figref#1{figure~\ref{#1}}
% Figure reference, capital. For start of sentence
\def\Figref#1{Figure~\ref{#1}}
\def\twofigref#1#2{figures \ref{#1} and \ref{#2}}
\def\quadfigref#1#2#3#4{figures \ref{#1}, \ref{#2}, \ref{#3} and \ref{#4}}
% Section reference, lower-case.
\def\secref#1{section~\ref{#1}}
% Section reference, capital.
\def\Secref#1{Section~\ref{#1}}
% Reference to two sections.
\def\twosecrefs#1#2{sections \ref{#1} and \ref{#2}}
% Reference to three sections.
\def\secrefs#1#2#3{sections \ref{#1}, \ref{#2} and \ref{#3}}
% Reference to an equation, lower-case.
\def\eqref#1{equation~\ref{#1}}
% Reference to an equation, upper case
\def\Eqref#1{Equation~\ref{#1}}
% A raw reference to an equation---avoid using if possible
\def\plaineqref#1{\ref{#1}}
% Reference to a chapter, lower-case.
\def\chapref#1{chapter~\ref{#1}}
% Reference to an equation, upper case.
\def\Chapref#1{Chapter~\ref{#1}}
% Reference to a range of chapters
\def\rangechapref#1#2{chapters\ref{#1}--\ref{#2}}
% Reference to an algorithm, lower-case.
\def\algref#1{algorithm~\ref{#1}}
% Reference to an algorithm, upper case.
\def\Algref#1{Algorithm~\ref{#1}}
\def\twoalgref#1#2{algorithms \ref{#1} and \ref{#2}}
\def\Twoalgref#1#2{Algorithms \ref{#1} and \ref{#2}}
% Reference to a part, lower case
\def\partref#1{part~\ref{#1}}
% Reference to a part, upper case
\def\Partref#1{Part~\ref{#1}}
\def\twopartref#1#2{parts \ref{#1} and \ref{#2}}

\def\ceil#1{\lceil #1 \rceil}
\def\floor#1{\lfloor #1 \rfloor}
\def\1{\bm{1}}
\newcommand{\train}{\mathcal{D}}
\newcommand{\valid}{\mathcal{D_{\mathrm{valid}}}}
\newcommand{\test}{\mathcal{D_{\mathrm{test}}}}

\def\eps{{\epsilon}}


% Random variables
\def\reta{{\textnormal{$\eta$}}}
\def\ra{{\textnormal{a}}}
\def\rb{{\textnormal{b}}}
\def\rc{{\textnormal{c}}}
\def\rd{{\textnormal{d}}}
\def\re{{\textnormal{e}}}
\def\rf{{\textnormal{f}}}
\def\rg{{\textnormal{g}}}
\def\rh{{\textnormal{h}}}
\def\ri{{\textnormal{i}}}
\def\rj{{\textnormal{j}}}
\def\rk{{\textnormal{k}}}
\def\rl{{\textnormal{l}}}
% rm is already a command, just don't name any random variables m
\def\rn{{\textnormal{n}}}
\def\ro{{\textnormal{o}}}
\def\rp{{\textnormal{p}}}
\def\rq{{\textnormal{q}}}
\def\rr{{\textnormal{r}}}
\def\rs{{\textnormal{s}}}
\def\rt{{\textnormal{t}}}
\def\ru{{\textnormal{u}}}
\def\rv{{\textnormal{v}}}
\def\rw{{\textnormal{w}}}
\def\rx{{\textnormal{x}}}
\def\ry{{\textnormal{y}}}
\def\rz{{\textnormal{z}}}

% Random vectors
\def\rvepsilon{{\mathbf{\epsilon}}}
\def\rvtheta{{\mathbf{\theta}}}
\def\rva{{\mathbf{a}}}
\def\rvb{{\mathbf{b}}}
\def\rvc{{\mathbf{c}}}
\def\rvd{{\mathbf{d}}}
\def\rve{{\mathbf{e}}}
\def\rvf{{\mathbf{f}}}
\def\rvg{{\mathbf{g}}}
\def\rvh{{\mathbf{h}}}
\def\rvu{{\mathbf{i}}}
\def\rvj{{\mathbf{j}}}
\def\rvk{{\mathbf{k}}}
\def\rvl{{\mathbf{l}}}
\def\rvm{{\mathbf{m}}}
\def\rvn{{\mathbf{n}}}
\def\rvo{{\mathbf{o}}}
\def\rvp{{\mathbf{p}}}
\def\rvq{{\mathbf{q}}}
\def\rvr{{\mathbf{r}}}
\def\rvs{{\mathbf{s}}}
\def\rvt{{\mathbf{t}}}
\def\rvu{{\mathbf{u}}}
\def\rvv{{\mathbf{v}}}
\def\rvw{{\mathbf{w}}}
\def\rvx{{\mathbf{x}}}
\def\rvy{{\mathbf{y}}}
\def\rvz{{\mathbf{z}}}

% Elements of random vectors
\def\erva{{\textnormal{a}}}
\def\ervb{{\textnormal{b}}}
\def\ervc{{\textnormal{c}}}
\def\ervd{{\textnormal{d}}}
\def\erve{{\textnormal{e}}}
\def\ervf{{\textnormal{f}}}
\def\ervg{{\textnormal{g}}}
\def\ervh{{\textnormal{h}}}
\def\ervi{{\textnormal{i}}}
\def\ervj{{\textnormal{j}}}
\def\ervk{{\textnormal{k}}}
\def\ervl{{\textnormal{l}}}
\def\ervm{{\textnormal{m}}}
\def\ervn{{\textnormal{n}}}
\def\ervo{{\textnormal{o}}}
\def\ervp{{\textnormal{p}}}
\def\ervq{{\textnormal{q}}}
\def\ervr{{\textnormal{r}}}
\def\ervs{{\textnormal{s}}}
\def\ervt{{\textnormal{t}}}
\def\ervu{{\textnormal{u}}}
\def\ervv{{\textnormal{v}}}
\def\ervw{{\textnormal{w}}}
\def\ervx{{\textnormal{x}}}
\def\ervy{{\textnormal{y}}}
\def\ervz{{\textnormal{z}}}

% Random matrices
\def\rmA{{\mathbf{A}}}
\def\rmB{{\mathbf{B}}}
\def\rmC{{\mathbf{C}}}
\def\rmD{{\mathbf{D}}}
\def\rmE{{\mathbf{E}}}
\def\rmF{{\mathbf{F}}}
\def\rmG{{\mathbf{G}}}
\def\rmH{{\mathbf{H}}}
\def\rmI{{\mathbf{I}}}
\def\rmJ{{\mathbf{J}}}
\def\rmK{{\mathbf{K}}}
\def\rmL{{\mathbf{L}}}
\def\rmM{{\mathbf{M}}}
\def\rmN{{\mathbf{N}}}
\def\rmO{{\mathbf{O}}}
\def\rmP{{\mathbf{P}}}
\def\rmQ{{\mathbf{Q}}}
\def\rmR{{\mathbf{R}}}
\def\rmS{{\mathbf{S}}}
\def\rmT{{\mathbf{T}}}
\def\rmU{{\mathbf{U}}}
\def\rmV{{\mathbf{V}}}
\def\rmW{{\mathbf{W}}}
\def\rmX{{\mathbf{X}}}
\def\rmY{{\mathbf{Y}}}
\def\rmZ{{\mathbf{Z}}}

% Elements of random matrices
\def\ermA{{\textnormal{A}}}
\def\ermB{{\textnormal{B}}}
\def\ermC{{\textnormal{C}}}
\def\ermD{{\textnormal{D}}}
\def\ermE{{\textnormal{E}}}
\def\ermF{{\textnormal{F}}}
\def\ermG{{\textnormal{G}}}
\def\ermH{{\textnormal{H}}}
\def\ermI{{\textnormal{I}}}
\def\ermJ{{\textnormal{J}}}
\def\ermK{{\textnormal{K}}}
\def\ermL{{\textnormal{L}}}
\def\ermM{{\textnormal{M}}}
\def\ermN{{\textnormal{N}}}
\def\ermO{{\textnormal{O}}}
\def\ermP{{\textnormal{P}}}
\def\ermQ{{\textnormal{Q}}}
\def\ermR{{\textnormal{R}}}
\def\ermS{{\textnormal{S}}}
\def\ermT{{\textnormal{T}}}
\def\ermU{{\textnormal{U}}}
\def\ermV{{\textnormal{V}}}
\def\ermW{{\textnormal{W}}}
\def\ermX{{\textnormal{X}}}
\def\ermY{{\textnormal{Y}}}
\def\ermZ{{\textnormal{Z}}}

% Vectors
\def\vzero{{\bm{0}}}
\def\vone{{\bm{1}}}
\def\vmu{{\bm{\mu}}}
\def\vtheta{{\bm{\theta}}}
\def\va{{\bm{a}}}
\def\vb{{\bm{b}}}
\def\vc{{\bm{c}}}
\def\vd{{\bm{d}}}
\def\ve{{\bm{e}}}
\def\vf{{\bm{f}}}
\def\vg{{\bm{g}}}
\def\vh{{\bm{h}}}
\def\vi{{\bm{i}}}
\def\vj{{\bm{j}}}
\def\vk{{\bm{k}}}
\def\vl{{\bm{l}}}
\def\vm{{\bm{m}}}
\def\vn{{\bm{n}}}
\def\vo{{\bm{o}}}
\def\vp{{\bm{p}}}
\def\vq{{\bm{q}}}
\def\vr{{\bm{r}}}
\def\vs{{\bm{s}}}
\def\vt{{\bm{t}}}
\def\vu{{\bm{u}}}
\def\vv{{\bm{v}}}
\def\vw{{\bm{w}}}
\def\vx{{\bm{x}}}
\def\vy{{\bm{y}}}
\def\vz{{\bm{z}}}

% Elements of vectors
\def\evalpha{{\alpha}}
\def\evbeta{{\beta}}
\def\evepsilon{{\epsilon}}
\def\evlambda{{\lambda}}
\def\evomega{{\omega}}
\def\evmu{{\mu}}
\def\evpsi{{\psi}}
\def\evsigma{{\sigma}}
\def\evtheta{{\theta}}
\def\eva{{a}}
\def\evb{{b}}
\def\evc{{c}}
\def\evd{{d}}
\def\eve{{e}}
\def\evf{{f}}
\def\evg{{g}}
\def\evh{{h}}
\def\evi{{i}}
\def\evj{{j}}
\def\evk{{k}}
\def\evl{{l}}
\def\evm{{m}}
\def\evn{{n}}
\def\evo{{o}}
\def\evp{{p}}
\def\evq{{q}}
\def\evr{{r}}
\def\evs{{s}}
\def\evt{{t}}
\def\evu{{u}}
\def\evv{{v}}
\def\evw{{w}}
\def\evx{{x}}
\def\evy{{y}}
\def\evz{{z}}

% Matrix
\def\mA{{\bm{A}}}
\def\mB{{\bm{B}}}
\def\mC{{\bm{C}}}
\def\mD{{\bm{D}}}
\def\mE{{\bm{E}}}
\def\mF{{\bm{F}}}
\def\mG{{\bm{G}}}
\def\mH{{\bm{H}}}
\def\mI{{\bm{I}}}
\def\mJ{{\bm{J}}}
\def\mK{{\bm{K}}}
\def\mL{{\bm{L}}}
\def\mM{{\bm{M}}}
\def\mN{{\bm{N}}}
\def\mO{{\bm{O}}}
\def\mP{{\bm{P}}}
\def\mQ{{\bm{Q}}}
\def\mR{{\bm{R}}}
\def\mS{{\bm{S}}}
\def\mT{{\bm{T}}}
\def\mU{{\bm{U}}}
\def\mV{{\bm{V}}}
\def\mW{{\bm{W}}}
\def\mX{{\bm{X}}}
\def\mY{{\bm{Y}}}
\def\mZ{{\bm{Z}}}
\def\mBeta{{\bm{\beta}}}
\def\mPhi{{\bm{\Phi}}}
\def\mLambda{{\bm{\Lambda}}}
\def\mSigma{{\bm{\Sigma}}}

% Tensor
\DeclareMathAlphabet{\mathsfit}{\encodingdefault}{\sfdefault}{m}{sl}
\SetMathAlphabet{\mathsfit}{bold}{\encodingdefault}{\sfdefault}{bx}{n}
\newcommand{\tens}[1]{\bm{\mathsfit{#1}}}
\def\tA{{\tens{A}}}
\def\tB{{\tens{B}}}
\def\tC{{\tens{C}}}
\def\tD{{\tens{D}}}
\def\tE{{\tens{E}}}
\def\tF{{\tens{F}}}
\def\tG{{\tens{G}}}
\def\tH{{\tens{H}}}
\def\tI{{\tens{I}}}
\def\tJ{{\tens{J}}}
\def\tK{{\tens{K}}}
\def\tL{{\tens{L}}}
\def\tM{{\tens{M}}}
\def\tN{{\tens{N}}}
\def\tO{{\tens{O}}}
\def\tP{{\tens{P}}}
\def\tQ{{\tens{Q}}}
\def\tR{{\tens{R}}}
\def\tS{{\tens{S}}}
\def\tT{{\tens{T}}}
\def\tU{{\tens{U}}}
\def\tV{{\tens{V}}}
\def\tW{{\tens{W}}}
\def\tX{{\tens{X}}}
\def\tY{{\tens{Y}}}
\def\tZ{{\tens{Z}}}


% Graph
\def\gA{{\mathcal{A}}}
\def\gB{{\mathcal{B}}}
\def\gC{{\mathcal{C}}}
\def\gD{{\mathcal{D}}}
\def\gE{{\mathcal{E}}}
\def\gF{{\mathcal{F}}}
\def\gG{{\mathcal{G}}}
\def\gH{{\mathcal{H}}}
\def\gI{{\mathcal{I}}}
\def\gJ{{\mathcal{J}}}
\def\gK{{\mathcal{K}}}
\def\gL{{\mathcal{L}}}
\def\gM{{\mathcal{M}}}
\def\gN{{\mathcal{N}}}
\def\gO{{\mathcal{O}}}
\def\gP{{\mathcal{P}}}
\def\gQ{{\mathcal{Q}}}
\def\gR{{\mathcal{R}}}
\def\gS{{\mathcal{S}}}
\def\gT{{\mathcal{T}}}
\def\gU{{\mathcal{U}}}
\def\gV{{\mathcal{V}}}
\def\gW{{\mathcal{W}}}
\def\gX{{\mathcal{X}}}
\def\gY{{\mathcal{Y}}}
\def\gZ{{\mathcal{Z}}}

% Sets
\def\sA{{\mathbb{A}}}
\def\sB{{\mathbb{B}}}
\def\sC{{\mathbb{C}}}
\def\sD{{\mathbb{D}}}
% Don't use a set called E, because this would be the same as our symbol
% for expectation.
\def\sF{{\mathbb{F}}}
\def\sG{{\mathbb{G}}}
\def\sH{{\mathbb{H}}}
\def\sI{{\mathbb{I}}}
\def\sJ{{\mathbb{J}}}
\def\sK{{\mathbb{K}}}
\def\sL{{\mathbb{L}}}
\def\sM{{\mathbb{M}}}
\def\sN{{\mathbb{N}}}
\def\sO{{\mathbb{O}}}
\def\sP{{\mathbb{P}}}
\def\sQ{{\mathbb{Q}}}
\def\sR{{\mathbb{R}}}
\def\sS{{\mathbb{S}}}
\def\sT{{\mathbb{T}}}
\def\sU{{\mathbb{U}}}
\def\sV{{\mathbb{V}}}
\def\sW{{\mathbb{W}}}
\def\sX{{\mathbb{X}}}
\def\sY{{\mathbb{Y}}}
\def\sZ{{\mathbb{Z}}}

% Entries of a matrix
\def\emLambda{{\Lambda}}
\def\emA{{A}}
\def\emB{{B}}
\def\emC{{C}}
\def\emD{{D}}
\def\emE{{E}}
\def\emF{{F}}
\def\emG{{G}}
\def\emH{{H}}
\def\emI{{I}}
\def\emJ{{J}}
\def\emK{{K}}
\def\emL{{L}}
\def\emM{{M}}
\def\emN{{N}}
\def\emO{{O}}
\def\emP{{P}}
\def\emQ{{Q}}
\def\emR{{R}}
\def\emS{{S}}
\def\emT{{T}}
\def\emU{{U}}
\def\emV{{V}}
\def\emW{{W}}
\def\emX{{X}}
\def\emY{{Y}}
\def\emZ{{Z}}
\def\emSigma{{\Sigma}}

% entries of a tensor
% Same font as tensor, without \bm wrapper
\newcommand{\etens}[1]{\mathsfit{#1}}
\def\etLambda{{\etens{\Lambda}}}
\def\etA{{\etens{A}}}
\def\etB{{\etens{B}}}
\def\etC{{\etens{C}}}
\def\etD{{\etens{D}}}
\def\etE{{\etens{E}}}
\def\etF{{\etens{F}}}
\def\etG{{\etens{G}}}
\def\etH{{\etens{H}}}
\def\etI{{\etens{I}}}
\def\etJ{{\etens{J}}}
\def\etK{{\etens{K}}}
\def\etL{{\etens{L}}}
\def\etM{{\etens{M}}}
\def\etN{{\etens{N}}}
\def\etO{{\etens{O}}}
\def\etP{{\etens{P}}}
\def\etQ{{\etens{Q}}}
\def\etR{{\etens{R}}}
\def\etS{{\etens{S}}}
\def\etT{{\etens{T}}}
\def\etU{{\etens{U}}}
\def\etV{{\etens{V}}}
\def\etW{{\etens{W}}}
\def\etX{{\etens{X}}}
\def\etY{{\etens{Y}}}
\def\etZ{{\etens{Z}}}

% The true underlying data generating distribution
\newcommand{\pdata}{p_{\rm{data}}}
% The empirical distribution defined by the training set
\newcommand{\ptrain}{\hat{p}_{\rm{data}}}
\newcommand{\Ptrain}{\hat{P}_{\rm{data}}}
% The model distribution
\newcommand{\pmodel}{p_{\rm{model}}}
\newcommand{\Pmodel}{P_{\rm{model}}}
\newcommand{\ptildemodel}{\tilde{p}_{\rm{model}}}
% Stochastic autoencoder distributions
\newcommand{\pencode}{p_{\rm{encoder}}}
\newcommand{\pdecode}{p_{\rm{decoder}}}
\newcommand{\precons}{p_{\rm{reconstruct}}}

\newcommand{\laplace}{\mathrm{Laplace}} % Laplace distribution

\newcommand{\E}{\mathbb{E}}
\newcommand{\Ls}{\mathcal{L}}
\newcommand{\R}{\mathbb{R}}
\newcommand{\emp}{\tilde{p}}
\newcommand{\lr}{\alpha}
\newcommand{\reg}{\lambda}
\newcommand{\rect}{\mathrm{rectifier}}
\newcommand{\softmax}{\mathrm{softmax}}
\newcommand{\sigmoid}{\sigma}
\newcommand{\softplus}{\zeta}
\newcommand{\KL}{D_{\mathrm{KL}}}
\newcommand{\Var}{\mathrm{Var}}
\newcommand{\standarderror}{\mathrm{SE}}
\newcommand{\Cov}{\mathrm{Cov}}
% Wolfram Mathworld says $L^2$ is for function spaces and $\ell^2$ is for vectors
% But then they seem to use $L^2$ for vectors throughout the site, and so does
% wikipedia.
\newcommand{\normlzero}{L^0}
\newcommand{\normlone}{L^1}
\newcommand{\normltwo}{L^2}
\newcommand{\normlp}{L^p}
\newcommand{\normmax}{L^\infty}

\newcommand{\parents}{Pa} % See usage in notation.tex. Chosen to match Daphne's book.

\DeclareMathOperator*{\argmax}{arg\,max}
\DeclareMathOperator*{\argmin}{arg\,min}

\DeclareMathOperator{\sign}{sign}
\DeclareMathOperator{\Tr}{Tr}
\let\ab\allowbreak

\usepackage[round]{natbib}
\renewcommand{\bibname}{References}
\renewcommand{\bibsection}{\subsubsection*{\bibname}}

\usepackage{hyperref}
\usepackage{url}

\usepackage{algorithm}
\usepackage{algorithmic}
\usepackage{amssymb}
\usepackage{booktabs,array,longtable}

\usepackage[utf8]{inputenc} % allow utf-8 input
\usepackage[T1]{fontenc}    % use 8-bit T1 fonts

\usepackage{url}            % simple URL typesetting
\usepackage{amsfonts}       % blackboard math symbols
\usepackage{nicefrac}       % compact symbols for 1/2, etc.
\usepackage{microtype}      % microtypography
\usepackage{xcolor}         % colors

\usepackage{etoc}


% The \author macro works with any number of authors. There are two commands
% used to separate the names and addresses of multiple authors: \And and \AND.
%
% Using \And between authors leaves it to LaTeX to determine where to break the
% lines. Using \AND forces a line break at that point. So, if LaTeX puts 3 of 4
% authors names on the first line, and the last on the second line, try using
% \AND instead of \And before the third author name.



% Recommended, but optional, packages for figures and better typesetting:
\usepackage{microtype}
\usepackage{graphicx}
\usepackage{subcaption}
\usepackage{booktabs}

\usepackage{hyperref}
\hypersetup{hidelinks}

\usepackage{amsmath}
\usepackage{amssymb}
\usepackage{mathtools}

\allowdisplaybreaks
\usepackage{wrapfig}
\usepackage{array,enumitem}
\usepackage{thmtools,thm-restate}
\usepackage{amsthm}
\usepackage{multirow}
% if you use cleveref..
%\usepackage[capitalize,noabbrev]{cleveref}

%%%%%%%%%%%%%%%%%%%%%%%%%%%%%%%%
% THEOREMS
%%%%%%%%%%%%%%%%%%%%%%%%%%%%%%%%
\theoremstyle{plain}
\newtheorem{theorem}{Theorem}[section]
\newtheorem{proposition}[theorem]{Proposition}
\newtheorem{lemma}[theorem]{Lemma}
\newtheorem{corollary}[theorem]{Corollary}
\theoremstyle{definition}
\newtheorem{definition}[theorem]{Definition}
\newtheorem{assumption}[theorem]{Assumption}
\theoremstyle{remark}
\newtheorem{remark}[theorem]{Remark}
%\newtheorem{proof}[theorem]{Proof}
\newtheorem{condition}[theorem]{Condition}
\newtheorem{example}[theorem]{Example}





% Make the "Part I" text invisible
\renewcommand \thepart{}
\renewcommand \partname{}


%% Red for convergence analysis
\newcommand{\red}[1]{\textcolor{red}{#1}}
%% brown for high-probability bound
\newcommand{\brown}[1]{\textcolor{brown}{#1}}
%% blue for others
\newcommand{\blue}[1]{\textcolor{blue}{#1}}
\newcommand{\Y}{Y}
\newcommand{\Yt}{\tilde{Y}}
\newcommand{\Yb}{\bar{Y}}
%\newcommand{\X}{X}
\newcommand{\HGamma}{\hat{\Gamma}}
\newcommand{\HUpsilon}{\hat{\Upsilon}}
%\newcommand{\Xt}{\tilde{X}}
%\newcommand{\Xb}{\bar{X}}
\newcommand{\tepsilon}{\tilde{\epsilon}}
\newcommand{\hL}{\hat{L}}
\newcommand{\Hambda}{\hat{\Lambda}}
\newcommand{\HK}{\hat{K}}
\newcommand{\htau}{\hat{\tau}}
\newcommand{\fL}{\mathcal{L}}
\newcommand{\fhL}{\hat{\mathcal{L}}}
\newcommand{\calG}{\mathcal{G}}
\newcommand{\calH}{\mathcal{H}}
\newcommand{\calhatE}{\mathcal{\hat{E}}}





\newcommand{\fix}{\marginpar{FIX}}
\newcommand{\new}{\marginpar{NEW}}

\begin{document}
\runningtitle{ Distributionally Robust Multi-Objective Optimization}

\twocolumn[

\aistatstitle{From Dual Tracking to Clipping: Provably Faster Distributionally Robust Multi-Objective Optimization}

\aistatsauthor{
  Yufeng Yang\textsuperscript{1} \quad
  Fangning Zhuo\textsuperscript{1} \quad
  Ziyi Chen\textsuperscript{2} \quad
  Heng Huang\textsuperscript{3} \quad
  Yi Zhou\textsuperscript{1}
}

\aistatsaddress{
  \textsuperscript{1}Texas A\&M University \quad
  \textsuperscript{2}A*STAR \quad
  \textsuperscript{3}University of Maryland
}
]

\begingroup
\renewcommand{\thefootnote}{\fnsymbol{footnote}}
\footnotetext[1]{Preprint Under Review}
\footnotetext[1]{Email: \texttt{yufeng.yang@tamu.edu}.}
\footnotetext[1]{Email: \texttt{yi.zhou@tamu.edu}.}
\endgroup




\begin{abstract}
  % Multi-objective optimization (MOO) has attracted growing attention across various fields, particularly in settings that require learning with multiple criteria or tasks. However, the standard MOO remain limited in considering distributional shifts over data. In this paper, we formulate a novel distributionally-robust multi-objective optimization (DR-MOO) framework by minimizing multiple objectives under their worst-case distributions. We propose new Pareto-type solution concepts for DR-MOO and design multi-gradient descent algorithms (MGDA). By exploiting a Lagrangian dual formulation of DR-MOO, we design a double-loop MGDA that uses an inner-loop to estimate the dual variables and achieves a total sample complexity $\mathcal{O}(\epsilon^{-12})$ to achieve an $\epsilon$-Pareto-stationary point. Moreover, by incorporating gradient clipping into the algorithm design effectively mitigates biased gradient challenge and removes the need for double sampling. We propose a single-loop double-clip MGDA that achieves significantly improved sample complexity $\mathcal{O}(\epsilon^{-4})$. Our analysis holds under relaxed assumptions, without requiring convexity or boundedness of the objectives or their gradients. Experiments demonstrate competitive performance of our algorithms against existing MGDA methods.
  Multi-objective optimization (MOO) has received growing attention in applications that require learning under multiple criteria. However, most existing MOO formulations do not explicitly account for distributional shifts in the data. We introduce distributionally robust multi-objective optimization (DR-MOO), which minimizes multiple objectives under their respective worst-case distributions. We propose Pareto-type solution concepts for DR-MOO and develop multi-gradient descent algorithms (MGDA) with provable guarantees. Leveraging a Lagrangian dual reformulation, we first design a double-loop MGDA that uses an inner loop to estimate dual variables and achieves a total sample complexity \(\mathcal{O}(\epsilon^{-8})\) for reaching an \(\epsilon\)-Pareto-stationary point. To further improve convergence, we combine large-batch sampling with gradient clipping to accommodate generalized smoothness and control bias in stochastic preference updates, eliminating the need for double sampling.  This yields a single-loop double-clip MGDA with substantially improved sample complexity \(\mathcal{O}(\epsilon^{-4})\). Our theory applies to nonconvex problems without requiring uniformly bounded gradients of the dual objectives. Experiments demonstrate that our methods are competitive with state-of-the-art MGDA baselines.
\end{abstract}

%% To be added in future version:
%1. Motivation on why DR-MOO is important?
%2.(Done) Table of Comparisons on existing MOO methods.
%3. Highlight disadvantage, and identify its solovability in practice (reflected on Ablation study)
%4.(Done) Add additonal related work section on considering robustness on MTL.
%5.(Done) Existing DRO results (List on More related work section.)
%6.(Important) Rewrite main theorem statements to include problem-dependent hyparameters.
%7. Brieflly mention limitations and future work on non-linear scalarization methods mentioned by reviewer. 

\section{Introduction}

%\textbf{Identify Challenge, Motivation and weakness (How to address) - checklist.
%0. Motivation on why studying DR-MOO? From perspective of Network structure. (Done, in introduction)
%1. Primal DRO includes distributionally uncertainty, this makes Pareto-concept no-longer apply
%2. Dual DRO introduces dual variable, making problem becomes a nested stochastic programming
%3. Dual DR-MOO’s loss is unbounded, the gradient noise is affine linear, this severely increase the confidencial parameters order: The high-order dependence comes from controlling affine-linear, unbounded stochastic noise under high probability.
%1. For algorithm1, theoretical limitation is high-inner-loop complexity, however, in practice, this is much better (link to ablation study, wall-clock time comparison and inner-step ablation)
%2. For algorithm2, theoretical limitation is batch size, link to ablation study showing that B=256 in practice is good enough to ensure stability and performance. At last, highlight potential extension of algorithm2 on introducing advanced normalized methods in MOO setting.}

%{Writting logic: 1.Introduce MOO and Solution Concepts; 2. Identify stochastic MOO hasn't consider distribution shifts in practical scenarios; 3. Impose first research question on Formulation DR-MOO; 4. Overview State-of-art MOO solvers; 5. Impose limitation and second research question: Algorithms for DR-MOO; 5. Answer two research question in contribution subsection.}
% Multi-objective optimization (MOO) has received intensive attentions in application fields including automatic driving~\citep{automativedriving}, AI4Science~\citep{ai4science}, healthcare~\cite{healthcare} and recommendation systems~\cite{recommendation}. Mathematically, MOO seeks to optimize $m \in \mathbf{N}$ potentially conflicting objectives simultaneously, as formulated in \eqref{eq: MOO}.
% \begin{align}
%    \min_{{\theta} \in \mathbf{R}^n} \Big\{\Phi({\theta}):=\big[\phi^1({\theta}),\phi^2({\theta}), ..., \phi^m({\theta})\big] \Big\},
%    \label{eq: MOO}
% \end{align}
% where the $i$-th objective usually takes the stochastic form $\phi^i({\theta}):=\mathbb{E}_{\xi\sim\mathbb{P}^i}[\ell^i({\theta}, \xi)]$ and ${\theta}$ denotes the model parameters. 
% %introduce Pareto-solution concepts
% However, directly solving the vector-valued objective~\eqref{eq: MOO} is challenging due to disparities in both gradient magnitudes and directions such that objectives with larger gradient magnitudes tend to dominate the update direction, which significantly degrade performance of less-favored objectives with smaller gradients. Thus, a widely-adopted target is find a Pareto-stationary point $\theta$ such that there is no common descent directions for all objectives. 
% While Pareto-type solution concepts effectively capture trade-offs among multiple criteria and admit tractable algorithmic solutions~\citep{1stempiricalMGDAwork,paretoMTL,stochasticMGDA}, existing works typically rely on the implicit assumption that data are drawn from stationary distributions. This assumption is often violated in practice and may lead to degraded performance under distributional shifts. For example, when training a multi-objective model~\eqref{eq: MOO} using data from multiple sources, shifts may naturally arise due to class imbalance, domain mismatch, or adversarial perturbations.

% To solve this problem, distributionally robust optimization (DRO)~\cite{DRO1st,DRO2nd,DRO3rd} has been extensively studied. DRO formulates a min–max optimization problem, seeking robust model parameters $\theta$ by minimizing the objective under worst-case data distribution within an ambiguity set, where deviations from a nominal distribution are quantified via information divergences such as $f$-divergence and Wasserstein distance. While these divergence-based formulations enable tractable reformulations and efficient algorithm design in single-objective settings, their extension to MOO~\eqref{eq: MOO}, along with Pareto-based solution concepts, remains unexplored. This naturally raises the question
% \begin{itemize}[leftmargin=*]
% \item \textit{Can we formulate distributionally robust multi-objective optimization (DR-MOO) in a manner that admits Pareto-type solution concepts?}
% \end{itemize}

Multi-objective optimization (MOO) has attracted increasing interest in applications such as autonomous driving \citep{automativedriving}, AI for science \citep{ai4science}, recommendation systems \citep{recommendation} and LLM Alignment~\citep{RACO}. Mathematically, MOO aims to optimize a set of \(m\) objectives simultaneously:
\begin{align}
    \min_{\theta \in \mathbf{R}^n}\ \Big\{\Phi(\theta):=\big[\phi^1(\theta),\phi^2(\theta),\ldots,\phi^m(\theta)\big]\Big\},\label{eq: MOO}
\end{align}
%\begin{restatable}{equation}{mooproblem}
%\text{(MOO):}~\min_{\theta \in \mathbf{R}^n}\ \Big\{\Phi(\theta):=\big[\phi^1(\theta), \phi^2(\theta), \ldots, \phi^m(\theta)\big]\Big\},\nonumber
%\end{restatable}
where each objective $i$ has a stochastic form \(\phi^i(\theta):=\mathbb{E}_{\xi\sim\mathbb{P}^i}[\ell^i(\theta,\xi)]\), and \(\theta\) denotes the model parameters.
A core challenge in solving MOO is that objectives may exhibit highly disparate gradient magnitudes and directions, so naive gradient-based updates can over-emphasize certain objectives and harm others. Consequently, much of the existing MOO literature focuses on Pareto-type solutions (e.g., Pareto stationary), which capture intrinsic trade-offs and admit tractable first-order algorithms \citep{1stempiricalMGDAwork,paretoMTL,stochasticMGDA}.

However, modern multi-task learning (MTL)~\citep{MTL-DL-survey} commonly adopts a shared representation with task-specific prediction heads. This representation-based structure is naturally fragile to task-dependent distribution shifts. Different tasks may suffer from adversarial perturbations, class imbalance or domain mismatch. Moreover, \citet{robustMTL5} showed that adversarially trained shared representations can theoretically improve generalization to new target tasks, suggesting that robust training is critical for improving reliability of MTL. From an optimization viewpoint, such distribution shifts challenge the standard Pareto optimality notion used in MOO.
In particular, a solution that is Pareto-optimal under the nominal task distributions may fail to remain Pareto-optimal when each objective is evaluated under its own shifted distribution, as illustrated by the toy example in Figure~\ref{fig: Toy-example Pareto-visualization}( Appendix~\ref{Appendix: Toy exmple visualization}). Distributionally robust optimization (DRO)~\citep{DRO1st,DRO2nd,DRO3rd} provides a principled framework for handling distribution shift by optimizing performance under the worst-case distribution within a divergence-based ambiguity set. While DRO is well understood for a single objective, extending it to MOO is challenging: one must model objective-wise distribution shifts and define meaningful, analytically tractable Pareto-type solutions under worst-case perturbations. This motivates the following questions:
\begin{itemize}[leftmargin = *, topsep = 0pt]
\item \textit{How should we formulate distributionally robust MOO (DR-MOO) to model objective-wise distribution shifts? What Pareto-type notions are both distributionally robust and algorithmically tractable?}
\end{itemize}
%On the algorithmic side, the MOO literature~\citep{gradient-balance-MTL} has largely been built around gradient-balancing methods, most notably MGDA \citep{1stempiricalMGDAwork,MGDA1stwork,stochasticMGDA} and its variants, with its convergence guarantees established for standard expected-risk objectives. However, standard MOO formulations typically optimize objectives under fixed nominal distributions \(\{\mathbb{P}^i\}_{i=1}^m\) and do not directly extend to DR-MOO. Distributional robustness introduces nested worst-case objectives whose dual reformulations require auxiliary variables and fundamentally alter the optimization geometry. Even when the underlying losses are uniformly Lipschitz continuous and smooth, the resulting dual objectives may satisfy neither property globally. Consequently, existing MGDA analyses and algorithmic designs do not readily apply, and naively combining MGDA with DRO subroutines can lead to biased gradients and poor convergence. These challenges motivate two algorithmic questions. First, MGDA must be adapted to the nested dual structure of DR-MOO to obtain valid convergence guarantees. Second, normalized and clipped gradient methods have proved effective for single-objective nonconvex DRO~\citep{jinnonconvexdro,qirevistfDRO}, raising the question of whether these designs can be extended to multi-objective gradient balancing:
%\begin{itemize}[leftmargin=*, topsep=0pt]
%\item \textit{How can MGDA be adapted to the dual structure of DR-MOO while retaining non-asymptotic convergence guarantees?}
%\item \textit{Can normalization and clipping techniques for single-objective DRO be extended to multi-objective gradient balancing to yield more efficient DR-MOO algorithms?}
%\end{itemize}
%We answer both questions affirmatively by developing two MGDA-type methods: a double-loop method that accommodates the dual DRO structure and a single-loop double-clip method that exploits its generalized-smooth geometry. To the best of our knowledge, this work is the first to systematically formulate nonconvex DR-MOO and establish non-asymptotic MGDA-type convergence guarantees. Table~\ref{Table: Comparison with existing methods} in Appendix~\ref{Appendix: comparison with other algorithms} compares our assumptions and guarantees with those of existing MOO methods.

On the algorithmic side, the MOO literature~\citep{gradient-balance-MTL} centers on gradient-balancing methods, particularly MGDA~\citep{1stempiricalMGDAwork,MGDA1stwork,stochasticMGDA}, whose guarantees largely concern standard expected-risk objectives under fixed nominal distributions. DR-MOO instead introduces nested worst-case objectives and auxiliary dual variables, fundamentally altering the optimization geometry. Even with uniformly Lipschitz continuous and smooth losses, the resulting dual objectives need not satisfy either property globally. Thus, existing MGDA analyses do not readily apply, while naive combinations of MGDA and DRO subroutines may yield biased gradients and poor convergence. Building on the effectiveness of normalized and clipped methods for single-objective nonconvex DRO~\citep{jinnonconvexdro,qirevistfDRO}, these challenges motivate two algorithmic questions:
\begin{itemize}[leftmargin=*, topsep=0pt]
\item \textit{How can MGDA be adapted to the dual structure of DR-MOO while retaining non-asymptotic convergence guarantees?}
\item \textit{Can normalization and clipping techniques for single-objective DRO be extended to multi-objective gradient balancing to yield more efficient DR-MOO algorithms?}
\end{itemize}
We answer both questions by developing a double-loop method that accommodates the dual DRO structure and a single-loop double-clip method that exploits its generalized-smooth geometry. To the best of our knowledge, this is the first systematic formulation of nonconvex DR-MOO with non-asymptotic MGDA-type guarantees. Appendix~\ref{Appendix: comparison with other algorithms} provides a detailed comparison with existing MOO methods in Table~\ref{Table: Comparison with existing methods}.

%%%%%% Highlight what challenges DR-MOO induce? and How we solve it?
\textbf{Our Contributions}
\begin{enumerate}[leftmargin = *, topsep = 0pt]
    \item We formulate $f$-divergence-regularized DR-MOO and introduce distributionally robust Pareto dominance and optimality. By leveraging strong duality to obtain a finite-dimensional reformulation, we translate these concepts into a form amenable to first-order optimization. This leads to an algorithmically tractable notion of distributionally robust Pareto stationarity, which serves as our convergence criterion.
    
    \item Leveraging this structure, we develop a double-loop MGDA in which an inner SGD tracks the optimal dual variables, controlling the bias in robust-objective gradients. The outer loop independently samples both inner-iterate indices and data batches to obtain unbiased stochastic Gram-matrix estimates for preference updates. We prove convergence to an \(\epsilon\)-distributionally robust Pareto-stationary point using \(\mathcal{O}(\epsilon^{-4})\) outer iterations, each with \(\mathcal{O}(\epsilon^{-4})\) inner steps, yielding \(\mathcal{O}(\epsilon^{-8})\) total sample complexity.
    
    \item To avoid costly dual tracking, we derive a sufficient surrogate stationarity criterion and propose a single-loop Double-Clip MGDA method. Our key innovation is to incorporate gradient clipping into both the primal--dual updates and the MGDA preference update, yielding a jointly clipped optimization and gradient-balancing scheme. This design controls the generalized-smooth geometry and the bias from stochastic gradients with affine-bounded variance, eliminating both inner-loop tracking and double sampling while achieving \(\mathcal{O}(\epsilon^{-4})\) sample complexity. Experiments validate our theory and demonstrate competitive performance against state-of-the-art MGDA baselines.

\end{enumerate}


% We show that the dual formulation of DR-MOO satisfies an $L$-smoothness property along certain primal–dual trajectories. This property enables us to characterize the relationship between gradient norms and function value gaps, thereby eliminating the need for bounded gradient assumptions. Building on this insight, we design a double-loop MGDA algorithm. By combining a simple inner-loop SGD, with stochastic MGDA and double sampling, we ensure that both the stochastic gradient estimates and the preference vector updates exhibit controllably small bias relative to the full gradients. Under this design, the algorithm attains a sample complexity of $\mathcal{O}(\epsilon^{-12})$ to reach an $\epsilon$-distributionally-robust-Pareto-stationary point.



% To further reduce the sample complexity, we exploit the specific gradient structure of the dual DR-MOO formulation~\eqref{eq: dual DR-MOO}. In particular, we show that obtaining an $\epsilon$-Pareto-stationary point for DR-MOO~\eqref{eq: dual DR-MOO} is equivalent to finding an $\epsilon$-Pareto-stationary point of a rescaled surrogate objective. This equivalence eliminates the need for accurately estimating the dual variables. Building on this insight, we propose a single-loop algorithm tailored to the structure of the rescaled surrogate objective. By incorporating dual variable update and gradient clipping into the MGDA-based updates of both the model parameters and the preference vector, the proposed method removes the need for double sampling and inner-loop SGD, and achieves a sample complexity of $\mathcal{O}(\epsilon^{-4})$ for attaining an $\epsilon$-Pareto-stationary point. We conduct numerical experiments to validate the proposed methods. The results demonstrate that our approaches are competitive with state-of-the-art MGDA baselines.




\section{Related Works}
%{writting Logic: 1. DRO: Focus on bottleneck of scalability and strong assumption: Convex, bounded loss; 2. MOO: Focus on bottleneck of strong assumptions.}

%Most DRO papers study ambiguity sets defined via information divergences, including \(f\)-divergences \citep{DRO3rd,duchiDRO,lamfDRO,large-scaleDRO,qiaaai}, KL divergence \citep{huKLDRO,qiKLDRO}, Wasserstein distance \citep{blanchetwasserstein,gaoregularizeWDRO,gaoWDRO}, and its regularized variants \citep{sinkhornDROwang,wangzihansinkhornDRO,yufengSinkhornDRO}.A central challenge in DRO is to solve the worst-case optimization over probability distributions into a tractable formulation that supports practical stochastic algorithms.  Existing approaches often incur substantial computational costs even under single-objective case. For example, \citet{DRO3rd,mini-maxWDRO} parameterize the adversarial distribution using finite samples and cast \(f\)-divergence DRO as a min-max problem, but the complexity scales linearly with the dataset size. \citet{yufengSinkhornDRO} formulates Sinkhorn DRO as a contextual bilevel problem but at the cost of a double-loop procedure to solve it. Although there are other more tractable formulations, several of them rely on restrictive assumptions on loss objectives. for example, \citet{large-scaleDRO,sinkhornDROwang} require the loss to be globally bounded and convex, whereas the nonconvex analysis of \citet{qiaaai} still assumes globally bounded losses. Under multi-objective case, tractable DRO formulation with provable algorithm analysis was less studied. Recently, \citet{DRMOOLatorre} introduce Pareto-optimality and stationarity notions for convex loss with DRO formulation. However, their results characterize solutions via existence arguments through the primal robust objectives, which still involve a generally intractable functional maximization over probability distributions, without an accompanying algorithm or non-asymptotic convergence analysis.  Notably, \citet{jinnonconvexdro,qirevistfDRO} revisit strong dual formulation of \(f\)-divergence DRO \citep{DROBook,duchiDRO,large-scaleDRO} and develop algorithms with convergence guarantees for nonconvex objectives with unbounded losses. These results motivate our focus on constructing pareto-optimality and stationary condition via primal-dual transformation, and develop practical algorithms for \(f\)-divergence regularized DRO in the multi-objective setting.

%Most DRO studies consider ambiguity sets defined by statistical divergences or probability metrics, including \(f\)-divergences~\citep{DRO3rd,duchiDRO,lamfDRO,large-scaleDRO,qiaaai}, KL divergence~\citep{huKLDRO,qiKLDRO}, Wasserstein distance~\citep{blanchetwasserstein,gaoregularizeWDRO,gaoWDRO}, and its regularized variants~\citep{sinkhornDROwang,wangzihansinkhornDRO,yufengSinkhornDRO}. A central challenge is to transform the resulting worst-case optimization over probability distributions into a tractable formulation that supports practical stochastic algorithm design. Even in the single-objective setting, existing approach can incur substantial computational costs. For example, \citet{DRO3rd,mini-maxWDRO} parameterize the adversarial distribution using finite samples and formulate \(f\)-divergence DRO as a min-max problem, resulting in per-iteration complexity that scales linearly with the dataset size. Similarly, \citet{yufengSinkhornDRO} formulate Sinkhorn DRO as a contextual bilevel problem that requires a computationally expensive double-loop procedure. Although more tractable reformulations have been developed, their analyses often rely on restrictive assumptions on the loss. In particular, \citet{large-scaleDRO,sinkhornDROwang} require globally bounded and convex losses, whereas the nonconvex analysis of \citet{qiaaai} retains the global boundedness assumption. Notably, \citet{jinnonconvexdro,qirevistfDRO} exploit the strong-duality-based reformulation of \(f\)-divergence DRO~\citep{DROBook,duchiDRO,large-scaleDRO} to develop algorithms with convergence guarantees for nonconvex and potentially unbounded losses. However, extending such tractable reformulations and provable algorithms to the multi-objective setting remains largely underexplored. Recently, \citet{DRMOOLatorre} introduced Pareto-optimality and stationarity notions for a DR-MOO formulation with convex losses. Their results, however, characterize solutions through existence arguments involving the primal robust objectives, whose evaluation still requires solving a generally intractable functional maximization over probability distributions; consequently, they do not provide a practical algorithm or a non-asymptotic convergence analysis. Motivated by the strong-duality approach in single-objective DRO, we transform the primal DR-MOO problem into a finite-dimensional dual formulation, characterize robust Pareto optimality and stationarity through the dual objectives, and develop practical algorithms with non-asymptotic guarantees for \(f\)-divergence-regularized DR-MOO.


\textbf{Distributionally robust optimization (DRO).}Most DRO studies consider ambiguity sets defined by \(f\)-divergences~\citep{DRO3rd,duchiDRO,lamfDRO,large-scaleDRO,qiaaai}, KL divergence~\citep{huKLDRO,qiKLDRO}, Wasserstein distance~\citep{blanchetwasserstein,gaoregularizeWDRO,gaoWDRO}, and its regularized variants~\citep{sinkhornDROwang,wangzihansinkhornDRO,yufengSinkhornDRO}. A central challenge is to transform the resulting worst-case optimization over probability distributions into a tractable formulation supporting practical stochastic algorithms. Even in the single-objective setting, \citet{DRO3rd,mini-maxWDRO} explicitly parameterize the adversarial distribution, leading to complexity that scales with the dataset size, while the bilevel formulation of \citet{yufengSinkhornDRO} requires a double-loop procedure. Other tractable formulations often rely on restrictive assumptions, \citet{large-scaleDRO,sinkhornDROwang} require globally bounded and convex losses, whereas the nonconvex analysis of \citet{qiaaai} retains global boundedness. Notably, \citet{jinnonconvexdro,qirevistfDRO} exploit strong-duality-based reformulations of \(f\)-divergence DRO~\citep{DROBook,duchiDRO,large-scaleDRO} to provide convergence guarantees for nonconvex, potentially unbounded losses. Extending such tractable formulations and provable algorithms to MOO, however, remains underexplored. Recently, \citet{DRMOOLatorre} introduced Pareto-optimality and stationarity notions for DR-MOO with convex losses, but characterized them through existence arguments on the primal robust objectives, which still involve a generally intractable functional maximization over distributions. Consequently, no practical algorithm or non-asymptotic convergence analysis was provided. Motivated by strong duality, we instead derive a finite-dimensional dual formulation, characterize robust Pareto optimality and stationarity through the dual objectives, and develop practical algorithms with non-asymptotic guarantees for \(f\)-divergence-regularized DR-MOO.


\textbf{Gradient-Balancing Methods in MOO.}
Gradient-balancing methods construct a common descent direction by adapting the model update and preference vector across objectives. Following the formulation of MTL as MOO~\citep{1stempiricalMGDAwork}, many methods build on MGDA~\citep{MGDA1stwork} and fall into two broad lines. The first develops stochastic MGDA variants with convergence guarantees. Stochastic MGDA~\citep{stochasticMGDA} requires increasing batch sizes, whereas MoDo~\citep{MODO} and SDMGrad~\citep{SDMkaiyiJi} use double sampling to control bias, and MoCo~\citep{MOCO} uses moving-average gradient tracking. Their analyses, however, typically assume uniformly bounded objective gradients, which may fail for dual DRO objectives whose gradients grow with the loss. The second line mitigates gradient conflicts through structured manipulations or interpretable optimization formulations, including CAGrad~\citep{CAGrad}, IMTL-Grad~\citep{IMTLGrad}, PCGrad~\citep{PCGrad}, GradDrop~\citep{GradDrop}, RotoGrad~\citep{Rotograd}, and NashMTL~\citep{nashMTL}. In practice, the LibMTL implementations of SDMGrad and NashMTL~\citep{LibMTL} normalize gradient Gram matrices to improve stability. Although effective empirically, these mechanisms remain insufficiently understood theoretically, particularly in stochastic settings with unbounded gradients. This theory-practice gap, together with the analysis of \citet{qimgda} that relaxes uniform boundedness by relating gradient norms to function-value gaps under generalized smoothness, motivates our provable gradient-clipping scheme tailored to dual DR-MOO objectives.


%{Writting Logic: Overview through developments of MGDA + variants, as well as heuristics desgin, identify their weakness.}
% Based on the MOO formulation, gradient-balanced methods aim to construct a common descent direction by dynamically adjusting model parameters and preference vectors to achieve balanced improvement across objectives. Notably, \citet{1stempiricalMGDAwork} view multi-task learning (MTL) as a MOO problem and show that it can be effectively solved using MGDA~\citep{MGDA1stwork}. Subsequent MGDA-type methods can be broadly divided into two streams. The first stream focuses on mitigating gradient conflicts through heuristic gradient manipulation, including CAGrad~\citep{CAGrad}, IMTL-Grad~\citep{IMTLGrad}, PCGrad~\citep{PCGrad}, GradDrop~\citep{GradDrop}, and RotoGrad~\citep{Rotograd}. While effective empirically, these methods generally lack convergence guarantees in stochastic settings. 

% The second stream emphasizes theoretical guarantees. \citet{stochasticMGDA} establish the first convergence guarantee for stochastic MGDA, albeit at the cost of requiring increasing batch sizes. To address this limitation, MoDo~\citep{MODO} and SDMGrad~\citep{SDMKaiyiJi} employ double-sampling strategies to maintain fixed batch sizes while ensuring unbiased stochastic gradient estimator. Later, MoCo~\citep{MOCO} and its variance-reduced extension~\citep{heshanSTORM} further improve balanced gradient estimation by introducing auxiliary tracking variables. Although these methods provide convergence guarantees under mild conditions, their analyses typically rely on bounded gradient assumptions. Notably, \citet{qimgda} relaxes this requirement by proving convergence of stochastic MGDA with double sampling under generalized smoothness conditions, motivating further investigation of MGDA-based methods for DR-MOO.
% \vspace{-0.8em}
% \section{Preliminaries}
% %%%%%%%%%%%%%%%%%%%% Compressed
% Recall that each objective in MOO formulation~\eqref{eq: MOO} takes the stochastic form
% \(\phi^i(\theta):=\mathbb{E}_{\xi\sim\mathbb{P}^i}[\ell^i(\theta,\xi)]\), and \(\theta\) denotes the model parameter. 
% In MOO, a solution is preferred if it improves at least one objective without worsening any other objective. This induces the standard notion of Pareto dominance. A solution is Pareto-optimal if it is not dominated by any other feasible solution. Equivalently, at a Pareto-optimal solution, improving one objective must necessarily degrade at least one other objective. Since finding Pareto-optimal solutions is generally difficult, convergence analyses in machine learning typically rely on the weaker notion of Pareto-stationary.
% \begin{definition}[Pareto-stationary]\label{def: pstationary}
%     Denote ${\mathcal{W}}$ as the probability simplex. A solution ${\theta}$ is called Pareto-stationary if there exists $w\in {\mathcal{W}}$ such that $\|\sum_{i=1}^{m} w^i \nabla \phi^i({\theta}) \|=0$.
%     ${\theta}$ is said to be $\epsilon$-Pareto-stationary if there exists $w\in {\mathcal{W}}$ such that $\|\sum_{i=1}^{m} w^i \nabla \phi^i({\theta}) \|\leq\epsilon$.
% \end{definition}
% Pareto-stationary is a necessary condition for Pareto-optimal. Moreover, if each objective $\ell^i(\cdot)$ is strictly convex for all $i\in [m]$, then Pareto-stationary is also sufficient for Pareto-optimal \citep{marucsciac1982fritz}.











%Recall the MOO formulation~\eqref{eq: MOO}, each objective $i$ has a stochastic form $\phi^i({\theta}):=\mathbb{E}_{\xi\sim\mathbb{P}^i}[\ell^i({\theta}, \xi)]$ and ${\theta}$ denotes the model parameters. 
%In MOO, a solution is naturally considered better than another if it achieves no higher loss on every objective, leading to the following Pareto-dominance condition.


%\begin{definition}[Pareto-dominance]
%For $\theta_1,\theta_2\in\mathbf{R}^{n}$, we say that $\theta_1$ Pareto-dominates $\theta_2$ if and only if $\phi^i(\theta_1)\le \phi^i(\theta_2)$ for all $i\in \{1,\dots,m\}$, and $\phi^j(\theta_1)<\phi^j(\theta_2)$ for some $j\in \{1,\dots,m\}$.

%\end{definition}


%\begin{definition}[Pareto-optimal]
%A solution ${\theta}^*$ is Pareto-optimal if no other solution dominates it. 
%\end{definition}
%In other words, for a Pareto-optimal solution $\theta^*$, any alternative solution $\theta$ that improves some objectives must necessarily worsen at least one other objective.  
%Finding a Pareto-optimal solution is in general difficult. Consequently, in machine learning it is common to analyze convergence using the following weaker notion of Pareto-stationary.



\section{Distributionally Robust MOO}

%In this section, we introduce distributionally-robust multi-objective optimization (DR-MOO), a new formulation for robust multi-objective optimization considering distribution shifts in the data. We then discuss Pareto-dominance, optimality and stationary conditions for DR-MOO.

%\subsection{Problem Formulation}
Distributionally robust optimization (DRO) has been widely adopted in machine learning to ensure robust model performance under distribution shifts. Specifically, consider a supervised learning problem with the nonconvex loss function denoted by $\ell(\theta;\xi)$, where $\theta\in \mathbf{R}^n$ denotes the collection of model parameters and $\xi$ corresponds to a data sample that follows an underlying data distribution $\mathbb{Q}$. We are interested in the following regularized DRO problem
\begin{align}
  &\text{(DRO)}: \quad \min_{\theta \in \mathbf{R}^n} \sup_{\mathbb{Q}} H({\theta},\mathbb{Q}), \nonumber\\
  &\quad
   \text{where} ~ H({\theta},\mathbb{Q}) = \mathbb{E}_{\xi \sim \mathbb{Q}} \bigr [ \ell(\theta;\xi)\bigr] - \lambda D_f(\mathbb{Q} | \mathbb{P}).
    \label{primal1}
\end{align}
Here, $\lambda>0$ is regularization hyperparameter and $D_f$ denotes the $f$-divergence that measures the 
discrepancy between the nominal data distribution $\mathbb{P}$ and the underlying data distribution $\mathbb{Q}$. Such minimax formulation $\min_{\theta}\sup_{\mathbb{Q}}$
seeks model parameters $\theta$ that perform well under the worst-case distribution, thereby improving robustness to distribution shifts away from $\mathbb{P}$.
To model robustness in the presence of multiple, potentially conflicting objectives, we extend DRO to the multi-objective setting. Specifically, we formulate distributionally robust multi-objective optimization as
\begin{align}
   &(\text{DR-MOO}):\nonumber\\
   &\quad~\min_{{\theta} \in \mathbf{R}^n} \sup_{\mathbb{Q}^1,\ldots,\mathbb{Q}^m} 
   \big[ H^1({\theta},\mathbb{Q}^1), \ldots, H^m({\theta},\mathbb{Q}^m)\big],
\label{eq: primal DR-MOO}
\end{align}
where $H^i({\theta}, \mathbb{Q}^i)
=\mathbb{E}_{\xi\sim\mathbb{Q}^i}[\ell^i({\theta},\xi)]
-\lambda D_f(\mathbb{Q}^i\mid \mathbb{P}^i),\forall i\in[m]$.
Unlike scalarized DRO, this formulation preserves vector-valued trade-offs: each task is evaluated under its own worst-case distribution, and robustness is enforced over Pareto trade-offs rather than a fixed weighted aggregation. Building on the MOO framework, DR-MOO naturally adopts Pareto-based solution concepts~\citep{marucsciac1982fritz}. However, distributional uncertainty causes objective values to vary across distributions, motivating an extension of Pareto dominance and optimality to characterize the robust Pareto front.
\begin{definition}[Distributionally robust Pareto-dominance]\label{Def: Pareto-dominance in DR-MOO}
In DR-MOO, for $\theta_1$, $\theta_2\in \mathbf{R}^n$, we say that ${\theta}_1$ distributionally robustly dominates ${\theta}_2$ if and only if $\sup_{\mathbb{Q}^i}H^i({\theta}_1,\mathbb{Q}^i)\leq \sup_{\mathbb{Q}^i}H^i({\theta}_2,\mathbb{Q}^i)$ for \textit{all} $i\in \{1,\dots,m\}$, and $\sup_{\mathbb{Q}^j}H^j({\theta}_1,\mathbb{Q}^j) < \sup_{\mathbb{Q}^j}H^j({\theta}_2,\mathbb{Q}^j)$ for \textit{some} $j\in \{1,\dots,m\}$.
% \begin{enumerate}[leftmargin = *, topsep = 0pt]
%  \item ${\theta}_1$ distributionally-robustly dominates ${\theta}_2$, denoted ${\theta}_1 \preceq {\theta}_2$, if and only if
%      $\sup_{\mathbb{Q}^i}H^i({\theta}_1,\mathbb{Q}^i)\leq \sup_{\mathbb{Q}^i}H^i({\theta}_2,\mathbb{Q}^i)$ for \textit{all} $i\in \{1,\dots,m\}$, and $\sup_{\mathbb{Q}^j}H^j({\theta}_1,\mathbb{Q}^j) < \sup_{\mathbb{Q}^j}H^j({\theta}_2,\mathbb{Q}^j)$ for \textit{some} $j\in \{1,\dots,m\}$.
     
%      \item ${\theta}_1$ strictly distributionally-robustly dominates ${\theta}_2$, denoted $\theta_1\prec\theta_2$, if and only if $\sup_{\mathbb{Q}^i}H^i({\theta}_1,\mathbb{Q}^i) < \sup_{\mathbb{Q}^i}H^i({\theta}_2,\mathbb{Q}^i)$ for all $i\in \{1,\dots,m\}$.
% \end{enumerate}
\end{definition}
\begin{definition}[Distributionally robust Pareto-optimal]\label{Def: Pareto-optimal under distribution uncertain}
A solution ${\theta}^*$ is distributionally robust Pareto-optimal if no other solution distributionally robustly dominates it.
\end{definition}
Intuitively, these definitions compare model parameters after adversarially perturbing each objective’s data distribution towards its worst-case distribution. When there is no distribution shift, they reduce to the standard Pareto-dominance and optimality condition.  Although extending Pareto notions to the distributionally robust setting is natural, directly optimizing the primal robust objectives
\(\sup_{\mathbb{Q}^i} H^i(\theta,\mathbb{Q}^i)\) remains challenging.
Even when the underlying losses \(\ell^i(\theta;\xi)\) are smooth, the induced robust objectives involve distributional maximization and may be nonsmooth.
These challenges motivate defining Pareto stationarity for DR-MOO through a primal--dual transformation, as detailed in the next subsection.

% In non-convex DR-MOO, it is natural to extend the Pareto-stationary condition in Definition \ref{def: pstationary} for MOO to the distributionally-robust setting. However, directly transplanting this definition to DR-MOO is challenging since the primal objective function $\sup_{\mathbb{Q}^i} H(\theta, \mathbb{Q}^i)$, which involves an untractable infinite-dimension optimization sub-problem and may not be differentiable. These difficulties motivate us to define Pareto-stationary condition in its dual formulation, as explained in the next subsection.

% On the other hand, to find $\theta^*$ satisfying Pareto-stationary condition, i.e.,$\min_{w\in W}\|\nabla\sup_{\mathbb{Q}^i} H(\theta,\mathbb{Q}^i)w\|^2=0$. It requires computation of $\nabla\sup_{\mathbb{Q}^i} H(\theta;\mathbb{Q}^i)$, which includes an untractable infinite-dimension optimization sub-problem. These difficulties motivate us to instead establish Pareto-stationary condition from dual formulation. Once strong duality holds and the dual objective admits well-structured dependence on variables over Euclidean space, establishing Pareto-stationary for dual function is sufficient to guarantee a distributionally robust Pareto-optimal solution, due to the equivalence of objective values between the primal and dual formulations. Therefore, we define the distributionally robust Pareto-stationary condition over dual function $\phi(\theta)$ as follows




    
\subsection{Dual Formulation and Distributionally Robust Pareto-Stationarity}

% Distributionally-robust Pareto-optimal is global solution concept, hard to achieve in nonconvex optimization. In the existing literature, a popular local solution concept for the standard MOO problem is Pareto-stationary. But, one cannot directly define it for DR-MOO since ???. In this subsection, we study the dual formulation of DR-MOO, based on which we propose a tractable notion of distributionally-robust Pareto-stationary.

% In this subsection, We leverage Lagrange duality theory \citep{DRO2nd,DRO3rd,large-scaleDRO}, rewriting DR-MOO through dual reformulation. To make the problem solvable using first-order methods. Throughout, we make the following standard assumptions on loss function $\ell(\cdot)$, which have been widely adopted in the existing DRO literature \citep{jinnonconvexdro,qirevistfDRO,yufengSinkhornDRO}. 
In this subsection, we leverage Lagrangian duality to reformulate DR-MOO into an equivalent dual form that facilitates gradient-based algorithm design. Throughout, we impose the following standard assumptions on the loss functions, which are widely adopted in the DRO literature \citep{jinnonconvexdro,qirevistfDRO}.
\begin{assumption}\label{assumption 1}
The loss functions $\ell^i(\cdot)$ and divergence function of DR-MOO in \eqref{eq: primal DR-MOO} satisfy the following conditions.
\begin{itemize}[leftmargin = *, topsep = 0pt]
    \item The sample loss $\ell^i(\theta,\xi)$ is finite $\mathbb{P}_i$-almost surely.
    \item There exists $G,L>0$ such that the loss function $\ell^{i}(\cdot, \xi)$ is $G$-Lipschitz continuous and $L$-smooth for all $\xi$ and all $i\in[m]$.
        \item For any ${\theta}$, the variance of $\ell^i(\cdot),~\forall i\in[m]$ with respect to sample $\xi~\sim \mathbb{P}^i$ is bounded by $\kappa^2$. %i.e., $
    %\mathbb{E}_{\xi\sim \mathbb{P}^i}\big[\big(\ell^i({\theta};\xi) - \mathbb{E}_{\xi\sim \mathbb{P}^i}[\ell^i({\theta};\xi)] \big)^2\big]\leq \kappa^2.$
   % where $\kappa$ is a finite constant.
    \item The divergence base function $f:[0,+\infty)\rightarrow (-\infty,+\infty]$ is convex and satisfies $f(1)=0, f(0)=\lim_{t\rightarrow 0^{+}}f(t)$. Additionally, its convex conjugate $f^*(\cdot)$ is $M$-smooth.
\end{itemize}    
\end{assumption}
{Several $f$-divergences satisfy the $M$-smoothness assumption, including the $\chi^2$-divergence, smoothed CVaR divergence, etc.} Under Assumption \ref{assumption 1}, \citet{DROBook,DRO3rd,large-scaleDRO} show that the regularized DRO problem in \eqref{primal1} admits the following equivalent dual formulation.
\begin{align}
  &\sup_{\mathbb{Q}} H({\theta},\mathbb{Q}) =
  \phi(\theta),\nonumber\\
  &\quad ~\text{and}~ \phi(\theta)=
  :\min_{\eta\in \mathbf{R}} \lambda \mathbb{E}_{\xi\sim\mathbb{P}}\big[f^*\big(\frac{\ell({\theta};\xi)-\eta}{\lambda}\big)\big]+\eta,\label{eq: single-objective dual DRO}
\end{align}
where $\eta$ denotes the dual variable. Leveraging these facts, primal DR-MOO \eqref{eq: primal DR-MOO} can be rewritten in the following equivalent dual form
\begin{align}
  &\text{(Dual of DR-MOO):}\nonumber\\
  &\quad\min_{\theta} \Big\{ \Phi({\theta}) := \big[\phi^1({\theta}), \phi^2({\theta}), ..., \phi^m({\theta})\big] \Big\},
 \label{eq: dual DR-MOO}
%\label{eq: dual-inner}
\end{align}
where $\phi^i({\theta}) = \min_{\eta^i\in \mathbf{R}} \lambda \mathbb{E}_{\xi\sim\mathbb{P}^i}\big[f^*\big(\frac{\ell^i({\theta};\xi)-\eta^i}{\lambda}\big)\big]+\eta^i$.
% When considering distributionally robust optimization in the context of multi-objective optimization, given multiple loss $\ell^i({\theta},\xi),\forall i\in[m]$ and underlying distribution $\xi^i\sim\mathbb{Q}^i, \forall i\in[m]$, to find Pareto-optimal point of primal formulations of DR-MOO \eqref{eq: primal DR-MOO}, definition \ref{Def: Pareto-optimal under distribution uncertain} indicates one can find the Pareto-optimality condition in terms of dual function $\phi({\theta})$. In the context of non-convex optimization, we start finding the Pareto-stationary point of $\phi({\theta})$. 
% Thus, we shift our interest of \eqref{eq: primal DR-MOO} into following strong duality formulation
%Following this conclusion, we conclude for $f$-divergence regularized DRO, one can characterize Pareto-stationary condition among $\Phi({\theta})$ and $\fL({\theta}, {\eta})$ by leveraging this structure and Lemma 2.6 proved in \cite{jinnonconvexdro} as follows.
%Therefore, DR-MOO can be rewritten as a standard MOO problem, in which each objective function takes the dual form of a DRO problem. In particular, note that this MOO problem has a nested structure, each $\phi^i({\theta})$ is in a minimization form over a dual variable $\eta^i$. Since each objective function $\phi^i({\theta})$ is differentiable, we can leverage the standard definition of Pareto-stationary to define Pareto-stationary for DR-MOO.
% Through this way, DR-MOO is solvable as standard MOO problems with $m$ objectives in their dual form. Each objective $\phi^i({\theta})$ is defined via a one-dimensional convex minimization over the dual variable $\eta^i$, yielding a tractable structure. Importantly, each dual objective $\phi^i({\theta})$ is continuously differentiable and eliminates the operator $\sup_{\mathbb{Q}}$ over infinite functional space, this enables first-order algorithm designs. and More importantly, allows us to leverage the Pareto-stationary in Definition~\ref{Def: Pareto-stationary under distribution uncertain} to define the following notion of distributionally-robust Pareto-stationary, which we adopt as the solution metric for DR-MOO.
In this dual form, DR-MOO can be treated as a standard MOO problem with \(m\) dual objectives. Moreover, for finite $\mathcal L^i(\theta,\eta^i)$ with nonempty
$\arg\min_{\eta^i}\mathcal L^i(\theta,\eta^i)$ at every $\theta$, \citet{jinnonconvexdro} showed that each \(\phi^i(\theta)\) is differentiable and its gradient admits the closed-form expression
\begin{align}
\nabla \phi^i(\theta)=\nabla_\theta \fL^i(\theta,\eta^{i,*}),
\eta^{i,*}\in \arg\min_{\eta\in\mathbf{R}} \fL^i(\theta,\eta^i),\nonumber
\end{align}
where
\(
\fL^i(\theta,\eta^i) :=\lambda \mathbb{E}_{\xi\sim\mathbb{P}}
[f^*(\frac{\ell^i(\theta;\xi)-\eta^i}{\lambda})]+\eta^i
\).
This dual representation replaces the infinite-dimensional maximization over
$\mathbb{Q}^i$ with a finite-dimensional problem and leads to the following
Pareto-stationarity condition for DR-MOO.
\begin{definition}[$\epsilon$-Distributionally robust Pareto-stationary]\label{def: dr-pstationary}
   Denote $\mathcal{W}$ as the probability simplex. We say that $\theta$ is an $\epsilon$-distributionally robust Pareto-stationary solution of DR-MOO if it is an $\epsilon$-Pareto-stationary solution of the dual formulation, i.e., there exists $w \in \mathcal{W}$ such that$
    \| \sum_{i=1}^{m} w^i \nabla \phi^i(\theta)\| \leq \epsilon$.
\end{definition}
Related notions of Pareto dominance, optimality, and stationarity under
distributional uncertainty were also studied by \citet{DRMOOLatorre}. Their framework
considers convex, possibly nonsmooth losses $\ell^i(\cdot)$ and characterizes
these notions through the existence of worst-case distributions
$\mathbb{Q}^{i,*}$ in the primal formulation. In contrast, we consider
nonconvex, $L$-smooth losses $\ell^i(\cdot)$ and define stationarity through the dual objectives
$\phi^i(\theta)$ induced by the $f$-divergence reformulation. Crucially, their
primal characterization requires identifying the worst-case distributions by
solving infinite-dimensional maximization problems, which limits its use in
gradient-based optimization. Our dual characterization avoids this obstacle and
enables the MGDA-type algorithms developed in the next two sections, together
with provable convergence to $\epsilon$-Pareto-stationary solutions.
\section{Double-Loop MGDA}\label{sec: Double-loop algorithm}
%In this section, we develop a double-loop multi-gradient descent algorithm (MGDA) to solve the dual DR-MOO problem in \eqref{eq: dual DR-MOO}, and establish its convergence guarantee in the nonconvex setting.
%\subsection{Algorithm Design}
%{Q: need notation $\xi = [\xi^1, \xi^2, ..., \xi^m]$.Writting Logic: 1. Introduce MGDA. 2. Why Double-loop: Mitigate estimation bias of true gradient $\nabla\phi^i(\theta)$.}
%{Add a short paragraph to summarize this section.}
% The standard MOO in \eqref{eq: MOO}, a central challenge is to handling conflicts among objective gradients. To address this, various types of multi-gradient descent algorithms (MGDA) \citep{MGDA1stwork,CAGrad,MOCO,MODO,SDMkaiyiJi,stochasticMGDA,qimgda,CAGrad} have been proposed. The key idea is to computing a balanced update direction \(d\) that maximizes the worst-case improvement across objectives, i.e.,
% $\max_{d\in\mathbf{R}^m}\frac{1}{\gamma} \min_{i\in[m]}\{\phi^i(\theta)-\phi^i(\theta-\gamma d) \}$,
% where $\gamma$ denotes the learning rate. This leads to the following MGDA update rules \citep{stochasticMGDA,qimgda}. 
% \begin{align}
% \text{(MGDA):} \quad d^* &=\nabla\Phi(\theta)w_{\rho}^*, \nonumber\\
%     \text{where} ~ w_{\rho}^* &=\arg\min_{w\in W}\frac{1}{2}\|\nabla \Phi(\theta)w\|^2+\frac{\rho}{2}\|w\|^2, \label{eq: MGDA surrogate}
% \end{align}
% and an \(\ell_2\)-regularization term is added to ensure that \(w_\rho^*\) is uniquely defined \citep{MODO,qimgda}, which in turn yields a tractable iterative update of the form $w_{t+1}=w_t -\nabla\Phi(\theta)^{\top}\nabla\Phi(\theta)w_t+\rho w_t$. 
For the standard MOO problem, a central challenge is handling conflicts among objective gradients. To address this issue, a variety of multi-gradient descent algorithms (MGDA) has been proposed \citep{MGDA1stwork,CAGrad,MOCO,MODO,SDMkaiyiJi,stochasticMGDA,qimgda}. The key idea is to compute an update direction \(d\) that maximizes the worst-case improvement across objectives, i.e.,
$
\max_{d\in\mathbf{R}^{{n}}}\frac{1}{\gamma}\min_{i\in[m]}\{\phi^i(\theta)-\phi^i(\theta-\gamma d)\},
$
where \(\gamma\) denotes the learning rate. 
This leads to the following MGDA update rule widely adopted in~\citet{stochasticMGDA,SDMkaiyiJi,qimgda,MODO}
\begin{align}
&\text{(MGDA):} \quad d^* =\nabla\Phi(\theta)w_{\rho}^*,\nonumber \\
&\quad\text{where} ~ w_{\rho}^* =\arg\min_{w\in W}\frac{1}{2}\|\nabla \Phi(\theta)w\|^2+\frac{\rho}{2}\|w\|^2, \label{eq: MGDA surrogate}
\end{align}
%{where $d^*$ is conflict-avoidant (CA) direction}, 
where \(\ell_2\)-regularization term ensures that \(w_\rho^\star\) is uniquely defined. To avoid solving~\eqref{eq: MGDA surrogate} exactly, we use the one-step preference update
\begin{align}
    w_{t+1}=\Pi_\mathcal{W}\big(w_t-(\nabla\Phi(\theta)^{\top}\nabla\Phi(\theta)w_t+\rho w_t)\big).
    \label{eq: MGDA one-step approximation}
\end{align}
However, applying MGDA~\eqref{eq: MGDA one-step approximation} to the dual formulation~\eqref{eq: dual DR-MOO} faces two challenges. First, computing the exact minimizers \(\eta^{i,*}\) is generally impractical. Second, estimating the Gram matrix \(\nabla\Phi(\theta)^\top\nabla\Phi(\theta)\) using sample-averaged stochastic gradients can introduce bias into the preference update.
%\begin{wrapfigure}{r}{0.5\textwidth}
%\vspace{-0.5em}
%\centering
%\begin{minipage}{0.5\textwidth}
%\hrule
%\vspace{0.3em}
%\captionsetup{type=algorithm,
%    labelfont=bf,
%    textfont=normalfont,
%    justification=raggedright,
%    singlelinecheck=false
%    }
%\caption{Double-loop MGDA for DR-MOO}
%\label{alg1}
%\vspace{-0.5em}
%\hrule
%\vspace{0.3em}
%\small
%\begin{algorithmic}[1]
%\STATE \textbf{Initialize} $\theta_0$, $w_0$, $\{\eta^i_{0,0}\}_{i=1}^m$,$\eta^i_{0,-1}=\eta_{0,0}^i$,$\rho$,
%and learning rates $\alpha,\beta,\gamma$.
%\FOR{$t=0,\ldots,T-1$}
%    \STATE Set $\eta_{t,0}^i = \eta^i_{t-1,d}$.
%    \FOR{$d=0,\ldots,D-1$ and $i\in[m]$}
%        \STATE Evaluate $V^i_d=\nabla_{\eta}L^i(\theta_t,\eta^i_{t,d};\xi^i_{t,d})$.
%        \STATE $\eta^i_{t,d+1}=\eta^i_{t,d}-\gamma V^i_d$.
%    \ENDFOR
%    \STATE Sample index $d,\bar d,\tilde d\overset{i.i.d.}{\sim}\text{Unif}\{0,\dots,D-1 \}$
%    \STATE Draw $\xi_t,\bar\xi_t,\tilde\xi_t$  independently.
%    \STATE Evaluate
    %$Y_t=\nabla_{\theta}\fL(\theta_t,\eta_{t,d};\xi_t)$.
%    \STATE Evaluate
%    $\bar Y_t=\nabla_{\theta}\fL(\theta_t,\eta_{t,\bar d};\bar{\xi}_t)$.
%    \STATE Evaluate
%    $\widetilde Y_t=\nabla_{\theta}\fL(\theta_t,\eta_{t,\tilde d};\widetilde{\xi}_t)$.
%    \STATE $\theta_{t+1}=\theta_t-\alpha Y_t w_t$.
%    \STATE $w_{t+1}=
%    \Pi_{\mathcal W}\!\big(
%    w_t-\beta(\bar Y_t^\top\widetilde Y_t w_t+\rho w_t)
%   \big)$.
%\ENDFOR
%\end{algorithmic}
%\vspace{0.3em}
%\hrule
%\end{minipage}
%\vspace{-1em}
%\end{wrapfigure}
\begin{algorithm}[t]
\hrule
\vspace{0.3em}

\caption{Double-loop MGDA for DR-MOO}
\label{alg1}

\vspace{-0.5em}
%\hrule
\vspace{0.3em}

\footnotesize
\begin{algorithmic}[1]
\STATE \textbf{Initialize} $\theta_0$, $w_0$,
$\{\eta^i_{0,0}\}_{i=1}^m$, $\rho$, and learning rates
$\alpha,\beta,\gamma$.

\FOR{$t=0,\ldots,T-1$}
    \STATE Run inner-SGD independently
    \FOR{$d=0,\ldots,D-1$ and $i\in[m]$}
        \STATE Evaluate
        $V^i_d=\nabla_{\eta}L^i(
        \theta_t,\eta^i_{t,d};\xi^i_{t,d})$.
        \STATE
        $\eta^i_{t,d+1}=\eta^i_{t,d}-\gamma V^i_d$.
        %\STATE In parallel, draw $\xi_{t,d},\xi_{t,\bar{d}},\xi_{t,\tilde{d}}$ and evaluate
        %\STATE  %$V^i_d=\nabla_{\eta}L^i(
        %\theta_t,\eta^i_{t,d};\xi^i_{t,d})$; $V^i_{\bar d}=\nabla_{\eta}L^i(
        %\theta_t,\eta^i_{t,\bar d};\xi^i_{t,\bar d})$ and $V^i_{\tilde{d}}=\nabla_{\eta}L^i(
        %\theta_t,\eta^i_{t,\tilde{d}};\xi^i_{t,\tilde{d}})$.
        %\STATE $\begin{aligned}
        %  \eta_{t,d+1}^i
        %    &=\eta_{t,d}^i-\gamma V_{d}^i,\\
          %\bar\eta_{t,\bar{d}+1}^i
            %&=\eta_{t,\bar{d}}^i-\gamma {V}_{\bar{d}}^i,\\
          %\eta_{t,\tilde{d}+1}^i
         %   &=\eta_{t,d}^i-\gamma {V}_{\tilde{d}}.
        %\end{aligned}$
    \ENDFOR

    \STATE Sample $d,\bar d,\tilde d
\overset{\mathrm{i.i.d.}}{\sim}\operatorname{Unif}\{0,\ldots,D-1\}$.

    \STATE Draw
    $\xi_t,\bar\xi_t,\tilde\xi_t$ independently.

    \STATE
    $Y_t=\nabla_{\theta}\fL(
    \theta_t,\eta_{t,d};\xi_t)$.

    \STATE
    $\bar Y_t=\nabla_{\theta}\fL(
    \theta_t,\eta_{t,\bar d};\bar\xi_t)$.

    \STATE
    $\widetilde Y_t=\nabla_{\theta}\fL(
    \theta_t,\eta_{t,\tilde d};\tilde\xi_t)$.

    \STATE
    $\theta_{t+1}=\theta_t-\alpha Y_t w_t$.

    \STATE
    $\displaystyle
    w_{t+1}=\Pi_{\mathcal W}\!\big(
    w_t-\beta(
    \bar Y_t^\top\widetilde Y_t w_t+\rho w_t
    )\big)$.
    \STATE $\eta_{t+1,0} = \eta_{t,d}$
\ENDFOR
\end{algorithmic}
\vspace{0.3em}
%\hrule
\end{algorithm}
To address these challenges, we propose double-loop MGDA (Algorithm~\ref{alg1}), which combines inner-loop dual tracking with independent sampling of data and dual variable indices. Its two loops play complementary roles in controlling the errors arising from approximate dual minimization and stochastic Gram-matrix estimation, as detailed below.
\begin{itemize}[leftmargin = *, topsep = 0pt]
    \item \textbf{Inner loop (dual tracking).} Since the exact minimizers \(\eta^{i,*}\) are unavailable, we run vanilla SGD \citep{lanSGD} (lines 4-7 in Algorithm~\ref{alg1}) to approximately solve \(\min_{\eta^i} \fL^i(\theta,\eta^i)\). This inner loop tracks the dual minimizers as \(\theta_t\) evolves, controlling the bias between \(\nabla_\theta \fL(\theta_t,\eta_t)\) and the desired \(\nabla \Phi(\theta_t)\).

    \item \textbf{Independent index and batch sampling} While dual tracking controls approximation bias, reusing the same stochastic gradient estimator with dependent randomness introduces a covariance term into the
    expected Gram-matrix estimate. We therefore adopt a double-sampling strategy \citep{SDMkaiyiJi,MODO}, independently sampling both inner-iterate indices and data batches. Specifically, we draw three independent index--batch pairs to
    construct conditionally independent estimators $Y_t$, $\bar Y_t$, and
    $\tilde Y_t$ such that
    $\mathbb{E}[\bar{Y}_t^{\top}\tilde{Y}_tw_t|\calH_t]=\mathbb E[\bar Y_t\mid\mathcal H_t]^\top
\mathbb E[\tilde Y_t\mid\mathcal H_t]w_t$. Thus, the cross-product $\bar Y_t^\top\tilde Y_t$ eliminates the covariance
    term arising from estimator reuse and provides an unbiased stochastic Gram-matrix estimate, yielding convergence without requiring dynamically growing minibatches~\citep{stochasticMGDA}.
\end{itemize}
%\begin{remark}
%    The inner subproblem \(\min_{\eta^i} \fL^i(\theta,\eta^i)\) is one-dimensional and convex. Hence, the inner loop often takes few iterations and little memory to reach a desired accuracy, and the overall computation overhead is much lower than the outer loop. {We compare the wall-clock time of each outer iteration of Double-Loop MGDA with single-loop baselines in deep learning settings. Figure~\ref{fig: Wall-clock} (Appendix~\ref{Appendix: inner-iteration and wall-clock}) shows that dual-variable computation is lightweight compared to updates of the model parameters $\theta_t$ and the preference vector $w_t$.}
%\end{remark}
%\begin{algorithm}[ht]
%\caption{Double-loop MGDA for DR-MOO}\label{alg1}
%\small
%\resizebox{0.9\linewidth}{!}{
%\begin{minipage}{0.9\linewidth}
%\begin{algorithmic}[1]
%\STATE Initialize ${\theta}_0$, ${w}_0$, $\{{\eta}_{0,0}^i \}_{i=1}^m$, $\rho$; learning rates $\alpha,\beta,\gamma$;
%\FOR{$t=0.....T-1$}
%        \FOR {$d=0....D-1$ and all $i=1,2,...,m$}
%            \STATE Evaluate $ V^i_d=\nabla_{\eta}L^i({\theta}_t, \eta^i_{t,d};\xi^i_{t,d})$
%             \STATE $\eta^i_{t,d+1}=\eta^i_{t,d}-\gamma V_d^i$
%    \ENDFOR
%    \STATE {Sample $d,\bar{d},\tilde{d}\sim \{0,...,D-1 \}$ independently}
%    \STATE {Draw $\xi_t, \bar{\xi}_t,\tilde{\xi}_t$ independently}
%    \STATE  Evaluate $\Y_{t}=\nabla_{{\theta}} \fL({\theta}_t, {\eta}_{{t,d}}; {\xi}_t)$
%    \STATE Evaluate $\Yb_t=\nabla_{{\theta}} \fL({\theta}_t, {\eta}_{{t,\bar{d}}}; \bar{\xi}_t)$
%    \STATE Evaluate $ \Yt_t=\nabla_{{\theta}} \fL({\theta}_t, {\eta}_{t,\tilde{d}}; \tilde{\xi}_t)$
%    \STATE  ${\theta}_{t+1}={\theta}_t-\alpha \Y_{t} {w}_t $
    
%    \STATE ${w}_{t+1}=\Pi_{\mathcal{W}}({w}_t-\beta\big( \Yb_t^{\top} \Yt_t{w}_t+\rho {w}_t)\big)$
%\ENDFOR
%\end{algorithmic}
%\end{minipage}
%}
%\end{algorithm}
\subsection{Convergence Analysis}
%We first establish a special $L$-smooth property for $\phi^i(\theta)$.
%\begin{restatable}[$L$-smooth of $\phi^i({\theta})$]{lemma}{smoothproperty}\label{lemma: smooth of phi}
%    Let Assumption \ref{assumption 1} hold. Denote $\eta_{{\theta}}^{i,*}\in \arg\min_{{\eta}\in \mathbf{R}}\fL^i({\theta},\eta^i)$. For any $\theta,\theta'$, we have$
%        \|\nabla \phi^i({\theta})-\nabla_{{\theta}} \fL^i({\theta'},\eta^{i,*}_{{\theta}}) \|\leq L_0\|{\theta}-{\theta}' \|$
%holds for $i\in[m]$, where $L_0=G^2M\lambda^{-1}+L$.
%\end{restatable}
%This lemma implies that the function geometry in a neighborhood of the primal–dual pair $(\theta, \eta_{\theta}^{i,*})$ is $L$-smooth. Intuitively, once an accurate estimate of the dual variable is obtained via the inner-loop SGD in Algorithm~\ref{alg1}, the convergence analysis becomes more tractable. Moreover, this smoothness condition yields a convenient descent inequality for $\phi^i(\theta)$.
%In particular, letting $\theta'=\theta-\frac{\nabla \phi^i(\theta)}{L_0(\|\nabla \phi^i(\theta) \|+1)^{1/2}}$ and applying the descent lemma gives 
%\begin{align}  
%\|\nabla\phi^i(\theta_t)\|^2 \leq 2L_0(\phi^i(\theta_t)-\phi^{i,*})(\|\nabla\phi^i(\theta_t)\|+ 1). \label{eq: boundgrad}
%\end{align}
%Therefore, whenever $\phi^{i}(\theta_t)-\phi^{i,*}\le F$, defining $\Lambda := \sup\{u\geq 0| u^2\leq {2L_0}F(u+1) \}$, we obtain $\|\nabla\phi^i(\theta_t)\|\leq \Lambda$.
%This establishes an explicit link between gradient boundedness and the function-value gap, allowing us to avoid a global bounded-gradient assumption on the dual DR-MOO formulation~\eqref{eq: dual DR-MOO} in the high-probability analysis. 
We first establish the local regularity of the dual objectives in
Appendix~\ref{Appendix: Relevant Properties of Phi}. Specifically,
Lemmas~\ref{lemma: gradient bound of phi} and~\ref{lemma: smooth of phi} show
that, at primal-dual pair $(\theta,\eta_{\theta}^{i,*})$, $\phi^i$ has locally bounded gradient,
$\|\nabla\phi^i(\theta)\|\leq G$ and locally $L_0$-smooth around
$\theta$.
Consequently, once the inner-loop SGD provides a sufficiently accurate estimate of $\eta_{\theta}^{i,*}$, the analysis follows techniques similar to those for double-sampling stochastic MGDA with nonconvex $L_0$-smooth objectives~\citep{qimgda}, while additionally controlling bias terms arising from inexact dual estimation. Building on this observation, we obtain the following
convergence result. The complete
statements, hyperparameter choices, and proof are provided in
Appendix~\ref{Appendix: proof of convergence algorithm 1}.
\begin{restatable}[Convergence of Algorithm~\ref{alg1}]{theorem}{convgdoubleloopmgda}\label{thm: convergence of algorithm 1}
    Let Assumption \ref{assumption 1} hold. Given $\epsilon$, denote $\Delta_{\theta_0}$ as initial outer-objective gap, $\overline\Delta_\eta$ as $\epsilon$-independent almost-sure upper bound on the warm-start gaps.
    Set $\rho=\mathcal{O}( \epsilon^2)$, $\beta =\mathcal{O}(\epsilon^2)$,
    $\alpha = \mathcal{O}(\epsilon^2)$, and $\gamma = \mathcal{O}(\epsilon^2)$ for Algorithm~\ref{alg1}. Then, after $T=\Theta(\max\{\Delta_{\theta_0}\alpha^{-1}\epsilon^{-2},\beta^{-1}\epsilon^{-2}\})$ outer iterations, each associated with $D=\Theta(\overline{\Delta}_{\eta}\gamma^{-1}G^4\epsilon^{-2})$ inner iterations, we have \[\frac{1}{T}\mathbb{E}\big[\sum_{t=0}^{T-1}\|\nabla \Phi(\theta_t)w_t\|^2\big]\leq 94\epsilon^2.\]
\end{restatable}

Theorem~\ref{thm: convergence of algorithm 1} provides a baseline guarantee for applying MGDA directly to the dual DR-MOO formulation~\eqref{eq: dual DR-MOO}. Its outer-loop complexity \(T=\mathcal{O}(\epsilon^{-4})\) matches the rates established for stochastic MGDA variants~\citep{SDMkaiyiJi,MOCO,MODO,qimgda}, whereas its higher sample complexity of \(\mathcal{O}(\epsilon^{-8})\) arises from the inner-loop accuracy needed to control bias due to inexact dual variables. Nevertheless, for fixed \(\theta\), each inner problem \(\min_{\eta^i}\fL^i(\theta,\eta^i)\) is one-dimensional and convex, and thus typically requires few iterations and little memory to reach sufficient accuracy. The ablation study in Figure~\ref{fig: batch-size and inner-step ablation} (Appendix~\ref{Appendix: inner-iteration and wall-clock}) shows that the number of inner steps has little effect on empirical convergence, while the wall-clock comparison in Figure~\ref{fig: Wall-clock} (Appendix~\ref{Appendix: inner-iteration and wall-clock}) indicates that dual-variable updates are lightweight relative to updates of \(\theta_t\) and \(w_t\) in deep-learning settings. Together with the constant \(\mathcal{O}(1)\) batch size sufficient for our convergence guarantee, these results support the practicality of double-loop MGDA in peak-memory-constrained settings. Next, we discuss the main technical novelties in our analysis.
%\begin{wrapfigure}{r}{0.5\textwidth}
%\vspace{-1em}
%\centering
%\begin{minipage}{0.5\textwidth}
%\hrule
%\vspace{0.2em}
%\captionsetup{
%    type=algorithm,
%    labelfont=bf,
%    textfont=normalfont,
%    justification=raggedright,
%    singlelinecheck=false
%}
%\caption{Double-Clip MGDA for DR-MOO}
%\label{alg1-3}

%\vspace{-0.4em}
%\hrule
%\vspace{0.25em}

%\small
%\renewcommand{\baselinestretch}{0.94}\selectfont
%\begin{algorithmic}[1]
%\STATE \textbf{Initialize} $\theta_0$, $\eta_0$, $w_0$, $\rho$, $\beta,\gamma$,
%\STATE Clipping Rule:
%$\alpha_t=\min\!\{c_1,\frac{c_2}{\|X_t w_t\|}\}$,
%$\mu_t=\min\!\{f_1,\frac{f_2}{\|Z_t w_t\|}\}$.
%\FOR{$t=0,\ldots,T-1$}
%    \STATE Evaluate
    %$Z_t=\nabla_{\eta}\fhL(\theta_t,\eta_t;\{\xi_t\}_B)$ with $B=N_2$.
    %\STATE $\eta_{t+1}=\eta_t-\gamma\mu_t Z_t w_t$.
    %\STATE Evaluate
    %$X_t=\nabla_{\theta}\fhL(\theta_t,\eta_{t+1};\{\bar{\xi}_t\}_B)$ with $B=N_1$.
    %\STATE $\theta_{t+1}=\theta_t-\gamma\alpha_t X_t w_t$.
    %\STATE $w_{t+1}=
    %\Pi_{\mathcal{W}}\big(w_t-\beta(
    %\alpha_t X_t^\top X_t w_t
    %+\mu_t Z_t^\top Z_t w_t
    %+\rho w_t
    %)\big)$.
%\ENDFOR
%\end{algorithmic}
%\vspace{0.25em}
%\hrule
%\end{minipage}
%\vspace{-1.5em}
%\end{wrapfigure}
\begin{algorithm}[t]
\hrule
\vspace{0.3em}

\caption{Double-Clip MGDA for DR-MOO}
\label{alg1-3}

\vspace{-0.5em}
\hrule
\vspace{0.3em}

\footnotesize
\begin{algorithmic}[1]
\STATE \textbf{Initialize} $\theta_0$, $\eta_0$, $w_0$, $\rho$, $\beta,\gamma$,
\STATE Clipping Rule:
$\alpha_t=\min\!\{c_1,\frac{c_2}{\|X_t w_t\|}\}$,
$\mu_t=\min\!\{f_1,\frac{f_2}{\|Z_t w_t\|}\}$.
\FOR{$t=0,\ldots,T-1$}
    \STATE Evaluate
    $Z_t=\nabla_{\eta}\fhL(\theta_t,\eta_t;\{\xi_t\}_B)$ with $B=N_2$.
    \STATE $\eta_{t+1}=\eta_t-\gamma\mu_t Z_t w_t$.
    \STATE Evaluate
    $X_t=\nabla_{\theta}\fhL(\theta_t,\eta_{t+1};\{\bar{\xi}_t\}_B)$ with $B=N_1$.
    \STATE $\theta_{t+1}=\theta_t-\gamma\alpha_t X_t w_t$.
    \STATE $w_{t+1}=
    \Pi_{\mathcal{W}}\big(w_t-\beta(
    \alpha_t X_t^\top X_t w_t
    +\mu_t Z_t^\top Z_t w_t
    +\rho w_t
    )\big)$.
\ENDFOR
\end{algorithmic}

\vspace{0.3em}
\hrule
\end{algorithm}
\begin{itemize}[leftmargin = *, topsep = 0pt]
\item \textbf{Bias from inexact dual tracking.} While double sampling removes the \emph{stochastic} bias in MGDA, it does not eliminate the \emph{optimization} bias induced by approximate dual solutions, as both \(\nabla_{\theta}\fL(\theta_t,\eta_{t,\bar d})\) and \(\nabla_{\theta}\fL(\theta_t,\eta_{t,\tilde d})\) may differ non-negligibly from \(\nabla\Phi(\theta_t)\). Prior work~\citep{qimgda} assumes access to unbiased stochastic gradients and can exploit martingale structure for terms like $\mathbb{E}[\sum_{t=0}^{\tau-1}(Y_{t}^i-\nabla_{\theta} \fL^i(\theta_t,\eta_{t,d}))w_t^i]$.  In contrast, because the exact dual minimizers $\eta^{i,*}$ are unavailable, our analysis involves bias sequences of the form $\mathbb{E}[\sum_{t=0}^{T-1}(\nabla_{\theta}\fL^i(\theta_t,\eta^i_{t,d})-\nabla\phi^i(\theta_t))w_t^i]$, whose errors may accumulate over time. Sampling \(d,\bar d,\tilde d\) independently  ensures that the corresponding dual-gradient biases have the same conditional mean given \(\calH_t\). This enables partial cancellation between the bias terms arising from the parameter and preference updates after taking conditional expectation (see~\eqref{eq: key insights of algorithm 1}). Together with Young's inequality, \(\mathcal{O}(\epsilon^2)\) dual-tracking accuracy bounds their combined contribution by \(\mathcal{O}(\|\nabla\Phi(\theta_t)w_t\|^2)+\mathcal{O}(\epsilon^2)\). Consequently, \(\mathcal{O}(\epsilon^{-4})\) inner iterations suffice, without requiring the stronger \(\mathcal{O}(\epsilon^4)\) mean-squared gradient approximation error.

%While double sampling removes the \emph{stochastic} bias in MGDA, it does not remove the \emph{optimization} bias induced by using an approximate dual solution. In particular, although
%$\mathbb{E}[\bar{Y}_t^{\top}\tilde{Y}_tw_t+\rho w_t]=\nabla_{\theta} \fL(\theta_t,\eta_{t,\bar{d}})^{\top}\nabla_{\theta} \fL(\theta_t,\eta_{t,\tilde{d}})w_t+\rho w_t$, for dual DR-MOO the bias between $\nabla_{\theta} \fL(\theta_t,\eta_{t,\bar{d}}),\nabla_{\theta} \fL(\theta_t,\eta_{t,\tilde{d}})$ and $\nabla \Phi(\theta_t)$ is non-negligible. Prior work such as \citet{qimgda} analyzes unbiased stochastic gradients and {can exploit martingale structure (via optional stopping) for terms like $\mathbb{E}[\sum_{t=0}^{\tau-1}(Y_{t}^i-\nabla_{\theta} L^i(\theta_t,\eta_{t,d}))w_t^i]$.} In contrast, in our setting the lack of exact minimizers $\eta^{i,*}$
% yields non-martingale error sequences of the form
%$\mathbb{E}[\sum_{t=0}^{\tau-1} (\nabla_{\theta} L^i(\theta_t,\eta^i_{t,d})-\nabla \phi^i(\theta_t))w_t^i]$, whose contribution can accumulate over time. Consequently, to ensure $\frac{1}{T}\sum_{t=0}^{T-1}\|\nabla\Phi(\theta_t)w_t\|^2=\mathcal{O}(\epsilon^2)$, we need sufficiently tight control of the inner-loop error, e.g., $\mathbb{E}_{\eta^i_{t,d}}[\|\nabla_{\theta} L^i(\theta_t,\eta^i_{t,d})-\nabla \phi^i(\theta_t)\|^2]=\mathcal{O}(\epsilon^4)$, which in turn drives up the required number of inner iterations (See Corollary~\ref{Corollary: estimation bias of inner-loop} for details).
%\item %\red{\textbf{Controlling unbounded gradients for dual of DR-MOO~\eqref{eq: dual DR-MOO}.} To establish convergence in MOO setting, existing works~\citep{MODO,MOCO,SDMkaiyiJi,heshanSTORM} typically assumes a uniform bound on the balanced gradient \(\|\nabla \Phi(\theta_t) w\|~\forall t \le T\). This requirement is restrictive and does not generally hold for the dual of DR-MOO objectives in~\eqref{eq: dual DR-MOO}, as stochastic gradients can be unbounded. To solve this, we leverage the gradient–function value relationship~\eqref{eq: boundgrad} with a stopping-time argument to control gradient growth without bounded-gradient assumptions. Specifically, we define \(\tau=\min\{\tau_1,\tau_2,\tau_3\}\), where \(\tau_2\) and \(\tau_3\) record the first time the stochastic approximation error becomes large (or \(T\) is reached), and \(\tau_1\) records the first time that there exists \(i\) such that \(\phi^i(\theta_{t+1})-\phi^\star \ge F/2\). Under the hyperparameter choices in Theorem~\ref{thm: convergence of algorithm 1}, we show \(\mathbb{P}(\tau=T)\ge 1-\delta/2\).  
%Finally, re-arranging inequality $\mathbb{E}[\Phi(\theta_\tau)w]-\Phi^*w\leq \frac{F\delta}{8} -\frac{\alpha}{2} \mathbb{E}[\sum_{t=0}^{\tau-1}\|\nabla \Phi(\theta_t)w_t \|^2]$ (Lemma~\ref{lemma: descent lemma of algorithm1}) gives desired high-probability bound stated in Theorem~\ref{thm: convergence of algorithm 1}.}
\end{itemize}


\section{Double-Clip MGDA for DR-MOO}\label{sec: single-loop algorithm}
%In this section, we develop a Double-Clip MGDA with to solve a equivalent surrogate of the dual DR-MOO problem~\eqref{eq: dual DR-MOO}, achieving improved sample complexity.
%\subsection{Algorithm Design}
%{Writting Logic: 1. why reformulation: Double-loop is time consuming; 2. Why clipping:1.clipping on $\theta$ is not enough, previous analysis imply, increasing batch size cannot improve complexity 2. Double-sampling is complicated, how to eliminate -> Use Double cliping. 3. Summarize how MGDA changes after inducing gradient clipping. }
% Having observed unavoidable bottlenecks on estimation bias of $\nabla \Phi(\theta)$, double-sampling and high-complexity induced from inner-loop. In this section, we seek a reformulated Dual of DR-MOO objective, which both admits unbiased stochastic gradient estimator in analysis and preserves equivalence in terms of distributionally-robust Pareto-stationary condition. Building on this reformulation, we propose Algorithm~\ref{alg1-3}, which eliminates the need for double sampling through gradient clipping, and enabling efficient stochastic MGDA updates of the preference vector with provable guarantees.
The preceding double-loop analysis shows that accurately estimating \(\nabla\Phi(\theta)\) requires tight dual tracking, resulting in conservative total sample complexity. Extending the rescaled surrogate condition of \citet{jinnonconvexdro} to multi-objective setting yields a sufficient Pareto-stationarity criterion, which underlies our single-loop method (Algorithm~\ref{alg1-3}).
\begin{restatable}[Convergence Criterion Reformulation]{lemma}{reformedcondition}\label{lemma: optimality condition}
    Let Assumption \ref{assumption 1} hold. Given $\epsilon$, if one can obtain $({\theta},{\eta})$ and preference vector $w$ such that
\begin{align}
    G\!\sum_{i=1}^{m}\!|w^i\nabla_{\eta^i} \fL^i({\theta}, \eta^i) | + \| \!\sum_{i=1}^{m} w^i \nabla_{{\theta}}\fL^i({\theta},\eta^i)  \|\leq \epsilon,
    \label{eq: epsilon Optimal condition}
\end{align}
then the corresponding ${\theta}$ satisfies
$\| \nabla\Phi({\theta})w \|\leq\epsilon$. Furthermore, the condition in \eqref{eq: epsilon Optimal condition} can be achieved by optimizing the rescaled function $\fhL(\theta,\eta) = \fL(\theta, G\sqrt{m}\eta)$ such that
$\|\nabla_{{\theta}, {\eta}} \fhL({\theta}, \eta)w \|\leq {\epsilon}/{\sqrt{2}}$.
\end{restatable} 
Compared with the single-objective construction of \citet{jinnonconvexdro}, Lemma~\ref{lemma: optimality condition} introduces an additional factor of \(\sqrt{m}\) to account for aggregating dual gradients across \(m\) objectives. This scaling ensures that stationarity of the weighted surrogate \(\fhL(\theta,\eta)\) implies Pareto stationarity of \(\Phi(\theta)\). The reformulation admits unbiased stochastic gradient estimators and eliminates the need to accurately track \(\eta^{i,*}\), enabling joint single-loop updates of \(\theta\), \(\eta\), and \(w\). To further control stochastic gradient errors without double sampling, we incorporate gradient clipping and propose Algorithm~\ref{alg1-3}, whose key features are summarized below.
\begin{itemize}[leftmargin = *, topsep = 0pt]
\item \textbf{Single-loop update.} %The convergence criterion reformulation avoids the need to accurately estimate $\nabla \Phi(\theta)$, and therefore eliminates the inner-loop used to compute the dual variables. Since The rescaled function $\fhL(\theta,\eta)$ satisfies generalized $(\hL_0,\hL_1)$-smooth condition~\citep{generalizedsmooth}in $\theta$, and is $\hL_2$-smooth in $\eta$. Applying additional one-step update on $\eta$ each iteration leads to acclerated convergence along balanced gradient direction.
%The reformulated condition in Lemma~\ref{lemma: optimality condition} avoids the need to accurately estimate $\nabla \Phi(\theta)$, thereby eliminating the inner loop for computing dual variables. 
Following the analysis of \citet{qirevistfDRO}, the rescaled function \(\fhL^i(\theta,\eta^i)\) is generalized \((\hL_0,\hL_1)\)-smooth in \(\theta\) and \(\hL_2\)-smooth in \(\eta\) (Lemma~\ref{lemma: generalized-smooth} in Appendix~\ref{sec: property of hat L}). Specifically, the effective smoothness of the \(\theta\)-gradient grows with the magnitude of \(|\nabla_{\eta}\fhL^i(\theta,\eta^i)|\). Therefore, performing an additional one-step balanced-gradient update on \(\eta\) at each iteration controls this coupling, stabilizes the optimization dynamics, and facilitates convergence along the joint balanced-gradient direction.
\item %\textbf{``Double clipping'' when updating $(\theta_t,\eta_t)$ and $w_t$}.
%Inspired by the efficiency of normalized methods~\citep{jinnonconvexdro,generalizedsmooth,whyadamisgood} under heavy-tailed noise and generalized smoothness, we adopt gradient clipping and large-batch sampling for stochastic gradients of the rescaled function $\fhL(\theta,\eta)$, which exhibits affine-linear variance (Lemma~\ref{lemma: affine bounded noise} in Appendix~\ref{sec: property of hat L}). Except clipping the balanced gradients $X_t w_t$ and $Z_t w_t$ when updating $(\theta_t,\eta_t)$, we introduce gradient clipping on gram matrix $X_{t}^{\top}X_t$, $Z_t^{\top}Z_t$ when updating $w_t$ to make their update magnitudes and variance consistent. As a result, double-clipping mechanism make each update of $(\theta_t, \eta_t)$ and $w_t$ on comparable scales while introducing a vanishing bias approaching pareto-stationary region, thereby eliminating the need for the double-sampling strategy commonly used for bias reduction in Algorithm~\ref{alg1}.
%To better exploit the one-step update on \(\eta\), we modify the stochastic MGDA update for \(w_t\) by adding an additional term \(\mu_t Z_t^{\top}Z_t w_t\). This term helps balance the two components \(\|\nabla_{\theta}\fhL(\theta_t,\eta_{t+1})w_t\|\) and \(\|\nabla_{\eta}\fhL(\theta_t,\eta_t)w_t\|\) in analysis, which is crucial for establishing convergence to a distributionally robust Pareto-stationary point measured by \(\|\nabla_{\theta,\eta}\fhL(\theta_t,\eta_{t+1})w_t\|\).
\textbf{``Double (gradient)-clipping'' when updating $(\theta_t,\eta_t)$ and $w_t$.}
Inspired by normalized methods for heavy-tailed noise and generalized
smoothness~\citep{jinnonconvexdro,generalizedsmooth,whyadamisgood}, we combine
large-batch sampling with gradient clipping to handle the affine-bounded
variance of the stochastic gradients of $\fhL(\theta,\eta)$
(Lemma~\ref{lemma: affine bounded noise} in
Appendix~\ref{sec: property of hat L}). Unlike clipping methods for single-objective optimization, which clip only the gradient \(\nabla_{\theta}\fhL(\theta,\eta),\nabla_{\eta} \fhL(\theta,\eta)\) used to update the optimization variables, our method also clips the stochastic Gram-matrix contributions to the preference-vector $w_t$'s update. Specifically, we first clip the balanced gradients \(X_tw_t\) and \(Z_tw_t\) when updating \(\theta_t\) and \(\eta_t\), and then reuse the clipping factors \(\alpha_t\) and \(\mu_t\) to control the magnitudes of \(X_t^\top X_tw_t\) and \(Z_t^\top Z_tw_t\) in the update of \(w_t\). The additional term \(\mu_tZ_t^\top Z_tw_t\) incorporates dual-gradient information, enabling \(w_t\) to balance both \(\|\nabla_{\theta}\fhL(\theta_t,\eta_{t+1})w_t\|\) and \(\|\nabla_{\eta}\fhL(\theta_t,\eta_t)w_t\|\). This preference-level clipping is crucial for controlling the magnitude and variance of the \(w_t\) update, keeping it commensurate with the primal–dual updates and eliminating the need for double sampling.
% Intuitively, once $|\nabla_{\eta^i}L^i(\theta,\eta^i)|$ vanishes to a small region, $(\hL_0,\hL_1)$-smooth reduced to $\hL_0$-smooth. For this purpose, in Algorithm \ref{alg1-3}, we add an additional extrapolation on $\eta=[\eta_i]_{i=1}^{m}$ (Line 4-5 in Algorithm \ref{alg1-3}) at each iteration, aiming to fasten the convergence along CA direction, $\| \nabla_{\eta}\fhL(\theta_t,\eta_t^i)w_t\|$.
% \citet{qirevistfDRO} showed that each $\fhL^i(\theta,\eta^i)$ satisfies so called $(\hL_0,\hL_1,\hL_2)$-smooth condition, where   $(\hL_0,\hL_1)$-smooth condition \citep{generalizedsmooth} holds in terms of $\theta$, $\|\nabla_{\theta}\fhL^i(\theta,\eta^i)-\nabla\fhL^i(\theta',\eta^i)\|\leq(\hL_0+\hL_1|\nabla_{\eta^i}\fhL^i(\theta,\eta^i)|)\|\theta-\theta'\|$, and $\hL_2$-smooth holds in terms of $\eta^i$, $\|\nabla_{\eta}\fhL^i(\theta, \eta^i)-\nabla_{\eta}\fhL^i(\theta,\eta^{i'})\|\leq L_2\|\eta^i-\eta^{i'}\|$. 
% Since generalized-smooth part critically depends on $|\nabla_{\eta^i}\fhL(\theta,\eta^i) |$, controlling $|\nabla_{\eta^i}\fhL(\theta,\eta^i)|$ to be bounded makes problem easier to address both in theory and practice. Intuitively, once $|\nabla_{\eta^i}L^i(\theta,\eta^i)|$ vanishes to a small region, $(\hL_0,\hL_1)$-smooth reduced to $\hL_0$-smooth. For this purpose, in Algorithm \ref{alg1-3}, we add an additional extrapolation on $\eta=[\eta_i]_{i=1}^{m}$ (Line 4-5 in Algorithm \ref{alg1-3}) at each iteration, aiming to fasten the convergence along CA direction, $\| \nabla_{\eta}\fhL(\theta_t,\eta_t^i)w_t\|$.
\end{itemize}
%%%%%%%%%%% un-compressed version %%%%%%%%%%%%
%\begin{algorithm}[t]
%\caption{Double-Clip MGDA for DR-MOO}
%\label{alg1-3}
%\small
%\resizebox{0.9\linewidth}{!}{
%\begin{minipage}{0.9\linewidth}
%\begin{algorithmic}[1]
%\STATE Initialize ${\theta}_0$, ${\eta}_0$, ${w}_0$ and $\rho$.
%\STATE Learning rates $\beta$,$\gamma$
%\STATE Clipping rule $\alpha_t=\min\{c_1,\frac{c_2}{\|X_tw_t\|}\}$,$\mu_t=\min\{ f_1,\frac{f_2}{\|Z_tw_t\|}\}$,
%\FOR{$t = 0, ..., T-1$}
%    \STATE Evaluate $Z_{t}=\nabla_{ {\eta}}\fhL({\theta}_t,{\eta}_t;\{{\xi}_{t}\}_{B})$ with $B=N_2$. 
%    \STATE ${\eta}_{t+1}={\eta}_t - {\gamma}{\mu_t}Z_t {w}_t$
%    \STATE Evaluate $X_{t}=\nabla_{{\theta}}\fhL({\theta}_t, {\eta}_{t+1};\{\bar{\xi}_{t}\}_{B})$ with $B=N_1$.
%    \STATE ${{\theta}_{t+1}={\theta}_t-\gamma {\alpha_t}X_{t}{w}_t }$
%    \STATE  ${w}_{t+1}\!=\!\Pi_{\mathcal{W}}\big({w}_t -\beta(\alpha_t {X}_{t}^{\top}\!{X}_{t} {w}_t \!+ \mu_t{Z}_t^{\top}\!{Z}_tw_t \!+ \rho {w}_t)\big)$
%\ENDFOR 
%\end{algorithmic}
%\end{minipage}
%}
%\end{algorithm}
%When optimizing single dual DRO objective \eqref{eq: single-objective dual DRO}, \citet{qirevistfDRO} showed that combining constant learning rate for update of $\eta$, and gradient-clipping for update $\theta$ attains $\mathcal{O}(\epsilon^{-4})$ sample complexity under mild assumptions. However, for dual DR-MOO formulation \eqref{eq: dual DR-MOO}, we find that such design becomes suboptimal. In particular, due to the structure of MGDA \eqref{eq: MGDA surrogate}, combination of constant and gradient clipping yields $\mathcal{O}(\epsilon^{-6})$ sample complexity for finding $(\theta,\eta)$ satisfies $\epsilon$-distributionally robust Pareto-stationary condition. We present our convergence analysis in Appendix~\ref{Appendix: Convergence Analysis of Algorithm 2}.
%{TL;DR: I want to highlight that, when introducing gradient clipping, we can control stochastic gradient magnitude in a small-level, instead of using double-sampling to ensure unbiasness.}



%To tackle this challenge, we propose Algorithm \ref{alg1-3}, inducing gradient clipping not only for problem-dependent parameter $\theta$ and $\eta$, but also for preference vector $w_t$ (Line 4-8 in Algorithm \ref{alg1-3}). 
%%% I'm not sure if we should add this
%And we observe the iterative update ${w}_{t+1}=\Pi_{\mathcal{W}}\big({w}_t-\beta( \alpha_t {X}_{t}^{\top}{X}_{t} {w}_t+ \mu_t{Z_t}^{\top} {Z}_tw_t+\rho {w}_t)\big)$ can be interpreted as a more general stochastic formulation than the MGDA surrogate formulation \eqref{eq: MGDA surrogate}. Due to space limitation, we leave the discussion for future work. 
\subsection{Convergence Analysis}
We now establish the convergence of Algorithm~\ref{alg1-3} under the generalized-smooth geometry of \(\fhL(\theta,\eta)\), without assuming uniformly bounded gradients as in prior analyses~\citep{MODO,MOCO,SDMkaiyiJi}. Instead, we use the relationship between gradient norms and function-value gaps~\citep{Haochuangensmooth} to jointly control stochastic gradients and function values with high probability. The full result, hyperparameter choices, and proof appear in Appendix~\ref{Appendix: Convergence of Algorithm 3}.
\begin{restatable}[Convergence of Algorithm~\ref{alg1-3}]{theorem}{convgalgorithm3}\label{thm: convergence of alg3}
     Let Assumption \ref{assumption 1} hold. Denote
    $\hat{\Delta}_{\theta_0,\eta_0}$ as the initial objective gap in terms of $\fhL(\theta,\eta)$. Given $\delta,\epsilon$ satisfy $\delta\epsilon\leq 1/m\cdot\min\{\mathcal{O}({1}/{\Delta_{\theta_0,\eta_0}^{1/2}}),\mathcal{O}((\beta/\gamma)^{1/2})\}$. Set the hyperparameters in Algorithm~\ref{alg1-3} as $c_1=f_1={1}/{2}$, $c_2=f_2=\delta\epsilon$, 
    $\rho=\mathcal{O}(\delta^2\epsilon^2)$, $\beta,\gamma = \mathcal{O}(1)$. Choose batch sizes $N_1,N_2=\mathcal{O}(\delta^{-3}\epsilon^{-2})$. Then, after
    $T
    =\max\{ \Theta(\Delta_{\theta_0,\eta_0}\gamma^{-1}\delta^{-2}\epsilon^{-2}),\Theta(\beta^{-1}\delta^{-2}\epsilon^{-2})\}$ iterations, with probability at least $1-\delta$, we have 
    \[\frac{1}{T}\sum_{t=0}^{T-1}\|\nabla_{\theta,\eta}\fhL(\theta_t, \eta_{t+1})w_t\|\leq 34\epsilon.\]
    \vspace{-1.5em}
\end{restatable}
Theorem~\ref{thm: convergence of alg3} establishes an \(\mathcal{O}(\epsilon^{-2})\) iteration complexity and an \(\mathcal{O}(\epsilon^{-4})\) total sample complexity for Algorithm~\ref{alg1-3}, matching the complexity of smooth nonconvex stochastic optimization~\citep{lowerbound}. Under our current convergence analysis framework targeting
\begin{align}
    \frac{1}{T}\mathbb{E}[\sum_{t=0}^{T-1}\gamma\|\nabla_{\theta,\eta}\fhL(\theta_t,\eta_{t+1})w_t\|^2]\leq\mathcal{O}(\gamma\epsilon^2),
    \label{eq: general analysis framework}
\end{align}
and the adopted learning-rate design, batch sizes \(N_1,N_2=\mathcal{O}(\epsilon^{-2})\) are needed to control the stochastic errors and preserve the stated iteration complexity. Despite this theoretical bottleneck, the batch-size ablation in Figure~\ref{fig: batch-size and inner-step ablation} (Appendix~\ref{Appendix: ablation on batch size and inner-steps}) shows that \(N_1=N_2=256\) provides stable performance across all experiments. We next explain the key technical novelties underlying our convergence guarantee.
\begin{itemize}[leftmargin = *, topsep = 0pt] 
\item \textbf{Bias control via double clipping and large-batch sampling.} To establish the target bound in~\eqref{eq: general analysis framework}, the
terms arising from the preference update, including
$\gamma\beta\mathbb{E}[\|\mu_tZ_t^\top Z_tw_t\|^2]$ and
$\gamma\beta\mathbb{E}[\|\alpha_tX_t^\top X_tw_t\|^2]$, must be controlled at
the $\mathcal{O}(\gamma\epsilon^2)$ level. Using the clipped factor $\mu_t$, we obtain
\begin{align}
\gamma\beta\mathbb{E}[\|\mu_tZ_t^\top Z_tw_t\|^2]
&\leq
2\gamma\beta\delta^2\epsilon^2
\mathbb{E}[\|Z_t-\nabla_{\eta}\fhL(\theta_t,\eta_t)\|_F^2]\nonumber\\
&\quad+
2\gamma\beta\delta^2\epsilon^2
\mathbb{E}[\|\nabla_{\eta}\fhL(\theta_t,\eta_t)\|_F^2].\nonumber
\end{align}
Here, the factor $\delta^2\epsilon^2$ follows from
$\mu_t\|Z_tw_t\|\leq\delta\epsilon$. An analogous bound holds for
$\gamma\beta\mathbb{E}[\|\alpha_tX_t^\top X_tw_t\|^2]$. Meanwhile, the descent
analysis contains the stochastic cross terms
$\gamma\mu_t\langle(Z_t-\nabla_{\eta}\fhL(\theta_t,\eta_t))w,Z_tw_t\rangle$
and
$\gamma\alpha_t\langle(X_t-\nabla_{\theta}\fhL(\theta_t,\eta_{t+1}))w,
X_tw_t\rangle$. After applying Cauchy--Schwarz and the clipping rules, these
terms are proportional to the stochastic estimation errors. Batch sizes
$N_1,N_2=\mathcal{O}(\epsilon^{-2})$ ensure that
$\mathbb{E}[\|(X_t-\nabla_{\theta}\fhL(\theta_t,\eta_{t+1}))w\|]$ and
$\mathbb{E}[\|(Z_t-\nabla_{\eta}\fhL(\theta_t,\eta_t))w\|]$ are both
$\mathcal{O}(\epsilon)$, keeping the corresponding contributions at the
$\mathcal{O}(\gamma\epsilon^2)$ level. Together, double clipping and large-batch
sampling remove the need to set $\beta,\gamma=\mathcal{O}(\epsilon^2)$ and
thereby improve the iteration complexity.
%The critical design of Algorithm~\ref{alg1-3} comes from Double-Clipping. As a result, we found it eliminates the need of double-sampling, since double gradient-clipping makes the update magnitudes of $w_t,\eta_t,\theta_t$ stay controllable. To elaborate, the last step for yielding convergence of MGDA takes the general form 
% \begin{align}
%     \frac{1}{T}\mathbb{E}[\sum_{t=0}^{T-1}\gamma\|\nabla \fhL(\theta_t,\eta_{t+1})w_t\|^2]\leq \mathcal{O}(\gamma \epsilon^2),\label{eq: general form of convergence}
% \end{align}
% where $\gamma$ is the learning rate for update optimization variables, i.e., $(\theta,\eta)$. This forces dominant oscillation terms induced by stochastic MGDA, $\gamma \beta\mathbb{E}\|\mu_tZ_t^{\top}Z_tw_t \|^2,\gamma\beta \mathbb{E}\|\alpha_t X_t^{\top}X_tw_t\|^2$ must stay consistently at $\mathcal{O}(\epsilon^2)$-level. When setting $\mu_t=\{\frac{1}{2},\frac{\delta\epsilon}{\| Z_tw_t\|} \}$, we found $\gamma \beta\mathbb{E}\|Z_t^{\top}Z_tw_t \|^2$ can be upper bounded as follows
% \begin{align}
%     &\gamma\beta\mathbb{E}[\|\mu_tZ_t^{\top}Z_tw_t\|^2]
%     \leq \nonumber\\
%     &\quad 2\gamma\beta \delta^2\epsilon^2\mathbb{E}[\| \HUpsilon_t\|_F^2]+2\gamma\beta\delta^2\epsilon^2\|\nabla_{\eta}\fhL(\theta_t,\eta_t)\|_F^2,\nonumber
% \end{align}
% where the associated $\delta^2\epsilon^2$ on RHS is induced from normalization upper bound of $\mu_t$ such that $\mu_t\leq {\delta\epsilon}/{\|Z_tw_t\|}$. Similar arguments also holds for $\gamma\beta\mathbb{E}[\|\alpha_tX_t^{\top}X_t\|^2]$. As a result, above inequality eliminates the need of $\beta=\mathcal{O}(\epsilon^2)$ to achieve~\eqref{eq: general form of convergence}, yielding an improved iteration complexity of $T$.
\item \textbf{Objective-independent clipping radii.}
Unlike single-objective methods~\citep{qirevistfDRO,revisitclip,generalizedsmooth}, whose clipping radii depend on problem-specific smoothness and variance constants, our method uses a universal radius of \(c_1=f_1=1/2\) for \((\theta_t,\eta_t)\) and \(w_t\), independent of these parameters. This is essential because $X_t$, $Z_t$, and
their Gram-matrix contributions generally have different smoothness and
variance scales. Using separate parameter-dependent radii would prevent their
error terms from being uniformly controlled at the
$\mathcal{O}(\gamma\epsilon^2)$ level in~\eqref{eq: general analysis framework}.
\item \textbf{Controlling potentially unbounded gradients.}
At the exact primal--dual pair
$(\theta,\eta^{i,*})$, we have $\|\nabla\phi^i(\theta)\|\leq G$. However, the gradients of
$\fhL(\theta_t,\eta_t)$ and $\fhL(\theta_t,\eta_{t+1})$ need not be uniformly
bounded at arbitrary dual iterates. Nevertheless,
$\|\nabla_{\theta}\fhL^i(\theta_t,\eta^i_{t+1})\|$ is coupled with
$\|\nabla_{\eta}\fhL^i(\theta_t,\eta^i_{t+1})\|$. We exploit this coupling through the
gradient--function-value relationship~\eqref{eq: boundgrad} and a stopping-time
argument. Specifically, we define $\htau$ to track abnormally large stochastic
gradient errors and function-value gaps and show that
$\mathbb{P}(\htau<T)<\delta/2$. We then rearrange the stopped descent inequality
\begin{align}
&\mathbb{E}[\fhL(\theta_{\htau},\eta_{\htau})w]-\fhL^*w\nonumber\\
&\quad\leq \frac{\bar F\delta}{8m}
-\frac{\gamma}{2}\mathbb{E}[\sum_{t=0}^{\htau-1}
\alpha_t\|X_tw_t\|^2]
-\frac{\gamma}{2}\mathbb{E}[\sum_{t=0}^{\htau-1}
\mu_t\|Z_tw_t\|^2],\nonumber
\end{align}
from Lemma~\ref{lemma: descent lemma of alg1-3} and combine it with the variance
bounds for $X_t^i$ and $Z_t^i$ and the one-step relationship between
$\nabla_{\eta^i}\fhL^i(\theta_t,\eta^i_{t+1})$ and
$\nabla_{\eta^i}\fhL^i(\theta_t,\eta^i_t)$ to obtain the desired result.
 
% $\mathbb{E}[\fhL(\theta_{\tau},\eta_\tau)w]-\hL^*w\leq\frac{\bar{F}\delta}{8} -\frac{\gamma}{2}\mathbb{E}[\sum_{t=0}^{\tau-1}\alpha_t\| X_tw_t\|^2]-\frac{\gamma}{2}\mathbb{E}[\sum_{t=0}^{\tau-1}\mu_t\|Z_tw_t\|^2]$ (inequality \eqref{eq: stopping-time bound alg3} in Appendix~\ref{Appendix: Convergence of Algorithm 3}), we show that $P(\tau=T)\geq 1-\frac{\delta}{2}$. Then, reorganize inequality~\eqref{eq: stopping-time bound alg3}, leverage variance $\mathbb{E}\|Z_t^i-\nabla_{\eta} \fhL^i(\theta_t,\eta^i_t)\|^2,\mathbb{E}\|X_t^i-\nabla_{\eta} \fhL^i(\theta_{t},\eta^i_{t+1})\|^2\leq O(\epsilon^2)$, and one-step gradient relationship between $\nabla_{\eta^i}\fhL^i(\theta_{t},\eta^i_{t+1})$ and $\nabla_{\eta^i}\fhL^i(\theta_{t},\eta^i_t)$ give desired result.
\end{itemize}
%%%%%%%%%%%%% Full Table %%%%%%%%%%%%%%%%%%%%
%\vspace{-1.5em}
\begin{table*}[!t]
\caption{Test Accuracy under FGSM attack}
\label{Table: Multi-mnist FGSM}
\centering
%\small
{
\scriptsize
\renewcommand{\arraystretch}{1.08}
\setlength{\tabcolsep}{1.5pt}


\begin{tabular*}{\textwidth}{
    @{\extracolsep{\fill}}
    lccccc|ccccc
    @{}
}
\toprule
\multirow{2}{*}{Methods/Attack Level (in \%)}
& \multicolumn{5}{c|}{Multi-MNIST 2-digits ($70$-epochs training)}
& \multicolumn{5}{c}{Multi-MNIST 3-digits ($100$-epochs training)} \\
\cmidrule(lr){2-6}\cmidrule(lr){7-11}
& $0.00$ & $0.01$ & $0.03$ & $0.05$ & $0.08$
& $0.00$ & $0.01$ & $0.03$ & $0.05$ & $0.08$ \\
\midrule
Double-Clip MGDA
& $\boldsymbol{95.66}$ & $\boldsymbol{83.48}$ & $\boldsymbol{65.95}$ & $\boldsymbol{60.40}$ & $\boldsymbol{57.13}$
& $\boldsymbol{98.76}$ & $\boldsymbol{97.59}$ & $\boldsymbol{94.40}$ & $\boldsymbol{91.05}$ & $\boldsymbol{86.65}$ \\

Double-loop MGDA
& 92.80 & 72.81 & 57.71 & 54.49 & $51.63$
& 97.49 & 95.38 & 89.99 & 85.25 & 79.88 \\

MoCo~\citep{MOCO}
& 94.49 & 77.69 & 61.74 & 58.63 & 56.43
& 98.27 & 96.62 & 92.75 & 88.75 & 83.43 \\

NashMTL~\citep{nashMTL}
& 91.21 & 62.58 & 51.67 & 49.54 & 47.09
& 96.17 & 92.31 & 84.53 & 78.91 & 73.48 \\

FAMO~\citep{FAMO}
& 89.05 & 61.04 & 50.82 & 48.66 & 46.48
& 95.90 & 91.86 & 84.29 & 78.94 & 73.52 \\

SDMGrad~\citep{SDMkaiyiJi}
& 89.59 & 64.02 & 52.00 & 49.91 & 47.31
& 96.46 & 92.89 & 85.06 & 79.37 & 73.50 \\

MoDo~\citep{MODO}
& 91.10 & 64.08 & 51.95 & 49.73 & 47.49
& 96.50 & 93.15 & 86.03 & 80.56 & 74.81 \\

MGDA~\citep{MGDA1stwork,qimgda}
& 89.44 & 62.61 & 52.19 & 50.81 & 48.30
& 96.34 & 92.85 & 85.42 & 80.26 & 75.18 \\
\bottomrule
\end{tabular*}}
\end{table*}
%\begin{table}[ht]
%\caption{Test Performance over $128\times128$ CelebA Dataset
%}
%\label{Table: Test on CelebA}
%\centering
%\small
%\resizebox{0.7\columnwidth}{!}{
%\begin{tabular}{lccccc}
%\toprule
%Methods/Evaluation metric(in \%) & Averaged Accuracy & Balanced Accuracy & AUC \\
%\midrule
%Double-Clip MGDA & $\boldsymbol{88.41}$\% & $\boldsymbol{89.55}$\%  & $\boldsymbol{94.63}$\% \\
%Double-loop MGDA & 86.61\% & 87.97\% & 93.28\%  \\
%MoCo~\citep{MOCO} & 87.27\% & 88.50\% & 93.31\% \\
%NashMTL~\citep{nashMTL} & 85.58\% & 87.00\% & 92.41\% \\
%FAMO~\citep{FAMO} & 85.32\% & 86.03\% & 91.41\%  \\
%SDMGrad~\citep{SDMkaiyiJi} & 85.89\% &  87.06\% & 92.45\%  \\
%MoDo~\citep{MODO} & 85.66\% & 87.08\% & 92.42\%  \\
%MGDA~\citep{MGDA1stwork,qimgda} & 85.53\% & 86.87\% & 92.58\% \\
%\bottomrule
%\end{tabular}}
%\label{table: accuracy table for CelebA}
%\end{table}
%%%%%%%%% Compressed Version %%%%%%%%%%
%\begin{wraptable}{r}{0.7\textwidth}
%\vspace{-1.5em}
%\centering
%\begin{minipage}{0.8\textwidth}
%\vspace{0.3em}
%\captionsetup{type=table}
%\caption{Test Performance against Label Imbalance}
%\label{table:accuracy-celeba}
%\vspace{-0.4em}

%\resizebox{\columnwidth}{!}{
%\begin{tabular}{lccc}
%\toprule
%Methods/Evaluation metric (in \%) & Averaged Accuracy & Balanced Accuracy & AUC \\
%\midrule
%Double-Clip MGDA & $\boldsymbol{88.41}$ & $\boldsymbol{89.55}$ & $\boldsymbol{94.63}$ \\
%Double-loop MGDA & 86.61 & 87.97 & 93.28 \\
%MoCo~\citep{MOCO} & 87.27 & 88.50 & 93.31 \\
%NashMTL~\citep{nashMTL} & 85.58 & 87.00 & 92.41 \\
%FAMO~\citep{FAMO} & 85.32 & 86.03 & 91.41 \\
%SDMGrad~\citep{SDMkaiyiJi} & 85.89 & 87.06 & 92.45 \\
%MoDo~\citep{MODO} & 85.66 & 87.08 & 92.42 \\
%MGDA~\citep{MGDA1stwork,qimgda} & 85.53 & 86.87 & 92.58 \\
%\bottomrule
%\end{tabular}
%}
%\vspace{0.3em}
%\end{minipage}
%\vspace{-1em}
%\end{wraptable}
\begin{table*}[!t]
\centering
%\vspace{-1em}
\caption{Test Performance against Label Imbalance}
\label{table:accuracy-celeba}
%\vspace{-0.4em}
{
\scriptsize
\renewcommand{\arraystretch}{1.08}
\setlength{\tabcolsep}{1.5pt}


\begin{tabular*}{\textwidth}{
    @{\extracolsep{\fill}}
    lccccc|ccccc
    @{}
}
\toprule
Methods/Evaluation metric (in \%) & Averaged Accuracy & Balanced Accuracy & AUC \\
\midrule
Double-Clip MGDA & $\boldsymbol{88.41}$ & $\boldsymbol{89.55}$ & $\boldsymbol{94.63}$ \\
Double-loop MGDA & 86.61 & 87.97 & 93.28 \\
MoCo~\citep{MOCO} & 87.27 & 88.50 & 93.31 \\
NashMTL~\citep{nashMTL} & 85.58 & 87.00 & 92.41 \\
FAMO~\citep{FAMO} & 85.32 & 86.03 & 91.41 \\
SDMGrad~\citep{SDMkaiyiJi} & 85.89 & 87.06 & 92.45 \\
MoDo~\citep{MODO} & 85.66 & 87.08 & 92.42 \\
MGDA~\citep{MGDA1stwork,qimgda} & 85.53 & 86.87 & 92.58 \\
\bottomrule
\end{tabular*}%
}
\vspace{-1em}
\end{table*}
\section{Experiments}
In this section, we evaluate Algorithms~\ref{alg1} and~\ref{alg1-3} for solving the dual formulation of DR-MOO in~\eqref{eq: dual DR-MOO} in deep learning settings. To illustrate the practical relevance of DR-MOO, we focus on three representative sources of distributional heterogeneity: adversarial perturbations on Multi-MNIST\citep{mnist}, label imbalance on CelebA~\citep{celebA}, and shifts across domains in office-31 dataset~\citep{officedata31}. In all experiments, we use a ResNet-18 encoder~\citep{ResNet} with multiple MLP classification heads, set $\ell(\cdot)$ to the cross-entropy loss, and choose $f^*(\cdot)$ as the convex conjugate associated with the $\chi^2$-divergence. We compare our algorithms against MGDA~\citep{MGDA1stwork,stochasticMGDA}, MoDo~\citep{MODO}, MoCo~\citep{MOCO}, SDMGrad~\citep{SDMkaiyiJi}, NashMTL~\citep{nashMTL}, and FAMO~\citep{FAMO}. For all baselines, we treat \((\theta,\eta)\) as optimization variables, perform joint updates, and fine-tune all hyperparameters. Due to space limitations, we defer the Office-31, linear and logistic regression experiments and ablation studies to Appendices~\ref{Appendix: Synthetic Experiments} and~\ref{Appendix: Ablation study}.

\vspace{-1em}
\subsection{Robustness against Adversarial Attacks}
%In this section, we further evaluate Algorithms~\ref{alg1} and~\ref{alg1-3} on dual DR-MOO problem in~\eqref{eq: dual DR-MOO} with the cross-entropy loss in deep neural network settings. Specifically, we train neural networks consisting of a ResNet18~\citep{ResNet} encoder and multiple MLP classification heads. Our goal is to evaluate the performance of the proposed algorithms on the Multi-MNIST 2-digit and 3-digit datasets~\citep{mnist} when facing adversarial samples at test-time. To further verify the effectiveness of our proposed methods on MTL benchmarks, we include additional baselines NashMTL~\citep{nashMTL} and FAMO~\citep{FAMO}, which are widely used in practice.
%At initialization, we load ResNet18 encoder pre-trained on ImageNet~\citep{imageNet}. For the MLP classification heads, we use Kaiming uniform initialization for each layer. We use a unified batch size of $B=256$, $\rho = 1e^{-4}$ and set the learning rate for the classification heads to $1e^{-2}$ for all methods. We train the classification heads using SGD~\citep{lanSGD}. For Double-Loop MGDA and SDMGrad~\citep{SDMkaiyiJi}, we set the number of inner iterations to $5$. For 2-digit Multi-MNIST, we resize input images to $64\times 64$ and train the proposed neural networks for $70$ epochs across all methods. For 3-digit Multi-MNIST, we resize the input images to $96\times96$ and train all methods for $100$ epochs. 
%\begin{table}[ht]
%\caption{Test Accuracy under FGSM attack}
%\label{Table: Multi-mnist FGSM}
%\centering
%\small
%\resizebox{\columnwidth}{!}{
%\begin{tabular}{lccccc}
%\toprule
%Methods/Attack Level & $0.00$ & $0.01$ & %$0.03$ & $0.05$ & $0.08$ \\
%\midrule
%\multicolumn{6}{c}{Averaged Accuracy after 70 Epochs Training ($64\times 64$ Multi-MNIST 2-digits)} \\
%\midrule
%Double-Clip MGDA & $\boldsymbol{95.66}\%$ & $\boldsymbol{83.48}\%$ & $\boldsymbol{65.95}\%$ & $\boldsymbol{60.40}\%$ & $\boldsymbol{57.13}\%$ \\
%Double-loop MGDA & 92.80\% & 72.81\% & 57.71\% & 54.49\% & $51.63\%$ \\
%MoCo~\citep{MOCO} & 94.49\% & 77.69\% & 61.74\% & 58.63\% & 56.43\% \\
%NashMTL~\citep{nashMTL} & 91.21\% & 62.58\% & 51.67\% & 49.54\% & 47.09\% \\
%FAMO~\citep{FAMO} & 89.05\% & 61.04\% & 50.82\% & 48.66\% & 46.48\% \\
%SDMGrad~\citep{SDMkaiyiJi} & 89.59\% &  64.02\% & 52.00\% & 49.91\% & 47.31\% \\
%MoDo~\citep{MODO} & 91.10\% & 64.08\% & 51.95\% & 49.73\% & 47.49\% \\
%MGDA~\citep{MGDA1stwork,qimgda} & 89.44\% & 62.61\% & 52.19\% & 50.81\% & 48.30\% \\
%\midrule
%\multicolumn{6}{c}{Averaged Task Accuracy after $100$ Epochs Training ($96\times 96$ Multi-MNIST 3-digits)} \\
%\midrule
%Double-Clip MGDA & $\boldsymbol{98.76\%}$ & $\boldsymbol{97.59\%}$ & $\boldsymbol{94.40\%}$ & $\boldsymbol{91.05\%}$ & $\boldsymbol{86.65\%}$ \\
%Double-loop MGDA & 97.49\% & 95.38\% &  89.99\%  & 85.25\% & 79.88\% \\
%MoCo~\citep{MOCO} & 98.27\% & 96.62\% & 92.75\%  & 88.75\% & 83.43\% \\
%NashMTL~\citep{nashMTL} & 96.17\% & 92.31\% & 84.53\%  & 78.91\% & 73.48\% \\
%FAMO~\citep{FAMO} & 95.90\% & 91.86\% & 84.29\%  & 78.94\% & 73.52\% \\
%SDMGrad~\citep{SDMkaiyiJi} & 96.46\% & 92.89\% & 85.06\% & 79.37\% & 73.50\% \\
%MoDo~\citep{MODO} & 96.50\% & 93.15\% & 86.03\% & 80.56\% & 74.81\% \\
%MGDA~\citep{MGDA1stwork,qimgda} & 96.34\% & 92.85\% & 85.42\% & 80.26\% & 75.18\% \\
%\bottomrule
%\end{tabular}}
%\label{table: accuracy table for multi-mnist}
%\end{table}
%%%%%%%%% Compressed Version %%%%%%%%%
We conduct experiments on Multi-MNIST 2-digit and 3-digit datasets~\citep{mnist}, evaluating robustness of proposed formulation and algorithms under adversarial perturbations. For training configurations and hyperparameters, we refer to Table~\ref{Table: backbone-training-config} and ~\ref{Table: Learning Rate Hyper-parameter} in Appendix~\ref{Appendix: Hyper-parameter of main experiments.} for more details. At test time, we load model parameters, compute task gradient summation and apply FGSM~\citep{FGSM} attack on test data. The corresponding test accuracies for all methods under different attack strengths are reported in Table~\ref{Table: Multi-mnist FGSM}. These results show that models trained using existing MGDA solvers are less robust under adversarial perturbations, whereas the proposed methods better preserve task performance under increasing attack strength.
%%%%%% Column Wise Table %%%%%%%%%
%For Double-Loop MGDA (Algorithm~\ref{alg1}), we set $\gamma=3e^{-3}$, $\alpha=5e^{-4}$, and $\beta=1e^{-5}$. For Single-Loop Double-Clip MGDA (Algorithm~\ref{alg1-3}), we set the learning rate to $\gamma=5e^{-3}$, $\beta=1e^{-2}$, and the clipping thresholds to $f_1=1.0$, $f_2=0.5$, $c_1=2.0$, and $c_2=1.0$. For MGDA, MoDo, MoCo, and SDMGrad, the learning rates for updating $(\theta,\eta)$ are set to $1e^{-4}$, $1e^{-4}$, $5e^{-4}$, and $1e^{-4}$, respectively. For the preference vector, we set the corresponding learning rates to $1e^{-6}$, $1e^{-6}$, $1e^{-4}$, and $1e^{-6}$, respectively. For FAMO, we set the learning rate for updating $(\theta,\eta)$ to $1e^{-4}$, the learning rate for the preference vector to $2.5e^{-3}$, and $\gamma=1e^{-2}$. The same learning rate for updating $(\theta,\eta)$ is also used for NashMTL. We use a mirror-descent algorithm to approximately solve the linear system, with $100$ iterations at each time step.

%At test time, we reload the model parameters stored at the last epoch and generate adversarial test samples using the FGSM attack~\citep{FGSM}. Table~\ref{table: accuracy table for multi-mnist} reports the classification accuracies under different attack levels. The results show that our proposed methods, Double-Loop MGDA and Double-Clip MGDA, achieve competitive performance compared with other baselines under adversarial distribution shifts.
\subsection{Robustness against Label Imbalance}
%In this section, we further scale our methods to CelebA~\citep{celebA} dataset, evaluating Double-Loop MGDA, Double-Clip MGDA, and other baselines for solving the dual DR-MOO problem~\eqref{eq: dual DR-MOO} under label imbalance. Specifically, we resize image to be $128\times 128$, and select the $8$ most imbalanced attributes, including ``\textit{Bald, Wearing Hat, Eyeglasses, Receding Hairline, Narrow Eyes, Blond Hair, Bags Under Eyes, and Big Nose}.'' We reuse the neural network backbone and hyper-parameters used for 3-digit Multi-MNIST experiment. During training, to address label imbalance, we simplify the nominal distribution $\mathbb{P}$ over labels as a Bernoulli distribution with $p={1}/{2}$. To correct label imbalance, we use an importance-sampling strategy when approximating the dual DR-MOO loss~\eqref{eq: dual DR-MOO}. In particular, for batched samples, we compute the empirical fraction of the binary labels, and denoted them as $\bar{p}$ and $1-\bar{p}$, and multiply the corresponding importance ratios ${1}/{2\bar{p}}$ and ${1}/{2(1-\bar{p})}$ when computing the finite-sum approximation of dual DR-MOO objective~\eqref{eq: dual DR-MOO}. We train the neural network backbone for $100$ epochs.
%\begin{figure}[ht]
%    %\centering
%    \begin{subfigure}{0.23\textwidth}
        %\includegraphics[width=1.05\linewidth]{Exp_Fig/mean_bal_acc.png}
    %\end{subfigure}
    %    \hfill
    %\begin{subfigure}{0.23\textwidth}
        %\includegraphics[width=1.05\linewidth]{Exp_Fig/mean_auc.png}
    %\end{subfigure}  
    %\caption{Balanced Accuracy (Left) and AUC score (Right) over Test Data for Each Method.}
    %\label{fig: Multi-CelebA}
%\end{figure}

%At test time, we evaluate the averaged AUC and balanced accuracy scores over the $8$ tasks for each method. The corresponding results are reported in Table~\ref{table: accuracy table for CelebA}. As shown in the figure, our methods are competitive in handling label imbalance. In particular, Double-Clip MGDA achieves a strong AUC score using a ResNet18 backbone, validating its robustness against label imbalance.
We further scale our methods to CelebA dataset to evaluate robust performance of proposed formulation and algorithms under label imbalance. We adopt the same neural network backbone as Multi-MNIST, detailed hyper-parameters are provided in Table~\ref{Table: Learning Rate Hyper-parameter CelebA}.

To address label imbalance, we employ an importance-sampling strategy when approximating the dual of DR-MOO objective~\eqref{eq: dual DR-MOO}. At test time, we report averaged accuracy, balanced accuracy and averaged AUC across tasks, with results summarized in Table~\ref{table:accuracy-celeba}. These results indicate that the proposed algorithms improve model robustness under task-wise label imbalance.


\section{Conclusion}
In this paper, we introduced DR-MOO as a framework for modeling objective-wise distribution shifts in multi-objective optimization. We characterized distributionally robust Pareto solution concepts and proposed a dual-based definition of robust Pareto stationarity that enables convergence analysis. Building on this foundation, we developed two MGDA-type algorithms: a double-loop method based on dual tracking and a Double-Clip method that provably improves sample complexity through large-batch sampling and coordinated clipping in the parameter and preference updates. Our results establish convergence guarantees under distributional uncertainty and highlight the importance of controlling relative update magnitudes in stochastic MOO. We hope this work motivates further theoretical investigation of gradient normalization in existing MTL methods and inspires extensions of optimizers such as Adam~\citep{kingma2014adam} and Muon~\citep{jordan2024muon} to multi-objective settings.



%\section*{Impact Statement}
%This paper aims to advance the field of multi-objective optimization. While the proposed methods may have a range of downstream societal impacts depending on the application domain, we do not identify any specific societal consequences that require additional discussion here.

% In the unusual situation where you want a paper to appear in the
% references without citing it in the main text, use \nocite
%\nocite{langley00}


\subsection*{AI Use Statement}
In this work, we used generative AI tools to assist with experimental design and to check the consistency and mathematical correctness of theoretical arguments. We did not use generative AI tools to conceive or develop the proposed algorithms, generate mathematical claims or proofs, implement the methods, generate synthetic data, process datasets, or interpret experimental results. Tasks involving surveys, interviews, transcription, and qualitative or thematic analysis were not applicable. Additionally, we used generative AI tools to improve the clarity, grammar, concision, and organization of author-written text.

All AI-assisted suggestions were critically reviewed by the authors. Theoretical feedback was independently checked against the stated assumptions and complete derivations, while experimental suggestions were evaluated, implemented, and validated by the authors. All algorithms, proofs, experiments, results, and conclusions were developed or verified by the authors. We take responsibility for the final content of this work, including all text, claims, and artifacts produced with the aid of generative AI.
%\subsection*{Ethics statement}

%(This section is \textbf{recommended} and does not count toward the page limit.)

%If authors feel that their paper submission raises questions regarding the Code
%of Ethics, they are encouraged to include a paragraph of Ethics Statement (at
%the end of the main text before references) to address potential concerns where
%appropriate. Topics include, but are not limited to, studies that involve human
%subjects, practices to data set releases, potentially harmful insights,
%methodologies and applications, potential conflicts of interest and sponsorship,
%discrimination/bias/fairness concerns, privacy and security issues, legal
%compliance, and research integrity issues (e.g., IRB, documentation, research
%ethics). This statement should not be more than 1 page.

%\subsection*{Reproducibility statement}

%(This section is \textbf{recommended} and does not count toward the page limit.)

%It is important that the work published in ICLR is reproducible. Authors are
%strongly encouraged to include a paragraph-long Reproducibility Statement at the
%end of the main text (before references) to discuss the efforts that have been
%made to ensure reproducibility. This paragraph should not itself describe
%details needed for reproducing the results, but rather reference the parts of
%the main paper, appendix, and supplemental materials that will help with
%reproducibility. For example, for novel models or algorithms, a link to an
%anonymous downloadable source code can be submitted as supplementary materials;
%for theoretical results, clear explanations of any assumptions and a complete
%proof of the claims can be included in the appendix; for any datasets used in
%the experiments, a complete description of the data processing steps can be
%provided in the supplementary materials. Each of the above are examples of
%things that can be referenced in the reproducibility statement.



%\subsubsection*{Acknowledgments}
%Use unnumbered third level headings for the acknowledgments. All
%acknowledgments, including those to funding agencies, go at the end of the paper.

\bibliographystyle{apalike}
\bibliography{main}
%\section*{Checklist}


% %%% BEGIN INSTRUCTIONS %%%
%The checklist follows the references. For each question, choose your answer from the three possible options: Yes, No, Not Applicable.  You are encouraged to include a justification to your answer, either by referencing the appropriate section of your paper or providing a brief inline description (1-2 sentences).
%Please do not modify the questions.  Note that the Checklist section does not count towards the page limit. Not including the checklist in the first submission won't result in desk rejection, although in such case we will ask you to upload it during the author response period and include it in camera ready (if accepted).

%\textbf{In your paper, please delete this instructions block and only keep the Checklist section heading above along with the questions/answers below.}
% %%% END INSTRUCTIONS %%%


%\begin{enumerate}

%  \item For all models and algorithms presented, check if you include:
%  \begin{enumerate}
%    \item A clear description of the mathematical setting, assumptions, algorithm, and/or model. [Yes]
%    \item An analysis of the properties and complexity (time, space, sample size) of any algorithm. [Yes]
%    \item (Optional) Anonymized source code, with specification of all dependencies, including external libraries. [No]
%  \end{enumerate}

%  \item For any theoretical claim, check if you include:
%  \begin{enumerate}
%    \item Statements of the full set of assumptions of all theoretical results. [Yes]
%    \item Complete proofs of all theoretical results. [Yes]
%    \item Clear explanations of any assumptions. [Yes]
%  \end{enumerate}

%  \item For all figures and tables that present empirical results, check if you include:
%  \begin{enumerate}
%    \item The code, data, and instructions needed to reproduce the main experimental results (either in the supplemental material or as a URL). [Yes]
%    \item All the training details (e.g., data splits, hyperparameters, how they were chosen). [Yes]
%    \item A clear definition of the specific measure or statistics and error bars (e.g., with respect to the random seed after running experiments multiple times). [Yes]
%    \item A description of the computing infrastructure used. (e.g., type of GPUs, internal cluster, or cloud provider). [Yes]
%  \end{enumerate}

%  \item If you are using existing assets (e.g., code, data, models) or curating/releasing new assets, check if you include:
%  \begin{enumerate}
%    \item Citations of the creator If your work uses existing assets. [Not Applicable]
%    \item The license information of the assets, if applicable. [Not Applicable]
%    \item New assets either in the supplemental material or as a URL, if applicable. [Not Applicable]
%    \item Information about consent from data providers/curators. [Not Applicable]
%    \item Discussion of sensible content if applicable, e.g., personally identifiable information or offensive content. [Not Applicable]
%  \end{enumerate}

%  \item If you used crowdsourcing or conducted research with human subjects, check if you include:
%  \begin{enumerate}
%    \item The full text of instructions given to participants and screenshots. [Not Applicable]
%    \item Descriptions of potential participant risks, with links to Institutional Review Board (IRB) approvals if applicable. [Not Applicable]
%    \item The estimated hourly wage paid to participants and the total amount spent on participant compensation. [Not Applicable]
%  \end{enumerate}

%\end{enumerate}
% make sure main.bib exists and is in the project
%\newpage
\clearpage
\appendix
\onecolumn

%################################################
\part*{Appendix}
\addcontentsline{toc}{section}{Appendix}

% Tag everything from here as appendix and configure tag depths
\etocdepthtag.toc{mtappendix}
\etocsettagdepth{mtchapter}{none}      % hide main paper sections
\etocsettagdepth{mtappendix}{subsection} % show appendix sections + subsections
%\subsection*{ Table of Contents}\vspace{0.5em}
\etocsettocstyle
  {}
  {\vspace{1em}}
\tableofcontents
\clearpage

%################################################
\allowdisplaybreaks
\section{Limitations and Future Work}
First, our high-probability guarantees for Double-Clip MGDA require batch sizes that depend strongly on the confidence parameter \(\delta\). This dependence arises from controlling stochastic-gradient errors and function-value gaps throughout the optimization trajectory and may partly reflect limitations of the current analytical framework~\eqref{eq: general analysis framework}. Future work could develop sharper concentration arguments through alternative Lyapunov analyses and explore algorithmic extensions, such as momentum or variance reduction, to reduce this dependence. Second, our methods seek Pareto-stationary solutions without explicitly enforcing user-specified trade-offs. An important extension is therefore to incorporate preference guidance, potentially through proximal or mirror-descent updates with preference-aware regularization.

\section{A Toy Example on Motivation of Study}\label{Appendix: Toy exmple visualization}
We visualize the effect of distribution shift on the objective geometry and Pareto frontier using a simple bi-objective toy problem. 
Given a parameter $\theta$, we evaluate two objectives
\[
f_1(\theta) = (\theta-x_1)^2+b_1,
\qquad
f_2(\theta) = (\theta-x_2)^2+b_2.
\]
We simulate the nominal and shifted data distributions by perturbing $x_1=0,x_2=2$ and $b_1,b_2=0$ with Gaussian noise. 
We approximate the Pareto frontier by retaining only the non-dominated sampled points. 
Figure~\ref{fig: Toy-example Pareto-visualization} visualizes the objective geometry (Left) and Pareto frontier (Right) before and after perturbation. 
As shown, the perturbation changes not only the objective geometry but also the induced Pareto frontier. 
In particular, we found the frontier is substantially distorted, and some parameter values that are Pareto-optimal under the nominal distribution become non-optimal after the perturbation.
\begin{figure}[ht]
    \centering
    \begin{subfigure}{0.23\linewidth}
        \includegraphics[width=\linewidth]{pareto_shift_demo_01_function_geometry.png}
    \end{subfigure}
    \hspace{0.02\linewidth}
    \begin{subfigure}{0.23\linewidth}
        \includegraphics[width=\linewidth]{pareto_shift_demo_02_pareto_comparison.png}  
    \end{subfigure}
    \caption{Effects of data distribution shift on function geometry (Left) and Pareto-frontier (Right).}
    \label{fig: Toy-example Pareto-visualization}
\end{figure}
\section{Summary of Comparison with Closely Related Works}\label{Appendix: comparison with other algorithms}
Table~\ref{Table: Comparison with existing methods} compares our algorithms' theoretical results with those of other related algorithms. All complexity results are compared in terms of convergence to $\epsilon$-Pareto stationary point.
\begin{table}[ht!]
\centering
\small
\caption{Comparison of convergence results on variants of gradient-balancing methods. (\textbf{Note: For the proposed algorithms, the smoothness and sample noise are defined with respect to the dual of DR-MOO formulation~\eqref{eq: dual DR-MOO}}. The corresponding properties are provided in Lemma~\ref{lemma: smooth of phi}, Corollary~\ref{Corollary: estimation error of true gradient}, Lemma~\ref{lemma: generalized-smooth}, and Lemma~\ref{lemma: affine bounded noise}, respectively.) Explanation on the upper footmarks: $1:$ Bounded Noise indicates the stochastic gradient variance is bounded by a constant. Affine Bounded Noise indicates stochastic gradient variance is affine bounded;
$\times$ indicates ``Does Not Apply" since NashMTL is analyzed under access to full-gradient; Bounded Output Noise indicates stochastic oracles output a bounded noisy gradient.
$2:$ (semi-LS) indicates function geometry is $L$-smooth along certain trajectory;
(LS) indicates $L$-smooth assumption on objectives while (GS) indicates objectives satisfies $(L_0,L_1)$ generalized-smooth~\citep{generalizedsmooth,Haochuangensmooth}; $3:$ (BF) denotes the bounded function value assumption and (BG) denotes the bounded gradient assumption; $5:$ Sample complexity is measured in terms of achieving $\epsilon$-accurate Pareto stationary point under stochastic settings. N/A denotes no additional assumptions.
}
\resizebox{\columnwidth}{!}{
\begin{tabular}{lccccc}
\toprule
Methods & Sample Noise\textsuperscript{1} & Smoothness\textsuperscript{2} & Other
Assumption\textsuperscript{3} & Batch Size\textsuperscript{4} & Sample Complexity\textsuperscript{5}  \\
\midrule
Double-Loop MGDA & Bounded Noise & semi-LS & N/A & $O(1)$ & $O(\epsilon^{-8})$ \\
Double-Clip MGDA & Affine Bounded Noise& GS & N/A & ${O}(\epsilon^{-2})$ & ${O}(\epsilon^{-4})$ \\
\midrule
MoCo~\citep{MOCO} & Bounded Noise&LS & BG and BF & ${O}(1)$ & $O(\epsilon^{-4})$ \\
NashMTL~\citep{nashMTL} & $\times$ & LS & N/A & N/A & asymptotic convergence\\
%FAMO~\citep{FAMO}  \\
SDMGrad~\citep{SDMkaiyiJi} & Bounded Noise & LS & BG &$O(1)$& $O(\epsilon^{-4})$.  \\
MoDo~\citep{MODO} & Bounded Output Noise & LS & BG & $O(1)$&  $O(\epsilon^{-4})$  \\
MGDA~\citep{qimgda} & Bounded Noise & GS & N/A & ${O}(1)$ & ${O}(\epsilon^{-4})$ \\
\bottomrule
\end{tabular}}
\label{Table: Comparison with existing methods}
\end{table}

\textbf{Robust multi-task learning.}
A line of work in robust MTL addresses noisy or outlier tasks through statistical modeling, structured parameterization, or decomposition. Early approaches~\citep{robustmtl1} use Gaussian-process variants to limit the influence of outlier tasks, while \citet{robustmtl2} decompose the task-parameter matrix into shared and task-specific components to identify such tasks explicitly. More recently, \citet{robustMTL3} develop adaptive frameworks with statistical guarantees that exploit task similarity while allowing task-specific deviations. Other studies consider robustness to label noise~\citep{robustMTL4} and adversarial perturbations~\citep{robustMTL5}. Despite these advances, the machine-learning literature on robust MTL has largely focused on robustness through task modeling, noise mitigation, or adversarial training. Task losses are typically aggregated or adaptively reweighted into a scalar objective, rather than modeled using objective-wise ambiguity sets over data distributions. Consequently, this literature does not formulate robust MTL as a vector-valued DRO problem or study distributionally robust Pareto solutions. In contrast, DR-MOO preserves the vector-valued robust objective and establishes convergence to distributionally robust Pareto-stationary solutions.

\textbf{On first-order algorithmic solutions for solving DRO}
Several works study single-objective DRO formulations amenable to first-order optimization. For the $f$-divergence--regularized risk $\min_{\theta}\phi(\theta)$, where $\phi(\theta):=\min_{\eta\in\mathbb{R}}\{\lambda\mathbb{E}_{\xi\sim\mathbb{P}}[f^*((\ell(\theta;\xi)-\eta)/\lambda)]+\eta\}$, \citet{jinnonconvexdro} develop normalized SGD with momentum and establish an $\mathcal{O}(\epsilon^{-4})$ gradient complexity. Under assumptions comparable to ours, \citet{qirevistfDRO} propose D-SGD-C and its variance-reduced counterpart D-Spider-C, achieving gradient complexities of $\mathcal{O}(\epsilon^{-4})$ and $\mathcal{O}(\epsilon^{-3})$, respectively. These results demonstrate that normalization and clipping effectively exploit the generalized-smooth geometry and affine stochastic noise of dual DRO objectives. However, they address only single-objective parameter updates and do not account for the gradient balancing and stochastic preference updates required in MOO. This motivates us to adopt clipping as a building block and extend it jointly to the optimization and preference updates in DR-MOO. Earlier, \citet{large-scaleDRO} developed stochastic methods for convex CVaR and $\chi^2$-DRO using biased batch-gradient and multilevel Monte Carlo estimators under stronger boundedness and regularity assumptions.

Beyond regularized $f$-divergence DRO, \citet{qiaaai} study constrained Cressie--Read DRO by combining stochastic-gradient and Frank--Wolfe updates, while \citet{qiKLDRO} develop constant-batch stochastic algorithms for constrained KL-DRO with complexity independent of the dataset size. For Sinkhorn DRO, \citet{sinkhornDROwang} employ stochastic mirror descent with biased gradient estimators and obtain near-optimal sample complexity under convexity and boundedness assumptions. Other work studies empirical Group-DRO: \citet{droacceleration} develop ball-oracle acceleration with variance-reduced inner solvers, whereas \citet{yu2024efficient} propose the ALEG stochastic primal--dual method by exploiting its two-level finite-sum structure. These methods optimize scalar robust risks under problem-specific assumptions but do not address vector-valued robust objectives or adaptive preference balancing in DR-MOO.

\section{Hyper-parameter Configuration of Main Experiments}\label{Appendix: Hyper-parameter of main experiments.}
We conduct our deep learning experiments using two primary configurations: an Intel-based workstation featuring a single NVIDIA RTX 4090 GPU and 32 GB RAM, and a server equipped with four NVIDIA RTX A6000 GPUs.

\subsection{Neural Network Backbone, Common hyper-parameters and Implementation Details}
Table~\ref{Table: backbone-training-config} summarizes the network backbone and common training settings for multi-task classification on Multi-MNIST, CelebA, and Office-31. For the ResNet18 encoder, we freeze the initial convolutional and batch normalization layers, together with the first residual stage, throughout training.

For CelebA, we select eight attributes: \textbf{Bald, Wearing Hat, Eyeglasses, Receding Hairline, Narrow Eyes, Blond Hair, Bags Under Eyes, and Big Nose}. For each attribute, we estimate the empirical label proportions within each batch, denoted by $\tilde{p}$ and $1-\tilde{p}$. When evaluating the dual DR-MOO objective~\eqref{eq: dual DR-MOO}, we assign samples from the two classes weights of $\frac{1}{2\tilde{p}}$ and $\frac{1}{2(1-\tilde{p})}$, respectively. This batchwise class reweighting approximates a nominal distribution $\mathbb{P}$ with balanced class probabilities, i.e., $p=\frac{1}{2}$.

In our implementation of Double-Loop MGDA, we sample three iterates, $\eta_d$, $\eta_{\bar d}$, and $\eta_{\tilde d}$, from a single inner-loop SGD trajectory and use them to construct the stochastic gradient estimators $Y_t$, $\bar Y_t$, and $\tilde Y_t$ for the model and preference updates. Reusing the trajectory simplifies implementation and avoids additional inner-solver executions.



\begin{table}[ht]
\centering
%\small
\caption{Neural network backbone and Common hyper-parameter Settings}
\label{Table: backbone-training-config}
{
\scriptsize
\renewcommand{\arraystretch}{1.08}
\setlength{\tabcolsep}{1.5pt}
\begin{tabular*}{\textwidth}{
    @{\extracolsep{\fill}}
    ll
    @{}
}
\toprule
\textbf{Component} & \textbf{Configuration} \\
\midrule
Shared encoder & ResNet18 pre-trained on ImageNet~\citep{imageNet} \\
Classification heads & MLP:  Two-layer linear head with intermediate
ReLU activation and dropout \\
Objective & Dual of DR-MOO~\eqref{eq: dual DR-MOO} for all methods  with $\ell(\cdot)$ to be cross-entropy loss\\
$f^*(\cdot)$ & Convex conjugate dual of $\chi^2$ divergence, i.e., $f^*(t)=\frac{1}{4}(t+2)_{+}^2-1$ \\
Head initialization & Kaiming uniform initialization for each layer \\
Optimizer for heads & SGD \\
Batch size & $256$ \\
Learning rate for heads & $1e^{-2}$ for all methods \\
$\rho$ & $1e^{-5}$ for multi-mnist; $5e^{-5}$ for CelebA\\
$\lambda$ in DR-MOO~\eqref{eq: dual DR-MOO} & $0.6$ for multi-mnist and CelebA\\
Inner iterations & $5$ for Double-Loop MGDA and SDMGrad; \\
\midrule
2-digit Multi-MNIST input size & $128 \times 128$ \\
2-digit Multi-MNIST training epochs & $70$ for all methods \\
3-digit Multi-MNIST input size & $96 \times 96$ \\
3-digit Multi-MNIST training epochs & $100$ for all methods \\
CelebA input size & $128\times128$ \\
CelebA training epochs & $100$ for all methods.\\
\bottomrule
\end{tabular*}}
\end{table}

\begin{table}[ht]
\centering
\caption{Hyperparameter settings for Multi-MNIST experiments.}
\label{Table: Learning Rate Hyper-parameter}
{
\scriptsize
\renewcommand{\arraystretch}{1.08}
\setlength{\tabcolsep}{1.5pt}
\begin{tabular*}{\textwidth}{
    @{\extracolsep{\fill}}
    lcccccccc
    @{}
}
\toprule
\textbf{Method} 
& \textbf{LR for $(\theta,\eta)$} 
& \textbf{LR for preference ($\beta$)} 
& $\boldsymbol{f_1,f_2}$ 
& $\boldsymbol{c_1,c_2}$ \\
\midrule
Double-Loop MGDA 
& $\alpha=5e^{-4},\gamma=3e^{-3}$ 
& $1e^{-5}$ 

& -- 
& -- \\

 Double-Clip MGDA 
& $\gamma=5e^{-3}$
& $1e^{-2}$ 
& $1.0$,\,$0.5$ 
& $1.0,\,0.5$ \\

MGDA~\citep{qimgda,MGDA1stwork} 
& $1e^{-4}$ 
& $1e^{-6}$ 
& -- 
& -- \\

MoDo~\citep{MODO} 
& $1e^{-4}$ 
& $1e^{-6}$ 
& -- 

& -- \\

MoCo~\citep{MOCO}
& $5e^{-4}$ 
& $1e^{-4}$ 
& -- 
& -- \\

SDMGrad~\citep{SDMkaiyiJi}
& $1e^{-4}$ 
& $1e^{-6}$ 
& -- 
& -- \\

FAMO~\citep{FAMO}
& $1e^{-4}$ 
& $2.5e^{-3}$ 
& -- 
& -- \\

NashMTL~\citep{nashMTL}
& $1e^{-4}$ 
& -- 
& -- 
& -- \\
\bottomrule
\end{tabular*}}
\end{table}

\begin{table}[ht]
\centering
\small
\caption{Hyperparameter settings for CelebA Experiments.}
\label{Table: Learning Rate Hyper-parameter CelebA}
{
\scriptsize
\renewcommand{\arraystretch}{1.08}
\setlength{\tabcolsep}{1.5pt}
\begin{tabular*}{\textwidth}{
    @{\extracolsep{\fill}}
    lcccccccc
    @{}
}
\toprule
\textbf{Method} 
& \textbf{LR for $(\theta,\eta)$} 
& \textbf{LR for preference ($\beta$)} 
& $\boldsymbol{f_1,f_2}$ 
& $\boldsymbol{c_1,c_2}$ \\
\midrule
Double-Loop MGDA 
& $\alpha=5e^{-4},\gamma=3e^{-4}$ 
& $5e^{-5}$ 

& -- 
& -- \\

Double-Clip MGDA 
& $\gamma=5e^{-3}$
& $2e^{-2}$ 
& $1.0$,\,$0.5$ 
& $1.0,\,0.5$ \\

MGDA~\citep{qimgda,MGDA1stwork} 
& $1e^{-4}$ 
& $5e^{-5}$ 
& -- 
& -- \\

MoDo~\citep{MODO} 
& $1e^{-4}$ 
& $5e^{-5}$ 
& -- 

& -- \\

MoCo~\citep{MOCO}
& $5e^{-4}$ 
& $5e^{-5}$ 
& -- 
& -- \\

SDMGrad~\citep{SDMkaiyiJi}
& $1e^{-4}$ 
& $1e^{-6}$ 
& -- 
& -- \\

FAMO~\citep{FAMO}
& $1e^{-4}$ 
& $5e^{-6}$ 
& -- 
& -- \\

NashMTL~\citep{nashMTL}
& $1e^{-4}$ 
& -- 
& -- 
& -- \\
\bottomrule
\end{tabular*}}
\end{table}

\begin{table}[t]
\centering
\caption{Hyperparameter settings for the Office-31 experiments.}
\label{tab:office31_hyperparameters}
{
\scriptsize
\renewcommand{\arraystretch}{1.08}
\setlength{\tabcolsep}{1.5pt}
\begin{tabular*}{\textwidth}{
    @{\extracolsep{\fill}}
    lcc
    @{}
}
\toprule
Method & LR for $(\theta,\eta)$ & LR for preference $(w)$ \\
\midrule
Double-Loop MGDA
& $\alpha=5e^{-4},\ \gamma=1e^{-3}$
& $5e^{-4}$ \\

Double-Clip MGDA
& $\gamma=1e^{-3}$
& $2e^{-3}$ \\

MGDA~\citep{MGDA1stwork}
& $1e^{-4}$
& $5e^{-5}$ \\

MoDo~\citep{MODO}
& $1e^{-4}$
& $5e^{-5}$ \\

MoCo~\citep{MOCO}
& $5e^{-4}$
& $5e^{-5}$ \\

SDMGrad~\citep{SDMkaiyiJi}
& $1e^{-4}$
& $5e^{-6}$ \\

FAMO~\citep{FAMO}
& $5e^{-5}$
& $1e^{-2}$ \\

NashMTL~\citep{nashMTL}
& $1e^{-4}$
& -- \\
\bottomrule
\end{tabular*}}
\end{table}

\subsection{Learning Rate Choices}
Table~\ref{Table: Learning Rate Hyper-parameter} and Table~\ref{Table: Learning Rate Hyper-parameter CelebA} summarize the learning rate choices for all methods, which are used when training the network backbone on the Multi-MNIST and CelebA datasets. For method-specific hyperparameters, we set $\gamma = 1e^{-2}$ for FAMO~\citep{FAMO}, and the exponential moving average parameter $\text{ema}=0.95$ for MoCo~\citep{MOCO}, we leverage mirror-descent solver to solve nonlinear system introduced in NashMTL~\citep{nashMTL} and run $100$ iterations each time step.


\section{Additional Experiments}\label{Appendix: Synthetic Experiments}
\subsection{Multi-Task Linear Regression over Synthetic Data}
We evaluate Algorithm~\ref{alg1} and \ref{alg1-3} on the dual of DR-MOO problem~\eqref{eq: dual DR-MOO} with the mean-squared error loss. We consider a synthetic linear regression setup with $m=3$ objectives, data dimension $n=10$, and a total of $6000$ samples. Inputs are drawn i.i.d. from $X \sim \mathcal{N}(0, I)$.
To induce partial conflicts across objectives, we generate ground-truth parameters as $
\theta^{1,*} \sim \mathcal{N}(0, I), \quad
\theta^{2,*} \sim \mathcal{N}(-0.2\theta_1^*, 0.04 I), \quad
\theta^{3,*} \sim \mathcal{N}( 0.5\theta_1^*, 0.25 I)
$. We then generate labels according to $y^{i}=X\theta^{*,i}+\varepsilon^i$, where the noise terms have different variances $\varepsilon^1 \sim \mathcal{N}(0, 0.04)$, $\varepsilon^2 \sim \mathcal{N}(0, 0.36)$, $
\varepsilon^3 \sim \mathcal{N}(0, 0.25)$. 

We compare all methods for $T={600}$ iterations with tuned hyperparameters. For all methods, we use a batch size of $256$ to ensure numerical stability.
For the double-loop MGDA (Algorithm~\ref{alg1}), we set the inner-loop length $D=20$, learning rates $\gamma=5e^{-3}$, $\alpha=\beta=5e^{-5}$, and regularization parameter $\rho=1e^{-5}$.
For the single-loop double-clip MGDA (Algorithm~\ref{alg1-3}), we set learning rate $\gamma=3e^{-2}$, $\beta = 1e^{-2}$, clipping thresholds $c_1=f_1=0.5$, $c_2=f_2=0.1$ and $\rho=1e^{-5}$. For stochastic MGDA, we set learning rate for both $(\theta_t,\eta_t)$ and $w_t$ to $1e^{-5}$, with $\rho=0$. For MoDo and MoCo, we adopt the same settings as stochastic MGDA, but set $\rho=1e^{-5}$. For SDMGrad, we set $D=10$, learning rates $1e^{-4}$ for $(\theta_t,\eta_t)$ and $5e^{-4}$ for $w_t$, and $\rho=1e^{-5}$.
Figure~\ref{fig: synthetic experiments} (Left) reports the balanced stochastic gradient norm versus iteration, together with total sample consumption. 
Both our proposed methods achieve competitive performance relative to all baselines, highlighting the benefit of explicitly exploiting the structure of DR-MOO. While SDMGrad attains a similar iteration-wise convergence trend to Double-Loop MGDA, it requires substantially more samples and relies on additional heuristic normalization for numerical stability.
%As shown, our proposed methods (Double-Loop and Double-Clip MGDA) achieve competitive performance compared with all baseline approaches, validating the necessity of exploiting the structural properties of DR-MOO. While SDMGrad exhibits comparable convergence as Double-Loop MGDA in terms of iterations, it incurs substantially higher sample complexity and requires additional heuristic normalization to maintain numerical stability.
\begin{figure}[ht]
    \centering
    \begin{subfigure}{0.23\textwidth}
        \includegraphics[width=\linewidth]{Exp_Fig/gradstep.png}
    \end{subfigure}
    \hspace{0.01\linewidth}
    \begin{subfigure}{0.23\textwidth}
        \includegraphics[width=\linewidth]{Exp_Fig/gradient.png}
    \end{subfigure}  
    \hspace{0.01\linewidth}
      \begin{subfigure}{0.23\textwidth}
        \includegraphics[width=\linewidth]{Exp_Fig/gradstep_logistic.png}
    \end{subfigure}
    \hspace{0.01\linewidth}
    \begin{subfigure}{0.23\textwidth}
        \includegraphics[width=\linewidth]{Exp_Fig/gradient_logistic.png}
    \end{subfigure}
    \caption{Balanced Gradient vs.Time Steps and Sample Consumption for Linear (Left) and Logistic Regression (Right)}
    \label{fig: synthetic experiments}
\end{figure}
\begin{table}[t]
    \centering
    
    \setlength{\tabcolsep}{5pt}
    \caption{Balanced-gradient norm over 100 random seeds. For each seed, the
    metric is averaged over the final 20 iterations. Lower values are better.
    The intervals are two-sided 90\% Student-$t$ confidence intervals.}
    \label{tab:regression-gradient-results}
    \resizebox{0.75\columnwidth}{!}{
    \begin{tabular}{llcc}
        \toprule
        \textbf{Experiment} & \textbf{Method}
        & \textbf{Mean final gradient} & \textbf{90\% CI} \\
        \midrule
        Synthetic linear regression
        & Double-Loop MGDA
        & 2.530248
        & $[2.514342,\;2.546153]$ \\

        Synthetic linear regression
        & Double-Clip MGDA
        & \textbf{1.224745}
        & $[1.205730,\;1.243761]$ \\

        UCI white-wine logistic regression
        & Double-Loop MGDA
        & 0.239733
        & $[0.237789,\;0.241677]$ \\

        UCI white-wine logistic regression
        & Double-Clip MGDA
        & \textbf{0.175293}
        & $[0.174285,\;0.176301]$ \\
        \bottomrule
    \end{tabular}}
\end{table}
\subsection{Multi-Task Logistic Regression over UCI White Wine Dataset}
We evaluate Algorithm~\ref{alg1} and \ref{alg1-3} on the dual of DR-MOO problem in~\eqref{eq: dual DR-MOO} with the cross-entropy loss. Specifically, we train $m=3$ logistic regression models on the UCI white wine dataset~\citep{uci-dataset}, corresponding to binary classification tasks based on wine quality, alcohol content, and residual sugars. To induce class imbalance, we generate binary labels using task-specific quantile thresholds $s^i$, i.e., $y^i=\mathbf{1}(y\geq s^i),\forall i\in \{1,2,3\}$. We set thresholds for quality, residual sugars and alcohol as $0.5$, $0.8$, $0.1$ respectively. 

We compare all methods for $T={1000}$ iterations. For all methods, we use a batch size of $256$. For the double-loop MGDA (Algorithm~\ref{alg1}), we set the inner-loop length $D=15$, learning rates $\gamma=5e^{-3}$, $\alpha=1.5e^{-3}$, $\beta = 6e^{-4}$, and $\rho=1e^{-6}$. For the single-loop double-clip MGDA (Algorithm~\ref{alg1-3}), we set learning rate $\gamma = 1e^{-2}$, clipping thresholds $c_1=f_1=0.5$, $c_2=f_2=0.1$ and $\rho=1e^{-5}$. For MGDA, MoDo and MoCo, we set learning rates for parameter $(\theta,\eta)$ as $1e^{-3}$, and learning rate for preference vector $w_t$ as $6e^{-4}$. For regularization parameter, we set $\rho=1e^{-6}$ for MoCo and MoDo, $\rho=0$ for MGDA, respectively. For SDMGrad, we set $D=15$, learning rates for parameter $(\theta,\eta)$ as $1e^{-3}$, for preference vector $w_t$ as $5e^{-4}$, and regularization parameter $\rho=1e^{-4}$.  Figure~\ref{fig: synthetic experiments} reports the comparison results, from which one can see that our proposed methods
achieve competitive performance relative to the baselines.
\subsection{Multi-Seed Experiments over Linear and Logistic Regression}
We repeat the linear and logistic regression experiments over 100 random seeds and report the balanced-gradient norm averaged over the final 20 iterations. Table~\ref{tab:regression-gradient-results} presents the mean and 90\% confidence interval across seeds. Double-Clip MGDA consistently achieves a lower final gradient norm than Double-Loop MGDA on both tasks, while the narrow confidence intervals indicate stable performance across different random initializations and stochastic samples.
\subsection{Multi-Domain Classification on Office-31}\label{Appendix: office-31 dataset experiment}
We further evaluate our methods on Office-31~\citep{officedata31}, which contains $31$ object classes from the Amazon, DSLR, and Webcam domains.  We treat each domain as a separate task and objective, reuse the neural network backbone from the preceding experiments and common hyper-parameters. Table~\ref{tab:office31_hyperparameters} presents fine-tuned learning rates of each method. To avoid repeatedly sampling the smaller domains solely to match the size of Amazon, we independently shuffle each domain and partition every domain loader into $31$ minibatches per pass. This results in batch sizes of $63$--$64$ for Amazon, $10$--$11$ for DSLR, and $17$--$18$ for Webcam. Thus, each training image is used once during a complete pass through its domain loader. At each training step, we concatenate one mini-batch from each domain for the forward pass and mask labels that do not belong to the domain whose objective is being computed.

For each domain, we optimize the same dual DR-MOO objective~\eqref{eq: dual DR-MOO}. We train each method for $70$ epochs. Table~\ref{table:accuracy-office31} reports classification accuracy on test set. Under this setting, Double-Clip MGDA achieves the highest observed worst-domain and mean-domain test accuracy. Double-Loop MGDA ranks behind MoCo but outperforms the remaining baselines on the worst-domain metric. These results demonstrate that the proposed algorithms are competitive when applied to the common DR-MOO objective on this multi-domain classification task.
\begin{table}[ht]
\centering
\small
\caption{Office-31 test accuracy (\%). ``Worst'' is the minimum accuracy across the three domains, and ``Mean'' weights the domains equally.}
\label{table:accuracy-office31}

\scriptsize
\renewcommand{\arraystretch}{1.08}
\setlength{\tabcolsep}{1.5pt}
{
\begin{tabular*}{\textwidth}{
    @{\extracolsep{\fill}}
    lccccc
    @{}
}
\toprule
\textbf{Method} & \textbf{Amazon} & \textbf{DSLR} & \textbf{Webcam} & \textbf{Worst} & \textbf{Mean} \\
\midrule
Double-Clip MGDA & $\boldsymbol{79.06}$ & $\boldsymbol{98.06}$ & $\boldsymbol{98.61}$ & $\boldsymbol{79.06}$ & $\boldsymbol{91.91}$ \\
Double-Loop MGDA & 78.62 & $\boldsymbol{98.06}$ & 97.22 & 78.62 & 91.30 \\
MoCo~\citep{MOCO} & 78.84 & $\boldsymbol{98.06}$ & 97.92 & 78.84 & 91.61 \\
SDMGrad~\citep{SDMkaiyiJi} & 78.17 & 97.09 & 97.92 & 78.17 & 91.06 \\
MoDo~\citep{MODO} & 78.17 & 97.09 & 97.92 & 78.17 & 91.06 \\
MGDA~\citep{MGDA1stwork,qimgda} & 78.17 & 97.09 & 97.92 & 78.17 & 91.06 \\
NashMTL~\citep{nashMTL} & 77.95 & 97.09 & 97.92 & 77.95 & 90.99 \\
FAMO~\citep{FAMO} & 75.72 & 96.12 & 97.22 & 75.72 & 89.69 \\
\bottomrule
\end{tabular*}}
\end{table}

\section{Ablation Studies}\label{Appendix: Ablation study}
%To provide a comprehensive understanding of the hyperparameters and practical performance of the proposed dual of DR-MOO~\eqref{eq: dual DR-MOO}, Double-Loop MGDA and Double-Clip MGDA, we conduct ablation studies on key components that may affect performance. These include determining feasible regularization hyper-parameter, learning rates, batch size selection for Double-Clip MGDA, the number of inner iterations for Double-Loop MGDA, and a comparison of iteration-wise wall-clock time against other baselines in deep learning experiments.
We examine sensitivity to the regularization parameter $\lambda$, learning rates, batch size, and the number of inner iterations, and compare the per-iteration computational cost of our methods with the baselines.

\subsection{Ablation on $\lambda$}
According to classical results on regularization hyperparameters~\citep{PRML}, from the perspective of the primal DR-MOO formulation~\eqref{eq: primal DR-MOO}, a smaller $\lambda$ permits larger deviations in the shifted distribution $\mathbb{Q}$, while a larger $\lambda$ enforces tighter control. However, our theoretical analysis reveals that several problem-dependent quantities scale with $\lambda^{-1}$. In particular, the smoothness parameters satisfy $\hL_0, \hL_2 = \mathcal{O}(\lambda^{-1})$, and the stochastic gradient variance satisfies $\hat{K}_0,\hat{K}_2 = \mathcal{O}(\lambda^{-2})$, both of which critically influence convergence behavior (See Appendix~\ref{Appendix: Relevant Properties of Phi}, ~\ref{sec: property of hat L} for more details).

To study the sensitivity to $\lambda$, we conduct an ablation experiment. We fix the optimizer to the proposed single-loop Double-Clip MGDA, adopt the same hyperparameters as in the Multi-MNIST experiments (see Table~\ref{Table: Multi-mnist FGSM}), and vary $\lambda \in \{0.4, 0.8, 1.6\}$. At test time, we evaluate robustness using FGSM attacks, and report classification accuracy under different attack levels in Table~\ref{Table: Lambda Ablation}. The results show that $\lambda=0.4$ yields slightly worse robustness compared to $\lambda=0.8$ and $1.6$, suggesting that small $\lambda$ may degrade optimization performance.
\begin{table}[t]
\caption{Ablation Comparison (Accuracy in \%) on $\lambda$}
\label{Table: Lambda Ablation}
\centering
{
\scriptsize
\renewcommand{\arraystretch}{1.08}
\setlength{\tabcolsep}{1.5pt}
\begin{tabular*}{\textwidth}{
    @{\extracolsep{\fill}}
    lccccc
    @{}
}
\toprule
Method/Attack Level & Clean & $0.010$ & $0.030$ & $0.050$ & $0.080$ \\
\midrule
$\lambda=0.4$ & 95.54 & 80.16 & 63.13 & 58.94 & 55.78 \\
$\lambda=0.8$ & 95.67 & 83.50 & 65.97 & 60.39 & 57.12 \\
%$\lambda=1.0$ & 0.9599 & 0.8356 & $\boldsymbol{0.6617}$ & $\boldsymbol{0.6133}$ & $\boldsymbol{0.5797}$ & $\boldsymbol{0.6726}$ \\
$\lambda=1.6$ & ${96.01}$ & ${83.57}$ & 65.96 & 60.77 & 57.38 \\
\bottomrule
\end{tabular*}}
\end{table}

\subsection{Ablation on Learning Rates}
During our experiments, we observed that the balanced gradient norm plays a critical role in training stability. Its magnitude also induces different feasible regions for effective learning rates between unclipped methods and Double-Clip MGDA. To select feasible learning rates and ensure convergence across control groups, we fix the batch size to $B=256$, $\rho = 1e^{-5}$, and $\gamma = 1e^{-3}$ for both linear and logistic regression. We then vary the learning rates for updating the model parameters $\theta$ and the preference vector $w$ over predefined candidate sets.

For linear regression, we perform a greedy search over the set $\{1e^{-5}, 5e^{-5}, 5e^{-4}, 1e^{-3}, 5e^{-3}, 1e^{-2}\}$ for all methods except Double-Clip MGDA. For Double-Clip MGDA, we restrict the search to $\{1e^{-3}, 5e^{-3}, 1e^{-2}, 5e^{-2}, 1e^{-1}\}$. For logistic regression, we use the candidate set $\{1e^{-5}, 5e^{-5}, 5e^{-4}, 1e^{-3}, 5e^{-3}, 1e^{-2}\}$ for all methods except Double-Clip MGDA, and $\{5e^{-3}, 1e^{-3}, 1e^{-2}, 5e^{-2}, 1e^{-1}\}$ for Double-Clip MGDA.
\begin{figure}[ht]
    \centering
    \begin{subfigure}{0.23\linewidth}
        \includegraphics[width=\linewidth]{Exp_Fig/regression_learning_rate_all_methods.png}
    \end{subfigure}
    \hspace{0.02\linewidth}
    \begin{subfigure}{0.23\linewidth}
        \includegraphics[width=\linewidth]{Exp_Fig/logistic_learning_rate_all_methods.png}  
    \end{subfigure}
    \caption{Ablation on effective learning rates for Linear(Left) and Logistic Regressions(Right)}
    \label{fig: ablation on LR}
\end{figure}

Figure~\ref{fig: ablation on LR} illustrates the relationship between the learning rate and the averaged balanced gradient norm over the last $20$ iterations. As shown, unclipped methods exhibit a different feasible region compared with Double-Clip MGDA. To achieve comparable convergence error, Double-Clip MGDA (green line) requires a larger learning rate, which is consistent with Theorem~\ref{thm: convergence of alg3} statements on choices of $\beta, \gamma$. 
%We emphasize that this experiment is intended to visualize the feasible learning-rate regions across methods. The hyperparameters used here differ slightly from those in the main experiments, where all parameters are further tuned to balance convergence speed and stability.


\subsection{Ablation on Batch Size }\label{Appendix: ablation on batch size and inner-steps}
\begin{figure}[ht]
    \centering
    \begin{subfigure}{0.23\linewidth}
        \includegraphics[width=\linewidth]{Exp_Fig/regression_single_batch_size_ablation.png}
        \caption{Linear}
    \end{subfigure}
    \hspace{0.02\linewidth}
    \begin{subfigure}{0.23\linewidth}
        \includegraphics[width=\linewidth]{Exp_Fig/logistic_single_batch_size_ablation.png}
        \caption{Logistic}
    \end{subfigure} 
        \begin{subfigure}{0.23\linewidth}
        \includegraphics[width=\linewidth]{Exp_Fig/regression_double_loop_inner_steps_ablation.png}
        \caption{Linear}
    \end{subfigure}
    \hspace{0.02\linewidth}
    \begin{subfigure}{0.23\linewidth}
        \includegraphics[width=\linewidth]{Exp_Fig/logistic_double_inner_steps_ablation.png}
        \caption{Logistic}
    \end{subfigure}
    \caption{Ablation on Effects of Batch Size (Left) and Inner-Step (Right)}
    
    \label{fig: batch-size and inner-step ablation}
\end{figure}


As suggested by Theorem~\ref{thm: convergence of alg3}, Double-Clip MGDA requires a batch size $N_1, N_2=\mathcal{O}(\epsilon^{-2})$ to guarantee convergence. In this section, we study the practical choice of batch size that balances computational efficiency and optimization stability. Specifically, for linear and logistic regression, we unify the learning rates across all control groups as $\beta, \gamma = 1e^{-2}$, and set $c_1=f_1=1.0$, $c_2=f_2=0.5$, and $\rho = 1e^{-5}$. We then vary the batch sizes $N_1, N_2$ over $\{32, 64, 128, 256, 512, 1024\}$.

The convergence behavior, measured by the balanced gradient, is shown in Figure~\ref{fig: batch-size and inner-step ablation} (a) and (b). As observed, when the batch size is small (e.g., $32$ or $64$), the gradients exhibit significant fluctuations during training, indicating instability. As the batch size increases beyond $256$, this instability becomes negligible. These results justify our practical choice of batch size $256$ as a sufficient trade-off between stability and efficiency.

\subsection{Effect of the Number of Inner Iterations}\label{Appendix: inner-iteration and wall-clock}
According to Theorem~\ref{thm: convergence of algorithm 1}, the primary convergence bottleneck of Double-Loop MGDA arises from the high complexity of the inner loop, which requires $D=\mathcal{O}(\epsilon^{-4})$ inner iterations per step to guarantee convergence. In this section, we study the practical choice of the number of inner iterations, aiming to balance overall computational cost and convergence performance.

Specifically, to ensure convergence across all control groups, for linear regression we unify the hyperparameters as $\alpha=\beta = 5e^{-5}$, $\gamma = 1e^{-3}$, and $\rho = 1e^{-5}$. For logistic regression, we set $\alpha, \gamma = 1e^{-3}$, $\beta = 6e^{-4}$, and $\rho = 1e^{-6}$ across all groups. We then vary the number of inner iterations over $\{5, 10, 15, 20, 25\}$. The convergence behavior, measured by the balanced gradient, is shown in Figure~\ref{fig: batch-size and inner-step ablation} (c) and (d). As observed, for linear regression, the convergence behavior is relatively insensitive to the number of inner steps under these hyperparameter settings. For logistic regression, although using $25$ inner steps leads to faster convergence, the final convergence results across all control groups show no significant differences. This suggests that, since the inner subproblem is convex, it is comparatively easier to solve than the outer loop responsible for updating model parameters and the preference vector.

In deep learning experiments, to reduce computational overhead, we further set the number of inner iterations to $5$ for Double-Loop MGDA. We compare its iteration-wise wall-clock time with Double-Clip MGDA and other baselines, with results shown in Figure~\ref{fig: Wall-clock}. Combined with the accuracy results reported in Table~\ref{Table: Multi-mnist FGSM} and Table~\ref{table:accuracy-celeba}, these findings indicate that setting the number of inner iterations to $5$ is sufficient to achieve performance comparable to other baselines. Moreover, with only $5$ inner iterations, the additional computational overhead becomes negligible relative to the cost of the network backbone and dataset size.


%\begin{figure}[ht]
%    \centering
%    \begin{subfigure}{0.3\linewidth}
        %\includegraphics[width=\linewidth]{Exp_Fig/regression_double_loop_inner_steps_ablation.png}
    %\end{subfigure}
    %\hspace{0.02\linewidth}
    %\begin{subfigure}{0.3\linewidth}
        %\includegraphics[width=\linewidth]{Exp_Fig/logistic_double_inner_steps_ablation.png}
    %\end{subfigure}  
    %\caption{Ablations on Effects of Inner-Loop Iteration for Linear(Left) and Logistic(Right) Regression}
    %\label{fig: Inner-Loop Ablation}
%\end{figure}

\begin{figure}[ht]
    \centering
    %\begin{subfigure}{0.23\textwidth}
    %\includegraphics[width=1.05\linewidth]{Exp_Fig/gradstep.png}
    %\end{subfigure}
        %\hfill
    %\begin{subfigure}{0.23\textwidth}
    \includegraphics[width=0.6\linewidth]{Exp_Fig/step_time_comparison.png}
    %\end{subfigure}  
    \caption{Iteration-wise Wall-clock Time Comparison for Each Method.}
    \label{fig: Wall-clock}
\end{figure}



%\section{Notations for Convergence Analysis}

%Throughout, when analyzing Double-loop MGDA algorithm \ref{alg1} in section \ref{sec: Double-loop algorithm}, we denote the gradient of vectorized multi-objective loss $\Phi({\theta})$ as $\nabla_{{\theta}} \Phi({\theta}) =\begin{bmatrix}
%    \nabla \phi^1({\theta}),
%    \dots, \nabla\phi^m({\theta})
%\end{bmatrix} = \begin{bmatrix}
%\nabla_{{\theta}}\fL^1({\theta},\eta_\theta^{1,*})\dots
%\nabla_{{\theta}}\fL^m({\theta}, \eta_\theta^{m,*})
%\end{bmatrix}\in \mathbf{R}^{n\times m}$. In practice, given $\eta=[\eta^1,\dots,\eta^m]$, we evaluate $\nabla_{{\theta}}\fL({\theta}, \eta^i)$ as  $[\nabla_{{\theta}} \fL^1({\theta},\eta^1),\dots,\nabla_{{\theta}} \fL^m({\theta}, \eta^m)]\in \mathbf{R}^{n\times m}$, and $\nabla_{{\eta}}\fL({\theta}, {\eta})$ as $ \text{diag}[\nabla_{\eta^1} \fL^1({\theta}, \eta^1)\dots\nabla_{\eta^m} \fL^m({\theta}, \eta^m)] \in \mathbf{R}^{m\times m} $.
%Combining $\nabla_{{\theta}}\fL({\theta}, {\eta})$ and $\nabla_{{\eta}}\fL({\theta}, \eta)$, the full matrix $\nabla_{{\theta}, {\eta}}\fL({\theta}, {\eta})\in\mathbf{R}^{(n+m)\times m}$ can be expressed as
%\begin{align}
%    \nabla_{{\theta}, {\eta}}\fL({\theta}, {\eta}) =  \begin{bmatrix}
%       \nabla_{{\theta}} \fL^1({\theta},\eta^1) &\dots  & \nabla_{{\theta}} \fL^m({\theta},\eta^m)  \\
%       \nabla_{\eta_1}\fL^1({\theta}, \eta^1) & 0  &0  \\
%        0& \nabla_{\eta_2}\fL^2({\theta}, \eta^2) & 0  \\
%        0& 0 & \nabla_{\eta_m}\fL^m({\theta}, \eta^m) \\
%    \end{bmatrix}\in \mathbf{R}^{(n+m)\times m}. \nonumber
%\end{align}

%When analyzing Double-loop MGDA Algorithm \ref{alg1}, denote $ V_d\in \mathbf{R}^{m\times m}$ as the stochastic gradient estimator of $\nabla_{\eta}\fL({\theta}_t, \eta_{t,d})$, where $t,d$ are iteration index of outer and inner-loop in Algorithm \ref{alg1} respectively; Denote $\Y_t, \Yt_t, \Yb_t$ as the (mutually) independent stochastic gradient estimators of $\nabla_{{\theta}}\fL({\theta}_t, {\eta}_{t,d}),\nabla_{\theta} \fL({\theta}_t, {\eta}_{t,\tilde{d}})$ and $\nabla_{\theta} \fL({\theta}_t, {\eta}_{t,\bar{d}})$. For estimation error of $\nabla_{\eta}\fL(\theta_t,\eta_{t,d})$,
%we denote it as $\Upsilon_{t,d}=[\Upsilon_{t,d}^1\dots\Upsilon_{t,d}^m]\in \mathbf{R}^{m\times m}$, where $\Upsilon^i_{t,d}:=V_d^i- \nabla_{\eta}\fL^i(\theta_t,\eta_{t,d})$. For approximation error of $\nabla_{\theta}\fL(\theta,\eta)$, we denote as ${\Gamma}_{t,2} = (\Gamma^{1}_{t,2},\dots,\Gamma^m_{t,2}), {\Gamma}_{t,3} = (\Gamma^{1}_{t,3},\dots,\Gamma^m_{t,3})\in \mathbf{R}^{n\times m}, {\Gamma}_{t,4} = (\Gamma^{1}_{t,4},\dots,\Gamma^m_{t,4})\in \mathbf{R}^{n\times m}$, where  $\Gamma^i_{t,2}=Y_t^i - \nabla_{\theta}\fL^i(\theta_t,\eta^i_{t,d}), \Gamma^i_{t,3}=\Yb^i - \nabla_{\theta}\fL^i(\theta_t,\eta^i_{t,\bar{d}})$ and $\Gamma^i_{t,4}=\Yt^i_t - \nabla_{\theta}\fL^i(\theta_t,\eta^i_{t,\tilde{d}})$. Similarly, for optimization error among $\nabla_{\theta} \fL(\theta_t,\eta_{t,d}),\nabla_{\theta}\fL(\theta_t,\eta_{t,\bar{d}}),\nabla_{\theta}\fL(\theta_t,\eta_{t,\tilde{d}})$ and $\nabla\Phi(\theta_t)$, we denote them as $A_{t,2},A_{t,3},A_{t,4}$ respectively. At the end, for approximation error among $Y_t,\bar{Y}_t,\tilde{Y}_t$ and $\nabla\Phi(\theta)$, we denote it as $E_{t,2},E_{t,3}, E_{t,4}\in \mathbf{R}^{n\times m}$ respectively, where $E_{t,2} =\Gamma_{t,2}+A_{t,2}= Y_t-\nabla\Phi(\theta),E_{t,3}=\Gamma_{t,3}+A_{t,3}=\bar{Y}_t-\nabla\Phi(\theta)$ and $E_{t,4}=\Gamma_{t,4}+A_{t,4}= \tilde{Y}_t-\nabla\Phi(\theta)$.

%When analyzing Single-loop Double-Clip Algorithm \ref{alg1-3}, for ease of notations, we denote our rescaled loss function as $ \fhL(\theta,\eta)$, i.e.,
%\begin{align}
%    \fhL(\theta,\eta):= \fL({\theta}, G\sqrt{m}{\eta}) = \begin{bmatrix}
% \lambda \mathbb{E}_{\xi\sim\mathbb{P}}\Big[f^*\Big(\frac{\ell^1({\theta};\xi)-G\sqrt{m}\eta^1}{\lambda}\Big)\Big]+G\sqrt{m}\eta^1\\
% \dots\\
% \lambda \mathbb{E}_{\xi\sim\mathbb{P}}\Big[f^*\Big(\frac{\ell^m({\theta};\xi)-G\sqrt{m}\eta^m}{\lambda}\Big)\Big]+G\sqrt{m}\eta^m\\
%\end{bmatrix}.\nonumber
%\end{align}
%Given $(\theta,\eta)$,
%we denote $\nabla_{\theta}\fhL(\theta,\eta)\in \mathbf{R}^{n\times m}, \nabla_{\eta}\fhL(\theta,\eta)\in \mathbf{R}^{m\times m}$ as gradient of $\fhL(\theta,\eta)$ with respect to $\theta$ and $\eta$. Due to the nature of update rule stated in Single-loop Double-clip Algorithm \ref{alg1-3}, denote $Z_t$ as stochastic gradient estimator for $\nabla_{\eta}\fL(\theta_t, \eta_{t})$; And $X_t$ as stochastic gradient estimator for $\nabla_{\theta}\fL(\theta_t,\eta_{t+1})$. As a result, denote  $\HUpsilon_t=[\HUpsilon_t^1\dots\HUpsilon_t^m]\in \mathbf{R}^{m\times m}$ as stochastic approximation error between $Z_t$ and $\nabla_{\eta}\fhL(\theta_t,\eta_t)$, and ${\HGamma_t} = [\HGamma_t^{1},\dots,\HGamma_t^m]$ as the stochastic approximation error between $X_t$ and $\nabla_{\theta} \fhL(\theta_t, \eta_{t+1})$. 

%For problem dependent parameters used in Double-loop MGDA Algorithm \ref{alg1}, we denote $L_0, L_2$ as $L$-smooth parameter for $\fL(\theta,\eta)$, $K_0,K_1,K_2$ are the variance upper bounds of stochastic approximation error $\Gamma_{t,2}^i,\Gamma_{t,3}^i,\Gamma_{t,4}^i, \text{ and }\Upsilon_{t,d}^i,\forall i\in[m]$ respectively. For problem dependent parameters used in Single-loop Double-Clip Algorithm \ref{alg1-3}, we denote $\hL_0, \hL_1,\hL_2$ as the $(L_0,L_1,L_2)$-smooth~\citep{qirevistfDRO} parameters for rescaled function $\fhL(\theta,\eta)$, and $\hat{K}_0,\hat{K}_1,\hat{K}_2$ as the variance upper bound of $\Gamma_t^i,\HUpsilon_t^i$ respectively.
%(and $\tilde{\Upsilon}_t^i$ used in algorithm \ref{alg1-2})

%Throughout, {we restrict our parameter spaces to be Euclidean space, and denote $\|\cdot\|$ as $\ell_2$-norm and $|\cdot|$ as $\ell_1$-norm over Euclidean space for simplicity}. $F,\bar{F}$ are the upper bound of maximal function value gap we constructed in Theorem \ref{thm: convergence of algorithm 1} and Theorem \ref{thm: convergence of alg3}. $\Lambda, \Lambda$ are the gradient upper bound of $\|\nabla\phi^i(\theta_t)\|$ and $\|\nabla_{\eta^i}\fhL^i(\theta_t,\eta_t^i)\|$ leveraging~\eqref{eq: boundgrad}~\footnote{A proof of the gradient--function value relationship~\eqref{eq: boundgrad} can be found in Lemma~5.1 of~\citet{Haochuangensmooth}. We omit the proof here, as the standard $L$-smoothness condition is a special case of the generalized smoothness framework. 
%}, where we focuses on showing that the constructed $F$, $\hat{F}$, and $\bar{F}$ are finite constants, and function value gaps exceeding $F$, $\hat{F}$, and $\bar{F}$ before algorithm terminates are tail events.

%During non-asymptotic convergence analysis of Algorithm~\ref{alg1}, Algorithm~\ref{alg1-3}, denote $0<\delta<1$ as tail-event probability upper bound, $0<\epsilon<1$ as target accuracy, and $\tepsilon=\min\{1,\Lambda^{-1}\}\delta\epsilon^2$ as rescaled-target accuracy when analyzing inner-loop convergence of Algorithm~\ref{alg1}.


\section{Notation for Convergence Analysis}

Table~\ref{tab:convergence-notation} summarizes the notation used
in the convergence analysis of Double-loop MGDA
(Algorithm~\ref{alg1}) and Single-loop Double-Clip
(Algorithm~\ref{alg1-3}).
All parameter spaces are real space.
Throughout the paper, we assume that our studied objective $\phi^i(\theta)$, $\fL^i(\theta, \eta),\fhL^i(\theta, \eta),\forall i\in [m]$ are well-defined, i.e., function value evaluated at any $(\theta, \eta)$ is finite. Its minimal value $\phi^{i,*}$, $\fL^{i,*}$ and $\fhL^{i,*}$ is finite. The corresponding minimizer $\theta^*, \eta_{\theta}^{i,*}$ exists for all $i\in [m]$.
\begingroup
\small
\setlength{\tabcolsep}{6pt}
\renewcommand{\arraystretch}{1.18}
\setlength{\LTleft}{0pt}
\setlength{\LTright}{0pt}

\begin{longtable}{
    @{}
    >{\raggedright\arraybackslash}p{0.27\textwidth}
    >{\raggedright\arraybackslash}
      p{\dimexpr0.73\textwidth-2\tabcolsep\relax}
    @{}
}
\caption{Notation used in the convergence analysis.}
\label{tab:convergence-notation}\\
\toprule
\textbf{Notation} & \textbf{Definition / description} \\
\midrule
\endfirsthead

\multicolumn{2}{@{}l}{%
    \textit{Table~\thetable\ (continued)}}\\
\toprule
\textbf{Notation} & \textbf{Definition / description} \\
\midrule
\endhead

\midrule
\multicolumn{2}{r@{}}{\textit{Continued on next page}}\\
\endfoot

\bottomrule
\endlastfoot

% ------------------------------------------------------------
\multicolumn{2}{@{}l}{\textbf{Basic notation and gradient matrices}}\\*
\addlinespace[2pt]

$n,\ m$
&
Dimension of $\theta$ and number of objectives, respectively;
$[m]=\{1,\ldots,m\}$.
\\

$\theta,\ \eta$
&
$\theta\in\mathbb{R}^{n}$ and
$\eta=(\eta^1,\ldots,\eta^m)\in\mathbb{R}^{m}$.
\\

$\eta_{\theta}^{i,*}$
&
Optimal inner variable for objective $i$ at $\theta$.
\\

$\nabla_{\theta}\Phi(\theta)$
&
Exact multi-objective gradient matrix in $\mathbb{R}^{n\times m}$.
Its $i$th column is
$\nabla\phi^i(\theta)
 =\nabla_{\theta}\fL^i(\theta,\eta_{\theta}^{i,*})$.
\\

$\nabla_{\theta}\fL(\theta,\eta)$
&
Gradient matrix in $\mathbb{R}^{n\times m}$ with $i$th column
$\nabla_{\theta}\fL^i(\theta,\eta^i)$.
\\

$\nabla_{\eta}\fL(\theta,\eta)$
&
\textbf{Diagonal} matrix in $\mathbb{R}^{m\times m}$ with $i$th
diagonal entry $\nabla_{\eta^i}\fL^i(\theta,\eta^i)$.
\\

$\nabla_{\theta,\eta}\fL(\theta,\eta)$
&
Joint gradient matrix:
$\begin{bmatrix}
\nabla_{\theta}\fL(\theta,\eta)\\
\nabla_{\eta}\fL(\theta,\eta)
\end{bmatrix}
\in\mathbb{R}^{(n+m)\times m}$.
\\

$\|\cdot\|,\ |\cdot|$
&
Euclidean $\ell_2$ norm and $\ell_1$ norm, respectively.
\\

% ------------------------------------------------------------
\midrule
\multicolumn{2}{@{}l}{\textbf{Double-loop MGDA}}\\*
\addlinespace[2pt]

$t,\ d$
&
Outer-loop and inner-loop iteration indices, respectively.
\\

$V_d$
&
Stochastic estimator of
$\nabla_{\eta}\fL(\theta_t,\eta_{t,d})$;
$V_d\in\mathbb{R}^{m\times m}$.
\\

$Y_t,\ \bar{Y}_t,\ \tilde{Y}_t$
&
Mutually independent stochastic estimators of the
$\theta$-gradient matrices evaluated at
$\eta_{t,d}$, $\eta_{t,\bar d}$, and
$\eta_{t,\tilde d}$, respectively, with $\theta=\theta_t$.
Each belongs to $\mathbb{R}^{n\times m}$.
\\

$\Upsilon_{t,d}$
&
Inner-gradient estimation error:
$\Upsilon_{t,d}
 =V_d-\nabla_{\eta}\fL(\theta_t,\eta_{t,d})
 \in\mathbb{R}^{m\times m}$.
\\

$\Gamma_{t,2}$
&
Stochastic gradient error:
$\Gamma_{t,2}
 =Y_t-\nabla_{\theta}\fL(\theta_t,\eta_{t,d})$.
\\

$\Gamma_{t,3}$
&
Stochastic gradient error:
$\Gamma_{t,3}
 =\bar{Y}_t
  -\nabla_{\theta}\fL(\theta_t,\eta_{t,\bar d})$.
\\

$\Gamma_{t,4}$
&
Stochastic gradient error:
$\Gamma_{t,4}
 =\tilde{Y}_t
  -\nabla_{\theta}\fL(\theta_t,\eta_{t,\tilde d})$.
\\

$A_{t,2},\ A_{t,3},\ A_{t,4}$
&
Inner-optimization errors:
\newline
$A_{t,2}
 =\nabla_{\theta}\fL(\theta_t,\eta_{t,d})
  -\nabla_{\theta}\Phi(\theta_t)$,
\newline
$A_{t,3}
 =\nabla_{\theta}\fL(\theta_t,\eta_{t,\bar d})
  -\nabla_{\theta}\Phi(\theta_t)$,
\newline
$A_{t,4}
 =\nabla_{\theta}\fL(\theta_t,\eta_{t,\tilde d})
  -\nabla_{\theta}\Phi(\theta_t)$.
\\

$E_{t,2},\ E_{t,3},\ E_{t,4}$
&
Total gradient errors, with
$E_{t,j}=\Gamma_{t,j}+A_{t,j}$ for $j\in\{2,3,4\}$:
\newline
$E_{t,2}=Y_t-\nabla_{\theta}\Phi(\theta_t)$,
\newline
$E_{t,3}=\bar{Y}_t-\nabla_{\theta}\Phi(\theta_t)$,
\newline
$E_{t,4}=\tilde{Y}_t-\nabla_{\theta}\Phi(\theta_t)$.
All $\Gamma_{t,j}$, $A_{t,j}$, and $E_{t,j}$
belong to $\mathbb{R}^{n\times m}$.
\\

$L_0,\ L_2$
&
Smoothness parameters for $\fL(\theta,\eta)$.
\\

$K_0,\ K_1,\ K_2$
&
Constants controlling the variance bounds for the
columnwise stochastic errors
$\Gamma_{t,2}^i$, $\Gamma_{t,3}^i$,
$\Gamma_{t,4}^i$, and $\Upsilon_{t,d}^i$,
$i\in[m]$.
\\
$\calH_t, \calG_t$
& $\sigma$-algebra collects history randomness including initialization and all random draws in outer iterations before $t$; $\calH_t \vee \{\text{$\sigma$-algebra of all inner-solver randomness and output selections at iteration t.} \} $
\\
% ------------------------------------------------------------
\midrule
\multicolumn{2}{@{}l}{\textbf{Single-loop Double-Clip MGDA}}\\*
\addlinespace[2pt]

$\fhL(\theta,\eta)$
&
Rescaled vector-valued loss:
$\fhL(\theta,\eta)=\fL(\theta,G\sqrt{m}\,\eta)$.
Its $i$th component is
\[
\begin{aligned}
\fhL^i(\theta,\eta)
={}&\lambda\mathbb{E}_{\xi\sim\mathbb{P}}
\big[
f^*\!\left(
\frac{\ell^i(\theta;\xi)-G\sqrt{m}\,\eta^i}
     {\lambda}
\right)
\big]+G\sqrt{m}\,\eta^i.
\end{aligned}
\]
\\

$\nabla_{\theta}\fhL,\ \nabla_{\eta}\fhL$
&
Gradient matrices of the rescaled loss with respect to
$\theta$ and $\eta$, belonging to
$\mathbb{R}^{n\times m}$ and
$\mathbb{R}^{m\times m}$, respectively.
\\

$Z_t$
&
Stochastic estimator of
$\nabla_{\eta}\fhL(\theta_t,\eta_t)$;
$Z_t\in\mathbb{R}^{m\times m}$.
\\

$X_t$
&
Stochastic estimator of
$\nabla_{\theta}\fhL(\theta_t,\eta_{t+1})$;
$X_t\in\mathbb{R}^{n\times m}$.
\\

$\HUpsilon_t$
&
Error relative to the rescaled $\eta$-gradient:
$\HUpsilon_t
 =Z_t-\nabla_{\eta}\fhL(\theta_t,\eta_t)
 \in\mathbb{R}^{m\times m}$.
\\

$\HGamma_t$
&
Error relative to the rescaled $\theta$-gradient:
$\HGamma_t
 =X_t-\nabla_{\theta}\fhL(\theta_t,\eta_{t+1})
 \in\mathbb{R}^{n\times m}$.
\\

$\hL_0,\ \hL_1,\ \hL_2$
&
Generalized smoothness parameters of $\fhL$
under the $(L_0,L_1,L_2)$-smoothness framework
of \citet{qirevistfDRO}.
\\

$\hat K_0,\ \hat K_1,\ \hat K_2$
&
Constants controlling the variance bounds for the
columnwise stochastic errors in the single-loop analysis.
\\
$\mathcal{F}_t$
& 
$\sigma$-algebra collects history randomness including initialization and all random draws before $t$.
\\
$\htau,\htau_1,\htau_2$
& Stopping time
\\
% ------------------------------------------------------------
\midrule
\multicolumn{2}{@{}l}{\textbf{Bounds and convergence parameters}}\\*
\addlinespace[2pt]

$\bar F$
&
Upper bounds on the maximal function-value gaps
constructed in Theorem~\ref{thm: convergence of alg3}, respectively.
\\


$\Lambda$
&
Upper bounds on
$\|\nabla_{\eta^i}\fhL^i(\theta_t,\eta_t^i)\|$, obtained using gradient-function value relationship \eqref{eq: boundgrad}.
\\

$\delta\in(0,1)$
&
Upper bound on the tail-event probability.
\\

$\epsilon\in(0,1)$
&
Target accuracy.
\\

$\tepsilon$
&
Rescaled target accuracy for the inner-loop analysis:
$\tepsilon
 =\min\{1,G^{-1}\}\epsilon$.
\\

\end{longtable}
\endgroup

We establish the function-value-gap gradient bound relationship following Lemma~5.1 from \citet{Haochuangensmooth}; 
For function satisfies $\hat{L}_2$-smooth, i.e., $
\|\nabla_{\eta} \fhL^i({\theta},\eta)-\nabla_{\eta} \fhL^i({\theta},\bar{\eta}^{i
}) \|\leq \hat{L}_2\|\eta^i-\bar{\eta}^i \|$
In particular, letting $\bar{\eta}^{i}=\eta^i-\frac{\nabla_{\eta} \fhL^i(\theta,\eta^i)}{\hat{L}_2(\|\nabla_{\eta} \fhL^i(\theta,\eta^i)\|+1)^{1/2}}$ and applying the descent lemma gives 
\begin{align}  
\|\nabla_{\eta} \fhL^i(\theta,\eta^{i})\|^2 \leq 2\hat{L}_2(\fhL^i(\theta,\eta^i)-\inf_{\eta^i}\fhL^i(\theta, \eta^{i}))(\|\nabla_{\eta} \fhL^i(\theta,\eta^{i})\|+ 1). \label{eq: boundgrad}
\end{align}
Therefore, if one can find a constant $\bar{F}$ such that $\fhL^i(\theta,\eta^i)-\fhL(\theta, \eta^{i,*})\leq \bar{F}$, defining $\Lambda := \sup\{u\geq 0| u^2\leq {2\hat{L}_2}\bar{F}(u+1) \}$, we obtain $\|\nabla_{\eta} \fhL^i(\theta_t,\eta^i)\|\leq \Lambda$.

\section{Relevant Properties of $\phi^i(\theta)$}\label{Appendix: Relevant Properties of Phi}
\begin{restatable}[$G$-Lipschitz of $\phi^i(\theta)$]{lemma}{Glipschitzphi}\label{lemma: gradient bound of phi}
    Let Assumption \ref{assumption 1} hold. Denote $\eta_{{\theta}}^{i,*}\in \arg\min_{{\eta}\in \mathbf{R}}\fL^i({\theta},\eta^i)$. For any $\theta,\theta'$, we have$
        \|\phi^i({\theta})-\phi^i(\theta') \|\leq G\|{\theta}-{\theta}' \|$
holds for $i\in[m]$. Consequently, at $(\theta, \eta_{\theta}^{i,*})$, we have $\|\nabla \phi^i(\theta)\|\leq G$.
\end{restatable}
\begin{proof}
    Strong duality between $H^i(\theta,\mathbb{Q}^*)$ and $\phi^i(\theta)$ implies given a fixed $\theta$, there is a corresponding worst-case distribution $\mathbb{Q}^*$ such that $H^i(\theta,\mathbb{Q}^*) = \phi^i(\theta)$. Moreover, at $\theta'$ $\mathbb{Q}^*$ is feasible thus $\phi^i(\theta')\geq H^i(\theta',\mathbb{Q}^*)$. This implies
    \begin{align}
        \phi^i(\theta) - \phi^i(\theta') &\leq H^i(\theta,\mathbb{Q}^*) - H^i(\theta',\mathbb{Q}^*) \nonumber\\
        &\leq \mathbb{E}_{\mathbb{Q}^*}[\ell^i(\theta;\xi)-{\ell}^i(\theta';\xi)]\nonumber\\
        &\leq G\|\theta-\theta'\|,
    \end{align}
    which completes proof. 
\end{proof}
    

\begin{restatable}[$L$-smooth of $\phi^i({\theta})$ along certain $(\theta, \eta)$-trajectory ]{lemma}{smoothproperty}\label{lemma: smooth of phi}
    Let Assumption \ref{assumption 1} hold. Denote $\eta_{{\theta}}^{i,*}\in \arg\min_{{\eta}\in \mathbf{R}}\fL^i({\theta},\eta^i)$. For any $\theta,\theta'$, we have$
        \|\nabla \phi^i({\theta})-\nabla_{{\theta}} \fL^i({\theta'},\eta^{i,*}_{{\theta}}) \|\leq L_0\|{\theta}-{\theta}' \|$
holds for $i\in[m]$, where $L_0=G^2M\lambda^{-1}+L$.
\end{restatable}
%\smoothproperty*
\begin{proof}
Define {matrices} $A$ and $B$ as follows
\begin{align}
    A&= \mathbb{E}_{\xi}\Big[f^{*'}(\frac{\ell^i({\theta};\xi)-\eta_{{\theta}}^{i,*}}{\lambda})\nabla \ell^i({\theta};\xi) - f^{*'}(\frac{\ell^i({\theta};\xi)-\eta_{{\theta}}^{i,*}}{\lambda})\nabla \ell^i({\theta}';\xi) \Big]\nonumber\\
    B&=\mathbb{E}_{\xi}\Big[f^{*'}(\frac{\ell^i({\theta};\xi)-\eta_{{\theta}}^{i,*}}{\lambda})\nabla \ell^i({\theta}';\xi)-f^{*'}(\frac{\ell^i({\theta'};\xi)-\eta_{{\theta}}^{i,*}}{\lambda})\nabla \ell^i({\theta'};\xi) \Big]. \nonumber
\end{align}
It's obvious that $\|A+B\| =\| \nabla \phi_i({\theta})-\nabla_{{\theta}} \fL^i({\theta'},\eta^{i,*}_{{\theta}})\|\leq \| A\|+\|B\|$. Next, we bound $\| A\|$ and $\| B\|$ separately.
For $\|A\|$, we have
\begin{align}
    &\Big\|\mathbb{E}_{\xi}\Big[f^{*'}(\frac{\ell^i({\theta};\xi)-\eta_{{\theta}}^{i,*}}{\lambda})\nabla \ell^i({\theta};\xi)-f^{*'}(\frac{\ell^i({\theta};\xi)-\eta_{{\theta}}^{i,*}}{\lambda})\nabla \ell^i({\theta}';\xi) \Big]\Big\|\nonumber\\
    &\leq \mathbb{E}_{\xi} \big\|f^{*'}(\frac{\ell^i({\theta};\xi)-\eta_{{\theta}}^{i,*}}{\lambda})\nabla \ell^i({\theta};\xi)-f^{*'}(\frac{\ell^i({\theta};\xi)-\eta_{{\theta}}^{i,*}}{\lambda})\nabla \ell^i({\theta}';\xi) \big\| \nonumber\\
    &\leq \mathbb{E}_{\xi}\Big[|{{f^*}'}(\frac{\ell^i({\theta};\xi)-\eta^{i,*}_{{\theta}'}}{\lambda})|\cdot \|\big(\nabla \ell^i({\theta};\xi)-\nabla \ell^i({\theta}';\xi) \big) \big\|\Big]\nonumber\\
    &\leq L \mathbb{E}_{\xi}|{{f^*}'}(\frac{\ell^i({\theta};\xi)-\eta^{i,*}_{{\theta}}}{\lambda}) |\|{\theta}-{\theta}'\|\nonumber\\
    &=L\|{\theta}-{\theta}'\|,
\end{align}
where the first inequality utilizes Jensen's inequality; the second inequality utilizes Cauchy-Schwarz inequality; the third inequality utilizes $L$-smooth of $\ell(\cdot)$ and the last equality utilizes maximization arguments such that $(f^{*})'(\cdot)\geq 0$ (recall that the primal variable is $r =\mathrm{d}\mathbb{Q}/\mathrm{d}\textit{Unif}$, which is random-nykodym derivative over p.d.f.) and $\nabla_{\eta_i}\fL^i(\theta,\eta_{\theta}^{i,*})=1-\mathbb{E}_\xi[f^{*'}(\frac{\ell^i({\theta};\xi)-\eta^{i,*}_{{\theta}}}{\lambda})]=0$.


Similarly, for $\|B\|$, we have
\begin{align}
    &\|\mathbb{E}_{\xi}\Big[f^{*'}(\frac{\ell^i({\theta};\xi)-\eta_{{\theta}}^{i,*}}{\lambda})\nabla \ell^i({\theta}';\xi) - f^{*'}(\frac{\ell^i({\theta}';\xi)-\eta_{{\theta}}^{i,*}}{\lambda})\nabla \ell^i({\theta}';\xi) \Big] \| \nonumber\\
    &\leq \mathbb{E}_{\xi}\|f^{*'}(\frac{\ell^i({\theta};\xi)-\eta_{{\theta}}^{i,*}}{\lambda})\nabla \ell^i({\theta}';\xi) - f^{*'}(\frac{\ell^i({\theta}';\xi)-\eta_{{\theta}}^{i,*}}{\lambda})\nabla \ell^i({\theta}';\xi) \| \nonumber \\
    &\leq \mathbb{E}_{\xi}\Big[\|f^{*'}(\frac{\ell^i({\theta};\xi)-\eta_{{\theta}}^{i,*}}{\lambda}) -f^{*'}(\frac{\ell^i({\theta}';\xi)-\eta_{{\theta}}^{i,*}}{\lambda}) \|\|\nabla \ell^i({\theta}';\xi) \|\Big] \nonumber\\
    &\leq G\mathbb{E}_{\xi}\Big[ \|f^{*'}(\frac{\ell^i({\theta};\xi)-\eta_{{\theta}}^{i,*}}{\lambda}) -f^{*'}(\frac{\ell^i({\theta}';\xi)-\eta_{{\theta}}^{i,*}}{\lambda}) \|\Big] \nonumber\\
    &\leq G^2M\lambda^{-1}\|{\theta}-{\theta}' \|,
\end{align}
where the first inequality applies Jensen's inequality; the second inequality utilizes Cauchy-Schwarz inequality; the third inequality utilizes $G$-Lipschitz of $\ell(\cdot)$; the {fourth} inequality utilizes $M$-smooth of $f^*(\cdot)$, $G$-lipschitz of $\ell(\cdot)$. Combine above two inequality gives desired result.
\end{proof}

%{The gradient estimation $Y_3, Y_4$ is a biased estimator of $\nabla \Phi(\theta)$, making the optinal stopping theorem applied in \cite{qimgda} no longer valid in analysis. We need to derive new martingale sequences.} 









%\begin{corollary}
%Let $\theta' = \theta - \frac{ \nabla_{{\theta}} \fL({\theta}, {\eta}){w}}{L^0+L^1 \|\nabla_{{\eta}}\fL({\theta},{\eta}) \|_F}$, for any $w$, above inequality has equivalent formulation 
%\begin{align}
    %\|\nabla_{{\theta}}\fL({\theta}, {\eta}){w} \|^2\leq 2\Big(\fL({\theta}, {\eta}){w}-\Phi({\theta}^*){w}\Big)\cdot\Big({L^0+L^1\|\nabla_{{\eta}}\fL({\theta}, {\eta}) \|_F}\Big),
%\end{align}
%where $\Phi({\theta}^*)=\min_{{\theta}\in \mathbf{R}^d}\Phi({\theta})$.
%\begin{proof}
%    Put $\theta'=\theta-\frac{\nabla_{{\theta}}\fL({\theta}, {\eta})}{L^0+L^1 \|\nabla_{{\eta}}\fL({\theta},{\eta}) \|_F}$ into above descent lemma \eqref{eq: moo descent lemma}, we have
%    \begin{align}
        %\Phi({\theta}^*)w\leq \fL({\theta}',{\eta})w \leq \fL({\theta}, {\eta})w-\frac{ \|\nabla_{{\theta}}\fL({\theta}, {\eta})w \|^2}{2(L^0+L^1\| \nabla_{{\eta}}\fL({\theta}, {\eta})\|_F)}.
%    \end{align}
%    Re-arrange above inequality gives the desired result.
%\end{proof}
    
%\end{corollary}


%%%%%% MOO variance is redundant, we only need to introduce constant $K_1$, $K_2$ and $K_0$.
%\begin{lemma}[Variance bound]
%Under above assumptions \ref{assumption 1}, \ref{assumption 2}, for ${\delta}_{j}w,j\in\{2,3\}$, $\nu w$, taking expectation over any $\xi\sim \mathbb{P}$, we have
%\begin{align}
    %&\mathbb{E}_{{\xi}}\|\Gamma_j \|_F^2\leq mK_0+K_1 \|\nabla_{{\eta}}\fL({\theta}, {\eta}) \|_F^2,\forall j\in\{2,3 \}, \text{ and } \mathbb{E}_{{\xi}} \| \Upsilon \|_F^2\leq mK_2,
    %\label{eq: variance bound}
%\end{align}
%where $K_0$, $K_1$, $K_2$ are same as above lemma.
%\end{lemma}

%\begin{proof}
%   Expanding $\Gamma_j$, we have 
%    \begin{align}
        %\mathbb{E}_{{\xi}}\|\Gamma_{j} \|^2&=\mathbb{E}_{{\xi}}\|\sum_{i=1}^{m}{\Gamma}^i_{j} \|^2 \overset{(i)}{\leq} \mathbb{E}_{{\xi}^i}\big[ \sum_{i=1}^{m} \|\Gamma^i_{j} \|^2\big] \nonumber\\
        %& \overset{(ii)} {\leq} mK_0 + K_1 \sum_{i=1}^{m} \|\nabla_{\eta^i} \fL({\theta}, \eta^i)\|^2 =K_0 + K_1\|\nabla_{{\eta}} \fL({\theta}, {\eta}) \|_F^2
    %\end{align}
    %where (i) applies Jensen's inequality $f(\sum_{i=1}^{m}w_ix_i)\leq \sum_{i=1}^{m}w_i f(x_i)$ holds for any convex function $f(\cdot)$; (ii) applies independent randomness of $\delta^{i}_{j}$ against $w$ and $w_i \leq 1$. Similarly, for $\Upsilon$, we have
    %\begin{align}
        %&\mathbb{E}_{{\xi}} \|{\Upsilon} \|_F^2=\mathbb{E}_{{\xi}} \|\sum_{i=1}^{m}\Upsilon^i \|_F^2  \overset{}{\leq}\mathbb{E}_{{\xi}} \sum_{i=1}^{m} \|\Upsilon^i_{j} \|^2 \overset{}{\leq} mK_2.
    %\end{align}
%\end{proof}

\subsection{Inner-loop Convergence of Algorithm \ref{alg1}}\label{Appendix: Proof of inner loop}
%Notice that in algorithm \ref{alg1}, we update $\eta_i$ for each $i\in[m]$ separately, where $\nabla_{\eta_i} \fL({\theta},\eta_i)$ is proved to satisfy $L^2$-smooth property as well as bounded variance property. By leveraging the analysis of vanilla SGD in \cite{lanSGD}, we can obtain the following convergence bound for each $\nabla_{\eta_i}\fL({\theta}, \eta_i)$ as well as $\|\nabla_{{\eta}} \fL({\theta}, {\eta}) \|_F$ stated in following corollary.
\begin{lemma}[Lemma 1 \citep{qirevistfDRO}]
Let assumptions \ref{assumption 1} hold, fix $\theta$, gradient $\nabla_{\eta^i}L^i(\theta,\eta^i)$ satisfy $L_2$-smooth property, i.e., for any $\eta^i,\bar{\eta}^{i}\in \mathbf{R}$, we have
\begin{align}
    \|\nabla_{\eta^i}\fL^i(\theta,\eta^i) - \nabla_{\eta^i}\fL^i(\theta,\bar{\eta}^i) \|\leq L_2\|\eta^i-\bar{\eta}^i\|,
\end{align}
where $L_2 = M\lambda^{-1}$.
\end{lemma}
\begin{lemma}[Lemma 3 \citep{qirevistfDRO}]
Let assumption \ref{assumption 1} hold, for each $\fL(\theta_t,\eta^i),i\in[m]$, variance of $\Upsilon_t^i, \Gamma^i_{t,2},\Gamma^i_{t,3},\Gamma^i_{t,4}$ can be upper bounded as
\begin{align}
    \mathbb{E}_{{\xi}^{i}_{t}\sim \mathbb{P}^i}[\|\Gamma^{i}_{t,2} \|^2], \mathbb{E}_{\bar{\xi}^{i}_{t}\sim \mathbb{P}^i}[\|\Gamma^{i}_{t,3} \|^2], &\mathbb{E}_{\tilde{\xi}^{i}_{t}\sim \mathbb{P}^i}[\|\Gamma^{i}_{t,4} \|^2] \leq K_0+K_1|\nabla_{\eta^i}\fL({\theta}_t,\eta^i)|^2, \nonumber \\
    &\text{ and }\mathbb{E}_{\xi^{i}_{t,d}\sim \mathbb{P}^i}[\|\Upsilon_{t,d}^{i} \|^2 ]\leq K_2,
    \label{eq: variance bound}
\end{align}
where $K_0=8G^2+10G^2M^2\lambda^{-2}\kappa^2$, $K_1 = 8G^2$, $K_2 = M^2 \lambda^{-2} \kappa^2$.
\end{lemma}
\begin{proof}
    The proof exactly follows \citet{qirevistfDRO}'s proof by changing $\fL(\theta,G\eta)$ into $\fL(\theta,\eta)$.
\end{proof}
\begin{restatable}[Inner-loop Convergence of Algorithm~\ref{alg1}]{corollary}{innerloopconvergence}
Denote $\calH_t$ as $\sigma$-algebra generated by outer-loop and initialization up to time $t$.
Denote the rescaled accuracy $\tepsilon = \min\{1,G^{-1}\}\epsilon$, and 
$\overline{\Delta}_{\eta}= \max_{0\leq t\leq T, i\in[m]}\big \{\fL^i({\theta}_t,\eta^i_{t,0})-\phi^i({\theta_t})\}<\infty$ as $\epsilon$-independent almost-sure upper bound on the warm-start gaps over all $i\in[m]$ and $t\in [T]$.
For every $\eta^i_{d}$, choosing $\gamma= \min\{\frac{\tepsilon^2}{2L_2K_2G^2},\frac{1}{L_2} \}$, 
after $D=\Theta(\overline{\Delta}_{\eta}K_2L_2^2G^4\tepsilon^{-4})$ iterations, at uniformly sampled index 
$d^i$, we have
\begin{align}
\mathbb{E}_{\eta^i_{t,d}}\big[\|\nabla_{\eta^i} \fL^i({\theta}_t, \eta^i_{t,d}) \|^2|\calH_t\big]\leq {\tepsilon^2/G^2},
\label{eq: inner loop stationary}
\end{align}
which further implies
$\mathbb{E}_{\eta_{t,d}}[\|\nabla_{\eta} \fL({\theta}_t, {\eta}_{t,{d}}) \|_F^2|\calH_t]\leq {m\tepsilon^2/{G^2}}$. Same upper bound also holds for $\mathbb{E}_{\eta^i_{t,\bar{d}}}\big[\|\nabla_{\eta^i} \fL^i({\theta}_t, \eta^i_{t,\bar{d}}) \|^2|\calH_t\big]$ and $\mathbb{E}_{\eta^i_{t,\tilde{d}}}\big[\|\nabla_{\eta^i} \fL^i({\theta}_t, \eta^i_{t,\tilde{d}}) \|^2|\calH_t\big]$.
\end{restatable}
\begin{proof}
    The proof is similar as vanilla SGD analysis~\citep{lanSGD} with minor changes, we present the full steps here for completeness. Denote $\calH_{t,d} = \calH_t \vee \sigma\{\xi_{t,k}^i, 0\leq k<d,i\in[m] \}$
    
    For each $\fL^i({\theta}_t, \eta^{i})$, it satisfies $L_2$-smoothness w.r.t $\eta^i$, following the descent lemma, we have
    \begin{align}
        \fL^i({\theta}_t, \eta^i_{t,d+1})\leq \fL^i({\theta}_t, \eta^i_{t,d}) + \langle\nabla_{\eta^i}\fL^i({\theta}_t, \eta^i_{t,d}),\eta^i_{t,d+1}-\eta^i_{t,d} \rangle + \frac{L_2}{2} \|\eta^i_{t,d+1}-\eta^i_{t,d} \|^2.\nonumber
    \end{align}
    Put the update rule, $\eta^i_{t,d+1}=\eta^i_{t,d} - \gamma V^i_d$ in, for each $i\in[m]$, we then have
    \begin{align}
        \fL^i({\theta}_t, \eta^i_{t,d+1})\leq \fL^i({\theta}_t, \eta^i_{t,d}) - \gamma\langle\nabla_{\eta^i}\fL^i({\theta}_t, \eta^i_{t,d}),V^i_{d} \rangle + \frac{L_2\gamma^2}{2} \|V^i_{d} \|^2.\nonumber
    \end{align}
Taking conditional expectation over $\xi^{i}_{t,d}\sim \mathbb{P}^i$ and leverage variance upper bound of $\Upsilon_{t,d}^i$~\eqref{eq: variance bound}, we have
\begin{align}
    \mathbb{E}_{\xi^{i}_{t,d}}\big[\fL^i({\theta}_t, \eta^i_{t,d+1})|\calH_{t,d}\big]\leq \mathbb{E}_{\xi^{i}_{t,d}}\big[ \fL^i({\theta}_t, \eta^i_{t,d})|\calH_{t,d}\big] - \gamma(1-\frac{\gamma L_2}{2}) \|\nabla_{\eta^i} \fL^i({\theta}_t, \eta^i_{t,d}) \|^2 + \frac{L_2\gamma^2K_2}{2}. \nonumber
\end{align}
For learning rate $\gamma$ satisfies $\gamma \leq \frac{1}{L_2}$, we have 
$\gamma(1-\frac{\gamma L_2}{2})\geq \frac{\gamma}{2}$,
Taking conditional expectation over $\calH_t$, reorganizing above inequality  further implies
\begin{align}
   \frac{\gamma}{2} \mathbb{E}[\|\nabla_{\eta^i}\fL^i({\theta}_t, \eta^i_{t,d}) \|^2|\calH_t] \leq \mathbb{E}_{\xi^i_{t,d}}\big[\fL({\theta}_t, \eta^i_{t,d}) - \fL^i({\theta}_t, \eta^i_{t,d+1}) |\calH_t\big] + \frac{L_2\gamma^2K_2}{2}.\nonumber
\end{align}
Summing above inequality through $0$ to $D-1$, multiplying both sides with $\frac{2}{D\gamma}$, we have
\begin{align}
    \frac{1}{D} \sum_{d=0}^{D-1}\mathbb{E}[\|\nabla_{\eta^i} \fL^i({\theta}_t, \eta^i_{t,d})\|^2|\calH_t] \leq \frac{2\overline{\Delta}_{\eta}}{D\gamma}+L_2\gamma K_2.\nonumber
\end{align}
Choosing $\gamma = \min\{\frac{1}{L_2},\frac{\tilde\epsilon^2}{2L_2K_2G^2}\}$, $D = \overline{\Delta}_{\eta}\max\{\frac{8L_2K_2G^4}{\tepsilon^{4}},\frac{4L_2G^2}{\tepsilon^2} \} = \Theta(\overline{\Delta}_{\eta}L_2K_2G^4\tepsilon^{-4})$, RHS of above inequality satisfies 
\begin{align}
\mathbb{E}_{\eta_{t,d}}[\|\nabla_{\eta^i} \fL^i({\theta}_t, {\eta}^i_{t,{d}}) \|^2|\calH_t]\leq {{\tepsilon^2}/{G^2}},
\end{align}
which completes proof.
%This further implies
%\begin{align}
    %\mathbb{E}_{{d}}\|\nabla_{{\eta}} \fL({\theta}_t, {\eta}_{{d}}) \|_2^2=\max_{i\in[m]}\mathbb{E}_{{d}}\|\nabla_{\eta_i}L^i({\theta}_t, \eta_{i,d}) \|^2\leq \epsilon^2/mG.
%\end{align}
\end{proof}

\begin{corollary}[Estimation bias of $\nabla\phi^i(\theta_t)$]\label{Corollary: estimation bias of inner-loop}
    Let Assumption \ref{assumption 1} hold, for $\eta_{t,d}$ generated from inner-loop of Algorithm~\ref{alg1}, we have     $\|\nabla_{\theta}\fL^i(\theta_t, \eta^i_{t,d}) - \nabla \phi^i(\theta_t) \|^2\leq G^2\|\nabla_{\eta^i}\fL(\theta_t,\eta_{t,d}^i)\|^2$.
    Furthermore, their expectation satisfies
    \begin{align}
        \mathbb{E}_{\eta_{t,d}}[\| A_{t,2}^i\|^2|\calH_t]=\mathbb{E}_{\eta_{t,d}}[\|\nabla_{\theta}\fL^i(\theta_t, \eta^i_{t,d}) - \nabla \phi^i(\theta_t) \|^2|\calH_t]\leq \tepsilon^2.
        \label{eq: approximation error bound }
    \end{align}
    Correspondingly, we also have $ \mathbb{E}_{\eta_{t,d}}[\|\nabla_{\theta}\fL(\theta_t, \eta_{t,d}) - \nabla \phi(\theta_t) \|_F^2|\calH_t]\leq m\tepsilon^2$. Similar bounds also hold for $A_{t,3}$, $A_{t,4}$.
\end{corollary}
\begin{proof}
   Expanding the expression of $\nabla_{\theta}\fL(\theta_t, \eta_{t,d})$ and $\nabla \phi(\theta_t)$, we have
    \begin{align}
        & \|\mathbb{E}_{{\xi}^i_t}\big[f^{*'}(\frac{\ell^i(\theta_t;{\xi}_t^i)-\eta^i_{t,d}}{\lambda})\nabla \ell^i(\theta_t;{\xi}_t^i)-f^{*'}(\frac{\ell^i(\theta_t;{\xi}_t^i)-\eta^{i*}}{\lambda})\nabla\ell^i(\theta_t;{\xi}_t^i)|\calH_t\big] \| \nonumber \\
        &\leq \mathbb{E}_{{\xi}^i_{t}} \big[\|f^{*'}(\frac{\ell^i(\theta_t;{\xi}^i_{t})-\eta^i_{t,d}}{\lambda})\nabla \ell^i(\theta_t;{\xi}^i_t)-f^{*'}(\frac{\ell^i(\theta_t;{\xi}^i_t)-\eta^{i*}}{\lambda})\nabla \ell^i(\theta_t;{\xi}^i_t) \||\calH_t\big] \nonumber \\
        &\leq G\|\mathbb{E}_{{\xi}_{t}^i}[f^{*'}(\frac{\ell^i(\theta_t;{\xi}^i_t)-\eta^i_{t,d}}{\lambda}) - f^{*'}(\frac{\ell(\theta_t;{\xi}_t^i)-\eta^{i*}}{\lambda})|\calH_t]  \| \nonumber\\
        &=G\|\nabla_{\eta^i} \fL^i(\theta_t,\eta^i_{t,d})\|.\nonumber
    \end{align}
    where the first inequality utilizes Jesen's inequality; the second inequality utilizes Cauchy-schwarz inequality, and $G$-Lipschitz assumption of $\ell^i(\cdot)$ holding for $\xi_t^i$ and Cauchy-Schwarz inequality; the last equality utilizes the fact $f^*(\cdot)$ is a monotone non-decreasing convex function, where the sign of
    $|f^{*'}(\frac{\ell^i(\theta_t;{\xi}^i_t)-\eta^i_{t,d}}{\lambda}) - f^{*'}(\frac{\ell^i(\theta_t;{\xi}_t)-\eta^{i*}}{\lambda})|$ merely depends on relative position between $\eta^i_{t,d}$ and $\eta^{i*}$, which is homogeneous among different ${\xi}^i_t$. Thus $|\mathbb{E}_{{\xi}_t^i}[f^{*'}(\frac{\ell^i(\theta_t;{\xi}^i_t)-\eta^i_{t,d}}{\lambda}) - f^{*'}(\frac{\ell^i(\theta_t;{\xi}^i_t)-\eta^{i*}}{\lambda})|\calH_t]|=|\nabla_{\eta^i}\fL^i(\theta_t,\eta^i_{t,d})-\nabla_{\eta^i}\fL^i(\theta,\eta^{i,*})| = |\nabla_{\eta^i}\fL^i(\theta_t,\eta_{t,d}^i)|$. 
    
    Taking square on both sides, and taking conditional expectation over $\mathbb{E}_{\eta_{t,d}^i}$ we have
    \begin{align}
        &\mathbb{E}_{\eta_{t,d}^i}\Big[\|\mathbb{E}_{{\xi}^i_t}\big[f^{*'}(\frac{\ell^i(\theta_t;{\xi}_t^i)-\eta^i_{t,d}}{\lambda})\nabla \ell^i(\theta_t;{\xi}_t^i)-f^{*'}(\frac{\ell^i(\theta_t;{\xi}_t^i)-\eta^{i*}}{\lambda})\nabla\ell^i(\theta_t;{\xi}_t^i)|\calH_t\big] \|^2|\calH_t\Big]\nonumber\\
        &\leq G^2\mathbb{E}_{\eta^i_{t,d}}[\|\nabla_{\eta^i} \fL^i(\theta_t,\eta^i_{t,d})\|^2|\calH_t]\leq \tepsilon^2,\nonumber
    \end{align}
    where the last inequality applies $\mathbb{E}_{\eta^i_{t,d}}[\|\nabla_{\eta^i} \fL^i(\theta_t,\eta^i_{t,d})\|^2|\calH_t]\leq \tepsilon^2/G^2$.
    
    For multi-objective loss $\nabla_{\theta}\fL(\theta_{t},\eta_{t,d})$ and $\Phi(\theta_t)$, we have
    \begin{align}
        \mathbb{E}_{\eta_{t,d}}[\|\nabla_{\theta}\fL(\theta_t, \eta_{t,d}) - \nabla \Phi(\theta_t) \|_F^2|\calH_t] = \sum_{i=1}^{m} \mathbb{E}_{\eta^i_{t,d}}[\|\nabla_{\theta}\fL^i(\theta_t,\eta^i_{t,d}) - \nabla \phi^i(\theta_t) \|^2|\calH_t]\leq m\tepsilon^2, \nonumber
    \end{align}
which completes the proof.
\end{proof}
\begin{corollary}[Stochastic estimation error of $\nabla \Phi(\theta_t)$]\label{Corollary: estimation error of true gradient}
    For estimation error $E_{t,2}:=Y_{t}-\nabla\Phi(\theta)$, we have
    \begin{align}
        &\mathbb{E}_{\eta_{t,d},{\xi}_t}[\|E_{t,2}w_t\|^2|\calH_t]\leq C_B^2,\text{ and }\mathbb{E}_{\eta_{t,d},{\xi}_t}[\|E_{t,2}\|_F^2|\calH_t]\leq {m C_B^2}.
        %=\mathbb{E}_{\eta_{t,d},{\xi}_t}\|\Y_t-\nabla \Phi(\theta_t) \|_F^2
        \label{eq: bias bound},
    \end{align}
    where $C_B^2:= {2K_0}+{2K_1G^{-2}}\tepsilon^2+2\tepsilon^2$. Same upper bound also holds for $E_{t,3}$, $E_{t,4}$. 

\end{corollary}
\begin{proof}
Notice that, for LHS of above inequality, we can upper bound as follows
\begin{align}
    &\mathbb{E}_{\eta_{t,d}^i,{\xi}_t}[\|\Y^i_t-\nabla \phi^i(\theta_t) \|^2|\calH_t]\nonumber\\
    &=\mathbb{E}_{\eta_{t,d},{\xi}^i_t}[\|\Y_t^i- \nabla_{\theta} \fL^i(\theta_t,\eta_{t,d}^i)+\nabla_{\theta} \fL^i(\theta,\eta_{t,d}^i)-\nabla\phi^i(\theta_t)\|^2|\calH_t] \nonumber\\
    &\leq {2}\mathbb{E}_{\eta_{t,d},{\xi}^i_t}[\|\Gamma_{t,2}^i \|^2|\calH_t] + 2\mathbb{E}_{\eta^i_{t,d}}[\|\nabla_{\theta}\fL^i(\theta_t, \eta_{t,d}^i)-\nabla \phi^i(\theta) \|^2|\calH_t]\nonumber\\
    &\leq  {2K_0}+{2K_1\mathbb{E}_{\eta_{t,d}^i}[\|\nabla_{\eta}\fL^i(\theta_t, \eta_{t,d}^i) \|^2}|\calH_t]+2\tepsilon^2 \nonumber\\
    &\leq \underbrace{{2K_0}+{2K_1G^{-2}}\tepsilon^2}_{\sigma^2}+2\tepsilon^2 = C_B^2, \nonumber
\end{align}
where the first inequality, we utilize $(a+b)^2\leq 2a^2+2b^2$, the second inequality utilizes variance bound  and \eqref{eq: approximation error bound }, respectively and the last inequality utilizes \eqref{eq: inner loop stationary} to upper bound $\mathbb{E}_{\eta_{t,d}^i}\|\nabla_{\eta}\fL^i(\theta_t, \eta_{t,d}^i) \|^2$ obtained from inner-loop convergence.

Similarly, for $\mathbb{E}_{\eta_{t,d},\hat{\xi}}[\|(\Y_t-\nabla \Phi(\theta_t))w_t\|^2|\calH_t]$, we conclude
\begin{align}
    \mathbb{E}_{\eta_{t,d},{\xi}_t}[\|(\Y_t-\nabla \Phi(\theta_t))w_t\|^2|\calH_t] &=\mathbb{E}_{\eta_{t,d},{\xi}_t}[\|\sum_{i=1}^{m}(Y_t^i - \nabla\phi^i(\theta_t))w_t^i\|^2|\calH_t]\nonumber\\
    &\leq \mathbb{E}_{\eta^i_{t,d},{\xi}_t^i}\big[\sum_{i=1}^{m}w_t^i \|Y_t^i-\nabla\phi^i(\theta_t)\|^2|\calH_t\big] \nonumber\\
    &= \sum_{i=1}^{m} w_t^i \mathbb{E}_{\eta_{t,d}^i,{\xi}_t^i}\big[\|Y_t^i - \nabla\phi^i(\theta) \|^2|\calH_t \big]
    \leq  C_B^2,\nonumber
\end{align}
where the inequality applies Jensen's inequality, and $w_t$ is $\calH_t$-measurable and independent with $\xi_t$ and $\eta_{t,d}$.
\end{proof}
\section{Convergence of Algorithm \ref{alg1}}\label{Appendix: proof of descent lemma algorithm 1}
\subsection{Formal Statement of Theorem \ref{thm: convergence of algorithm 1} and Proof}\label{Appendix: proof of convergence algorithm 1}
Given problem dependent parameter $m$,$G$, $L_0=G^2M\lambda^{-1}+L$, $L_2=M\lambda^{-1}$, $K_0=8G^2+10G^2M^2\lambda^{-2}\kappa^2$, $K_1=8G^2$, $K_2=M^2\lambda^{-2}\kappa^2$, denote $\Delta_{\theta_0} = \max_{i\in[m]}\{\phi^i(\theta_0)-\phi^{i,*}\}<\infty$ as finite initial outer-objective-gap.
Set the hyper-parameters in Algorithm \ref{alg1} as follows
\begin{align}
    \sigma^2 &= {2K_0}+{2K_1G^{-2}}\min\{1,G^{-2}\}\epsilon^2 =\mathcal{O}(K_0)=\mathcal{O}(G^2M^2\lambda^{-2}\kappa^2)\nonumber\\
    C_B^2&=\sigma^2+2\min\{1,G^{-2}\}\epsilon^2 = \mathcal{O}(G^2M^2\lambda^{-2}\kappa^2)\nonumber\\
    \gamma&= \min\{\frac{\epsilon^2}{2L_2K_2G^4},\frac{\epsilon^2}{2L_2K_2G^2},\frac{1}{L_2} \}=\mathcal{O}(\frac{\lambda^3\epsilon^2}{M^3G^4\kappa^2})
    \nonumber\\
    \rho&= \frac{1}{16}\epsilon^2=\mathcal{O}(\epsilon^2) \nonumber\\
    \alpha&= \min\Big\{\frac{1}{2L_0},\frac{\epsilon^2}{8L_0C_B^2}\Big\} = \mathcal{O}(\frac{\lambda^{3}\epsilon^2}{G^4M^3\kappa^2}) \nonumber\\
    \beta&= \min \Big\{\frac{\epsilon^2}{32(mG^4+2mG^2C_B^2+mC_B^4)},\frac{2}{4\rho}\Big\}= \mathcal{O}(\frac{\epsilon^2}{m(G^2+G^2M^2\lambda^{-2}\kappa^2)^2})
    \nonumber\\
    D&= \overline{\Delta}_{\eta}\max\{4G^2\gamma^{-1}\epsilon^{-2}, 4G^4\gamma^{-1}\epsilon^{-2}, \} = \mathcal{O}( \overline{\Delta}_{\eta}G^4\gamma^{-1}\epsilon^{-2})\nonumber\\
    T&=\max\Big\{ {8}\beta^{-1}\epsilon^{-2}, 8\Delta_{\theta_0}\alpha^{-1}\epsilon^{-2}\Big\}=\max\Big\{\Theta(\Delta_{\theta_0}\alpha^{-1}\epsilon^{-2}), \Theta(\beta^{-1}\epsilon^{-2})\Big\}=\Theta(\epsilon^{-4}),
\end{align}
where
$\overline{\Delta}_{\eta}= \max_{0\leq t\leq T, i\in[m]}\big \{\fL^i({\theta}_t,\eta^i_{t,0})-\phi^i({\theta_t})\}<\infty$ is $\epsilon$-independent almost-sure upper bound on the warm-start gaps over all $i\in[m]$ and $t\in [T]$.
Then, we then have the following convergence statement on Double-Loop Algorithm \ref{alg1}.

\begin{restatable}[Formal Statement of Theorem~\ref{thm: convergence of algorithm 1}]{theorem}{convegencefirstalg}\label{restatethm: convergence of algorithm 1}
    Let Assumption \ref{assumption 1} hold. Given $0<\epsilon<1$, Denote $\Delta_{\theta_0} = \max_{i\in[m]}\{\phi^i(\theta_0)-\phi^{i,*}\}$ and 
    $\overline{\Delta}_{\eta}= \max_{0\leq t\leq T, i\in[m]}\big \{\fL^i({\theta}_t,\eta^i_{t,0})-\phi^i({\theta_t})\}<\infty$ as $\epsilon$-independent almost-sure upper bound on the warm-start gaps over all $i\in[m]$ and $t\in [T]$.
    Set $\rho=\mathcal{O}(\epsilon^2)$, $\beta=\mathcal{O}(\frac{\epsilon^2}{m(G^2+G^2M^2\lambda^{-2}\kappa^2)^2})$,
    $\alpha = \mathcal{O}(\frac{\lambda^{3}\epsilon^2}{G^4M^3\kappa^2})$, and $\gamma = \mathcal{O}(\frac{\lambda^3\epsilon^2}{M^3G^2\kappa^2})$ for Algorithm \ref{alg1}. Then, after $T=\max\big\{\Theta(\Delta_{\theta_0}\alpha^{-1}\epsilon^{-2}), \Theta(\beta^{-1}\epsilon^{-2})\big\}$ outer iterations, each associated with $D=\Theta(\overline{\Delta}_{\eta}G^4\gamma^{-1}\epsilon^{-2})$ inner-iterations, we have 
    \begin{align}
        \frac{1}{T}\sum_{t=0}^{T-1}\mathbb{E}[\|\nabla \Phi(\theta_t)w_t\|^2]\leq 94\epsilon^2.
    \end{align}
\end{restatable}
%\begin{restatable}[Descent Lemma of Algorithm \ref{alg1}]{lemma}{descentlemmafirstalgorithm}\label{lemma: descent lemma of algorithm1}
%    Under the same parameter choices stated in Theorem \ref{thm: convergence of algorithm 1}, for any $w\in \mathcal{W}$, we have
%    \begin{align}
%        \mathbb{E}[\Phi(\theta_\tau)w]-\Phi^*w\leq \frac{F\delta}{8} -\frac{\alpha}{2} \mathbb{E}[\sum_{t=0}^{\tau-1}\|\nabla \Phi(\theta_t)w_t \|^2],
%    \end{align}
%holds. 
%\end{restatable}
\begin{proof}
According to Lemma~\ref{lemma: gradient bound of phi}, we know $\|\nabla \phi^i(\theta_t)\|\leq G$. Put the update rule $\theta_{t+1} = \theta_t - \alpha Y_t{w}_t$ in the descent inequality of $L_0$-smooth function, we have
\begin{align}
    \Phi(\theta_{t+1})w&\leq \fL({\theta}_{t+1}, {\eta}_{\theta_t}^*){w} \leq \Phi(\theta_t){w} -\alpha\langle\nabla \Phi({\theta}_t){w},\Y_{t}{w}_t\rangle  +\frac{L_0}{2} \alpha^2 \|\Y_t{w}_t \|^2\nonumber\\
    &= \Phi(\theta_t){w} -\alpha\langle \nabla \Phi(\theta_t)w,\nabla \Phi(\theta_t) {w}_t\rangle -\alpha \langle \nabla \Phi(\theta_t)w,(\Y_{t}-\nabla\Phi(\theta_t))w_t \rangle  +\frac{L_0}{2} \alpha^2 \|\Y_{t}{w}_t \|^2 \nonumber\\
    &=\Phi(\theta_t){w} - \alpha\langle \nabla \Phi(\theta_t)w,\nabla \Phi(\theta_t) {w}_t\rangle - \alpha{\langle \nabla \Phi(\theta_t)w, (Y_t - \nabla_{\theta} \fL(\theta_t, \eta_{t,d}))w_t\rangle} \nonumber\\
    &\quad -\alpha{\langle \nabla \Phi(\theta_t)w, (\nabla_{\theta} \fL(\theta_t, \eta_{t,d})-\nabla \Phi(\theta_t))w_t\rangle }+\frac{L_0}{2}\alpha^2 \|Y_tw_t \|^2,
    %&= \Phi(\theta_t)w - \alpha \langle  \nabla \Phi(\theta_t)w,\nabla \Phi(\theta_t) {w}_t\rangle -\alpha \langle \nabla \Phi(\theta_t){w}, \Gamma_2{w}_t \rangle
    %\nonumber\\
    %&\quad - \alpha \langle\nabla\Phi(\theta_t)w, (\nabla_{\theta}\fL(\theta_t, \eta_{t,d})-\nabla\Phi(\theta_t))w_t \rangle  +\frac{L_0}{2} \alpha^2 \|Y_{2}{w}_t \|^2
    \label{eq: intermediate descent lemma alg1}
\end{align}
where the first equality  $\Y_tw_t = (\Y_t-\nabla \Phi(\theta_t)+\Phi(\theta_t))w_t$ and the second inequality further decompose $Y_t - \nabla \Phi(\theta_t)= Y_t-\nabla_{\theta} \fL(\theta_t, \eta_{t,d})+\nabla_{\theta}\fL(\theta_t, \eta_{t,d})-\nabla \Phi(\theta_t)$.
%the second equality further expands $Y_2-\nabla\Phi(\theta_t)=Y_2 - \nabla_{\theta} \fL(\theta_{t},\eta_{t,d})+\nabla_{\theta} \fL(\theta_t, \eta_{t,d})-\nabla\Phi(\theta_t)$.


Next, we upper bound the term $\langle \nabla \Phi({\theta}_t)w, \nabla \Phi({\theta}_t){w}_t \rangle$ from analyzing $\|w_{t+1}-w \|^2$.
\begin{align}
    \|{w}_{t+1} - {w} \|^2
    &= \| \Pi_{\mathcal{W}} \big({w}_t - \beta \big[\Yb_t^{\top}\Yt_t{w}_t+\rho {w}_t \big]\big) - {w} \|^2\leq \|{w}_t - \beta \big[\Yb_t^{\top}\Yt_t{w}_t+\rho {w}_t \big]- {w} \|^2 \nonumber\\
    &= \|{w}_t - {w} \|^2 -2\beta \langle {w}_t-{w}, \Yb_t^{\top}\Yt_t{w}_t+\rho {w}_t \rangle + \beta^2\|\Yb_t^{\top}\Yt_t{w}_t+\rho {w}_t\|^2,
\end{align}
where the first inequality is due to non-expansiveness of projection over probability simplex.

Re-arrange above inequality, we have
\begin{align}
    &2\beta \langle {w}_t-{w},\Yb_t^{\top}\Yt_t{w}_t \rangle \nonumber\\
    &\leq \big(\|{w}_t -{w}\|^2-\|{w}_{t+1} -{w} \|^2 \big) -2\beta\rho\langle {w}_t-{w}, {w}_t\rangle + \beta^2 \|\Yb_t^{\top}\Yt_t{w}_t+\rho {w}_t\|^2\nonumber\\
    &{\leq} \big(\|{w}_t -{w}\|^2-\|{w}_{t+1} -{w} \|^2 \big)+4\beta \rho + 2\beta^2\rho^2+ 2\beta^2\|\Yb_t^{\top}\Yt_t{w}_t \|^2,
    \label{eq: wtsquare}
\end{align}
where the first inequality applies $\langle w_t-w,w_t \rangle= \|w_t\|^2-\langle w, w_t\rangle\leq 2$, and $(a+b)^2\leq 2a^2+2b^2$. 

For $\Yb_t^{\top}\Yt_t$, one can establish the relationship as follows
\begin{align}
    \Yb_t^{\top}\Yt_t &=
    (\underbrace{\Yb_t-\nabla\Phi(\theta_t)}_{E_{t,3}}+\nabla \Phi(\theta_t))^{\top}(\underbrace{\Yt_t-\nabla\Phi(\theta_t)}_{E_{t,4}}+\nabla \Phi(\theta_t)) \nonumber \\
    & = E_{t,3}^{\top}E_{t,4} + E_{t,3}^{\top}\nabla\Phi(\theta_t)+\nabla \Phi(\theta_t)^{\top}E_{t,4}+\nabla \Phi(\theta_t)^{\top}\nabla\Phi(\theta_t),
    \label{eq: Y3Y4 reltion}
    %&= (\Gamma_3 +E_t+ \nabla\Phi(\theta_t))^{\top}(\Gamma_4+ E_t+\nabla \Phi(\theta_t))\nonumber\\
    %&= \Gamma_3^{\top}\Gamma_4+\Gamma_3^{\top} E_t+\Gamma_4^{\top}E_t +  E_t^{\top}E_t+\Gamma_3^{\top} \nabla \Phi(\theta_t)+\Gamma_4^{\top} \nabla \Phi(\theta_t)+2E_t^{\top}\nabla \Phi(\theta_t)+\nabla\Phi(\theta_t)^{\top}\Phi(\theta_t) \nonumber
\end{align}
Put above equality into \eqref{eq: wtsquare}, we have 
\begin{align}
    &2\beta \langle {w}_t-{w},\Yb_t^{\top}\Yt_t{w}_t \rangle \nonumber\\
    & \overset{(i)}{\leq} \big(\|{w}_t -{w}\|^2-\|{w}_{t+1} -{w} \|^2 \big)+4\beta \rho + 2\beta^2\rho^2\nonumber \nonumber\\
    &\quad + 2\beta^2 \|E_{t,3}^{\top}E_{t,4}w_t + E_{t,3}^{\top}\nabla\Phi(\theta_t)w_t+\nabla \Phi(\theta_t)^{\top}E_{t,4}w_t+\nabla \Phi(\theta_t)^{\top}\nabla \Phi(\theta_t)w_t \|^2\nonumber\\
    &\overset{(ii)}{\leq} \big(\|{w}_t -{w}\|^2-\|{w}_{t+1} -{w} \|^2 \big)+4\beta \rho + 2\beta^2\rho^2\nonumber\\
    & \quad+8\beta^2 \|E_{t,3}^{\top}E_{t,4}{w}_t \|^2 + 8\beta^2\|\nabla \Phi(\theta_t)^{\top}\nabla \Phi(\theta_t){w}_t \|^2 \nonumber\\
    &\quad +8\beta^2\|E_{t,3}^{\top}\nabla \Phi(\theta_t){w}_t \|^2+ 8\beta^2\|\nabla \Phi(\theta_t)^{\top}E_{t,4}{w}_t \|^2,
\end{align}
where (i) utilizes inequality \eqref{eq: Y3Y4 reltion}; (ii) utilizes the fact $(a+b+c+d)^2\leq 4a^2+4b^2+4c^2+4d^2$.

%In proof of Theorem \ref{thm: convergence of algorithm 1}, we constructed stopping time $\tau$ such that $\forall i\in[m], \phi^i(\theta_0)-\phi^{i,*}\leq F$ holds for any $t\leq \tau$. For the event $\tau=T$, this further implies 
For $\|\nabla \Phi(\theta_t)\|$ and $\|\nabla \Phi(\theta_t)w_t\|$, we have
$\|\nabla \Phi(\theta_t) \|_F\leq \sqrt{m}G$ and $\|\nabla \Phi(\theta_t)w_t \| \leq \sum_{i=1}^{m}w_t^i\| \nabla \phi^i(\theta_t)\|\leq G$ holds for $t\leq T$. Thus, for $\|\nabla\Phi(\theta_t)^{\top}\nabla \Phi(\theta_t){w}_t \|^2$, we can upper bound it as follows
\begin{align}
    &\|\nabla\Phi(\theta_t)^{\top}\nabla \Phi(\theta_t){w}_t \|^2\overset{(i)}{\leq} \|\nabla\Phi(\theta_t)\|_F^2\|\nabla\Phi(\theta_t)w_t \|^2\leq {m G^4},
    %&\overset{(vi)}{\leq} \Big(2\cdot\sum_{i=1}^{m}\Delta_i\cdot \big(L^0+L^1|\nabla_{\eta_{i,d}}\fL({\theta}_t, \eta_{i,d})|\big)\Big)^2\nonumber\\
    %&\overset{(vii)}{=}\Big(2M\Delta_{\max}L^0+2\Delta_{\max}L^1\sum_{i}^{m}|\nabla_{\eta_{i,d}}\fL({\theta}_t, \eta_{i,d})| \Big)^2\nonumber\\
    %&\overset{(viii)}{\leq}4M^2\Delta_{\max}^2L^0^2+4\Delta_{\max}^2L^1^2M(\sum_{i=1}^{m}|\nabla_{\eta_{i,d}}\fL({\theta}_t, \eta_{i,d})|)^2\nonumber \\
    %&\overset{(viiii)}{=}4M^2\Delta_{\max}^2L^0^2+4\Delta_{\max}^2L^1^2M\|\nabla_{\eta_i}\fL({\theta}_t, {\eta}_{i,d})\|_F^2
    %&\overset{(vii)}{\leq} \Big(\sum_{i=1}^{m} 2\Delta_i^2 + 2L^0^2+2L^1^2|\nabla_{\eta_i}\fL({\theta}_t, \eta_{i,D}) |^2 \Big)^2\nonumber\\
    %&\overset{(viii)}{\leq} 6M^2 \Delta_{\max}^2+6M^2L^0^4+6L^1^4\|\nabla_{\eta}\fL({\theta}_t, {\eta}_D)\|_F^4,
\end{align}
where (i) applies Cauchy-Schwarz inequality, sub-multiplicative property of Frobenius norm.
%(vi) applies fact \eqref{eq: induced single-objective descent lemma}; (vii), (viii) and (viiii) applies fact $(a+b)^2\leq 2a^2+2b^2$, definition of Frobenius norm, Cauchy-Schwarz inequality and denote $\Delta_{\max} = \max_{i\in [m]}\Delta_i$.

Similarly, for $ \|E_{t,3}^{\top}\nabla \Phi(\theta_t){w}_t \|^2$,$\|E_{t,4}^{\top}\nabla \Phi(\theta_t){w}_t \|^2$, we have
\begin{align}
    \| E_{t,3}^{\top}\nabla \Phi(\theta_t){w}_t \|^2
    \leq \|\nabla\Phi(\theta_t){w}_t\|^2 \|E_3 \|_F^2\leq {G^2 \|E_{t,3} \|_F^2}.\nonumber
    %&\leq \Big(2ML^0\Delta_{\max}+2L^1\sum_{i=1}^{m}|\nabla_{\eta_{i,d}}\fL({\theta}_t, \eta_{i,d})| \Big)\|{\delta}_{t,4}^{{\theta}_t}{w}_t\|^2 \nonumber\\
    %&\leq \Big(2ML^0\Delta_{\max}+2L^1\sqrt{\sum_{i=1}^{m}|\nabla_{\eta_{i,d}}\fL({\theta}_t, \eta_{i,d})|^2} \Big)\|{\delta}_{t,4}^{{\theta}_t}{w}_t\|^2 \nonumber \\
    %&= \Big(2ML^0\Delta_{\max}+2L^1\|\nabla_{{\eta}}\fL({\theta}_t, {\eta}_d) \|_F\Big)\|{\delta}_{t,4}^{{\theta}_t}{w}_t\|^2
\end{align}
Same arguments applies to $\|\nabla\Phi(\theta_t)^{\top}E_{t,4}{w}_t \|^2\leq mG^2 \|E_{t,4}{w}_t\|^2$.

Combine above inequalities, we have
\begin{align}
    &2\beta \langle {w}_t-{w},\Yb_t^{\top}\Yt_t{w}_t \rangle\nonumber\\
    &\leq \big(\|{w}_t -{w}\|^2-\|{w}_{t+1} -{w} \|^2 \big)+4\beta \rho + 2\beta^2\rho^2+8\beta^2 \|E_{t,3}^{\top}E_{t,4}{w}_t\|^2\nonumber\\
    &\quad + 8\beta^2 mG^4 + 8\beta^2 G^2(\|E_{t,3}\|_F^2 +m\|E_{t,4}w_t\|_F^2 ).
    %&\quad+ 32\beta^2M^2\Delta_{\max}^2L^0^2+32\beta^2\Delta_{\max}^2L^1^2M\|\nabla_{\eta_i}\fL({\theta}_t, {\eta}_{i,d})\|_F^2\nonumber\\
    %&\quad + \big(16\beta^2ML^0\Delta_{\max}+16\beta^2L^1\|\nabla_{{\eta}}\fL({\theta}_t, {\eta}_d) \|_F\big)\|\big({\delta}_{t,3}^{{\theta}_t}{w}_t\|^2+\|{\delta}_{t,4}^{{\theta}_t}{w}_t\|^2 \big)
\end{align}
Put equation $\Yb_t^{\top}\Yt_t=(\nabla\Phi({\theta}_t)+E_{t,3})^{\top}(\nabla \Phi(\theta_t)+E_{t,4})$ in LHS and re-arrange above inequality, we have
\begin{align}
    &2\beta\langle {w}_t-{w},\nabla \Phi({\theta}_t)^{\top}\nabla \Phi({\theta}_t){w}_t \rangle \nonumber\\
    &\leq \big(\|{w}_t -{w}\|^2-\|{w}_{t+1} -{w} \|^2 \big)+4\beta \rho + 2\beta^2\rho^2+8\beta^2 \|E_{t,3}^{\top} E_{t,4}{w}_t\|^2\nonumber\\
    & \quad -2\beta \langle w_t-w, E_{t,3}^{\top}E_{t,4}{w}_t+\nabla \Phi({\theta}_t)^{\top}E_{t,4}w_t+E_{t,3}^{\top}\nabla \Phi({\theta}_t)w_t\rangle\nonumber\\
    & \quad + 8\beta^2 mG^4 + 8\beta^2G^2(\|E_{t,3}\|_F^2 +m\|E_{t,4}w_t\|_2^2 ).
    %&\quad+ 32\beta^2M^2\Delta_{\max}^2L^0^2+32\beta^2\Delta_{\max}^2L^1^2M\|\nabla_{\eta_i}\fL({\theta}_t, {\eta}_{i,d})\|_F^2\nonumber\\
    %&\quad + \big(16\beta^2ML^0\Delta_{\max}+16\beta^2L^1\|\nabla_{{\eta}}\fL({\theta}_t, {\eta}_d) \|_F\big)\|\big({\delta}_{t,3}^{{\theta}_t}{w}_t\|^2+\|{\delta}_{t,4}^{{\theta}_t}{w}_t\|^2 \big)
    \label{eq: intermediate step on wt}
\end{align}
Combine   \eqref{eq: intermediate descent lemma alg1} and \eqref{eq: intermediate step on wt}, we have
\begin{align}
    &\Phi(\theta_{t+1})w - \Phi(\theta_t)w\nonumber\\
    &\leq \fL({\theta}_{t+1},{\eta}_{\theta_t}^*){w} - \Phi(\theta_t){w} \nonumber\\
    &\overset{(i)}{\leq} -\alpha \|\nabla \Phi(\theta_t){w}_t \|^2 -\alpha {\langle \nabla\Phi(\theta_t){w},(\Y_t - \nabla_{\theta} \fL(\theta_t,\eta_{t,d})) {w}_t\rangle} -\alpha {\langle \nabla\Phi(\theta_t)w, (\nabla_{\theta} \fL(\theta_t, \eta_{t,d})-\nabla\Phi(\theta_t))w_t\rangle}\nonumber\\
    & \quad + \alpha^2L_0\|\nabla \Phi({\theta}_t){w}_t \|^2 + \alpha^2L_0\|E_{t,2}{w}_t\|^2 \nonumber\\
    &\quad +\frac{\alpha}{2\beta}\Big(\big(\|{w}_t -{w}\|^2-\|{w}_{t+1} -{w} \|^2 \big)+4\beta \rho + 2\beta^2\rho^2+8\beta^2 \|E_3^{\top} E_4{w}_t\|^2\nonumber\\
    & \quad -2\beta \langle {w}_t-{w}, E_{t,3}^{\top}E_{t,4}{w}_t+\nabla \Phi({\theta}_t)^{\top}E_{t,4}w_t+E_{t,3}^{\top}\nabla \Phi({\theta}_t){w}_t\rangle\nonumber\\
    &\quad + 8\beta^2 mG^4 + 8\beta^2G^2(\|E_{t,3} \|_F^2 +m\|E_{t,4} w_t\|^2 ) \Big),
    \label{eq: one step descent inequality}
\end{align}
where (i) utilizes $(a+b)^2\leq 2a^2+2b^2$ to further decompose $\frac{\alpha^2 L_0}{2}\|Y_tw_t \|^2=\frac{\alpha^2L_0}{2}\|\nabla\Phi(\theta_t)w_t+E_{t,2}w_t \|^2$, and upper bound this term by $\alpha^2L_0 \|\nabla \Phi(\theta_t)w_t\|^2+\alpha^2L_0\|E_{t,2}w_t\|^2$ and replace $\langle \nabla \Phi(\theta_t)w, \nabla \Phi(\theta_t)w_t\rangle$ by \eqref{eq: intermediate step on wt}.


Sum above equations from $0$ to $T-1$, taking expectation over all randomness including $\xi_t, \bar{\xi}_t, \tilde{\xi}_t,\theta_t,\eta_{t,d},\eta_{t,\bar{d}},\eta_{t,\tilde{d}}, w_t$. For ${\alpha\leq \frac{1}{2L_0}}$, we upper bound $-\alpha \|\nabla\Phi(\theta_t)w_t\|^2+\alpha^2 L_0\|\nabla\Phi(\theta_t)w_t\|^2$ by $-\frac{\alpha}{2}\|\nabla \Phi(\theta_t)w_t\|^2$ and have following inequality
\begin{align}
    &\mathbb{E}[\Phi(\theta_{T}){w}] - \mathbb{E}[\Phi({\theta}_0){w}] \nonumber\\
    &\leq -\frac{\alpha}{2} \mathbb{E}\big[\sum_{t=0}^{T-1}\|\nabla \Phi({\theta}_t){w}_t \|^2\big] -\alpha {\mathbb{E}\big[ \sum_{t=0}^{T-1}\langle \nabla \Phi({\theta}_t){w},\underbrace{(Y_{t}-\nabla_{\theta}\fL(\theta_t, \eta_{t,d}))}_{\Gamma_{t,2}}{w_t}\rangle \big]}\nonumber\\
    &\quad -\alpha {\mathbb{E}\big[\sum_{t=0}^{T-1}\langle \nabla \Phi(\theta_t)w, \underbrace{(\nabla \fL(\theta_t,\eta_{t,d}) - \nabla \Phi(\theta_t))}_{A_{t,2}} w_t \big]} \nonumber\\
    &\quad -\alpha { \mathbb{E}\big[\sum_{t=0}^{T-1} \langle w_t - w, E_{t,3}^{\top} E_{t,4}w_t \rangle+\langle w_t -w, \nabla \Phi(\theta_t)^{\top}E_{t,4}w_t\rangle+\langle w_t -w, E_{t,3}^{\top}\nabla \Phi(\theta_t)w_t \rangle \big]} \nonumber\\
    &\quad +\frac{\alpha}{2\beta}\big(\|w_{0}-w \|^2-\|w_{\tau}-w\|^2\big)+ 2\alpha \rho T + \alpha\beta\rho^2T + 4\alpha\beta m G^4 T+\alpha^2 L_0 \mathbb{E}\big[ \sum_{t=0}^{T-1} \|E_{t,2}w_t\|^2\big] \nonumber\\
    &\quad +4\alpha\beta G^2 \mathbb{E}\big [\sum_{t=0}^{T-1} \|E_{t,3} \|_F^2+m\|{E}_{t,4}w_t\|^2 \big] + 4\alpha\beta \mathbb{E}\big[\sum_{t=0}^{T-1} \|E_{t,3}^{\top}E_{t,4}w_t\|^2 \big] \nonumber\\
    &\leq -\frac{\alpha}{2} \mathbb{E}\big[\sum_{t=0}^{T-1}\|\nabla \Phi({\theta}_t){w}_t \|^2\big] -\alpha {\mathbb{E}\big[ \sum_{t=0}^{T-1}\langle \nabla \Phi({\theta}_t){w},\Gamma_{t,2}{w_t}\rangle \big]}-\alpha {\mathbb{E}\big[\sum_{t=0}^{T-1}\langle \nabla \Phi(\theta_t)w, A_{t,2} w_t \rangle\big]} \nonumber\\
    &\quad +\alpha {\mathbb{E}\big[\sum_{t=0}^{T-1} \langle w-w_t, \Gamma_{t,3}^{\top} (\Gamma_{t,4}+A_{t,4})w_t \rangle +\langle w-w_t, A_{t,3}^{\top} (\Gamma_{t,4}+ A_{t,4})w_t \rangle \big]} \nonumber\\
    &\quad +\alpha{\mathbb{E}\big[\sum_{t=0}^{T-1}\langle w -w_t, \nabla \Phi(\theta_t)^{\top}\Gamma_{t,4}w_t\rangle + \langle w -w_t, \nabla \Phi(\theta_t)^{\top}A_{t,4}w_t\rangle
    \big]} \nonumber\\
    &\quad +\alpha{\mathbb{E}\big[\sum_{t=0}^{T-1}\langle w -w_t, \Gamma_{t,3}^{\top}\nabla \Phi(\theta_t)w_t \rangle + \langle w -w_t, A_{t,3}^{\top}\nabla \Phi(\theta_t)w_t \rangle\big]} \nonumber\\
    &\quad +\frac{\alpha}{\beta}+ 2\alpha \rho T + \alpha\beta\rho^2T + 4\alpha\beta m G^4 T+\alpha^2 L_0 \mathbb{E}\big[ \sum_{t=0}^{T-1} \|E_{t,2}w_t\|^2\big] \nonumber\\
    &\quad +4\alpha\beta G^2 \mathbb{E}\big [\sum_{t=0}^{T-1} \|E_{t,3} \|_F^2+m\|{E}_{t,4}w_t\|^2 \big] + 4\alpha\beta \mathbb{E}\big[\sum_{t=0}^{T-1} \|E_{t,3}^{\top}E_{t,4}w_t\|^2 \big],
    \label{eq: intermediate descent inequality 2}
\end{align}
where we decompose error $E_{t,j}= \Gamma_{t,j}+A_{t,j}$ for $j\in\{2,3,4 \}$ and rewrite $E_{t,3}^{\top}E_{t,4}=(\Gamma_{t,3}+A_{t,3})^{\top}(\Gamma_{t,4}+A_{t,4})=\Gamma_{t,3}^{\top}(\Gamma_{t,4}+A_{t,4})+A_{t,3}^{\top}(\Gamma_{t,4}+A_{t,4})$, and utilizing the fact $\| w_0-w_{\tau}\|^2\leq 2$ holds for any $w_0, w_{\tau}$.

We first bound $-\alpha \mathbb{E}\big[\sum_{t=0}^{T-1}\langle \nabla \Phi(\theta_t)w, \Gamma_{t,2}w_t\rangle \big]$. 
Recall, $\calG_t$ denotes $\sigma$-algebra collecting outer loop randomness up to $t-1$ and inner-loop randomness at current iteration $t$. Conditioned on $\calG_t$, randomness of $\Gamma_{t,2} = Y_{t,2}-\nabla_{\theta}\fL(\theta_t,\eta_{t,d})$ only depends on $\xi_{t}$, we have
\begin{align}
    \mathbb{E}[\langle \nabla \Phi(\theta_t)w, \Gamma_{2}w_t\rangle|\calG_t] = \langle \nabla \Phi(\theta_t)w,  \mathbb{E}[\Gamma_{t,2}|\calG_t]w_t \rangle =0.
    \label{eq: inner product bound 1}
\end{align}
Taking expectation over $\calG_t$, tower property further implies $\mathbb{E}[\langle \nabla \Phi(\theta_t)w, \Gamma_{2}w_t\rangle]=0$. Thus we have $ -\alpha\mathbb{E}\big[\sum_{t=0}^{T-1}\langle \nabla \Phi(\theta_t)w, \Gamma_{t,2}w_t\rangle \big]=0$.
Similar arguments also apply on $\alpha\mathbb{E}\big[\sum_{t=0}^{T-1}\langle w -w_t, \Gamma_{t,3}^{\top}\nabla \Phi(\theta_t)w_t \rangle\big]$ and $\alpha\mathbb{E}\big[\sum_{t=0}^{T-1}\langle w -w_t, \Gamma_{t,4}^{\top}\nabla \Phi(\theta_t)w_t \rangle\big]$. 

For $\alpha\mathbb{E}[\sum_{t=0}^{T-1} \langle w-w_t, \Gamma_{t,3}^{\top} (\Gamma_{t,4}+A_{t,4})w_t \rangle] $, we decompose it as
\begin{align}
    &\alpha \mathbb{E}\big[\sum_{t=0}^{T-1} \langle w-w_t, \Gamma_{t,3}^{\top} (\Gamma_{t,4}+A_{t,4})w_t \rangle \big] \nonumber\\
    &\qquad = \alpha \mathbb{E}\big [\sum_{t=0}^{T-1}\langle w-w_t, \Gamma_{t,3}^{\top}\Gamma_{t,4}w_t\rangle \big] + \alpha \mathbb{E}\big[\sum_{t=0}^{T-1}\langle w-w_t, \Gamma_{t,3}^{\top} A_{t,4}w_t \rangle \big] \nonumber.
\end{align}
For $\alpha\sum_{t=0}^{T-1}\langle w-w_t, \Gamma_{t,3}^{\top}\Gamma_{t,4}w_t \rangle$, conditioned on $\calG_t$, $\Gamma_{t,3}$ and $\Gamma_{t,4}$ shares independent randomness. Thus, we have
\begin{align}
    \mathbb{E}[\langle w-w_t, \Gamma_{t,3}^{\top}\Gamma_{t,4}w_t \rangle|\calG_t] = \mathbb{E}[ \langle w-w_t, \mathbb{E}[\Gamma_{t,3}|\calG_t]^{\top}\mathbb{E}[\Gamma_{t,4}|\calG_t]w_t \rangle ] = 0.
    \label{eq: inner product bound 2}
\end{align}
Taking expectation over $\calG_t$, applying tower property implies $\mathbb{E}[\langle w-w_t, \Gamma_{t,3}^{\top}\Gamma_{t,4}w_t \rangle] = 0$. Thus $\alpha\mathbb{E}[\sum_{t=0}^{T-1}\langle w-w_t, \Gamma_{t,3}\Gamma_{t,4}w_t\rangle]=0$

Similarly, for $\alpha \sum_{t=0}^{T-1}\langle w-w_t, \Gamma_{t,3}^{\top} A_{t,4}w_t \rangle$, conditioned on $\calG_t$, $\Gamma_{t,3}$ randomness depends on sample $\bar{\xi}_t$, whereas $A_{t,4}$ randomness conditioned on $\calG_t$ is fixed. Thus, we have
\begin{align}
    \mathbb{E}[\langle w-w_t, \Gamma_{t,3}^{\top} A_{t,4}w_t \rangle|\calG_t]= \langle w-w_t, \mathbb{E}[\Gamma_{t,3}^{\top}|\calG_t] A_{t,4}w_t \rangle = 0.
    \label{eq: inner product bound 3}
\end{align}
Taking expectation over $\calG_t$, tower property further implies $\alpha \mathbb{E}[\sum_{t=0}^{T-1}\langle w-w_t, \Gamma_{t,3}^{\top} A_{t,4}w_t \rangle] = 0$.



For $\alpha \mathbb{E}\big[\sum_{t=0}^{T-1}\langle w-w_t, A_{t,3}^{\top} (\Gamma_{t,4}+A_{t,4}) w_t \rangle \big]$, we decompose it as follows
\begin{align}
    &\alpha \mathbb{E}[\sum_{t=0}^{T-1}\langle w-w_t, A_{t,3}^{\top} (\Gamma_{t,4}+A_{t,4}) w_t \rangle] \nonumber\\
    &\quad = \alpha \mathbb{E}\big[\sum_{t=0}^{T-1}\langle w-w_t, A_{t,3}^{\top} \Gamma_{t,4} w_t \rangle\big] + \alpha \mathbb{E}\big[\sum_{t=0}^{T-1}\langle w-w_t, A_{t,3}^{\top}A_{t,4}w_t \rangle \big]. \nonumber
\end{align}
Similarly, conditioned on $\calG_t$, we have 
\begin{align}
    \mathbb{E}[\langle w-w_t, A_{t,3}^{\top}\Gamma_{t,4} w_t \rangle |\calG_t] = \langle w-w_t, A_{t,3}^{\top}\mathbb{E}[\Gamma_{t,4}|\calG_t] w_t \rangle = 0.
    \label{eq: inner product bound 4}
\end{align}
Taking expectation over $\calG_t$, tower property implies $ \alpha \mathbb{E}[\sum_{t=0}^{T-1}\langle w-w_t, A_{t,3}^{\top} \Gamma_{t,4} w_t \rangle] = 0$.

For $\alpha \sum_{t=0}^{T-1}\langle w-w_t, A_{t,3}^{\top} A_{t,4} w_t \rangle$, conditioned on $\calH_t$, we have
\begin{align}
    \alpha\sum_{t=0}^{T-1}  \mathbb{E}[\langle w-w_t, A_{t,3}^{\top} A_{t,4} w_t \rangle|\calH_t] &= \sum_{t=0}^{T-1} \alpha \mathbb{E}[\langle A_{t,3}(w-w_t), A_{t,4} w_t \rangle|\calH_t] \nonumber\\
    %&\leq \alpha \sum_{t=0}^{T-1} \mathbb{E}[\|A_{t,3}(w-w_t)\||\calH_t]\mathbb{E}[\|A_{t,4}w_t\||\calH_t] \nonumber\\
    &\leq \alpha \sum_{t=0}^{T-1}\sqrt{\mathbb{E}[\|A_{t,3}(w-w_t)\|^2|\calH_t]}\sqrt{\mathbb{E}[\|A_{t,4}w_t\|^2|\calH_t]}\nonumber\\
    &\leq  2\alpha T \tepsilon^2,
    \label{eq: inner product bound 5}
\end{align}
where the first inequality utilizes conditional Cauchy–Schwarz inequality; the second inequality utilizes Jensen's inequality; the last inequality utilizes Corollary~\ref{Corollary: estimation bias of inner-loop}. Taking expectation over $\calH_t$, tower property of expectation gives $\alpha \mathbb{E}[\sum_{t=0}^{T-1}\langle w-w_t, A_{t,3}^{\top} A_{t,4} w_t \rangle]\leq 2\alpha T\tepsilon^2$.


For 
$$\alpha\mathbb{E}[\sum_{t=0}^{T-1}\langle w-w_t, A_{t,3}^{\top}\nabla \Phi(\theta_t)w_t  \rangle]+\alpha \mathbb{E}[\sum_{t=0}^{T-1}\langle w-w_t, \nabla \Phi(\theta_t)^{\top}A_{t,4}w_t \rangle] -\alpha \mathbb{E}\big[\sum_{t=0}^{T-1}\langle \nabla \Phi(\theta_t)w, A_{t,2}w_t \rangle \big],$$ conditioned on $\sigma$-algebra of outer iterations up to $\calH_t$, $A_{t,2}$, $A_{t,3}$ and $A_{t,4}$'s randomness is dependent with $\eta_{t,d}, \eta_{t, \bar{d}}, \eta_{t,\tilde{d}}$ respectively.  Notice that Conditional on $\calH_t$, the estimates $\eta_{t,d}$, $\eta_{t,\bar d}$, and $\eta_{t,\tilde d}$ are independent with each other and generated by the same inner-loop algorithm, using the same $\calH_t$-measurable initialization, learning rate, number of iterations. Each estimate is selected uniformly random. Consequently, these estimates have the same conditional distribution given $\calH_t$. Since $A_{t,2}$, $A_{t,3}$, and $A_{t,4}$ are obtained by applying the same bias mapping to their respective estimates, they also have the same conditional distribution and the same conditional mean. We therefore denote $\mathcal{A}_t=\mathbb{E}_{\eta_{t,d}}[A_{t,2}|\calH_t] = \mathbb{E}_{\eta_{t,\bar{d}}}[A_{t,3}|\calH_t] = \mathbb{E}_{\eta_{t,\tilde{d}}}[A_{t,4}|\calH_t]$. Thus, we further have
\begin{align}
    &\mathbb{E}\big[ \langle w-w_t, A_{t,3}^{\top}\nabla \Phi(\theta_t)w_t  \rangle+ \langle w-w_t, \nabla \Phi(\theta_t)^{\top}A_{t,4}w_t  \rangle \nonumber\\ 
    &\quad - \langle \nabla \Phi(\theta_t)w, A_{t,2}w_t \rangle |\calH_t\big]\nonumber\\
    & = \langle w-w_t, \mathbb{E}[A_{t,3}^{\top}|\calH_t]\nabla \Phi(\theta_t)w_t  \rangle + \langle w-w_t, \nabla \Phi(\theta_t)^{\top} \mathbb{E}[A_{t,4}|\calH_t]w_t \rangle \nonumber \\
    &\quad -\langle\nabla \Phi(\theta_t)w,\mathbb{E}[A_{t,2}|\calH_t]w_t \rangle \nonumber\\
    & = \langle w-w_t, \mathcal{A}_t^{\top}\nabla \Phi(\theta_t)w_t  \rangle + \langle w-w_t, \nabla \Phi(\theta_t)^{\top} \mathcal{A}_tw_t \rangle -\langle\nabla \Phi(\theta_t)w,\mathcal{A}_tw_t \rangle \nonumber\\
    & = \langle \mathcal{A}_t(w-2w_t), \nabla \Phi(\theta_t)w_t \rangle\nonumber\\
    & \leq \frac{1}{4}\|\nabla \Phi(\theta_t)w_t \|^2 + \|\mathcal{A}_t(w-2w_t)\|^2\nonumber\\
    & \leq  \frac{1}{4}\|\nabla \Phi(\theta_t)w_t \|^2+(\|\mathcal{A}_tw\| + 2\|\mathcal{A}_tw_t\|)^2\nonumber\\
    & \leq \frac{1}{4}\|\nabla \Phi(\theta_t)w_t \|^2 + 9\tepsilon^2,
    \label{eq: key insights of algorithm 1}
\end{align}
where the first inequality applies young's inequality, $\langle a,b\rangle\leq \frac{1}{4}\| a\|^2+\|b\|^2$; the second inequality utilizes triangle-inequality $\|a-b\|\leq \|a\|+\|b\|$; the last inequality applies Corollary~\ref{Corollary: estimation bias of inner-loop}. This further implies
\begin{align}
    &\alpha\mathbb{E}[\sum_{t=0}^{T-1}\langle w-w_t, A_{t,3}^{\top}\nabla \Phi(\theta_t)w_t  \rangle]+\alpha \mathbb{E}[\sum_{t=0}^{T-1}\langle w-w_t, \nabla \Phi(\theta_t)^{\top}A_{t,4}w_t \rangle] -\alpha \mathbb{E}\big[\sum_{t=0}^{T-1}\langle \nabla \Phi(\theta_t)w, A_{t,2}w_t \rangle \big] \nonumber\\
    &\leq \frac{\alpha}{4}\sum_{t=0}^{T-1}\mathbb{E}[\|\nabla \Phi(\theta_t)w_t\|^2] + 9\alpha T \tepsilon^2.
    \label{eq: inner product bound 6}
\end{align}
% Loose bound, no longer used
%\begin{align}
%    \alpha\mathbb{E}[\sum_{t=0}^{T-1}\langle w-w_t, A_{t,3}^{\top}\nabla \Phi(\theta_t)w_t  \rangle |\calH_t] &=\alpha \sum_{t=0}^{T-1}\mathbb{E}[\langle A_{t,3}(w-w_t), \nabla \Phi(\theta_t)w_t  \rangle|\calH_t]\nonumber\\
%    &\leq \alpha \sum_{t=0}^{T-1}\mathbb{E}[\|A_{t,3}(w-w_t) \|\|\nabla \Phi(\theta_t)w_t \||\calH_t]\nonumber\\
%    &\leq \alpha G \sum_{t=0}^{T-1}\mathbb{E}[\|A_{t,3}(w-w_t) \||\calG_t] \\
%    &\leq 2\alpha T G \tepsilon,
%    \label{eq: inner product bound 6}
%\end{align}
%where the first inequality utilizes Cauchy-Schwarz inequality; the second inequality utilizes $\|\nabla \phi^i(\theta)\|\leq G$ the last inequality is due to $\|A_{t,3}(w-w_t)\|\leq \|A_{t,3}w\|+\|A_{t,3}w_t\|$ and Corollary~\ref{Corollary: estimation bias of inner-loop}. Taking expectation over $\calH_t$ and tower property implies $\alpha\mathbb{E}[\sum_{t=0}^{T-1}\langle w-w_t, A_{t,3}^{\top}\nabla \Phi(\theta_t)w_t  \rangle ] \leq 2\alpha TG\tepsilon$. Same bound also applies to $\alpha \mathbb{E}[\sum_{t=0}^{\tau-1}\langle w-w_t, \nabla \Phi(\theta_t)^{\top}A_{t,4}w_t \rangle]$.
%Leveraging similar arguments as ~\eqref{eq: inner product bound 6}, we conclude
%\begin{align}
%    -\alpha \mathbb{E}\big[\sum_{t=0}^{\tau-1}\langle \nabla \Phi(\theta_t)w, A_{t,2}w_t \rangle \big]
%    &\leq \alpha TG\tepsilon,
%    \label{eq: inner product bound 7}
%\end{align}
Since $w_t$ is $\calH_t$-measurable, $E_{t, 3}, E_{t, 4}$ are conditionally independent given $\mathcal{H}_t$, thus we have
\begin{align}
\mathbb{E}[\|E_{t, 3}\|_F^2\|E_{t, 4} w_t\|^2]  &=\mathbb{E}[\mathbb{E}[\|E_{t, 3}\|_F^2\|E_{t, 4} w_t\|^2 \mid \mathcal{H}_t]] \nonumber\\
& =\mathbb{E}[\mathbb{E}[\|E_{t, 3}\|_F^2 \mid \mathcal{H}_t] \mathbb{E}[\left\|E_{t, 4} w_t\right\|^2 \mid \mathcal{H}_t]] \leq m C_B^4.
\label{eq: quadratic bound 7}
\end{align}
Combining \eqref{eq: inner product bound 1},\eqref{eq: inner product bound 2},\eqref{eq: inner product bound 3},\eqref{eq: inner product bound 4},\eqref{eq: inner product bound 5},\eqref{eq: inner product bound 6},\eqref{eq: quadratic bound 7} we further have
\begin{align}
    &\mathbb{E}[\Phi(\theta_{T}){w}] - \mathbb{E}[\Phi({\theta}_0){w}] \nonumber\\
    &\leq -\frac{\alpha}{4} \mathbb{E}\big[\sum_{t=0}^{T-1}\|\nabla \Phi({\theta}_t){w}_t \|^2\big] +{11\alpha T \tepsilon^2} +\frac{\alpha}{\beta}+ 2\alpha \rho T + \alpha\beta\rho^2T + 4\alpha\beta mG^4 T\nonumber\\
    &+\alpha^2 L_0 \mathbb{E}\big[ \sum_{t=0}^{T-1} \|E_{t,2}w_t\|^2\big] +4\alpha\beta G^2 \mathbb{E}\big [\sum_{t=0}^{T-1} \|E_{t,3} \|_F^2+m\|{E}_{t,4}w_t\|_2^2 \big] + 4\alpha\beta \mathbb{E}\big[\sum_{t=0}^{T-1} \|E_{t,3}^{\top}E_{t,4}w_t\|^2 \big] \nonumber\\
    &\leq -\frac{\alpha}{4} \mathbb{E}\big[\sum_{t=0}^{T-1}\|\nabla \Phi({\theta}_t){w}_t \|^2\big] +{11\alpha T \tepsilon^2} +\frac{\alpha}{\beta}+ 2\alpha \rho T + \alpha\beta\rho^2T + 4\alpha\beta m G^4 T\nonumber\\
    &+\alpha^2 L_0 \mathbb{E}\big[ \sum_{t=0}^{T-1} \|E_{t,2}w_t\|^2\big] +4\alpha\beta  G^2 \mathbb{E}\big [\sum_{t=0}^{T-1} \|E_{t,3} \|_F^2+m\|E_{t,4}w_t\|_2^2 \big] + 4\alpha\beta \mathbb{E}\big[\sum_{t=0}^{T-1} \|E_{t,3}\|^2_F \|E_{t,4}w_t\|^2 \big] \nonumber\\
    &\leq -\frac{\alpha}{4} \mathbb{E}\big[\sum_{t=0}^{T-1}\|\nabla \Phi({\theta}_t){w}_t \|^2\big] +\frac{\alpha}{\beta}+{11\alpha T \tepsilon^2} + 2\alpha \rho T + \alpha\beta\rho^2T + 4\alpha\beta m G^4 T\nonumber\\
    &+\alpha^2 L_0 C_B^2T+8\alpha\beta m G^2C_B^2T + 4\alpha\beta mT C_B^4 
    %&\leq -\frac{\alpha}{2} \mathbb{E}\big[\sum_{t=0}^{T-1}\|\nabla \Phi({\theta}_t){w}_t \|^2\big] +\frac{\alpha}{\beta}+{5\alpha T\delta\epsilon^2+2\alpha T \delta^2\epsilon^4} \nonumber\\
    %&+ 2\alpha \rho T + \alpha\beta\rho^2T + 4\alpha\beta m \Lambda^4 T+\alpha^2 L_0 C_B^2T+8\alpha\beta m G^2C_B^2T + 4\alpha\beta mT C_B^4,
\end{align}
where the second and third inequality applies sub-multiplicative property of Frobenius-norm and Corollary~\ref{Corollary: estimation error of true gradient}. 
%{the last inequality replace scaled accuracy $\tepsilon$ by $\tepsilon =\min\{\Lambda^{-1},1\}\delta\epsilon^2$, which further implies $\Lambda\tepsilon\leq \delta\epsilon^2$ and $\tepsilon^2\leq \delta^2\epsilon^4$}. 
Re-organize above inequality, divide both sides with $\frac{4}{\alpha T}$, we have
\begin{align}
    &\frac{1}{T}\sum_{t=0}^{T-1}\mathbb{E}[\|\nabla \Phi(\theta_t)w_t \|^2] \nonumber\\
    &\leq \frac{4}{\alpha}\cdot \frac{2}{T}\Big(\mathbb{E}[\Phi(\theta_0)w-\Phi^*w]+\frac{\alpha}{\beta}+{11\alpha T \tepsilon^2}+{2\alpha \rho T}+ \alpha\beta\rho^2T+4\alpha\beta m G^4 T\nonumber\\
    &\quad+ {\alpha^2 L_0 C_B^2T}+{8\alpha\beta mG^2C_B^2T + 4\alpha\beta m T C_B^4}\Big) \nonumber\\
    &= \frac{8(\mathbb{E}[\Phi(\theta_0)w-\Phi^*w])}{\alpha T}+\frac{8}{\beta T}+88\tepsilon^2
    +16\rho+8\beta\rho^2+32\beta m G^4+8\alpha L_0C_B^2+64\beta m G^2C_B^2\nonumber\\
    &\quad +32\beta m C_B^4\nonumber\\
    &\leq 94\epsilon^2,
\end{align}
where the last inequality is due to our parameter choices.
%% Outdated bound, no-longer use
%\begin{align}
%    \mathbb{E}[\Phi(\theta_{\tau}){w}] - \Phi^*w&\leq \mathbb{E}[\Phi(\theta_0)w]-\Phi^*w-\frac{\alpha}{2} \mathbb{E}\big[\sum_{t=0}^{\tau-1}\|\nabla \Phi({\theta}_t){w}_t \|^2\big] + \underbrace{\frac{\alpha}{\beta}}_{\leq c_1} \nonumber \\
%    &+\underbrace{{5\alpha T \delta\epsilon^2+3\alpha T\delta^2\epsilon^4}}_{\leq c_2}+\underbrace{{5\alpha \Lambda \sigma+3\alpha\delta^2\epsilon^4+6\alpha\sigma^2}}_{\leq c_3} \nonumber\\
%    &+ \underbrace{2\alpha \rho T}_{\leq c_4} + \underbrace{\alpha\beta\rho^2T}_{\leq c_5} + \underbrace{\alpha^2 L_0 C_B^2T}_{\leq c_6}+ \underbrace{4\alpha\beta m \Lambda^4 T+8\alpha\beta m\Lambda^2C_B^2T + 4\alpha\beta m T C_B^4}_{\leq c_7}\nonumber\\
%    &\leq \frac{F\delta}{8}- \frac{\alpha}{2}\mathbb{E}[\sum_{t=0}^{\tau-1}\| \nabla \Phi(\theta_t)w_t\|^2],
%    \label{eq: MOO descent up to stopping time}
%\end{align}
%which completes the proof. 
\end{proof}
%Given problem dependent parameter $m$,$G$, $L_0=G^2M\lambda^{-1}+L$, $L_2=M\lambda^{-1}$, $K_0=8G^2+10G^2M^2\lambda^{-2}\kappa^2$, $K_1=8G^2$, $K_2=M^2\lambda^{-2}\kappa^2$, denote constant $C_0, C_1\geq 0$, $F$, $c_1...c_7\geq 0$ be some finite constant such that
%\begin{align}
%    \frac{F}{8}\geq \frac{\Delta_{\theta} + c_1+c_2+\dots+c_7}{\delta},
%\end{align}
%where $\Delta_{\theta_0} = \max_{i\in[m]}\{\phi^i(\theta_0)-\phi^{i,*}\}$. Define $0<\delta<1$ are pre-chosen tail-event probability upper bound.
%set the hyper-parameters in Algorithm \ref{alg1} as follows
%\begin{align}
%    \sigma^2 &= {2K_0}+{2K_1G^{-1}}\min\{1,G^{-2}\}\delta^2\epsilon^4 =\mathcal{O}(K_0)=\mathcal{O}(G^2M^2\lambda^{-2}\kappa^2)\nonumber\\
%    C_B&=\sigma^2+2\min\{1,G^{-2}\}\delta^2\epsilon^4 = \mathcal{O}(G^2M^2\lambda^{-2}\kappa^2)\nonumber\\
%    \gamma&= \min\{\frac{\delta^2\epsilon^4}{2L_2K_2G^2\Lambda^2},\frac{\delta^2\epsilon^4}{2L_2K_2G^2},\frac{1}{L_2} \}=\mathcal{O}(\frac{\lambda^3\delta^2\epsilon^4}{M^3G^2\Lambda^2\kappa^2})
%    \nonumber\\
%    \rho&= \frac{1}{8}\delta^2\epsilon^2=\mathcal{O}(\delta^2\epsilon^2) \nonumber\\%\min\Big\{\frac{C_0^2\delta^2\epsilon^2}{96mC_B^2\Delta_{\theta_0}},\frac{C_1^2\delta^2\epsilon^2}{32m^2C_B^4\Delta_{\theta_0}} \Big\}=\mathcal{O}(\delta^2\epsilon^2)\Big\} \nonumber\\
%    \alpha&= \min\Big\{\frac{1}{2L_0},c_1\beta ,\frac{c_2}{(5\delta\epsilon^2+3\delta^2\epsilon^4)T},\frac{c_3}{5\Lambda\delta+3\delta^2\epsilon^4+6\sigma^2},\frac{c_4}{2\rho T},\frac{c_5}{\beta\rho^2T},\sqrt{\frac{c_6}{L_0C_B^2T}},\frac{\delta\epsilon^2}{4L_0C_B^2},\nonumber\\
%    &\quad \frac{c_7}{(4m\Lambda^4+8m\Lambda^2C_B^2+4mC_B^4)\beta T},
%    \min\big\{\frac{b_1^2}{(C_0\Lambda+2\sqrt{m}C_1+4\Lambda\sqrt{m}C_0)^2}, \frac{b_2}{C_0^2L_0}\big\}\rho \nonumber\\
%    &\quad \frac{C_0\delta}{24mC_B^2\rho T},\frac{C_1^2\delta}{8m^2C_B^4\rho T}\Big\} \nonumber\\
%    &= \min\Big\{ \mathcal{O}(\beta),\mathcal{O}(\frac{1}{\delta\epsilon^2T}),\mathcal{O}(\frac{1}{(m\Lambda^4+mG^8M^8\lambda^{-8}\kappa^8)\beta T}),
%    \mathcal{O}(\frac{\rho}{\sqrt{m}\Lambda}),\mathcal{O}(\frac{\lambda\rho}{G^2M}),\nonumber\\
%    &\mathcal{O}(\frac{\lambda^8\delta}{m^2G^8M^8\kappa^8\rho T})\Big\},\nonumber\\
%    \beta&= \min \Big\{\frac{\delta\epsilon^2}{16(m\Lambda^4+2m\Lambda^2C_B^2+mC_B^4)},\frac{\delta\epsilon^2}{4\rho}, \frac{b_3}{4mC_1^2+8mC_0^2\Lambda^2}\rho\Big\}= \mathcal{O}(\frac{\rho}{m\Lambda^2})
%    \nonumber\\
%D&= 4\gamma^{-1}\Delta_{\eta}G^2 \max\{1,\Lambda^{2}\}\delta^{-2}\lambda^{-4}=\Theta(\gamma^{-1}\Delta_{\eta}G^2\Lambda^2\delta^{-2}\epsilon^{-4})\nonumber\\
%T&=\max\Big\{(20\Lambda \sigma+24\sigma^2+12\delta^2\epsilon^4)\delta^{-1}\epsilon^{-2}, {4}\delta^{-1}\epsilon^{-2}\beta^{-1}, 4\Delta_{\theta_0}\delta^{-1}\epsilon^{-2}\alpha^{-1}\Big\}\nonumber\\
%&=\max\Big\{\Theta(\Delta_{\theta_0}\alpha^{-1}\delta^{-1}\epsilon^{-2}), \Theta(\beta^{-1}\delta^{-1}\epsilon^{-2})\Big\}=\Theta(\delta^{-3}\epsilon^{-4}),
%\end{align}
%where ${\Delta}_{\eta}=\max_{t\in T, i\in[m]}\big \{\fL^i({\theta}_t,\eta^i_{0})-\fL^i({\theta_t},\eta^{i*})\}$
%Denote $b_1,\dots b_3\geq 0$ be some constant satisfying
%\begin{align}
%    \frac{F}{2}\geq     b_1+b_2+b_3+c_1+c_4+c_5+c_7 .\nonumber
%\end{align}
%Then, we then have the following convergence statement on Double-Loop Algorithm \ref{alg1}.
%\begin{restatable}[Formal Statement of Theorem~\ref{thm: convergence of algorithm 1}]{theorem}{convegencefirstalg}\label{restatethm: convergence of algorithm 1}
%    Let Assumption \ref{assumption 1} hold. Denote $\Delta_{\theta_0} = \max_{i\in[m]}\{\phi^i(\theta_0)-\phi^{i,*}\}$ and ${\Delta}_{\eta}=\max_{t\in T, i\in[m]}\big \{\fL^i({\theta}_t,\eta^i_{0})-\fL^i({\theta_t},\eta^{i,*})\}$. Given $0<\epsilon,\delta<1$, 
%    set $\rho=\mathcal{O}(\delta^2 \epsilon^2)$, $\beta =\mathcal{O}(\frac{\rho}{m\Lambda^2})$,
%    $\alpha = \min\big\{ \mathcal{O}(\beta),\mathcal{O}(\frac{1}{\delta\epsilon^2T}),\mathcal{O}(\frac{1}{(m\Lambda^4+mG^8M^8\lambda^{-8}\kappa^8)\beta T}),
%    \mathcal{O}(\frac{\rho}{\sqrt{m}\Lambda}),\mathcal{O}(\frac{\lambda\rho}{G^2M}),\mathcal{O}(\frac{\lambda^8\delta}{m^2G^8M^8\kappa^8\rho T})\big\}$, and $\gamma = \mathcal{O}(\frac{\lambda^3\delta^2\epsilon^4}{M^3G^2\kappa^2})$ for Algorithm \ref{alg1}. Then, after $T=\max\big\{\Theta(\Delta_{\theta_0}\alpha^{-1}\delta^{-1}\epsilon^{-2}), \Theta(\beta^{-1}\delta^{-1}\epsilon^{-2})\big\}$ outer iterations, each associated with $D=\Theta(\gamma^{-1}\Delta_{\eta}G^2\Lambda^2\delta^{-2}\epsilon^{-4})$ inner-iterations, we have 
%    \begin{align}
%        \frac{1}{T}\sum_{t=0}^{T-1}\|\nabla \Phi(\theta_t)w_t\|^2\leq 78\epsilon^2,
%    \end{align}
%    holds with probability at least $1-\delta$.
%\end{restatable}
%\begin{proof}
%\textbf{Part I: Tail Event $\{\tau<T\}$}.
%Define stopping time $\tau_1, \tau_2, \tau_3$ and $\tau$ as follows
%\begin{align}
%    \tau_1 &=\min\{t|\exists i \in[m],\phi^i(\theta_{t+1})-\phi^{i,*}>F \}\wedge T\nonumber\\
%    \tau_2 &=\min\{t|\exists i\in[m], j\in[2,3,4], \|E_{t,j}^i \|\geq \frac{C_0}{\sqrt{\alpha\rho}}\}\wedge T\nonumber\\
%    \tau_3 &={\min\{t|\exists i,k\in[m], \|E_{t,3}^i \|\| E_{t,4}^k\|\geq \frac{C_1}{\sqrt{\alpha\rho}}\} \wedge T } \nonumber\\ 
%    \tau &= \min\{\tau_1, \tau_2, \tau_3\}, \nonumber
%\end{align}
%where $C_0, C_1>0$ represent placeholder constants. We show that $\tau <T$ is a tail event, namely $P(\tau<T)\leq \delta$.

%For event $\tau<T$, it is equivalent to the following events
%\begin{align}
%    \{\tau<T \}=\{\tau_2 <T\} \cup \{\tau_3<T \}\cup \{\tau_1<T \}. \nonumber
%\end{align}
%For any $i\in[m]$, from Markov's inequality, we have
%\begin{align}
%    P(\|E_{t,j}^i\|\geq \frac{C_0}{\sqrt{\alpha\rho}})=P(\|E_{t,j}^i\|^2\geq \frac{C^2_0}{{\alpha\rho}})\leq \frac{\alpha \rho\mathbb{E}\|E_{t,j}^i \|^2}{C_0^2}, \forall j \in \{2,3,4\} \nonumber
%\end{align}
%Then, for event $\{\tau=\tau_2<T\}$, its probability can be upper bounded as
%\begin{align}
%    P(\tau_2<T)&\leq \sum_{t=0}^{T-1}\sum_{i=1}^{m}\sum_{j=2}^{4} P(\|E_{t,j}^i\|\geq \frac{C_0}{\sqrt{\alpha\rho}})\leq \sum_{t=0}^{T-1}\sum_{i=1}^{m}\sum_{j=2}^{4}\frac{\alpha\rho\mathbb{E}\|{E}_{t,j}^i\|^2}{C_0^2} \nonumber\\
%    &\leq \frac{3\alpha\rho mC_B^2T}{C_0^2}\leq  \frac{\delta}{8},
%\end{align}
%where the last two inequalities are due to the fact $\mathbb{E}[\|E_{t,j}\|_F^2]=\mathbb{E}[\sum_{i=1}^m \|E_{t,j}^i\|^2]\leq mC_B^2$ and ${\alpha\leq \frac{C_0^2\delta}{24C_B^2 m \rho T}}$.


%Similarly, for any $i\in[m]$, by Markov's inequality, we have
%\begin{align}
%    P(\|E_{t,3}^i\|\| E_{t,4}^k\|\geq \frac{C_1}{\sqrt{\alpha\rho}}) &= P(\|E_{t,3}^i\|^2\| E_{t,4}^k\|^2\geq \frac{C_1^2}{{\alpha\rho}}) \leq \frac{{\alpha\rho}\mathbb{E}[\|E_{t,3}^i\|^2\| E_{t,4}^k\|^2]}{C_1^2}. \nonumber
%\end{align}

%Then, for any $i\in[m]$, if the event $\{\tau=\tau_3<T\}$ happens, its probability can be upper bounded by
%\begin{align}
%P(\tau_3<T)&\leq\sum_{t=0}^{T-1}\sum_{i=1}^{m}\sum_{k=1}^{m}P(\|E_{t,3}^i\|\| E_{t,4}^k\|\geq \frac{C_1}{\sqrt{\alpha\rho}})\leq \sum_{t=0}^{T-1}\sum_{i=1}^m\sum_{k=1}^{m}\frac{\alpha\rho\mathbb{E}[\|E_{t,3}^i\|^2\|E_{t,4}^j \|^2]}{C_1^2} \nonumber\\
%    & =\sum_{t=0}^{T-1}\sum_{i=1}^m\sum_{k=1}^{m}\frac{\alpha\rho\mathbb{E}\big[\mathbb{E}_{\eta_{t,\bar{d}},\bar{\xi}_{t}}[\|E_{t,3}^i\|^2|\theta_t]\mathbb{E}_{\eta_{t,\tilde{d}},\tilde{\xi}_t}[\|E_{t,4}^j \|^2|\theta_t]\big]}{C_1^2}\nonumber\\
%    &\leq \frac{m^2\alpha\rho T C_B^4}{C_1^2}\leq \frac{\delta}{8},\nonumber
%\end{align}
%where the last two inequality applies the fact $E_{t,3},E_{t,4}$ shares independent randomness against each other and $\mathbb{E}_{\eta_{t,\bar{d}},\bar{\xi}_t}[\|E_{t,3}^i\|^2],\mathbb{E}_{\eta_{t,\tilde{d}},\tilde{\xi}_t}[\|E_{t,4}^i\|^2]\leq {C_B^2}$, and ${\alpha\leq \frac{C_1^2\delta}{8 m^2 \rho TC_B^4}}$.

%For event $\{\tau_1<T \}$, we know at $\tau+1$, there exists at least one $ i\in[m]$ such that $\phi^i(\theta_{\tau+1})-\phi^{i,*}\geq F$. Since $\tau_2, \tau_3\geq \tau_1$, we know for all $i\in[m], t\leq T$, $\|E_{t,j}^i\|\leq \frac{C_0}{\sqrt{\alpha\rho}}$ and $\|E_{t,3}^i\|\|E_{t,4}^k\|\leq \frac{C_1}{\sqrt{\alpha\rho}}$. Then, we have
%\begin{align}
%    \|E_{t,j}w_t\|&\leq \sum_{i=1}^{m}w_t^i\|E^i_{t,j}\|\leq 
%    \frac{C_0}{\sqrt{\alpha\rho}}, \forall j\in \{2,3,4 \},\label{eq: concentration bound 1}\\
%    \|E_{t,j}\|_F&\leq \frac{\sqrt{m}C_0}{\sqrt{\alpha\rho}}, \forall j\in\{2,3,4 \},\label{eq: concentration bound 2}\\
%    \|E_{t,3}^{\top}E_{t,4}w_t\|&\leq \sum_{k=1}^{m}w_t^k \sqrt{\sum_{i=1}^{m} \langle E_{t,3}^{i}, E_{t,4}^{k}\rangle^2 } \leq \sum_{k=1}^{m}w_t^i\sqrt{\sum_{i=1}^{m} \| E_{t,3}^i\|^2 \|E_{t,4}^k \|^2}\leq \frac{\sqrt{m}C_1}{\sqrt{\alpha\rho}}.
%    \label{eq: concentration bound 3} 
%\end{align}
 

%Let $w_i=1, w_{m/\{i\}}=0$, at time $t=\tau$, \eqref{eq: one step descent inequality} reduces to 
%\begin{align}
%    &\phi^i(\theta_{\tau+1}) - \phi^i(\theta_\tau)\nonumber \\
%    &\leq -\alpha \|\nabla \Phi(\theta_\tau)w_\tau \|^2 -{\alpha \langle \nabla\phi^i(\theta_\tau),({\Y_\tau - \nabla_{\theta} \fL(\theta_t, \eta_{t,d})}) {w}_\tau\rangle - \alpha \langle \nabla \phi^i(\theta_\tau),\nabla_{\theta} \fL(\theta_t, \eta_{t,d}) -\Phi(\theta_\tau) \rangle} \nonumber\\
%    & \quad + \alpha^2L_0\|\nabla \Phi({\theta}_\tau){w}_\tau \|^2 + \alpha^2L_0\|E_{\tau,2}{w}_\tau\|^2 \nonumber\\
%    &\quad +\frac{\alpha}{2\beta}\Big(\big(\|{w}_\tau -{w}\|^2-\|{w}_{\tau+1} -{w} \|^2 \big)+4\beta \rho + 2\beta^2\rho^2+8\beta^2 \|E_3^{\top} E_4{w}_\tau\|^2\nonumber\\
%    &\quad -2\beta \langle {w}_\tau-{w}, E_{\tau,3}^{\top}E_{\tau,4}{w}_\tau+E_{\tau,4}^{\top}\nabla \Phi({\theta}_t)w_\tau+E_{\tau,3}^{\top}\nabla \Phi({\theta}_\tau){w}_\tau\rangle\nonumber\\
%    & \quad + 8\beta^2 m\Lambda^4 + 8\beta^2 \Lambda^2(\|E_{\tau,3} \|_F^2 +\|E_{\tau,4} \|_F^2 ) \Big) \nonumber\\
%    &{=} -\alpha \|\nabla \Phi(\theta_\tau)w_\tau \|^2 -{\alpha \langle \nabla\phi^i(\theta_\tau),{E_{\tau,2}} {w}_\tau\rangle} + \alpha^2L_0\|\nabla \Phi({\theta}_\tau){w}_\tau \|^2 + \alpha^2L_0\|E_{\tau,2}{w}_\tau\|^2 \nonumber\\
%    &\quad +\frac{\alpha}{2\beta}\Big(\big(\|{w}_\tau -{w}\|^2-\|{w}_{\tau+1} -{w} \|^2 \big)+4\beta \rho + 2\beta^2\rho^2+8\beta^2 \|E_{\tau,3}^{\top} E_{\tau,4}{w}_\tau\|^2\nonumber\\
%    &\quad -2\beta \langle {w}_\tau-{w}, E_{\tau,3}^{\top}E_{\tau,4}{w}_\tau+E_{\tau,4}^{\top}\nabla \Phi({\theta}_t)w_\tau+E_{\tau,3}^{\top}\nabla \Phi({\theta}_\tau){w}_\tau\rangle\nonumber\\
%    &\quad + 8\beta^2 m\Lambda^4 + 8\beta^2 \Lambda^2(\|E_{\tau,3} \|_F^2 +\|E_{\tau,4} \|_F^2 ) \Big) \nonumber\\
%    &\overset{(i)}{\leq} -\frac{\alpha}{2} \|\nabla \Phi(\theta_\tau) w_\tau \|^2 + \alpha \| \nabla \phi^i(\theta_\tau)\| \|E_{\tau,2}w_\tau\| +\alpha^2 L_0\|E_{\tau,2}w_\tau \|^2+\frac{\alpha}{\beta} + 2\alpha\rho+\alpha\beta\rho^2 \nonumber\\
%    & \quad +4\alpha\beta\|E_{\tau,3}^{\top}E_{\tau,4}w_\tau \|^2 +2\alpha\| E_{\tau,3}^{\top}E_{\tau,4}w_\tau \|+2\alpha \|E_{\tau,4}^{\top}\nabla\Phi(\theta_t)w_\tau \|+ 2\alpha \|E_{\tau,3}^{\top}\nabla\Phi(\theta_\tau)w_\tau \|\nonumber\\
%    &\quad +4\alpha\beta m \Lambda^4+4\alpha\beta  \Lambda^2(\|E_{\tau,3}\|_F^2+\|E_{\tau,4}\|_F^2) \nonumber\\
%    &\overset{(ii)}{\leq} \alpha \Lambda \|E_{\tau,2}w_\tau \|+\alpha^2L_0\| E_{\tau,2}w_\tau\|^2 +\frac{\alpha}{\beta}+2\alpha\rho+\alpha\beta\rho^2+4\alpha\beta \|E_{\tau,3}^{\top}E_{\tau,4}w_\tau\|^2+2\alpha \|E_{\tau,3}^{\top} E_{\tau,4}w_\tau \|\nonumber\\
%    &\quad +2\alpha\Lambda \|E_{\tau,4} \|_F+2\alpha\Lambda\|E_{\tau,3} \|_F+4\alpha\beta m\Lambda^4+4\alpha\beta \Lambda^2(\|E_{\tau,3}\|_F^2+\|E_{\tau,4}\|_F^2)\nonumber\\
%    &\overset{(iii)}{\leq} \alpha\Lambda \frac{C_0}{\sqrt{\alpha\rho}}+\alpha^2L_0\frac{C_0^2}{\alpha\rho}+\frac{\alpha}{\beta}+2\alpha\rho+\alpha\beta\rho^2+4\alpha\beta \frac{{m}C_1^2}{\alpha\rho}+2\alpha\frac{\sqrt{m}C_1}{\sqrt{\alpha\rho}}\nonumber\\
%    &\quad +2\alpha\Lambda \frac{\sqrt{m}C_0}{\sqrt{\alpha\rho}}+2\alpha\Lambda \frac{\sqrt{m}C_0}{\sqrt{\alpha\rho}}+4\alpha\beta m\Lambda^4+8\alpha\beta\Lambda^2(\frac{mC_0^2}{\alpha\rho})\nonumber\\
%    &=\frac{C_0\Lambda\sqrt{\alpha}}{\sqrt{\rho}}+\frac{\alpha C_0^2L_0}{\rho}+\frac{\alpha}{\beta}+2\alpha\rho+\alpha\beta\rho^2+\frac{4mC_1^2\beta}{\rho}+\frac{2\sqrt{m\alpha}C_1}{\sqrt{\rho}}+\frac{2\Lambda\sqrt{m\alpha}C_0}{\sqrt{\rho}}+\frac{2\sqrt{\alpha m}C_0\Lambda}{\sqrt{\rho}}\nonumber\\
%    &\quad +4\alpha\beta m\Lambda^4+\frac{8mC_0^2\beta\Lambda^2}{\rho},
%    \label{eq: stopping time descent}
%\end{align}
%where $(i)$ applies the fact ${\alpha\leq \frac{1}{2L_0}}$, $\|w_{\tau}-w\|\leq 2$ and Cauchy-Schwarz inequality, and (ii) applies the fact at time $\tau$, $\|\nabla\phi^i(\theta_t)\|\leq \Lambda$ still holds. (iii) applies the fact \eqref{eq: concentration bound 1}, \eqref{eq: concentration bound 2} and \eqref{eq: concentration bound 3} to upper bound $\|E_{\tau,2}w_{\tau}\|$, $\|E_{\tau,2}w_{\tau} \|^2$, $\|E_{\tau,3}^{\top}E_{\tau,4}w_{\tau}\|$, and $\|E_{\tau,3}^{\top}E_{\tau,4}w_{\tau}\|^2$, respectively.



%Since, for $\frac{\alpha}{\beta}, 2\alpha\rho$, $\alpha\beta\rho^2$, $4\alpha\beta m \Lambda^4$, and $0<\delta<1$, we have
%\begin{align}
%    \frac{\alpha}{\beta}\leq c_1, 2\alpha\rho<c_4, \alpha\beta\rho^2\leq c_5, 4\alpha\beta m \Lambda^4 <c_7 \nonumber
%\end{align}
%Thus, for $b_1>0, b_2>0, b_3>0$ and $\frac{F}{2}$, we have
%$b_1+b_2+b_3+\frac{\alpha}{\beta}+2\alpha\rho+\alpha\beta\rho^2+4\alpha\beta m \Lambda^4\leq \frac{F}{2}$
%By setting
%\begin{align}
%    {\frac{\alpha}{\rho}}\leq \min\big\{\frac{b_1^2}{(C_0\Lambda+2\sqrt{m}C_1+4\Lambda\sqrt{m}C_0)^2}, \frac{b_2}{C_0^2L_0}\big\},\quad
%    {\frac{\beta}{\rho}}\leq \frac{b_3}{4m C_1^2+8m C_0^2\Lambda^2},\nonumber
%\end{align}
%\eqref{eq: stopping time descent} reduces to
%\begin{align}
%    \phi^i(\theta_{\tau+1})-\phi^{i}(\theta_{\tau})\leq b_1+b_2+b_3+\frac{\alpha}{\beta}+2\alpha\rho+2\alpha\beta\rho^2+4\alpha\beta m \Lambda^4\leq \frac{F}{2}.
%    \label{eq: stopping time descent 1} 
%\end{align}
%However, after stopping time $\tau$, we know this specific $\tilde{i}$ satisfying $\phi^{\tilde{i}}(\theta_{\tau+1})-\phi^{\tilde{i},*}>F$, for these specific $\tilde{i}$, combing \eqref{eq: stopping time descent 1}, we know 
%\begin{align}
%    \phi^{\tilde{i}}(\theta_{\tau})-\phi^{\tilde{i},*}>\frac{F}{2}.\nonumber
%\end{align}
%From \eqref{eq: MOO descent up to stopping time}, we know $\mathbb{E}[\phi^i(\theta_\tau)-\phi^{i,*}]\leq \frac{\delta F}{8}$.
%Using Markov inequality, we upper bound the probability as 
%$P(\phi^{i}(\theta_\tau)-\phi^{i,*}\geq \frac{F}{2})\leq \frac{\delta}{4}$. This inequality further suggests, for the event $\tau<T$, we have
%\begin{align}
%    P(\tau<T)\leq P(\tau_1<T)+P(\tau_2<T)+P(\tau_3<T)\leq \frac{\delta}{2},\nonumber
%\end{align}
%which is a tail event.

%\textbf{Part II: Convergence of $\frac{1}{T}\sum_{t=0}^{T-1}\mathbb{E}[\|\nabla \Phi(\theta_t)w_t \|^2|\tau=T]$ }.

%For metric $\frac{1}{T}\sum_{t=0}^{T-1}\mathbb{E}[\|\nabla \Phi(\theta_t)w_t \|^2|\tau=T]$, following \eqref{eq: MOO descent up to stopping time}, we have
%\begin{align}
%    &\frac{1}{T}\sum_{t=0}^{T-1}\mathbb{E}[\|\nabla \Phi(\theta_t)w_t \|^2|\tau=T] \nonumber\\
%    &\leq \frac{1}{T}\frac{1}{P(\tau=T)}\sum_{t=0}^{T-1}\mathbb{E}[\|\nabla \Phi(\theta_t)w_t \|^2]\nonumber\\
%    &\leq \frac{2}{\alpha}\cdot \frac{2}{T}\Big(\mathbb{E}[\Phi(\theta_0)w-\Phi^*w]+\frac{\alpha}{\beta}+{5\alpha\Lambda\sigma+5\alpha T\delta\epsilon^2+3\alpha T \delta^2\epsilon^4+3\alpha \delta^2\epsilon^4+6\alpha\sigma^2}\nonumber\\
%    &\quad +{2\alpha \rho T} + \alpha\beta\rho^2T + {\alpha^2 L_0 C_B^2T}+{4\alpha\beta m \Lambda^4 T+8\alpha\beta m\Lambda^2C_B^2T + 4\alpha\beta m T C_B^4}\Big) \nonumber\\
%    &= \frac{4(\mathbb{E}[\Phi(\theta_0)w-\Phi^*w])+4\alpha/\beta}{\alpha T} + {\frac{20 \Lambda \sigma}{T}+20\delta\epsilon^2+12\delta^2\epsilon^4+\frac{24\sigma^2}{T}+\frac{12\delta^2\epsilon^4}{T}}
%    +8\rho+4\beta\rho^2\nonumber\\
%    &\quad +4\alpha L_0C_B^2+16\beta m \Lambda^4+32\beta m \Lambda^2C_B^2+16\beta m C_B^4\leq 39\delta\epsilon^2,
%\end{align}
%where the last inequality is due to the choice of parameters $\alpha,\beta, \rho$ and $T$. Applying Markov's inequality,we have
%\begin{align}
%    P(\frac{1}{T}\sum_{t=0}^{T-1}\|\nabla \Phi(\theta_t)w_t \|^2\geq 78\epsilon^2 \mid \tau=T)\leq \frac{1/T\sum_{t=0}^{T-1}\mathbb{E}\big[\|\nabla \Phi(\theta_t)w_t \|^2\big|\tau=T]}{78\epsilon^2}\leq \frac{\delta}{2}.\nonumber
%\end{align}
%Combine all probability of aforementioned events, we have
%\begin{align}
%    &P(\frac{1}{T}\sum_{t=0}^{T-1}\|\nabla \Phi(\theta_t)w_t\|^2\leq 78\epsilon^2) \nonumber\\
%    &= 1-P(\tau<T)-P(\frac{1}{T}\sum_{t=0}^{T-1}\|\nabla \Phi(\theta_t)w_t \|^2>78\epsilon^2\mid \tau=T)\cdot P(\tau=T) \nonumber\\
%    &\geq 1-\frac{\delta}{2}-\frac{\delta}{2}=1-\delta.
%\end{align}
%Thus, we conclude Algorithm \ref{alg1} converge to $\epsilon$-Pareto stationary in the average sense,
%\begin{align}
%    \frac{1}{T}\sum_{t=0}^{T-1}\|\nabla \Phi(\theta_t) \|^2\leq 78\epsilon^2 \text{ w.h.p }, \nonumber
%\end{align}
%which completes the proof.

%\end{proof}







\section{Reformulation of Pareto-stationary condition}
\subsection{Optimality condition Reformulation}\label{Appendix: proof of reformulated optimality condition}
\reformedcondition*
\begin{proof}
Expanding $\| \nabla \Phi({\theta})w \|$, we have
\begin{align}
    &\|\nabla \Phi(\theta)w \| \nonumber\\
    &\overset{(i)}{\leq} 
    \|\nabla \Phi(\theta)w - \nabla_{\theta}\fL(\theta,\eta)w\| + \| \nabla_{\theta} \fL(\theta,\eta)w\|
    \nonumber\\
    %&= \|\sum_{i=1}^{m}w_i(\nabla_{{\theta}} L^i({\theta},\eta_i^*)-\nabla_{{\theta}} L^i({\theta}, \eta_i))\| + \| \sum_{i=1}^{m} w_i \nabla_{{\theta}}L^i({\theta},\eta_i)\| \nonumber\\
    &=\Big\|\sum_{i=1}^{m} w^i\Big(\mathbb{E}_{\xi}\Big[\big({f^*}'(\frac{\ell^i({\theta};\xi)-\eta^i}{\lambda})-{f^*}'(\frac{\ell^i({\theta};\xi)-\eta^{i,*}}{\lambda})\big)\nabla \ell^i({\theta},\xi)\Big]\Big)\Big \| + \Big\|\sum_{i=1}^{m}w^i\nabla_{{\theta}}\fL^i({\theta},\eta^i) \Big\| \nonumber \\
    &\overset{(ii)}{\leq} \sum_{i=1}^{m}w^i\Big\|\mathbb{E}_{\xi}\Big[\Big((f^*)'(\frac{\ell^i({\theta};\xi)-\eta^i}{\lambda})-(f^*)'(\frac{\ell^i({\theta};\xi)-\eta^{i,*}}{\lambda})\Big)\nabla \ell^i({\theta},\xi)\Big] \Big\| + \Big\| \sum_{i=1}^{m} w^i \nabla_{{\theta}}\fL^i({\theta},\eta^i)  \Big\|\nonumber\\
    &\overset{(iii)}{\leq} \sum_{i=1}^{m} w^i\mathbb{E}_{\xi}\Big[ \Big\| \Big((f^*)'(\frac{\ell^i({\theta};\xi)-\eta^i}{\lambda}) - (f^*)'(\frac{\ell^i({\theta};\xi)-\eta^{i,*}}{\lambda})\Big) \nabla \ell^i({\theta},\xi) \Big\|\Big] +\Big\|\sum_{i=1}^{m}w^i\nabla_{{\theta}} \fL^i({\theta}, \eta^i) \Big\|\nonumber \\
    & \overset{(iv)}{\leq} \sum_{i=1}^{m} w^i \mathbb{E}_{\xi}\Big[\Big\| (f^*)'(\frac{\ell^i({\theta};\xi)-\eta^i}{\lambda}) - (f^*)'(\frac{\ell^i({\theta};\xi)-\eta^{i,*}}{\lambda}) \Big\| \|\nabla \ell^{i}(\theta,\xi) \Big\| \Big]+ \Big \|\sum_{i=1}^{m}w^i\nabla_{{\theta}} \fL^i({\theta}, \eta^i) \Big \|\nonumber \\
    & \overset{(v)}{\leq} G \sum_{i=1}^{m} w^i \mathbb{E}_{\xi}|(f^*)'(\frac{\ell^i({\theta}, \xi)-\eta^i}{\lambda}) - (f^*)'(\frac{\ell^i({\theta}, \xi)-\eta^{i,*}}{\lambda}) |+ \|\sum_{i=1}^{m}w^i\nabla_{{\theta}} \fL^i({\theta},\eta^i) \|\nonumber \\
    & = G \sum_{i=1}^{m} w^i \mathbb{E}_{\xi}\|\nabla_{\eta^i}\fL^{i}({\theta}, \eta^i;\xi) - \nabla_{\eta^i}\fL^{i}({\theta}, \eta^{i,*};\xi) \|+\| \sum_{i=1}^{m} w^i \nabla_{{\theta}}\fL^i({\theta}, \eta_i)\| \nonumber \\
    & = G \sum_{i=1}^{m} w^i \big|\mathbb{E}_{\xi}\big[\nabla_{\eta_i} \fL^i({\theta},\eta^i;\xi)- \nabla_{\eta^i}\fL^i({\theta}, \eta^{i,*};\xi) \big] \big| + \|\sum_{i=1}^{m}w^i \nabla_{{\theta}} \fL^i({\theta}, \eta^i) \| \nonumber\\
    %&= G \sum_{i=1}^{m}w^i \big|\nabla_{\eta^i} L^i({\theta}, \eta^i) \big|+ \| \sum_{i=1}^{m} w^i \nabla_{{\theta}}L^i({\theta},\eta^i)  \|
    &= G  \big|\nabla_{\eta} \fL({\theta}, \eta) w\big|+ \|  \nabla_{{\theta}}\fL({\theta},\eta)w  \|,
\end{align}
where (i) and (ii) applies triangle inequality, $\|a+b \|\leq \| a\|+\| b\|$; (iii) applies Jensen's inequality inequality; (iv) applies Cauchy-Schwarz inequality; (v) applies $G$-Lipschitz continuity assumption for each $\ell^i(\cdot)$. Notice $(f^*)'$ is monotone non-decreasing, the sign of $\nabla_{\eta^i}\fL(\theta,\eta^i;\xi)-\nabla_{\eta^i}\fL(\theta,\eta^{i,*};\xi)$ depends on relative position between $\eta^i$ and $\eta^{i,*}$, this enables to put $\mathbb{E}_{\xi}$ into $|\cdot|$ without increasing its value.

Next we show condition \eqref{eq: epsilon Optimal condition} can be achieved by optimizing rescaled function $\fhL(\theta, \eta)=\fL({\theta}, G{\sqrt{m}}{\eta})$ such that $\|\nabla \fhL(\theta,\eta)w\|\leq \frac{\epsilon}{\sqrt{2}}$. Expanding $\|\nabla_{{\theta}, {\eta}}\fL({\theta}, G{\sqrt{m}}{\eta})w \|^2$, we have
\begin{align}
    \|\nabla_{{\theta}, {\eta}}\fhL({\theta}, {\eta}){w}\|^2=
    &\|\nabla_{{\theta}, {\eta}}\fL({\theta}, G{\sqrt{m}}{\eta}){w} \|^2\nonumber \\
    =&\Big\| \begin{bmatrix}
        \nabla_{{\theta}}\fL({\theta}, G\sqrt{m} {\eta})\\
        G\sqrt{m}\cdot\nabla_{{\eta}} \fL({\theta}, G\sqrt{m} {\eta})\\
    \end{bmatrix} {w} \Big\|^2 \nonumber\\
    %=&\Bigg\| \begin{bmatrix}
    %   \nabla_{{\theta}} L^1({\theta},G\sqrt{m}\eta^1) &...  & \nabla_{{\theta}} L^m({\theta},G\sqrt{m}\eta^m)  \\
    %   \nabla_{\eta^1}L^1({\theta}, \sqrt{m}\eta^1) & 0  &0  \\
    %    0& \nabla_{\eta^2}L^2({\theta}, G\sqrt{m}\eta^2) & 0  \\
    %    0& 0 & \nabla_{\eta^m}L^m({\theta}, G\sqrt{m}\eta^m) \nonumber \\ 
%\end{bmatrix} \begin{bmatrix}
    %w^1 \\\dots \\ w^m
%\end{bmatrix}\Bigg \|^2 \nonumber \\
%    =& \Bigg\|\begin{bmatrix}
    %\sum_{i=1}^{m}\nabla_{{\theta}}L^i({\theta},G\sqrt{m}\eta^i)w^i\\
    %G\sqrt{m}\nabla_{\eta_1}L^1({\theta},G\sqrt{m}\eta^1)w^1\\
    %\dots
    %\\G\sqrt{m} \nabla_{\eta^m}L^m({\theta}, G\sqrt{m}\eta^m)w^m
%\end{bmatrix}\Bigg\|^2
%\nonumber\\
    =&\|\sum_{i=1}^{m}\nabla_{{\theta}}\fL^i({\theta},G\sqrt{m}\eta^i)w^i\|^2 +G^2m\cdot\sum_{i=1}^{m}|\nabla_{\eta^i}\fL^i({\theta}, G\sqrt{m} \eta^i){w^i}|^2\nonumber \\
    =&\|\sum_{i=1}^{m}\nabla_{{\theta}}\fL^i({\theta},G\sqrt{m}\eta^i)w^i\|^2 +G^2m\cdot\sum_{i=1}^{m}|\nabla_{\eta_i}\fL^i({\theta}, G\sqrt{m} \eta^i){w^i}|^2\nonumber \\
     \overset{(i)}{\geq}& \| \sum_{i=1}^{m}\nabla_{{\theta}}\fL^i({\theta},G\sqrt{m}\eta^i)w^i\|^2+{G^2}\big(\sum_{i=1}^{m}|\nabla_{\eta^i}\fL^i({\theta}, G\sqrt{m}\eta^i)w^i|\big)^2\nonumber\\
    \overset{(ii)}{\geq}& \frac{1}{2}\Big(\| \sum_{i=1}^{m}w^i \nabla_{{\theta}}\fL^i({\theta}, G\sqrt{m}\eta^i)\|+{G}\sum_{i=1}^{m}| \nabla_{\eta^i}\fL^i({\theta},G\sqrt{m}\eta^i)w^i| \Big)^2\nonumber\\
    =&\frac{1}{2}\Big(\| \nabla_{{\theta}}\fL({\theta}, G\sqrt{m}\eta)w\|+{G}| \nabla_{\eta}\fL({\theta}, G\sqrt{m}\eta)w| \Big)^2,
\end{align}
where the first inequality applies $(a_1+....+a_m)^2\leq m\cdot \sum_{i=1}^{m}a_i^2$; the second inequality applies $\|a\|^2+\|b\|^2\geq \frac{1}{2}(a+b)^2$ holds for all $a,b\geq 0$. Taking square-root on both sides, we have
\begin{align}
    \|\nabla_{{\theta}, {\eta}}\fhL({\theta}, {\eta}){w} \|\geq \frac{1}{\sqrt{2}}\Big(\| \nabla_{{\theta}}\fL({\theta}, G\sqrt{m}\eta)w\|+{G} |\nabla_{\eta}\fL({\theta}, G\sqrt{m}\eta)w| \Big),\nonumber
\end{align}
which completes the proof.
\end{proof}

%\begin{remark}[Single-loop improvement]
%As showed in Lemma \ref{lemma: optimality condition}, one can alternatively optimize a rescaled function $\fL({\theta}, G\sqrt{m}\cdot {\eta})$ to reach $\epsilon$-level Pareto-stationary. However, noticed that such re-scaling is dependent with the number of objectives. When $m$ increases, it brings challenges in optimizing ${\theta}$ and $\eta$ simultaneously and may increase the number of iterations. Thus, there exists a trade-off between accuracy and algorithm complexity. The detailed algorithm is stated in algorithm \ref{alg1-2}.  
    
%\end{remark}




\section{Relavant Properties of Rescaled function $\fhL(\theta,\eta)$}\label{sec: property of hat L}
%In \cite{qirevistfDRO}, they have proved partial-smoothness property for single-objective, i.e., $\nabla_{\theta} L^i(\theta, \eta_i), \nabla_{\eta_i} L^i(\theta, \eta_i), \forall i\in[m]$ and the resultant descent lemma stated as follows
\begin{lemma}[Lemma 1 \citep{qirevistfDRO}]\label{lemma: generalized-smooth}
Let assumption \ref{assumption 1} hold, for each loss $\fhL^i(\theta, \eta^i)$, it satisfies $(\hL_0, \hL_1,\hL_2)$-smooth condition, i.e.,
\begin{align}
&\|\nabla_{\theta} {\fhL}^i({\theta}, \eta^i)-\nabla_{\theta} {\fhL}^i(\bar{\theta}, \eta^i)\|  \leq (\hL_0+\hL_1|\nabla_{\eta^i} {\fhL}^i({\theta}, \eta^i)|)\|{\theta}-\bar{\theta}\|, \nonumber\\
&|\nabla_{\eta^i} {\fhL}^i({\theta}, \eta^i)-\nabla_{\eta^i} {\fhL}^i({\theta}, \bar{\eta}^{i})| \leq \hL_2|\eta^i-\bar{\eta}^{i}|,
\end{align}
where $\hL_0=L+{G^2M}{\lambda}^{-1}, \hL_1={L}{(G\sqrt{m})^{-1}}, \hL_2 = {G^2Mm}{\lambda}^{-1}$
\end{lemma}
\begin{proof}
    The Proof follows \citet{qirevistfDRO} by changing $\fL(\theta, G\eta)$ into $\fhL(\theta,\eta)$. 
\end{proof}
\begin{remark}
The descent lemma with respect to $\theta$ is
\begin{align}
    {\fhL}^i({\theta}', \eta^i) \leq {\fhL}^i({\theta}, \eta^i)+\langle\nabla_{{\theta}} {\fhL}^i({\theta}, \eta^i), {\theta}'-{\theta}\rangle+\frac{\hL_0+\hL_1|\nabla_{\eta^i} {\fhL^i}({\theta}, \eta^i)|}{2}\|{\theta}'-{\theta}\|^2.
    \label{eq: single-objective descent lemma}
\end{align}
\end{remark}
\begin{lemma}[Lemma 3 \citep{qirevistfDRO}]\label{lemma: affine bounded noise}
Let assumption \ref{assumption 1} hold, for each objective $\fhL(\theta_t,\eta_t^i),i\in[m]$, the variance of $\HUpsilon_t^i, \HGamma_t^i$ can be upper bounded as
\begin{align}
    \mathbb{E}_{\bar{\xi}^{i}_{t}\sim \mathbb{P}^i}[\|\HGamma_t^{i} \|^2],\leq \hat{K_0}+\hat{K_1}|\nabla_{\eta^i}\fhL^i({\theta}_t, \eta_{t+1}^i)|^2,
    &\text{ and }\mathbb{E}_{\xi^{i}_{t}\sim \mathbb{P}^i}[\|\HUpsilon_{t}^{i} \|^2 ]\leq \hat{K}_2,
    \label{eq: variance bound for rescaled function}
\end{align}
where $\hat{K_0}=8G^2+10G^2M^2\lambda^{-2}\kappa^2$, $\hat{K_1} = 8/m$, $\hat{K_2} = mG^2 M^2 \lambda^{-2} \kappa^2$.
\end{lemma}
\begin{proof}
    The proof follows \citet{qirevistfDRO} by changing $\fL(\theta,G\eta)$ into $\fhL(\theta,\eta)$.
\end{proof}

%and by applying ${\theta}' = {\theta} - \frac{\nabla_{\theta}\fL({\theta},\eta_i)}{L^0+L^1| \nabla_{\eta_i}\fL({\theta}, \eta_i)|}$, \eqref{eq: single-objective descent lemma} has following facts
%\begin{align}
   % \|\nabla_{{\theta}}L^i({\theta}, \eta_i)\|^2\leq 2\Delta_i( L^0+L^1|\nabla_{\eta_i}L^i({\theta}, \eta_i)|).
    %\label{eq: induced single-objective descent lemma}
%\end{align}
%where $\Delta_i = \fL({\theta}, \eta_i)-\min_{{\theta},\eta}\fL({\theta}, \eta_i)$.

%Notice that, \cite{jinnonconvexdro} proved $(L^0, L^1)$-generalized smooth condition of $L^i({\theta},\eta_i)$ in terms of $\nabla_{{\theta}, \eta_i}\fL({\theta},\eta_i)$, namely
%\begin{align}
    %&\|\nabla_{{\theta},\eta_i}L^i({\theta},\eta_i)-\nabla_{{\theta},\eta_i}L^i({\theta}',\eta_i') \|\nonumber\\
    %&\leq (L+(G^2+1)\lambda^{-1}M+L\|\nabla_{{\theta},\eta_i}L^i({\theta},\eta_i) \|)\|({\theta},\eta_i)-({\theta}',\eta_i') \|,\forall i\in[m]
%\end{align}
%The RHS of above equation satisfies sub-linear property since $L+(G^2+1)\lambda^{-1}M+L\| \nabla_{\theta,\eta_i}\fL({\theta}, \eta_i)\|\leq  \mathcal{O}(\|\nabla_{\theta, \eta_i}\fL({\theta}, \eta) \|^2)$. Thus for each $L^i, i\in[m]$, Corollary 3.6 in \cite{li2023convex} still applies, we present the corollary here for completeness
%{\begin{corollary}[Corollary 3.6 in \cite{li2023convex}]
%    Suppose function $L^i({\theta},\eta_i)$ satisfies above $(L^0, L^1)$-smooth property. If $L^i({\theta}, \eta_{i})-L^*\leq F$ for some ${\theta}_t,\eta_i$ and $F\geq 0$. Then, for $\Lambda=\sup\{u|u^2\leq (2L^0+4L^1u)\cdot F) \}$, we have $ \|\nabla_{{\theta},\eta_i}\fL({\theta},\eta_i) \|\leq \Lambda$.
%\end{corollary}}
%Since $\|\nabla_{\eta_i}L^i({\theta}, \eta_i) \|\leq \|\nabla_{{\theta},\eta_i}L^i({\theta}, \eta_i) \|$ and $\|\nabla_{{\theta}}L^i({\theta}, \eta_i) \|\leq \|\nabla_{{\theta},\eta_i}L^i({\theta}, \eta_i) \|$ , we must have $\|\nabla_{\eta_i}\fL({\theta},\eta_i) \|\leq G$ as long as there exists some $({\theta},\eta_i)$ such that $\fL({\theta},\eta_i)-L^*\leq F$.

\begin{restatable}[Descent Lemma]{lemma}{singleloopdescentlemma}
    Let Assumption \ref{assumption 1} hold, define $\hL_0 = L+G^2M\lambda^{-1}, \hL_1 = \fL(G\sqrt{m})^{-1}$. Then, for rescaled multi-objective function $\fhL({\theta}, {\eta})$, given any preference vector $w$,
    we have the following descent lemma
    \begin{align}
        \fhL({\theta}', {\eta}){w} &\leq \fhL({\theta}, {\eta}){w} + \langle\nabla_{{\theta}}\fhL({\theta}, {\eta}){w}, {\theta}'-{\theta} \rangle +\frac{\hL_0+\hL_1 |\nabla_{\eta}\fhL(\theta,\eta)w|}{2}\|{\theta}'-{\theta} \|^2,
        \label{eq: moo descent lemma theta}
\end{align}
holds for any $\theta,\theta'\in \mathbf{R}^n$ and fixed $\eta$. Additionally, we have
\begin{align}
        \fhL({\theta}, {\eta}'){w}\leq \fhL({\theta}, {\eta}){w} +\langle\nabla_{\eta}\fhL(\theta,\eta)w,\eta'-\eta \rangle + \frac{\hL_2}{2}\|\eta'-\eta\|^2,
        \label{eq: moo descent lemma eta}
\end{align}
holds for any $\eta,\eta'\in\mathbf{R}^m$ given $\theta$.

\end{restatable}
\begin{proof}
Based on the descent lemma \eqref{eq: single-objective descent lemma} for each $\fhL^i(\theta, \eta)$, since $\fhL({\theta}', {\eta}){w} = \sum_{i=1}^{m} \fhL^i({\theta}',{\eta}_i)w^i$, we have
\begin{align}
    &\fhL({\theta}', {\eta}){w} \nonumber\\
    %&= \sum_{i=1}^{m} \fhL^i({\theta},{\eta}^i)w^i \nonumber\\
    &\overset{(i)}{\leq} \sum_{i=1}^{m}\Big( {\fhL}^i({\theta}, \eta^i)+\langle\nabla_{{\theta}} {\fhL}^i({\theta}, \eta^i), {\theta}'-{\theta}\rangle+\frac{\hL_0+\hL_1|\nabla_{\eta^i} {\fhL}({\theta}, \eta^i)|}{2}\|{\theta}-{\theta}'\|^2\Big)w^i \nonumber\\
    %& = \sum_{i=1}^{m} \fhL^i({\theta}, \eta^i) w^i + \sum_{i=1}^{m}\langle\nabla_{{\theta}} {\fhL}^i({\theta}, \eta^i), {\theta}'-{\theta}\rangle w^i + \frac{L_0}{2}\|{\theta}'-{\theta}\|^2+ \frac{L_1\sum_{i=1}^{m} |\nabla_{\eta^i} \fhL({\theta}, \eta^i) |w^i}{2}\|{\theta}'-{\theta} \|^2 \nonumber\\
    & \overset{(ii)}{=} \fhL({\theta}, {\eta}){w}+ \langle  \nabla_{\theta}^{\top}\fhL(\theta,\eta)(\theta'-\theta), w \rangle + \frac{\hL_0}{2}\|{\theta}'-{\theta} \|^2+ \frac{\hL_1\sum_{i=1}^{m} | \nabla_{\eta^i} \fhL({\theta}, {\eta}^i)| w^i}{2}\|{\theta}'-{\theta} \|^2 \nonumber\\
    %& \overset{}{=} \fhL({\theta}, {\eta}){w} + \langle \nabla_{\theta}\fhL({\theta}, {\eta})w,{\theta}'-{\theta} \rangle + \frac{L_0}{2}\|{\theta}'-{\theta} \|^2+ \frac{L_1\sum_{i=1}^{m} | \nabla_{\eta^i} \fhL({\theta}, {\eta}^i)| w^i}{2}\|{\theta}'-{\theta} \|^2  \nonumber\\
    &= \fhL({\theta}, {\eta}){w} + \langle \nabla_{\theta}\fhL({\theta}, {\eta})w,{\theta}'-{\theta} \rangle + \frac{\hL_0}{2}\|{\theta}'-{\theta} \|^2+ \frac{\hL_1 | \nabla_{\eta} \fhL({\theta}, {\eta})w| }{2}\|{\theta}'-{\theta} \|^2,  \nonumber
    %& \overset{(iv)}{=} \fL({\theta}, {\eta}){w} + \langle \nabla_{{\theta}}\fL({\theta}, {\eta}){w}, {\theta}'-{\theta}\rangle + \frac{L^0}{2}\|{\theta}'-{\theta} \|^2+ \frac{L^1}{2}\|{\theta}'-{\theta} \|^2 \langle |\nabla_{{\eta}}\fL({\theta}, {\eta})|, w \rangle \nonumber\\
    %& \overset{(v)}{\leq} \fL({\theta}, {\eta}){w} + \langle \nabla_{{\theta}}\fL({\theta}, {\eta}){w}, {\theta}'-{\theta}\rangle + \frac{L^0}{2}\|{\theta}'-{\theta} \|^2 + \frac{L^1}{2}\|{\theta}'-{\theta} \|^2 \| w\|\|\nabla_{{\eta}}\fL({\theta}, {\eta}) \|_2\nonumber\\
    %&\overset{(vi)}{=} \fL({\theta},{\eta}){w} + \langle\nabla_{{\theta}}\fL({\theta}, {\eta}){w}, {\theta}'-{\theta} \rangle +\frac{L^0+L^1 \|\nabla_{{\eta}}\fL({\theta}, {\eta})\|_2}{2}\|{\theta}'-{\theta} \|^2\nonumber\\
    %&\leq \fL({\theta},{\eta}){w} + \langle\nabla_{{\theta}}\fL({\theta}, {\eta}){w}, {\theta}'-{\theta} \rangle +\frac{L^0+L^1\Lambda}{2}\|{\theta}'-{\theta} \|^2
\end{align}
where (i) applies descent lemma \eqref{eq: single-objective descent lemma} for each $\fhL^i({\theta}, \eta^i)$; (ii) rewrites $\sum_{i=1}^{m}\langle\nabla_{{\theta}} {\fhL}^i({\theta}, \eta^i), {\theta}'-{\theta}\rangle w^i$ as $\langle \nabla_{\theta}^{\top}\fhL(\theta, \eta)(\theta'-\theta),w\rangle$. 

Then, for descent lemma \eqref{eq: moo descent lemma eta}, we have
    \begin{align}
         &\fhL(\theta,\eta')w-\fhL(\theta, \eta)w-\langle \nabla_{\eta}\fhL(\theta,\eta)w, \eta'-\eta \rangle \nonumber \\ 
         %=&\sum_{i=1}^{m}\Big(\fhL^i(\theta, \eta^{i'})-\fhL^i(\theta, \eta^i)-\langle \nabla_{\eta^i} \fhL^i(\theta,\eta^i), \eta^{i'}-\eta^i \rangle\Big) w^i\nonumber\\
        =& \sum_{i=1}^{m}\Big(\int_{0}^{1}\langle \nabla_{\eta^i} \fhL^i(\theta, \eta^i_{{\iota^i}}), {\eta^{i'}}-\eta^i\rangle \mathrm{d} {\iota}^i - \int_{0}^{1} \langle \nabla_{\eta^i} \fhL^i(\theta, \eta^i), \eta^{i'}-\eta^i\rangle \mathrm{d}{\iota}^i\Big)w^i  \nonumber \\
        %=& \sum_{i=1}^{m}\Big(\int_{0}^{1}\langle \nabla_{\eta^i}\fhL^i(\theta_t, \eta^i_{{\iota}^i})-\nabla_{\eta^i}\fhL(\theta_t,\eta^i), \eta^{i'}-\eta^i\rangle \mathrm{d} {\iota}^i\Big)w^i \nonumber\\
        \overset{(i)}{\leq}& \sum_{i=1}^{m}
        \Big(\int_{0}^{1}  \|\nabla_{\eta^i} \fhL^i(\theta, \eta^i_{\iota^i})-\nabla_{\eta^i} \fhL(\theta,\eta^i ) \| \|\eta^{i'}-\eta^i \|\mathrm{d}\iota^i \Big) w^i\nonumber \\
        %=&\sum_{i=1}^{m}\Big(\int_{0}^{1}  |\nabla_{\eta^i} \fhL^i(\theta_t, \eta^i_{\iota^i})-\nabla_{\eta^i} \fhL(\theta,\eta^i )| |\eta^{i'}-\eta^i |\mathrm{d}\iota^i \Big) w^i\nonumber\\
        \overset{(ii)}{\leq}& \sum_{i=1}^{m} \Big(\int_{0}^{1} \iota^i \hL_2|\eta^{i'}-\eta^i|^2 \mathrm{d}\iota^i \Big) w^i=\frac{1}{2}\hL_2\Big(\sum_{i=1}^{m}|\eta^{i'}-\eta^i|^2w^i\Big) \leq \frac{\hL_2}{2}\|\eta'-\eta\|^2,\nonumber
    \end{align}
    where $\iota^i\in[0,1]$, $\eta^i_{\iota^i}=(1-\iota^i)\eta^i+\iota^i\eta^{i'}$; the first inequality utilizes Cauchy-Schwarz inequality in terms of $L_2$-norm ($\nabla_{\eta^i}\fhL^i(\theta, \eta^i), \eta^i\in\mathbf{R}$), and the second inequality is due to $\hL_2$-smooth property for each $\nabla_{\eta^i}\fhL^i(\theta, \eta^i)$ and the last inequality applies $w^i\leq 1,\forall i\in[m]$.
\end{proof}


\begin{lemma}[Relationship between $\nabla_{\theta}\fhL^i(\theta,\eta^i)$ and $\nabla_{\eta}\fhL^i(\theta, \eta^i)$]
    Given $(\theta,\eta^i)$, for $\nabla_{\theta}\fhL^i(\theta,\eta^i)$ and $\nabla_{\eta}\fhL^i(\theta,\eta^i)$, we have the following relationship
    \begin{align}
        \|\nabla_{\theta}\fhL^i(\theta, \eta^i)\|\leq G+|\nabla_{\eta^i}\fhL^i(\theta,\eta^i)|.
        \label{eq: gradient relationship}
    \end{align}
\end{lemma}
\begin{proof}
    Expanding $\nabla_{\theta}\fhL^i(\theta,\eta^i)$, we have
    \begin{align}
        \|\nabla_{\theta}\fhL^i(\theta, \eta^i)\|&=\|\mathbb{E}_{\xi}\Big[(f^*)'\big(\frac{\ell^i(\theta;\xi^i)-G\sqrt{m}\eta^i}{\lambda} \big)\nabla\ell^i(\theta;\xi^i)\Big]\|\nonumber\\
        &\leq \mathbb{E}_{\xi}\|(f^*)'\big(\frac{\ell^i(\theta;\xi^i)-G\sqrt{m}\eta^i}{\lambda}\big)\nabla \ell^i(\theta;\xi^i)\|\nonumber\\
        %&=\mathbb{E}_{\xi}\Big[|(f^*)'\big(\frac{\ell^i(\theta;\xi^i)-G\sqrt{m}\eta^i}{\lambda}\big)|\|\nabla \ell^i(\theta;\xi^i)\|\Big]\nonumber\\
        &\leq G \mathbb{E}_{\xi}\Big[ |(f^*)'\big( \frac{\ell^i(\theta;\xi^i)-G\sqrt{m}\eta^i}{\lambda}\big) |\Big]\nonumber\\
        &\overset{(i)}{=} G|\mathbb{E}_{\xi}\big[(f^*)'\big(\frac{\ell^i(\theta;\xi^i)-G\sqrt{m}\eta^i}{\lambda}\big) \big] |\nonumber\\
        &=G |1-\frac{\nabla_{\eta^i}\fhL^i(\theta,\eta^i)}{G\sqrt{m}}|\leq G+\frac{|\nabla_{\eta^i}\fhL^i(\theta,\eta^i)|}{\sqrt{m}}\leq G+|\nabla_{\eta^i}\fhL^i(\theta,\eta^i)|,
    \end{align}
where the first inequality applies Jensen's inequality; the second inequality applies $G$-Lipschitz assumption of $\ell^i(\theta;\xi^i)$; (i) utilizes the maximization arguments of convex conjugate function, the primal variable $r={\mathrm{d}\mathbb{Q}}/{\mathrm{d}\mathbb{P}}\geq 0$ thus implies ${f^*}'(s)\geq0$; the last inequality applies $|a-b|\leq |a|+|b|$ and $m\geq 1$.
\end{proof}

\begin{corollary}[Gradient Upper bound of $\nabla_{\eta}\fhL(\theta_t, \eta_{t+1})$]
    Given update rule $\eta_{t+1} = \eta_t - \gamma\mu_t Z_tw_t$, if there exists $\Lambda$ such that $|\nabla_{\eta^i}\fhL(\theta_t, \eta_{t}^i)|<\Lambda$~\footnote{Notice, $\nabla_{\eta^i}\fhL(\theta_t,\eta_t^i)\in \mathbf{R}$ is a scalar, this implies $|\nabla_{\eta^i}\fhL(\theta_t, \eta_{t}^i)|=\|\nabla_{\eta^i}\fhL(\theta_t, \eta_{t}^i)\|$.} holds for all $t\leq T$. Then, for $|\nabla_{\eta}\fhL(\theta_t, \eta_{t+1})w|$, we have
    \begin{align}
        |\nabla_{\eta} \fhL(\theta_t, \eta_{t+1})w|\leq \hL_2\gamma\mu_t\|Z_tw_t\| + |\nabla_{\eta}\fhL(\theta_t, \eta_t)w|\leq\hL_2\gamma\mu_t\|Z_tw_t\|+\Lambda,
    \end{align}
    holds for all $t\leq T$.
\end{corollary}
\begin{proof}
Expanding $|\nabla_{\eta}\fhL(\theta_t,\eta_{t+1})w|$, we have
\begin{align}
    |\nabla_{\eta} \fhL(\theta_t, \eta_{t+1})w|&=\sum_{i=1}^{m}w^i|\nabla_{\eta^i}\fhL^i(\theta_t, \eta_{t+1}^i) |\nonumber\\
    &\leq \hL_2\langle w, |\eta_{t+1}-\eta_t|\rangle + |\nabla_{\eta} \fhL(\theta_t, \eta_{t})w|\nonumber\\
    &=\hL_2\gamma\mu_t\langle w, |Z_tw_t| \rangle + |\nabla_{\eta} \fhL(\theta_t, \eta_{t})w|\nonumber\\
    &\leq \hL_2\gamma\mu_t\|Z_tw_t\| + |\nabla_{\eta}\fhL(\theta_t, \eta_t)w|\nonumber\\
    &\leq\hL_2\gamma\mu_t\|Z_tw_t\|+\Lambda,
    \label{eq: partial-smooth property}
\end{align}
where the first inequality applies $\hL_2$-smooth w.r.t. $|\nabla_{\eta}\fhL^i(\theta, \eta^i)|$; the second inequality applies Cauchy-Schwarz inequality and $\|Z_tw_t\| = \sqrt{\sum_{i=1}^{m}(Z_t^iw_t^i)^2}=\| | Z_tw_t| \|$.
\end{proof}


\begin{corollary}
If there exists $\Lambda$ such that $|\nabla_{\eta^i}\fhL(\theta_t, \eta_{t}^i)|<\Lambda$. For $\nabla_{\theta}\fhL(\theta_{t},\eta_{t+1})$ and $\nabla_{\eta} \fhL(\theta_t, \eta_{t})$, we have the following upper bound
    \begin{align}
        \|\nabla_{\theta}\fhL(\theta_t,\eta_{t+1})\|_F^2\leq 3mG^2 + 3\gamma^2\mu_t^2\hL_2^2 \|Z_tw_t\|^2+3m\Lambda^2.
        \label{eq: gradient relationship 2}
    \end{align}
\end{corollary}
\begin{proof}
    Expanding $\|\nabla_{\theta}\fhL(\theta_t, \eta_{t+1})\|^2$, we have
    \begin{align}
        \|\nabla_{\theta}\fhL(\theta_t, \eta_{t+1})\|_F^2&=\sum_{i=1}^{m}\|\nabla_{\theta}\fhL^i(\theta_t, \eta_{t+1}^i) \|^2\nonumber\\
        &\leq \sum_{i=1}^{m}(G+|\nabla_{\eta^i}\fhL^i(\theta_t, \eta_{t+1}^i)|)^2\nonumber\\
        &\leq \sum_{i=1}^{m}(G+\gamma\mu_t\hL_2|Z_t^iw_t^i|+\Lambda)^2\nonumber\\
        &\leq 3mG^2+3\gamma^2\mu_t^2\hL_2^2\sum_{i=1}^{m}|Z_t^iw_t^i|^2 + 3m\Lambda^2 \nonumber\\
        &= 3mG^2 + 3\gamma^2\mu_t^2\hL_2^2 \|Z_tw_t\|^2+3m\Lambda^2,
    \end{align}
    where the first inequality leverages \eqref{eq: gradient relationship} for each $\nabla_{\theta}\fhL^i(\theta_t, \eta_{t+1})$ and $\nabla_{\eta^i}\fhL^i(\theta_t, \eta_{t+1})$; the second inequality leverages the $\hL_2$-smooth of $\nabla_{\eta^i}\fhL^i(\theta_t, \eta_{t+1}^i)$, such that $|\nabla_{\eta^i}\fhL^i(\theta_t, \eta_{t+1}^i)|\leq \hL_2|\eta_{t+1}^i-\eta_{t}^i|+|\nabla_{\eta^i}\fhL^i(\theta_t, \eta_{t}^i)|\leq \gamma\mu_t\hL_2|Z_t^iw_t^i|+\Lambda$; the third inequality leverages $(a+b+c)^2\leq 3a^2+3b^2+3c^2$.
\end{proof}



\begin{corollary}
    For $\HGamma_{t}={X}_t-\nabla_{\theta}\fhL(\theta_t, \eta_{t+1})$, if there exists $\Lambda<\infty$ such that
    $\|\nabla_{\eta^i}\fhL^i(\theta_t, \eta^i_{t})\|\leq \Lambda$ holds~\footnote{Notice, $\nabla_{\eta^i}\fhL(\theta_t,\eta_t^i)\in \mathbf{R}$ is a scalar, this implies $|\nabla_{\eta^i}\fhL(\theta_t, \eta_{t}^i)|=\|\nabla_{\eta^i}\fhL(\theta_t, \eta_{t}^i)\|$.}. Then, we have following upper bound
    \begin{align}
        \mathbb{E}\big[\|\HGamma_{t}w\|\big]\leq \Xi_3'=\sqrt{\frac{\hat{K}_0}{N_1}}+\sqrt{\frac{2\hat{K}_1\hL_2^2\gamma^2f_2^2}{N_1}}+\sqrt{\frac{2\hat{K}_1\Lambda^2}{N_1}},
        %\sqrt{\frac{\hat{K}_0}{N_1}}+\sqrt{\frac{4\hat{K}_1\hat{K}_2\hL_2^2\gamma^2}{N_1^2}}+\sqrt{\frac{4\hat{K}_1\hat{L}_2^2\gamma^2\Lambda^2}{N_1}}+\sqrt{\frac{2\hat{K_1}\Lambda^2}{N_1}}.
        \label{eq: estimation error bound thm2}
    \end{align}
\end{corollary}
\begin{proof}
    Expanding $\mathbb{E}\big[\|\HGamma_{t}w\|\big]$, we have
    \begin{align}
        \mathbb{E}\big[\|\HGamma_{t}w\|\big]&=\mathbb{E}\big[\mathbb{E}_{\xi_t,\bar{\xi}_t,\eta_{t+1}}[\| \HGamma_tw\||\theta_t,\eta_t,w_t,t]\big]\leq \mathbb{E}\big[\sqrt{\mathbb{E}_{\xi_t,\bar{\xi}_t,\eta_{t+1}}\big[\|\HGamma_{t}w\|^2 |\theta_t,\eta_t,w_t,t\big]}\big]\nonumber\\
        %&=\mathbb{E}\Big[\sqrt{\mathbb{E}_{\xi_t,\bar{\xi}_t,\eta_{t+1}}\big[\|\sum_{i=1}^{m}{\HGamma_{t}^i}w^i\|^2 |\theta_t,\eta_t,w_t,t\big]}\Big]\nonumber\\
        &\overset{(i)}{\leq} \mathbb{E}\Big[\sqrt{\mathbb{E}_{\xi_t,\bar{\xi}_t,\eta_{t+1}}\big[\sum_{i=1}^{m}w^i\|\HGamma_{t}^i\|^2|\theta_t,\eta_t,w_t,t\big]}\Big]\nonumber\\
        %&=\mathbb{E}\Big[\sqrt{\sum_{i=1}^{m}w^i\mathbb{E}_{{\xi_t,\bar{\xi}_t,\eta_{t+1}}}[\|\HGamma_{t}^i\|^2|\theta_t,\eta_t,w_t,t]}\Big]\nonumber\\
        &\overset{(ii)}{\leq}\mathbb{E}\Big[ \sqrt{\mathbb{E}_{\xi_t,\eta_{t+1}}\big[\sum_{i=1}^{m}w^i(\frac{\hat{K}_0}{N_1}+\frac{\hat{K}_1}{N_1}|\nabla_{\eta^i}\fhL^i(\theta_t, \eta_{t+1}^i)|^2|\theta_t,\eta_t,w_t,t\big]}\Big]\nonumber\\
        &\overset{(iii)}{\leq} \mathbb{E}\Big[\sqrt{\mathbb{E}\big[\sum_{i=1}^{m}w^i(\frac{\hat{K}_0}{N_1}+\frac{\hat{K}_1}{N_1}(\hL_2\gamma\mu_t|Z_t^iw_t^i|+\Lambda)^2)|\theta_t,\eta_t,w_t,t\big]}\Big]\nonumber\\
        &\overset{}{\leq} \mathbb{E}\Big[\sqrt{\mathbb{E}_{\xi_t}\big[\frac{\hat{K}_0}{N_1}+\frac{2\hat{K}_1\hL_2^2\gamma^2\mu_t^2}{N_1}\|Z_tw_t\|^2+\frac{2\hat{K}_1}{N_1}\Lambda^2|\theta_t,\eta_t,w_t,t\big]}\Big]\nonumber\\
        &\overset{(iv)}{\leq} \sqrt{\frac{\hat{K}_0}{N_1}+\frac{2\hat{K}_1\hL_2^2\gamma^2f_2^2}{N_1}+\frac{2\hat{K}_1\Lambda^2}{N_1}}\nonumber\\
        &\overset{(v)}{\leq} \underbrace{\sqrt{\frac{\hat{K}_0}{N_1}}+\sqrt{\frac{2\hat{K}_1\hL_2^2\gamma^2f_2^2}{N_1}}+\sqrt{\frac{2\hat{K}_1\Lambda^2}{N_1}}}_{\Xi'_3},
        %%%%%%%%%% old
        %&\overset{(iv)}{\leq}\mathbb{E}\Big[ \sqrt{\mathbb{E}_{\xi_t}\big[\frac{\hat{K}_0}{N_1}+\frac{4\hat{K}_1\hL_2^2\gamma^2\mu_t^2}{N_1}\|\HUpsilon_t w_t\|^2+\frac{4\hat{K}_1\hL_2^2\gamma^2\mu_t^2}{N_1}\|\nabla_{\eta}\fhL(\theta_t, \eta_t)w_t\|^2+\frac{2\hat{K}_1\Lambda^2}{N_1}|\theta_t,\eta_t,w_t,t\big]}\Big]\nonumber\\
        %&\overset{}{\leq}\sqrt{\frac{\hat{K}_0}{N_1}+\frac{4\hat{K}_1\hat{K}_2\hL_2^2\gamma^2\mu_t^2}{N_2N_1}+\frac{4\hat{K}_1\hat{L}_2^2\gamma^2\mu_t^2}{N_1}\Lambda^2+\frac{2\hat{K}_1}{N_1}\Lambda^2}\nonumber\\
        %&\overset{(v)}{\leq}\underbrace{\sqrt{\frac{\hat{K}_0}{N_1}}+\sqrt{\frac{4\hat{K}_1\hat{K}_2\hL_2^2\gamma^2\mu_t^2}{N_2N_1}}+\sqrt{\frac{4\hat{K}_1\hat{L}_2^2\gamma^2\mu_t^2\Lambda^2}{N_1}}+\sqrt{\frac{2\hat{K_1}\Lambda^2}{N_1}}}_{\Xi_3},
        \label{eq: inner product bound thm 2}
    \end{align}
    where (i) applies Jensen's inequality; (ii) applies the fact that $\mathbb{E}_{\bar{\xi}_{t}^i}\|\HGamma_{t}^i\|^2\leq \hat{K}_0/N_1+(\hat{K}_1/N_1)\cdot\|\nabla_{\eta^i}\fhL^i(\theta_t,\eta^{i}_{t+1})\|^2$; (iii) applies the fact $|\nabla_{\eta^i}\fhL^i(\theta_t,\eta^{i}_{t+1})|^2\leq (\hL_2\mu_t\gamma|Z_t^iw_t^i|+\|\nabla_{\eta^i}\fhL^i(\theta_t, \eta_{t}^i)\|)^2\leq 2\hL_2^2\mu_t^2\gamma^2 |Z_t^iw_t^i|^2+2\Lambda^2$;
    (iv) utilizes gradient clipping rule, $\mu_t\leq \frac{f_2}{\|Z_tw_t\|}$; the last inequality leverages $\sqrt{a+b+c}\leq \sqrt{a}+\sqrt{b}+\sqrt{c}$.
    %(iv) applies the fact $\|Z_tw_t\|^2\leq 2\|\HUpsilon_tw_t\|^2+2\|\nabla_{\eta}\fhL(\theta_t, \eta_t)w_t\|^2\leq 2\|\HUpsilon_tw_t\|^2+2\Lambda^2$ and $\mathbb{E}_{\xi_t}[|\HUpsilon_t^i|^2]\leq \hat{K}_2/N_2$; (v) applies the fact $\sqrt{a+b+c+d}\leq \sqrt{a}+\sqrt{b}+\sqrt{c}+\sqrt{d}$.
\end{proof}
    

\section{Convergence Analysis of Algorithm~\ref{alg1-3}}
\subsection{Formal Statement of Theorem \ref{thm: convergence of alg3} and Proof}\label{Appendix: Convergence of Algorithm 3}
Given well-defined problem dependent parameters $m$,$G$,$\hL_0=L+G^2M\lambda^{-1}$,$\hL_1=\fL(G\sqrt{m})^{-1}$,$\hL_2=G^2Mm\lambda^{-1}$,$\hat{K}_0=8G^2+10G^2M^2\lambda^{-2}\kappa^2$, $\hat{K}_1=8/m$, $\hat{K}_2=mG^2M^2\lambda^{-2}\kappa^2$,$\Lambda=\sup\{u\geq 0|u^2\leq2\hL_2\bar{F}(u+1) \}$, and place-hold constant $C_0$.
Let $d_1'\dots d_8'$, $e_1\dots e_3$ be constant. For $\bar{F}<\infty$ be a finite constant such that
\begin{align}
    \frac{\bar{F}}{8}\geq \max\Big\{ \frac{m(\Delta_{\theta_0,\eta_0}+d'_1+\dots+d'_8)}{\delta},e_1+e_2+ \sum_{j=3}^{8}d_j' \Big\} .
\end{align}
Set parameters of Algorithm \ref{alg1-3} as follows
\begin{align}
        c_2&=f_2 = \delta\epsilon ,\nonumber\\
        \alpha_t &= \min\{\frac{1}{2},\frac{\delta\epsilon}{\|X_tw_t\|} \},\mu_t = \min\{\frac{1}{2},\frac{\delta\epsilon}{\| Z_tw_t\|} \},\nonumber\\
       N_1 & = \max\big\{9\hat{K}_0\delta^{-3}\epsilon^{-2}, 18\hat{K}_1\delta^{-1}, 18\Lambda^2\hat{K}_1\delta^{-3}\epsilon^{-2}\}= \mathcal{O}(\max\{G^2M^2\lambda^{-2}\kappa^2,{\Lambda^2}/{m}\}{\delta^{-3}}\epsilon^{-2}) ,\nonumber\\
        N_2& = \max\big\{ {18}\hat{K_2}\delta^{-3}\epsilon^{-2}\big\}=\mathcal{O}(mG^2M^2\lambda^{-2}\kappa^2\delta^{-3}\epsilon^{-2}),\nonumber\\
        \Xi_1' &= mG^2+\gamma^2\hL_2^2\delta^2\epsilon^2+m\Lambda^2\leq mG^2+m\Lambda^2+1,\nonumber\\
        \Xi_2' &= \frac{m\hat{K}_0}{N_1}+\frac{2\hL_2^2\gamma^2\delta^2\epsilon^2\hat{K_1}}{N_1}+\frac{2m\Lambda^2\hat{K}_1}{N_1}\leq\frac{m\delta^3\epsilon^2}{3}, \nonumber\\
        \Xi_3' &=\sqrt{\frac{\hat{K}_0}{N_1}}+\sqrt{\frac{2\hat{K}_1\hL_2^2\gamma^2\delta^2\epsilon^2}{N_1}}+\sqrt{\frac{2\hat{K}_1\Lambda^2}{N_1}}\leq \delta^{3/2}\epsilon,\nonumber
        \\
        %\gamma &\leq \min\Big\{{\frac{2}{(\hL_0+\hL_1\Lambda)},\frac{1}{\hL_2+2\vartheta\hL_1^2\hL_2^2}, (\frac{16\vartheta}{10\delta^2\epsilon^2})^{1/3}} ,\blue{\frac{d_1}{T\delta^{5/2}\epsilon^2},\blue{\frac{d_2}{T\delta^{5/2}\epsilon^2},\blue{\beta d_3},(\frac{16\vartheta d_8}{\epsilon^4\delta^4T})^{1/4}}},\brown{ \frac{1}{192m^2T\delta^4\epsilon^4},\frac{e_1^2}{3m}}\Big\}\nonumber\\
        \gamma &=\min\Big\{{\frac{2}{(\hL_0+\hL_1\Lambda)},\frac{1}{\hL_2+2\vartheta\hL_1^2\hL_2^2}, (\frac{16\vartheta}{10\delta^2\epsilon^2})^{1/3}},{\beta d_3'},{\frac{9me_1^2}{C_0^2}},\frac{\min\{d_1',d_2'\}}{T\delta^{5/2}\epsilon^2},\nonumber\\
        &\frac{d_6'}{(3\Xi_2'+9\Xi_1')\delta^2\epsilon^2T\beta}, \frac{d_7'}{3m(\Lambda^2+\delta^3\epsilon^2)\delta^2\epsilon^2\beta T},(\frac{16\vartheta d_8'}{\delta^4\epsilon^4T})^{1/4}
        \Big\}\nonumber\\
        &=\min\Big\{ \mathcal{O}(\frac{1}{\hL_0+\hL_1\Lambda}), \mathcal{O}(\frac{1}{\hL_2+\vartheta\hL_1^2\hL_2^2}),\mathcal{O}(\beta),\mathcal{O}(\frac{1}{(mG^2+m\Lambda^2)\beta T\delta^2\epsilon^2}) \Big\},\nonumber\\
        \beta &= \min\Big\{{ \frac{4\rho^{-1}}{3},\frac{1}{90(mG^2+m\Lambda^2+1)},\frac{1}{30m\delta^3\epsilon^2},\frac{1}{30m\Lambda^2}},\frac{6e_2}{C_0^2},{\frac{4d_5'}{3d_4'\rho}}\Big\}=\mathcal{O}(\frac{1}{mG^2+m\Lambda^2}),\nonumber\\
        \rho &= \min\Big\{{\frac{1}{20}\delta^2\epsilon^2,{\frac{d_4'}{2\gamma T}}}\Big\}=\min\{\mathcal{O}(\delta^2\epsilon^2),\mathcal{O}(\frac{1}{\gamma T})\},\nonumber\\
        T&= {\max \Big\{ {10\hat{\Delta}_{\theta_0,\eta_0}\gamma^{-1}, 10\beta^{-1}} \Big\}}\delta^{-2}\epsilon^{-2}=\Theta(\max\{\hat{\Delta}_{\theta_0,\eta_0}\gamma^{-1},\beta^{-1}\}\delta^{-2}\epsilon^{-2})\nonumber\\
        \delta\epsilon&\leq \sqrt{\min\{\frac{C_0^2}{560m^2\hat{\Delta}_{\theta_0,\eta_0}},\frac{C_0^2\beta}{560m^2\gamma}}\}=\min\{\mathcal{O}({1}/{m\sqrt{\hat{\Delta}_{\theta_0,\eta_0}}}),\mathcal{O}(\sqrt{{\beta}/{\gamma}})\}
\end{align}

Then we are ready to present the following Convergence result.



Under the same parameter choices stated in Lemma~\ref{lemma: descent lemma of alg1-3}, we have the following high-probability convergence result.
\begin{restatable}[Formal Statement of Theorem~\ref{thm: convergence of alg3}]{theorem}{mainthmthree}\label{restatethm: convergence of alg3}
     Let Assumption \ref{assumption 1} hold. Denote the initial objective gap as $\Delta_{\theta_0,\eta_0}
    := \max_{i\in[m]}\{\fhL^i(\theta_0,\eta_0^i)-\fhL^{i,} \}$,  $\hL_0=L+G^2M\lambda^{-1}$,$\hL_1=\fL(G\sqrt{m})^{-1}$,$\hL_2=G^2Mm\lambda^{-1}$.
     Let $\delta,\epsilon$ satisfy $\delta\epsilon\leq 1/m\cdot\min\{\mathcal{O}({1}/{\Delta_{\theta_0,\eta_0}^{1/2}}),\mathcal{O}((\beta/\gamma)^{1/2}\}$. Set the hyperparameters in Algorithm~\ref{alg1-3} as $c_1,f_1={1}/{2}$, $c_2,f_2=\delta\epsilon$, 
    $\rho=\min\{\mathcal{O}(\delta^2\epsilon^2),\mathcal{O}(\frac{1}{\gamma T})\}$, $\beta=\mathcal{O}(\frac{1}{mG^2+m\Lambda^2})$ and $\gamma = \min\big\{ \mathcal{O}(\frac{1}{\hL_0+\hL_1\Lambda}), \mathcal{O}(\frac{1}{\hL_2+\vartheta\hL_1^2\hL_2^2}),\mathcal{O}(\beta),\mathcal{O}(\frac{1}{(mG^2+m\Lambda^2)\beta T\delta^2\epsilon^2}) \big\}$.
    Choose batch sizes $N_1=\mathcal{O}(\max\{G^2M^2\lambda^{-2}\kappa^2,\Lambda^2/m \}\delta^{-3}\epsilon^{-2})$, $N_2= \mathcal{O}(mG^2M^2\lambda^{-2}\kappa^2\delta^{-3}\epsilon^{-2})$. Then, after
    $T
    =\max\{ \Theta(\Delta_{\theta_0,\eta_0}\gamma^{-1}\delta^{-2}\epsilon^{-2}),\Theta(\beta^{-1}\delta^{-2}\epsilon^{-2})\}$ iterations, we have 
    \begin{align}
        \frac{1}{T}\sum_{t=0}^{T-1}\|\nabla_{\theta,\eta}\fhL(\theta_t, \eta_{t+1})w_t\|\leq 34\epsilon
        ,\nonumber
    \end{align}
    holds with probability at least $1-\delta$. 
\end{restatable}
\begin{proof} \textbf{Part I: Tail Event $\{\htau<T \}$}. Denote $\mathcal{F}_t$ as $\sigma$-algebra of initialization, and batch samples before time $t$, i.e., $\sigma\big(\theta_0, \eta_0, w_0, \{\xi_s,\bar{\xi}_s\}_{s<t} \big)$. Define event $\calhatE_t = \{\forall i\in[m], \fhL(\theta_t,\eta_t^i) - \hL^{i,*}\leq \bar{F} \}$, and stopping time
    \begin{align}
        \htau_1 &= \min\{t|\exists i\in[m],\|\HGamma^i_{t}\|\text{ or }\|\HUpsilon_t^i\|\geq \frac{C_0}{6\delta\epsilon\sqrt{\gamma m}} \} \wedge T \nonumber\\
        \htau_2 &= \min\{ t|\exists i\in[m], \fhL^i(\theta_{t+1},\eta_{t+1}^i)-\fhL^{i,*}>\bar{F}\}\wedge T\nonumber\\
        \htau &= \min\{\htau_1,\htau_2\}.
    \end{align}
We partition event $\{\htau<T \}$ to be $A= \{\htau_1\leq \htau_2, \htau_1< T \}$ and $B = \{\htau_2< \htau_1, \htau_2<T \}$. 
For $\mathbb{E}[\mathbf{1}(\calhatE_t)\|\HGamma_t^i\|^2]$ 
\begin{align}
    \mathbb{E}\Big[ \mathbf{1}(\calhatE_t) \|\hat{\Gamma}_{t}^i\|^2\Big] &=\mathbb{E}[\mathbb{E}[\mathbf{1}(\calhatE_t) \|\hat{\Gamma}_{t}^i\|^2|\mathcal{F}_t]] \overset{}{\leq} \Big(\frac{\hat{K}_0}{N_1}+\frac{2\hL_2^2\gamma^2f_2^2\hat{K_1}}{N_1}+\frac{2\Lambda^2\hat{K}_1}{N_1}\Big),
\end{align}
where the derivation follows ~\eqref{eq: variance intermediate bound alg3}. Similar arguments gives $\mathbb{E}[\mathbf{1}(\calhatE_t)\|\HUpsilon^i_t\|^2]\leq \frac{\hat{K_2}}{N_2}$. Then, for event $A$, since $\htau_1\leq \htau_2$, we have $\|\nabla_{\eta}\fhL^i(\theta_t,\eta_t^i)\|\leq \Lambda,\forall i\in[m]$ still holds. By Markov inequality, we know
\begin{align}
    P(\calhatE_t \cap\{\|\HGamma_{t}^i\|\geq \frac{C_0}{6\delta\epsilon\sqrt{\gamma   m}}\}) &= P(\calhatE_t\cap \{ \|\HGamma_{t}^i\|^2\geq \frac{{C_0^2}}{36m\gamma \delta^2\epsilon^2 }\})\nonumber\\
    &\leq \frac{36m\gamma\delta^2\epsilon^2}{{C_0^2}}(\frac{\hat{K}_0}{N_1}+\frac{2\hat{K}_1\gamma^2\hL_2^2\delta^2\epsilon^2}{N_1}+\frac{2\hat{K}_1\Lambda^2}{N_1}),\nonumber\\
    P(\calhatE_t \cap \{\|\HUpsilon_t^i\|\geq \frac{C_0}{6\delta\epsilon\sqrt{\gamma m}}\})&= P(\calhatE_t\cap \{ \|\HUpsilon_t^i\|^2\geq \frac{{C_0^2}}{36\delta^2\epsilon^2\gamma m}\})\leq \frac{36 m\gamma\delta^2\epsilon^2\hat{K}_2}{{C_0^2}N_2},
\end{align}
Then for event $A$, its probability $P(A)$ can be bounded as
\begin{align}
    P(A)&\leq P\Big(\bigcup_{t=0}^{T-1} \bigcup_{i=1}^m\Big[\calhatE_t \cap\big(\{\|\HGamma_t^i\| \geq \frac{C_0}{6\delta\epsilon\sqrt{\gamma m}}\} \cup\{\|\HUpsilon_t^i\| \geq \frac{C_0}{6\delta\epsilon\sqrt{\gamma m}}\}\}\big)\Big]\Big)\nonumber\\
    &\leq \sum_{i=1}^{m}\sum_{t=0}^{T-1}\big( P(\calhatE_t\cap \|\HGamma_{t}^i\|\geq \frac{C_0}{6\delta\epsilon\sqrt{\gamma   m}})+P(\calhatE_t\cap \|\HUpsilon_t^i\|\geq \frac{C_0}{6\delta\epsilon\sqrt{\gamma m}})\big)\nonumber\\
    &\leq \frac{36m^2\gamma T\delta^2\epsilon^2}{{C_0^2}}(\frac{\hat{K}_0}{N_1}+\frac{2\hat{K}_1\gamma^2\hL_2^2\delta^2\epsilon^2}{N_1}+\frac{2\hat{K}_1\Lambda^2}{N_1}+\frac{\hat{K}_2}{N_2})\leq \frac{\delta}{4},
\end{align}
since the pre-chosen $\delta,\epsilon$ satisfies ${56m^2\delta^4\epsilon^4\gamma T\leq{C_0^2}}$.

For event $B$, we know $\htau=\htau_2<T$, by setting index $i$ with  $w^i=1, w_{m/[i]}=0$, evaluate \eqref{eq: one-step moo descent alg3} at $t=\htau$,
we have
\begin{align}
&\fhL^i(\theta_{\htau+1},\eta_{\htau+1})\nonumber\\
&\leq \fhL^i(\theta_{\htau},\eta_{\htau}) +\gamma \mu_t\|\HUpsilon_t^i\|\| Z_tw_t\|+\gamma\alpha_t\| \HGamma_t^i\|\|X_tw_t\|+\frac{\gamma}{\beta}+\frac{3\gamma\beta\rho^2}{2}+2\gamma\rho+ 3\gamma\beta\delta^2\epsilon^2\|\HGamma_t\|_F^2\nonumber\\
&+9\gamma\beta\Xi'_1\delta^2\epsilon^2+3\beta\gamma\delta^2\epsilon^2\|\HUpsilon_t\|_F^2+3\gamma\beta\delta^2\epsilon^2m\Lambda^2+\frac{\gamma^4\delta^4\epsilon^4}{16\vartheta}\nonumber\\
&\leq \fhL^i(\theta_{\htau},\eta_{\htau})+\gamma \delta\epsilon\|\HUpsilon_t^i\|+\gamma\delta\epsilon\|\HGamma_t^i\|+\frac{\gamma}{\beta}+\frac{3\gamma\beta\rho^2}{2}+2\gamma\rho+ 3\gamma\beta\delta^2\epsilon^2\|\HGamma_t\|_F^2\nonumber\\
&+9\gamma\beta\Xi_1'\delta^2\epsilon^2+3\beta\gamma\delta^2\epsilon^2\|\HUpsilon_t\|_F^2+3\gamma\beta\delta^2\epsilon^2m\Lambda^2+\frac{\gamma^4\delta^4\epsilon^4}{16\vartheta}\nonumber\\
&\leq \fhL^i(\theta_{\htau},\eta_{\htau})+\frac{C_0}{3}\sqrt{\frac{\gamma}{m}}+\frac{\gamma}{\beta}+\frac{3\gamma\beta\rho^2}{2}+2\gamma\rho+\frac{C_0^2\beta}{6}+9\gamma\beta\Xi_1'\delta^2\epsilon^2+3\gamma\beta\delta^2\epsilon^2m\Lambda^2+\frac{\gamma^4\delta^4\epsilon^4}{16\vartheta},
\end{align}
where the first inequality leverages Cauchy-schwarz inequality to upper bound $\gamma\mu_t\langle\HUpsilon_tw,Z_tw_t\rangle$ and $\gamma\alpha_t\langle\HGamma_tw,X_tw_t \rangle$; the second inequality leverages clipping rule $\alpha_t<\frac{\delta\epsilon}{\|X_tw_t\|},\mu_t\leq \frac{\delta\epsilon}{\|Z_tw_t\|}$; the last inequality leverages fact $\|\HGamma_t^i\|,\|\HUpsilon_t^i\|\leq \frac{C_0}{6\delta\epsilon\sqrt{\gamma m}}$ since $\htau=\htau_2<\htau_1$.

Since we know, for ${0<\delta<1}$, we have
\begin{align}
    \frac{\gamma}{\beta}+\frac{3\gamma\beta\rho^2}{2}+2\gamma\rho+9\gamma\beta\Xi_1'\delta^2\epsilon^2+3\gamma\beta\delta^2\epsilon^2m\Lambda^2+\frac{\gamma^4\delta^4\epsilon^4}{16\vartheta}\leq d_3'+\dots+d_8'\leq\frac{\delta \bar{F}}{8}<\frac{\bar{F}}{2}.
\end{align}
By setting $\gamma,\beta$ satisfying${\gamma\leq \frac{9me_1^2}{C_0^2},\beta\leq \frac{6e_2}{C_0^2}}$,
we conclude $ \fhL^i(\theta_{\tau+1},\eta_{\tau+1})-\fhL^i(\theta_\tau,\eta_\tau)\leq \frac{\bar{F}}{2}$ since $e_1+e_2+\frac{\gamma}{\beta}+2\gamma\rho+9\gamma\beta\Xi_1\delta^2\epsilon^2+3\gamma\beta\delta^2\epsilon^2m\Lambda^2+\frac{3\gamma\beta\rho^2}{2}+\frac{\gamma^4\delta^4\epsilon^4}{16\vartheta}\leq \frac{F}{2}$.

However, we have $\max_{i\in[m]}\fhL^i(\theta_{\htau+1},\eta^i_{\htau+1})-\hL^{i,*}\geq \bar{F}$, thus, for this task, we have $\max_{i\in [m]}\fhL(\theta_{\htau},\eta_{\htau}^i)-\hL^{i,*}>\frac{\bar{F}}{2}$.
Leveraging Lemma~\ref{lemma: descent lemma of alg1-3} and Markov inequality, this further implies
\begin{align}
    P(B)&\leq P(\max_{i\in [m]}\fhL^i(\theta_{\tau},\eta_{\tau})-\fhL^{i,*}>\frac{\bar{F}}{2})\nonumber\\
    &\leq P(\bigcup_{i=1}^{m} \fhL^i(\theta_{\tau},\eta_{\tau})-\fhL^{i,*}>\frac{\bar{F}}{2}) \leq \sum_{i=1}^{m} P(\fhL^i(\theta_{\tau},\eta_{\tau})-\fhL^{i,*}>\frac{\bar{F}}{2})
    \leq m\cdot\frac{\bar{F}\delta/8m}{\bar{F}/2}= \frac{\delta}{4}.\nonumber
\end{align}
Thus, we have $P(\htau < T)\leq P(A)+P(B)\leq \frac{\delta}{2}$. 

\textbf{Part II: Convergence of $\frac{\gamma}{2T}\mathbb{E}[(\sum_{t=0}^{T-1}\alpha_t\|X_tw_t\|^2+\sum_{t=0}^{T-1}\mu_t\|Z_tw_t\|^2)\mathbf{1}(\htau = T)]$}.
Reorganize and divide both sides with $T$ on ~\eqref{eq: descent lemma alg1-3} in Lemma~\ref{lemma: descent lemma of alg1-3}, leveraging specified choices of parameter $\gamma$,$\beta$,$\rho$, and $N_1, N_2$. We have
%For $\frac{\gamma}{2T}\mathbb{E}[\sum_{t=0}^{T-1}\alpha_t\|X_tw_t\|^2+\sum_{t=0}^{T-1}\mu_t\|Z_tw_t\|^2|\htau = T]$, we have 
\begin{align}
    &\frac{\gamma}{2T}\mathbb{E}\big[\big(\sum_{t=0}^{T-1}\alpha_t\|X_tw_t\|^2+\sum_{t=0}^{T-1}\mu_t\|Z_tw_t\|^2\big)\mathbf{1}(\htau = T)\big] \nonumber\\
    & \leq \frac{\gamma}{2T}\mathbb{E}\big[\big(\sum_{t=0}^{\htau-1}\alpha_t\|X_tw_t\|^2+\sum_{t=0}^{\htau-1}\mu_t\|Z_tw_t\|^2\big)\big] \nonumber\\
    %& = \frac{\gamma}{2T}{P(\tau = T)} \mathbb{E}\big[\sum_{t=0}^{\tau-1}\alpha_t\|X_tw_t\|^2+\sum_{t=0}^{\tau-1}\mu_t\|Z_tw_t\|^2\big]\nonumber\\
    &\leq \frac{\hat{\Delta}_{\theta_0,\eta_0}}{T} +\gamma \delta\epsilon\sqrt{\hat{K}_2/N_2}+\gamma \delta\epsilon\Xi_3'+\frac{\gamma}{\beta T}+\frac{3\gamma\beta\rho^2}{2} +2\gamma\rho\nonumber\\
    &
    + 3\gamma\beta \delta^2\epsilon^2\Xi_2' + 9\gamma\beta\Xi_1'\delta^2\epsilon^2 + 3m\beta\gamma{\delta^2\epsilon^2}(\hat{K}_2/N_2)+3m\gamma\beta{\delta^2\epsilon^2}\Lambda^2 + \frac{\gamma^4\delta^4\epsilon^4}{16\vartheta}
    \nonumber\\
    &\leq (9\cdot \frac{1}{10}+2)\gamma\delta^2\epsilon^2 \leq 3\gamma\delta^2\epsilon^2,
\end{align}
 

This further implies
\begin{align}
    \frac{1}{T}\mathbb{E}[\sum_{t=0}^{T-1}\alpha_t\|X_tw_t\|^2\mathbf{1}(\htau =T)]\leq 6\delta^2\epsilon^2, \text{ and } \frac{1}{T}\mathbb{E}[\sum_{t=0}^{T-1}\mu_t{\|Z_tw_t\|^2}\mathbf{1}(\htau=T)]\leq 6 \delta^2\epsilon^2,\nonumber
\end{align}
Leveraging $\min\{\frac{a^2}{2},a\}\geq a-\frac{1}{2}$, we have
\begin{align}
    \frac{1}{T}\mathbb{E}[\sum_{t=0}^{T-1}\alpha_t\|X_tw_t\|^2\mathbf{1}(\htau = T)] &=\frac{1}{T}\mathbb{E}\big[\sum_{t=0}^{T-1}\delta^2\epsilon^2 \min\{\frac{\|X_tw_t\|^2}{2\delta^2\epsilon^2},\frac{\|X_tw_t\|}{\delta\epsilon} \} \mathbf{1}(\htau = T)\big]\nonumber\\
    &\geq\frac{\delta\epsilon}{T}\mathbb{E}[\sum_{t=0}^{T-1}\|X_tw_t\|\mathbf{1}(\htau = T)]-\frac{\delta^2\epsilon^2}{2} P(\tau= T)\nonumber,
\end{align}
Similar lower bound also applies to $\frac{1}{T}\mathbb{E}[\sum_{t=0}^{T-1}\mu_t\|Z_tw_t\|^2\mathbf{1}(\htau=T)]$.
This implies 
\begin{align}
    \frac{1}{T}\mathbb{E}[\sum_{t=0}^{T-1}\|X_tw_t\|\mathbf{1}(\htau=T)] \text{ and } \frac{1}{T}\mathbb{E}[\sum_{t=0}^{T-1}\|Z_tw_t\|\mathbf{1}(\htau = T)]\leq (6\delta\epsilon+\frac{\delta\epsilon}{2}P(\tau= T))\leq 7\delta\epsilon.\nonumber
\end{align}
Further more, for $\mathbb{E}[\|\HUpsilon_tw_t\|^2\mathbf{1}(\htau = T)]$, leveraging Lemma~\ref{lemma: conditioned GammaXi}, we have
\begin{align}
   \mathbb{E}[\|\HUpsilon_tw_t\|\mathbf{1}(\htau = T)]\leq \delta\epsilon, \text{ and } \mathbb{E}[\|\HGamma_tw_t\|\mathbf{1}(\htau = T)]\leq \delta\epsilon.\nonumber
\end{align}
%where we leverage variance bound of $\HUpsilon_t^i$ stated in ~\eqref{eq: variance bound for rescaled function} and choice of $N_2=18\delta^{-3}\epsilon^{-2}$.
This further implies
\begin{align}
\frac{1}{T}\mathbb{E}[\sum_{t=0}^{T-1}\| \nabla_{\theta}\fhL(\theta_t,\eta_{t+1})w_t\|\mathbf{1}(\htau = T)]\leq \frac{1}{T}\mathbb{E}[\sum_{t=0}^{T-1}\| \HGamma_tw_t\|\mathbf{1}(\htau = T)] + \frac{1}{T}\mathbb{E}[\sum_{t=0}^{T-1}\|X_tw_t\|\mathbf{1}(\htau = T)]\leq 8\delta\epsilon,\nonumber
\end{align}
and
\begin{align}
     \frac{1}{T}\mathbb{E}[\sum_{t=0}^{T-1}\|\nabla_{\eta}\fhL(\theta_t,\eta_t)w_t\|\mathbf{1}(\htau=T)] \leq \frac{1}{T}\mathbb{E}[\sum_{t=0}^{T-1}\| \HUpsilon_tw_t\|\mathbf{1}(\htau = T)] + \frac{1}{T}\mathbb{E}[\sum_{t=0}^{T-1}\|Z_tw_t\|\mathbf{1}(\htau = T)]\leq 8\delta\epsilon.
\end{align}
Additionally, we know
\begin{align}
    &\frac{1}{T}\mathbb{E}[\sum_{t=0}^{T-1}\|\nabla_{\eta}\fhL(\theta_t,\eta_{t+1})w_t\|\mathbf{1}(\htau = T)]\nonumber\\
    &\leq \frac{1}{T}\mathbb{E}[ \sum_{t=0}^{T-1}\|\nabla_{\eta}\fhL(\theta_t,\eta_{t+1})w_t - \nabla_{\eta}\fhL(\theta_t,\eta_t)w_t\|\mathbf{1}(\htau = T)]+\frac{1}{T}\mathbb{E}[\sum_{t=0}^{T-1}\|\nabla_{\eta}\fhL(\theta_t,\eta_t)w_t\|\mathbf{1}(\htau = T)]\nonumber\\
    &\leq\frac{1}{T}\mathbb{E}[\sum_{t=0}^{T-1}\sqrt{\sum_{i=1}^{m}\hL_2^2\gamma^2\mu_t^2(Z_t^iw_t^i)^2}\mathbf{1}(\htau = T)] + \frac{1}{T}\mathbb{E}[\sum_{t=0}^{T-1}\|\nabla_{\eta}\fhL(\theta_t,\eta_t)w_t\|\mathbf{1}(\htau= T)]\nonumber\\
    &\overset{(i)}{=} \frac{1}{T}\mathbb{E}[\sum_{t=0}^{T-1}\hL_2\gamma\mu_t\|Z_tw_t\|\mathbf{1}(\htau = T)]+ \frac{1}{T}\mathbb{E}[\sum_{t=0}^{T-1}\|\nabla_{\eta}\fhL(\theta_t,\eta_t)w_t\|\mathbf{1}(\htau= T)]\nonumber\\
    &\overset{(ii)}{\leq} \delta\epsilon+8\delta\epsilon = 9\delta\epsilon,
\end{align}
where (i) utilizes $\hL_2$-smooth property of $\nabla_{\eta^i}\fhL(\theta_t,\eta_t)$, and update equation $\eta_{t+1}^i-\eta_t^i=-\gamma \mu_tZ_t^iw_t^i$; (ii) utilizes $\mu_t\leq \frac{\delta\epsilon}{\|Z_tw_t \|}$ and $\gamma\leq \frac{1}{\hL_2}$.
Thus, in conclusion, we have
\begin{align}
    \frac{1}{T}\mathbb{E}[(\sum_{t=0}^{T-1}\|\nabla_{\theta}\fhL(\theta_t,\eta_{t+1})w_t\|+\|\nabla_{\eta}\fhL(\theta_t,\eta_{t+1})w_t \|)\mathbf{1}(\htau = T)]\leq 17\delta\epsilon,
\end{align}
By Markov inequality, this implies
\begin{align}
    P(
    \frac{1}{T}[\sum_{t=0}^{T-1}\|\nabla_{\theta}\fhL(\theta_t,\eta_{t+1})w_t\|+\|\nabla_{\eta}\fhL(\theta_t,\eta_{t+1})w_t \|]>34{\epsilon},\htau = T)\leq \frac{\delta}{2}.
\end{align}
Thus we have
\begin{align}
    &P(\frac{1}{T}[\sum_{t=0}^{T-1}\|\nabla_{\theta}\fhL(\theta_t,\eta_{t+1})w_t\|+\|\nabla_{\eta}\fhL(\theta_t,\eta_{t+1})w_t \|]{>} 34{\epsilon})\nonumber\\
    &\leq P(\tau<T) + P(
    \frac{1}{T}[\sum_{t=0}^{T-1}\|\nabla_{\theta}\fhL(\theta_t,\eta_{t+1})w_t\|+\|\nabla_{\eta}\fhL(\theta_t,\eta_{t+1})w_t \|]>34{\epsilon},\htau = T)\nonumber\\
    &\leq \frac{\delta}{2} + \frac{\delta}{2}= \delta,
\end{align}
This implies $P(\frac{1}{T}[\sum_{t=0}^{T-1}\|\nabla_{\theta,\eta}\fhL(\theta_t,\eta_{t+1})w_t\|]\leq 34{\epsilon})\geq 1-\delta$, which completes proof.
\end{proof}

\subsection{Descent Lemma of Algorithm~\ref{alg1-3}}\label{Appendix: Descent Lemma of Algorithm 1-3}

\begin{restatable}[Descent Lemma of Algorithm \ref{alg1-3}]{lemma}{maindescentlemmathmtwo}\label{lemma: descent lemma of alg1-3}
Under the same hyper-parameters choices stated in Theorem \ref{thm: convergence of alg3}, then for fixed $w\in \mathcal{W}$, we have
    \begin{align}
        &\mathbb{E}\big[\fhL(\theta_{\htau}, \eta_{\htau})w\big]-\fhL^*w \leq \frac{\bar{F}\delta}{8m}-\mathbb{E}\Big[\frac{\gamma}{2}\sum_{t=0}^{\htau-1}\mu_t\|Z_tw_t\|^2+\frac{\gamma}{2}\sum_{t=0}^{\htau-1}\alpha_t\|X_tw_t\|^2 \Big],
        \label{eq: descent lemma alg1-3}
    \end{align}
holds.
\end{restatable}



\begin{proof} 
%Starting with MGDA update rule described in Algorithm~\ref{alg1-3},  
Denote $\mathcal{F}_t$ as $\sigma$-algebra of initialization, and batch samples before time $t$, i.e., $\sigma\big(\theta_0, \eta_0, w_0, \{\xi_s,\bar{\xi}_s\}_{s<t} \big)$. Define event $\calhatE_t = \{\forall i\in[m], \fhL(\theta_t,\eta_t^i) - \hL^{i,*}\leq \bar{F} \}$. Since $\calhatE_t \in \mathcal{F}_t$ and $\hat{\tau}_2$ detects an excessive objective gap at the next iterate, we have $\{t<\hat{\tau}\} \subseteq \calhatE_t$. Thus, on $\calhatE_t$, $\|\nabla_{\eta^i}\fhL(\theta,\eta^i)\|$ is bounded by $\Lambda$.

For term ${-\alpha_t \langle {X}_tw, X_tw_t \rangle-\mu_t\langle Z_tw,Z_tw_t\rangle}$, 
Expanding $\|w_{t+1}-w\|^2$, we have
\begin{align}
    \|{w}_{t+1} - {w} \|^2&= \| \Pi_{\mathcal{W}} \big({w}_t - \beta \big[{\alpha_t{X}_t^{\top}X_t{w}_t}+
   \mu_t{{Z}_t^{\top}Z_t{w}_t}+\rho {w}_t \big]\big) - {w} \|^2\nonumber\\
    &\leq \|{w}_t - \beta \big[{\alpha_t{X}_t^{\top}X_t{w}_t}+\mu_t{{Z}_t^{\top}Z_t{w}_t}+\rho {w}_t \big]- {w} \|^2 \nonumber\\
    &= \|{w}_t - {w} \|^2 -2\beta \langle {w}_t-{w}, {\alpha_t{X_t}^{\top}X_t{w}_t}+\rho {w}_t \rangle -2\beta\langle w_t-w,\mu_t{{Z}_t^{\top}Z_t{w}_t} \rangle \nonumber\\
    &\quad+ \beta^2\|\alpha_t{{X}_t^{\top}X_t{w}_t}+\rho {w}_t + \mu_t{{Z}_t^{\top}Z_t{w}_t} \|^2\nonumber \\
    &\leq \|{w}_t - {w} \|^2 -2\beta{\alpha_t} \langle {w}_t-{w}, {{X_t}^{\top}X_t{w}_t} \rangle -2\beta\mu_t\langle w_t-w,{{Z}_t^{\top}Z_t{w}_t} \rangle \nonumber\\
    &\quad + 3\beta^2\alpha_t^2 \|{{X}_t^{\top} X_t w_t}\|^2+{3\beta^2\mu_t^2}\|{Z}_t^{\top}Z_tw_t\|^2+3\beta^2 \rho^2 + {4}\beta\rho,
    \label{eq: coro MGDA from 1st principle}
\end{align}
where the first inequality is due to non-expansiveness of projection over probability simplex; the second inequality applies $(a+b+c)^2\leq 3a^2+3b^2+3c^2$ and $\|w\|\|w_t\|\leq 1$.

For $\|{{X}_t^{\top}X_t{w}_t}\|^2$, we further decompose it as
\begin{align}
    &\|{{X}_t^{\top}X_t{w}_t}\|^2 \nonumber\\
    &= \|{\big(\underbrace{{X}_t-\nabla_{\theta}\fhL(\theta_t, \eta_{t+1})}_{\hat{\Gamma}_{t}}+\nabla_{\theta}\fhL(\theta_t, \eta_{t+1})\big)^{\top}X_t{w}_t\|^2}\nonumber\\
    &\leq  2\|\HGamma_{t}\|^2_F\|X_tw_t \|^2 + 2\|\nabla_{\theta}\fhL(\theta_t, \eta_{t+1})\|_F^2 \|X_tw_t\|^2\nonumber\\
    & \leq 2\|\HGamma_{t}\|^2_F\|X_tw_t \|^2 + (6mG^2+6\gamma^2\mu_t^2\hL_2^2\|Z_tw_t\|^2+6m\Lambda^2)\|X_tw_t\|^2\nonumber\\
    &\leq2\|\HGamma_{t}\|^2_F\|X_tw_t \|^2 + (6mG^2+6\gamma^2\hL_2^2{f_2}^2+6m\Lambda^2)\|X_tw_t\|^2,
    \label{eq: upper bound MGDA pre-conditioning}
\end{align}
where the last three inequalities leverages Cauchy-Schwarz inequality, Multiplicative property of Frobenius norm, $(a+b)^2\leq 2a^2+2b^2$, and \eqref{eq: gradient relationship 2} respectively. 

Similarly, $\|{Z}_t^{\top}Z_tw_t\|^2$ can be upper bounded as
\begin{align}
    \|{Z}_t^{\top}Z_tw_t\|^2&=\|({\HUpsilon}_t+\nabla_{\eta}\fhL(\theta_t,\eta_t))^{\top}Z_tw_t\|^2\nonumber\\
    &\leq 2\|{\HUpsilon}_t^{\top}Z_tw_t\|^2+2\|\nabla_{\eta}\fhL(\theta_t,\eta_t)^{\top}Z_tw_t\|^2\nonumber\\
    &\leq 2\| \HUpsilon_t\|^2_F\|Z_tw_t\|^2+2m\Lambda^2\|Z_tw_t\|^2,
    \label{eq: upper bound MGDA pre-conditioning 2}
\end{align}
where above inequalities leverages $(a+b)^2\leq 2a^2+2b^2$; Cauchy-Schwarz inequality, sub-multiplicative property of Frobenius norm and $\|\nabla_{\eta}\fhL(\theta_t,\eta_t)\|_F\leq \sqrt{m}\Lambda$, respectively.

Merge \eqref{eq: upper bound MGDA pre-conditioning}, \eqref{eq: upper bound MGDA pre-conditioning 2} into \eqref{eq: coro MGDA from 1st principle}, re-organize it and multiply both sides with $\gamma/2\beta$, we have
\begin{align}
    &-\gamma\alpha_t\langle X_t{w},X_t{w}_t \rangle -\gamma{\mu_t}\langle{Z}_tw, Z_tw_t \rangle\nonumber \\
    &\leq -\gamma\alpha_t \|X_tw_t\|^2-\gamma\mu_t\| Z_tw_t\|^2+\frac{\gamma}{2\beta}(\|w_t-w\|^2-\|w_{t+1}-w\|^2)+\frac{3\gamma\beta\rho^2}{2}+2\gamma\rho\nonumber\\
    &\quad +3\beta\gamma\alpha_t^2 \|\HGamma_{t}\|_F^2 \|X_t w_t\|^2+9\beta\gamma\alpha_t^2\underbrace{(mG^2+\gamma^2\hL_2^2{f_2^2}+m\Lambda^2)}_{\Xi'_1}\|X_tw_t\|^2\nonumber\\
    &\quad+ 3\beta\gamma\mu_t^2\| {\HUpsilon}_t\|^2_F\|Z_tw_t\|^2+3\beta\gamma\mu_t^2m\Lambda^2\|Z_tw_t\|^2\nonumber\\
    &\leq-\gamma\alpha_t \|X_tw_t\|^2-\gamma\mu_t\| Z_tw_t\|^2+\frac{\gamma}{2\beta}(\|w_t-w\|^2-\|w_{t+1}-w\|^2)+\frac{3\gamma\beta\rho^2}{2}+2\gamma\rho \nonumber\\
    &\quad+ 3\gamma\beta c_2^2\|\HGamma_t\|_F^2 + 9\gamma\beta\Xi_1'c_2^2 + 3\beta\gamma{f_2^2}\|\HUpsilon_t\|_F^2+3\beta\gamma{f_2^2}m\Lambda^2,
    \label{eq: coro MDGA}
\end{align}
where the last inequality leverages clipping structure $\alpha_t = \min\{c_1,\frac{c_2}{\|X_tw_t\|} \} \text{ and }\mu_t = \min\{f_1,\frac{f_2}{\|Z_tw_t\|} \}$.

Next, for descent lemma with respect to $\nabla_{\eta}\fhL(\theta_t, \eta_t)$~\eqref{eq: moo descent lemma eta}, Merge $\eta_{t+1}=\eta_t-\gamma\mu_tZ_tw_t$, we have
\begin{align}
    \fhL(\theta_t,\eta_{t+1})w
    &\leq \fhL(\theta_t, \eta_t)w - {\gamma\mu_t} \langle \nabla_{\eta}\fhL(\theta_t, \eta_t)w,Z_tw_t \rangle + \frac{\hL_2{\gamma^2\mu_t^2}}{2}\|Z_tw_t\|^2\nonumber\\
    &=\fhL(\theta_t, \eta_t)w +{\gamma\mu_t} \langle{\HUpsilon}_tw,Z_tw_t\rangle-{\gamma\mu_t}\langle{Z}_tw,Z_tw_t \rangle + \frac{\hL_2\gamma^2\mu_t^2}{2}\|Z_tw_t\|^2,
    \label{eq: coro descent 1}
\end{align}
where the last equality decomoposes $\nabla_{\eta}\fhL(\theta_t,\eta_t)w=(\nabla_{\eta}\fhL(\theta_t,\eta_t)-Z_t+Z_t)w = (-\HUpsilon_t+Z_t)w$.

For descent lemma with respect to $\nabla_{\theta}\fhL(\theta_t, \eta_{t+1})$~\eqref{eq: moo descent lemma theta}, Merge update rule $\theta_{t+1} = \theta_{t}-\gamma\alpha_tX_tw_t$ into descent lemma \eqref{eq: moo descent lemma theta}, we have
\begin{align}
    &\fhL({\theta}_{t+1}, {\eta}_{t+1}){w}\nonumber\\
    &\leq \fhL({\theta}_{t}, {\eta}_{t+1}){w} + \langle \nabla_{\theta}\fhL({\theta}_{t}, {\eta}_{t+1})w,{\theta}_{t+1}-{\theta}_t \rangle + \frac{\hL_0}{2}\|{\theta}_{t+1}-{\theta}_t \|^2+ \frac{\hL_1 | \nabla_{\eta} \fhL({\theta}_t, {\eta}_{t+1})w|}{2}\|{\theta}_{t+1}-{\theta}_t \|^2  \nonumber \\
    &=  \fhL({\theta}_t, {\eta}_{t+1}){w} -\gamma\alpha_t \langle \nabla_{\theta}\fhL({\theta}_t, {\eta}_{t+1})w,X_tw_t \rangle + \frac{\gamma^2\alpha_t^2\hL_0}{2}\|X_tw_t\|^2 \nonumber\\
    &\quad + \frac{\gamma^2\alpha_t^2\hL_1 | \nabla_{\eta} \fhL({\theta}_t, {\eta}_{t+1})w|}{2} \|X_tw_t\|^2\nonumber\\
    &\leq \fhL({\theta}_t, {\eta}_{t+1}){w} -\gamma\alpha_t \langle (\underbrace{\nabla_{\theta}\fhL({\theta}_t, {\eta}_{t+1})-{{X}_t}}_{-\HGamma_t})w,X_tw_t \rangle - \gamma\alpha_t \langle {{X}_tw},X_tw_t \rangle+ \frac{\hL_0}{2}\alpha_t^2\gamma^2\|X_tw_t\|^2  \nonumber\\
    &\quad  + \frac{\hL_1\Lambda+{\gamma\mu_t\hL_1\hL_2\|Z_tw_t\|}}{2}\alpha_t^2\gamma^2\| X_tw_t \|^2,
    \label{eq: coro descent 2}
\end{align}
where the last inequality utilizes  $|\nabla_{\eta} \fhL(\theta_t, \eta_{t+1})w|\leq \hL_2\gamma\mu_t\|Z_tw_t\| + |\nabla_{\eta}\fhL(\theta_t, \eta_t)w|\leq\hL_2\gamma\mu_t\|Z_tw_t\|+\Lambda$ via gradient-clipping update,
$\eta_{t+1}=\eta_t-\gamma\mu_tZ_tw_t$.

Merge \eqref{eq: coro MDGA},\eqref{eq: coro descent 1}, \eqref{eq: coro descent 2}, we have
\begin{align}
    &\fhL(\theta_{t+1},\eta_{t+1})w\nonumber\\
    &\leq \fhL(\theta_t, \eta_t)w +{\gamma\mu_t} \langle{\HUpsilon}_tw,Z_tw_t\rangle+\gamma\alpha_t \langle\HGamma_tw,X_tw_t \rangle-{\gamma\mu_t}\langle{Z}_tw,Z_tw_t \rangle - \gamma\alpha_t \langle {{X}_tw},X_tw_t \rangle \nonumber\\
    &\quad + \frac{\hL_0}{2}\alpha_t^2\gamma^2\|X_tw_t\|^2+ \frac{\hL_2\gamma^2\mu_t^2}{2}\|Z_tw_t\|^2+ \frac{\hL_1\Lambda+{\gamma\mu_t \hL_1\hL_2\|Z_tw_t\|}}{2}\alpha_t^2\gamma^2\| X_tw_t \|^2\nonumber\\
    &\leq  \fhL(\theta_t, \eta_t)w +{\gamma\mu_t} \langle{\HUpsilon}_tw,Z_tw_t\rangle+\gamma\alpha_t \langle\HGamma_tw,X_tw_t \rangle -\gamma\alpha_t \|X_tw_t\|^2-\gamma\mu_t\| Z_tw_t\|^2 \nonumber\\
    &\quad +\frac{\gamma}{2\beta}(\|w_t-w\|^2-\|w_{t+1}-w\|^2)+\frac{3\gamma\beta\rho^2}{2}+2\gamma\rho + 3\gamma\beta c_2^2\|\HGamma_t\|_F^2 + 9\gamma\beta\Xi_1'c_2^2+ 3\beta\gamma{f_2^2}\|\HUpsilon_t\|_F^2\nonumber\\
    &\quad +3\beta\gamma{f_2^2}m\Lambda^2+ \frac{\hL_0}{2}\alpha_t^2\gamma^2\|X_tw_t\|^2+ \frac{\hL_2\gamma^2\mu_t^2}{2}\|Z_tw_t\|^2+ \frac{\hL_1\Lambda+{\gamma\mu_t \hL_1\hL_2\|Z_tw_t\|}}{2}\alpha_t^2\gamma^2\| X_tw_t \|^2\nonumber\\
    &\leq \fhL(\theta_t, \eta_t)w +{\gamma\mu_t} \langle{\HUpsilon}_tw,Z_tw_t\rangle+\gamma\alpha_t \langle\HGamma_tw,X_tw_t \rangle -\gamma\alpha_t \|X_tw_t\|^2-\gamma\mu_t\| Z_tw_t\|^2 \nonumber\\
    &\quad +\frac{\gamma}{2\beta}(\|w_t-w\|^2-\|w_{t+1}-w\|^2)+\frac{3\gamma\beta\rho^2}{2}+2\gamma\rho + 3\gamma\beta c_2^2\|\HGamma_t\|_F^2 + 9\gamma\beta\Xi_1'c_2^2+ 3\beta\gamma{f_2^2}\|\HUpsilon_t\|_F^2\nonumber\\
    &\quad +3\gamma\beta{f_2^2}m\Lambda^2+ \frac{\hL_0}{2}\alpha_t^2\gamma^2\|X_tw_t\|^2+ \frac{\hL_2\gamma^2\mu_t^2}{2}\|Z_tw_t\|^2+ \frac{\hL_1\Lambda}{2}\alpha_t^2\gamma^2\| X_tw_t \|^2 \nonumber\\
    &\quad+\vartheta\gamma^2\mu_t^2\hL_1^2\hL_2^2\|Z_tw_t\|^2+\frac{\gamma^4c_2^4}{16\vartheta}\nonumber\\
    &=\fhL(\theta_t,\eta_t)w+{\gamma\mu_t} \langle{\HUpsilon}_tw,Z_tw_t\rangle+\gamma\alpha_t \langle\HGamma_tw,X_tw_t \rangle-\gamma\alpha_t(1-\frac{\hL_0+\hL_1\Lambda}{2}\gamma\alpha_t)\|X_tw_t\|^2\nonumber\\
    &\quad - \gamma\mu_t(1-\frac{\hL_2+2\vartheta\hL_1^2\hL_2^2}{2}\gamma\mu_t)\|Z_tw_t\|^2+\frac{\gamma}{2\beta}(\|w_t-w\|^2-\|w_{t+1}-w\|^2)+\frac{3\gamma\beta\rho^2}{2}+2\gamma\rho  \nonumber\\
    &\quad+ 3\gamma\beta c_2^2\|\HGamma_t\|_F^2+ 9\gamma\beta\Xi_1'c_2^2+ 3\beta\gamma{f_2^2}\|\HUpsilon_t\|_F^2+3\gamma\beta{f_2^2}m\Lambda^2 + \frac{\gamma^4c_2^4}{16\vartheta},
    \label{eq: moo descent alg3 intermediate}
\end{align}
where we apply young's inequality $\frac{ab}{2}\leq  \vartheta a^2+\frac{b^2}{16\vartheta}$ on $\frac{\hL_1\hL_2\gamma\mu_t\|Z_tw_t\|\cdot\gamma^2\alpha_t^2\|X_tw_t\|^2}{2}$,  and clipping rule ${\alpha_t = \min\{c_1, \frac{c_2}{\|X_tw_t\|} \}}\leq c_1 \text{ and }\frac{c_2}{\|X_tw_t\|}$ and $\mu_t$.

For $\gamma\mu_t\langle \HUpsilon_tw,Z_tw_t\rangle+\gamma\alpha_t\langle \HGamma_tw, X_tw_t\rangle$, we have
\begin{align}
    \gamma\mu_t\langle \HUpsilon_tw,Z_tw_t\rangle+\gamma\alpha_t\langle \HGamma_tw, X_tw_t\rangle&\leq \gamma\mu_t\| \HUpsilon_tw\| \|Z_tw_t\| + \gamma\alpha_t\|\HGamma_tw\|\|X_tw_t\|\nonumber\\
    &\leq \gamma {f_2}\|\HUpsilon_tw\|+\gamma c_2\| \HGamma_tw\|,
    \label{eq: non-martingale-sequence alg3}
\end{align}
where the first inequality utilizes Cauchy-schawarz inequality and the second inequality utilizes $\alpha_t\leq \frac{c_2}{\|X_tw_t\|}$ and $\mu_t\leq \frac{f_2}{\| Z_tw_t\|}$. This further reduces \eqref{eq: moo descent alg3 intermediate} into
\begin{align}
    &\fhL(\theta_{t+1},\eta_{t+1})w\nonumber\\
    &\leq \fhL(\theta_t,\eta_t)w+\gamma {f_2}\|\HUpsilon_tw\|+\gamma c_2\| \HGamma_tw\|-\gamma\alpha_t(1-\frac{\hL_0+\hL_1\Lambda}{2}\gamma\alpha_t)\|X_tw_t\|^2\nonumber\\
    &\quad - \gamma\mu_t(1-\frac{\hL_2+2\vartheta\hL_1^2\hL_2^2}{2}\gamma\mu_t)\|Z_tw_t\|^2+\frac{\gamma}{2\beta}(\|w_t-w\|^2-\|w_{t+1}-w\|^2)+\frac{3\gamma\beta\rho^2}{2}+2\gamma\rho  \nonumber\\
    &\quad+ 3\gamma\beta c_2^2\|\HGamma_t\|_F^2+ 9\gamma\beta\Xi_1'c_2^2+ 3\beta\gamma{f_2^2}\|\HUpsilon_t\|_F^2+3\gamma\beta{f_2^2}m\Lambda^2 + \frac{\gamma^4c_2^4}{16\vartheta}.
    \label{eq: one-step moo descent alg3}
\end{align}
Taking conditional expectation over $\xi_t$, for $\|\HUpsilon_tw\|$, we have
\begin{align}
    \mathbf{1}(\calhatE_t)\mathbb{E}_{\xi_t}[\|\HUpsilon_tw\||\mathcal{F}_t]\leq \mathbf{1}(\calhatE_t) \sqrt{\mathbb{E}_{\xi_t}[\|\HUpsilon_tw\|^2|\mathcal{F}_t]}\leq \mathbf{1}(\calhatE_t)\sqrt{\hat{K}_2/N_2},
    \label{eq: variance bound thm 3}
\end{align}
where the last inequality follows from variance upper bound of $\HUpsilon_t^i,\forall i\in[m]$~\eqref{eq: variance bound for rescaled function}.
%Following proof logic of ~\eqref{eq: estimation error bound thm2}, taking conditional expectation over $\bar{\xi}_t$, we have
%\begin{align}
    %\mathbb{E}\big[\mathbb{E}_{\bar{\xi}_t,\xi_t,\eta_{t+1}}[\|\HGamma_tw\||\theta_t,\eta_{t},w_t,t]\big]&\leq \mathbb{E}\big[\sqrt{\mathbb{E}_{\bar{\xi}_t,\xi_t}[\|\HGamma_tw\|^2|\theta_t,\eta_{t},w_t,t]}\big]\nonumber\\
    %&\leq \mathbb{E}\Big[\sqrt{\mathbb{E}_{\xi_t}\big[\frac{\hat{K}_0}{N_1}+\frac{2\hat{K}_1\hL_2^2\gamma^2\mu_t^2}{N_1}\|Z_tw_t\|^2+\frac{2\hat{K}_1}{N_1}\Lambda^2|\theta_t,\eta_t,w_t,t\big]}\Big]\nonumber\\
    %&\leq \sqrt{\frac{\hat{K}_0}{N_1}+\frac{2\hat{K}_1\hL_2^2\gamma^2f_2^2}{N_1}+\frac{2\hat{K}_1\Lambda^2}{N_1}}\nonumber\\
    %&\leq \underbrace{\sqrt{\frac{\hat{K}_0}{N_1}}+\sqrt{\frac{2\hat{K}_1\hL_2^2\gamma^2f_2^2}{N_1}}+\sqrt{\frac{2\hat{K}_1\Lambda^2}{N_1}}}_{\Xi'_3},
    %\label{eq: variance bound of HGamma alg3}
%\end{align}
%where the first inequality leverages variance upper bound of $\HGamma_t^i,\forall i\in[m]$~\eqref{eq: variance bound for rescaled function}, derivation of \eqref{eq: estimation error bound thm2}; the second inequality utilizes gradient clipping rule, $\mu_t\leq \frac{f_2}{\|Z_tw_t\|}$; the last inequality leverages $\sqrt{a+b+c}\leq \sqrt{a}+\sqrt{b}+\sqrt{c}$.

Following similar proof logic as~\eqref{eq: estimation error bound thm2},  taking conditional expectation over ${\xi}_t,\bar{\xi}_t$, $ \|\hat{\Gamma}_{t}\|_F^2$ becomes
\begin{align}
    &\mathbf{1}(\calhatE_t)\mathbb{E}_{{\xi}_t,\bar{\xi}_t,\eta_{t+1}}\Big[  \|\hat{\Gamma}_{t}\|_F^2|\mathcal{F}_t\Big]\nonumber\\
    &\overset{(i)}{\leq} \mathbf{1}(\calhatE_t)\mathbb{E}_{{\xi}_t,\bar{\xi}_t,\eta_{t+1}}\Big[\frac{m\hat{K}_0}{N_1}+\frac{\hat{K}_1}{N_1}\|\nabla_{\eta}\fhL(\theta_t, \eta_{t+1})\|_F^2|\mathcal{F}_t\Big]\nonumber\\
    &\overset{}{\leq} \mathbf{1}(\calhatE_t) \mathbb{E}_{{\xi}_t}\Big[ \frac{m\hat{K}_0}{N_1}+\frac{\hat{K}_1}{N_1}(2\hL_2^2\gamma^2\mu_t^2\|Z_tw_t\|^2+2m\Lambda^2))|\mathcal{F}_t\Big]\nonumber\\
    &\overset{(iii)}{\leq} \mathbf{1}(\calhatE_t) \mathbb{E}[
    \frac{m\hat{K}_0}{N_1}+\frac{\hat{K}_1}{N_1}(2\hL_2^2\gamma^2\mu_t^2\|Z_tw_t\|^2+2m\Lambda^2))|\mathcal{F}_t]\nonumber\\
    &\overset{(iv)}{\leq}  \mathbf{1}(\calhatE_t)\big(\underbrace{\frac{m\hat{K}_0}{N_1}+\frac{2\hL_2^2\gamma^2f_2^2\hat{K_1}}{N_1}+\frac{2m\Lambda^2\hat{K}_1}{N_1}}_{\Xi'_2}\big),
    \label{eq: variance intermediate bound alg3}
\end{align}
where the (i) applies the fact $\mathbb{E}_{\bar{\xi}_t^i}[\|\HGamma_{t}^i\|^2|\mathcal{F}_t,\xi_t]\leq (\hat{K}_0+\hat{K}_1\|\nabla_{\eta^i}\fhL(\theta_t,\eta^i_{t+1})\|^2)/N_1$; (iii) further upper bound $\|\nabla_{\eta}\fhL(\theta_t, \eta_{t+1})\|^2_F=\sum_{i=1}^{m}\|\nabla_{\eta^i}\fhL(\theta_t, \eta_{t+1}^i)\|^2\leq \sum_{i=1}^{m}(\hL_2\gamma\mu_t \|Z_t^iw_t^i\|+\Lambda)^2\leq 2\hL_2^2\gamma^2\mu_t^2\|Z_tw_t\|^2+2m\Lambda^2$; (iii) leverages $\mu_t\leq \frac{f_2}{\|Z_tw_t\|}$. 

Similarly, repeating above conditional on $\mathcal F_t$, and using the gradient bound on $\calhatE_t$, we obtain, for any fixed \(w\in\mathcal W\),
$$
\mathbf1({\calhatE_t})
\mathbb E[\|\widehat\Gamma_t w\|\mid\mathcal F_t]
\le \mathbf1({\widehat{\mathcal E}_t})\Xi'_3.
$$

Combining \eqref{eq: variance bound thm 3},\eqref{eq: estimation error bound thm2},\eqref{eq: variance intermediate bound alg3} into \eqref{eq: one-step moo descent alg3}, sum it from $0$ to $\htau-1$. Since $\{t< \htau\}\subseteq \calhatE_t$, $\mathbf{1}(\calhatE_t)=1$, {take expectation on both sides} and utilizing $\|w_0-w\|^2\leq 2$, we have
\begin{align}
    &\mathbb{E}[\fhL(\theta_{\htau},\eta_{\htau})w]-\fhL^*w\nonumber\\
    &\leq \fhL(\theta_0,\eta_0)w-\fhL^*w-\mathbb{E}[\sum_{t=0}^{\htau-1}\gamma\alpha_t(1-\frac{\hL_0+\hL_1\Lambda}{2}\gamma\alpha_t)\|X_tw_t\|^2]\nonumber\\
    & -\mathbb{E}[\sum_{t=0}^{\htau-1}\gamma\mu_t(1-\frac{\hL_2+2\vartheta\hL_1^2\hL_2^2}{2}\gamma\mu_t)\|Z_tw_t\|^2] +\mathbb{E}[\sum_{t=0}^{\htau-1}\gamma{f_2}\|\HUpsilon_tw\|]+\mathbb{E}[\sum_{t=0}^{\htau-1}\gamma  c_2\| \HGamma_tw\|]\nonumber\\
    &+\frac{\gamma}{\beta}+\frac{3\gamma\beta\rho^2T}{2}
    +2\gamma\rho T + \mathbb{E}[\sum_{t=0}^{\htau-1}3\gamma\beta c_2^2\|\HGamma_t\|_F^2] + \mathbb{E}[\sum_{t=0}^{\htau-1}[9\gamma\beta\Xi_1'c_2^2]+ \mathbb{E}[\sum_{t=0}^{\htau-1}3\beta\gamma{f_2^2}\|\HUpsilon_t\|_F^2]\nonumber\\
    &+\mathbb{E}[\sum_{t=0}^{\htau-1}3\beta\gamma{f_2^2}m\Lambda^2]+ \frac{\gamma^4c_2^4T}{16\vartheta}\nonumber\\
    &\leq \fhL(\theta_0,\eta_0)w-\fhL^*w-\mathbb{E}[\sum_{t=0}^{\htau-1}\gamma\alpha_t(1-\frac{\hL_0+\hL_1\Lambda}{2}\gamma\alpha_t)\|X_tw_t\|^2] \nonumber\\
    &-\mathbb{E}[\sum_{t=0}^{\htau-1}\gamma\mu_t(1-\frac{\hL_2+2\vartheta\hL_1^2\hL_2^2}{2}\gamma\mu_t)\|Z_tw_t\|^2]+\gamma T{f_2}\sqrt{\hat{K}_2/N_2}+\gamma c_2\Xi_3'T+\frac{\gamma}{\beta}+\frac{3\gamma\beta\rho^2T}{2} +2\gamma\rho T\nonumber\\
    &
    + 3\gamma\beta c_2^2\Xi_2' T + 9\gamma\beta\Xi_1'c_2^2 T+ 3m\beta\gamma{f_2^2}(\hat{K}_2/N_2)T+3m\beta\gamma{f_2^2}\Lambda^2T + \frac{\gamma^4c_2^4T}{16\vartheta}\nonumber\\
    &\leq \fhL(\theta_0,\eta_0)w-\fhL^*w-\mathbb{E}[\frac{\gamma}{2}\sum_{t=0}^{\htau-1}\alpha_t\|X_tw_t\|^2] -\mathbb{E}[\frac{\gamma}{2}\sum_{t=0}^{\htau-1}\mu_t\|Z_tw_t\|^2] +\underbrace{\gamma T{f_2}\sqrt{\hat{K}_2/N_2}}_{\leq d'_1}+\underbrace{\gamma c_2\Xi_3'T}_{\leq d'_2}\nonumber\\
    &+\underbrace{\frac{\gamma}{\beta}}_{\leq d'_3}+\underbrace{2\gamma\rho T}_{\leq d'_4}+\underbrace{\frac{3\gamma\beta\rho^2T}{2}}_{\leq d'_5} 
    + \underbrace{3\gamma\beta c_2^2\Xi_2' T + 9\gamma\beta\Xi_1'c_2^2 T}_{\leq d'_6}+ \underbrace{3m\beta\gamma{f_2^2}(\hat{K}_2/N_2)T+3m\beta\gamma{f_2^2}\Lambda^2T}_{\leq d'_7} \nonumber\\
    &+ \underbrace{\frac{\gamma^4c_2^4T}{16\vartheta}}_{\leq d'_{8}}\nonumber\\
    &\leq \Delta_{\theta_0, \eta_0} -\mathbb{E}[\frac{\gamma}{2}\sum_{t=0}^{\htau-1}\alpha_t\|X_tw_t\|^2] -\mathbb{E}[\frac{\gamma}{2}\sum_{t=0}^{\htau-1}\mu_t\|Z_tw_t\|^2] + d_1'+\dots d_8'    \nonumber\\
    &\leq \frac{\bar{F}\delta}{8m} -\frac{\gamma}{2}\mathbb{E}[\sum_{t=0}^{\htau-1}\alpha_t\| X_tw_t\|^2]-\frac{\gamma}{2}\mathbb{E}[\sum_{t=0}^{\htau-1}\mu_t\|Z_tw_t\|^2],
    \label{eq: stopping-time bound alg3}
\end{align}
where the first inequality $\mathbb{E}[\sum_{t=0}^{\htau-1}{\|\hat{\Gamma}}\|_F^2] \leq \mathbb{E}[\sum_{t=0}^{T-1}{\mathbf{1}(\calhatE_t)\|\hat{\Gamma}}\|_F^2] \leq \sum_{t=0}^{T-1}\mathbb{E}[\mathbf{1}(\calhatE_t)\|\HGamma\|_F^2]\leq \sum_{t=0}^{T-1}\mathbb{E}[\mathbf{1}(\calhatE_t)\mathbb{E}[\|\HGamma\|_F^2|\mathcal{F}_t]]\leq T\Xi_2'$, similar arguments also applies to upper bound $\mathbb{E}[\sum_{t=0}^{\htau-1}\|\HGamma_tw\|]\leq T\Xi_3'$,  $\mathbb{E}[\sum_{t=0}^{\htau-1}\|\HUpsilon_t\|_F^2]\leq Tm\hat{K}_2/N_2$, $\mathbb{E}[\sum_{t=0}^{\htau-1}\|\HUpsilon_tw\|]\leq T\sqrt{\hat{K}_2/N_2}$; The last inequality follows the parameter choice of $\gamma,\beta, \rho, \alpha_t,\mu_t$, $N_1$, $N_2$, construction of $\bar{F}$, and $\htau \leq T$.
\end{proof}

\begin{lemma}\label{lemma: conditioned GammaXi}
   For event $\{\htau = T \}$, under the parameter choices as Theorem~\ref{thm: convergence of alg3}, we have
   \begin{align}
        \mathbb{E}[\|\HGamma_tw_t\|\mathbf{1}(\htau = T)]\leq \delta\epsilon \text{ and }\mathbb{E}[\|\HUpsilon_tw_t\|\mathbf{1}(\htau = T)]\leq \delta\epsilon.
   \end{align}
\end{lemma}
\begin{proof}
    Notice $\theta_t, \eta_t$ is $\mathcal{F}_t$-measurable. Moreover, since $\calhatE_t$ only depends on $\theta_t,\eta_t$, we have $\calhatE_t \in \mathcal{F}_t$.
    At $\htau = T$, every state encountered before $T$ is safe. Therefore we conclude $\{\htau = T \}\subseteq \calhatE_t$ and $\mathbf{1}(\htau = T)\leq \mathbf{1}(\calhatE_t)$. Thus, applying tower property, we have
    \begin{align}
        \mathbb{E}[\|\HGamma_tw_t\|^2\mathbf{1}(\htau = T)] &\leq \mathbb{E}[\|\HGamma_tw_t\|^2\mathbf{1}(\calhatE_t)]\nonumber\\
        &= \mathbb{E}[\mathbb{E}[\|\HGamma_tw_t\|^2\mathbf{1}(\calhatE_t)|\mathcal{F}_t]]\nonumber\\
        & = \mathbb{E}[\mathbf{1}(\calhatE_t)\mathbb{E}[\|\HGamma_tw_t\|^2|\mathcal{F}_t]]\nonumber\\
        & =  \mathbb{E}[\mathbf{1}(\calhatE_t)\mathbb{E}[\mathbb{E}[\|\HGamma_tw_t\|^2|\mathcal{F}_t,\xi_t]|\mathcal{F}_t]],
    \end{align}
    For $\mathbb{E}[\|\HGamma_tw_t\|^2|\mathcal{F}_t,\xi_t]$, leveraging variance bound Lemma~\ref{lemma: affine bounded noise}, we know
    \begin{align}
        \mathbb{E}[\|\HGamma_t^i\|^2 |\mathcal{F}_t, \xi_t] \leq \frac{\hat{K}_0+\hat{K}_1|\nabla_{\eta^i} \fhL^i(\theta_t, \eta_{t+1}^i)|^2}{N_1}\leq \frac{\hat{K}_0+2\hat{K}_1( \Lambda^2+{\hL}_2^2 \gamma^2 f_2^2) }{N_1} \text{ on } \calhatE_t,
    \end{align}
    where the last inequality leverages the $\hL_2$-smooth in terms of $\nabla_{\eta}\fhL^i(\theta_t, \eta_{t})$ such that $|\nabla_{\eta^i} \widehat{L}^i\left(\theta_t, \eta_{t+1}^i\right)| \leq \Lambda+\hL_2 \|\eta_{t+1}-\eta_t\|$ and clipping step-size design such that $\|\eta_{t+1}-\eta_t\|\leq \gamma \mu_t \|Z_tw_t\|\leq \gamma f_2$.
    Substitute parameter choice of $N_1 = \max\big\{9\hat{K}_0\delta^{-3}\epsilon^{-2}, 18\hat{K}_1\delta^{-1}, 18\Lambda^2\hat{K}_1\delta^{-3}\epsilon^{-2}\}$ gives 
    \begin{align}
        \mathbb{E}[\|\HGamma_tw_t\|^2\mathbf{1}(\htau=T)]\leq \frac{\hat{K}_0+2\hat{K}_1( \Lambda^2+{\hL}_2^2 \gamma^2 f_2^2) }{N_1} P(\calhatE_t)\leq \frac{\delta^3\epsilon^2}{3}.
    \end{align}
    Leveraging Cauchy-schwarz inequality, we have
    \begin{align}
        \mathbb{E}[\|\HGamma_tw_t\|\mathbf{1}(\htau=T)] 
        &\leq \sqrt{\mathbb{E}[(\|\HGamma_tw_t\|^2\mathbf{1}(\htau = T))P(\htau=T)}\nonumber\\
        &\leq \frac{\delta^{3/2}\epsilon}{\sqrt{3}}\leq \delta\epsilon,
    \end{align}
    which completes proof. Similar arguments also applies to $\mathbb{E}[\|\HUpsilon_tw_t\|\mathbf{1}(\tau = T)]$.
\end{proof}





%
%



%%% END INSTRUCTIONS %%%






\end{document}
